% =========================================================================
% 《泛函分析（第 5–8 章）期末闭卷冲刺与作业详解全集 v2》
%
% 编译方式（需 MacTeX / TeX Live，执行三遍以收敛目录、书签与总页数）：
%   xelatex -interaction=nonstopmode 泛函分析_期末闭卷冲刺与作业详解_v2.tex
%   xelatex -interaction=nonstopmode 泛函分析_期末闭卷冲刺与作业详解_v2.tex
%   xelatex -interaction=nonstopmode 泛函分析_期末闭卷冲刺与作业详解_v2.tex
% =========================================================================

\documentclass[11pt,a4paper]{article}

% ------------------------- 页面几何与留白 -------------------------
\usepackage[
  a4paper,
  left=21mm,
  right=21mm,
  top=23mm,
  bottom=22mm,
  headheight=18pt,
  headsep=6mm,
  footskip=10.5mm
]{geometry}

% ------------------------- 核心宏包 -------------------------
\usepackage{amsmath,amssymb,mathtools}
\usepackage{calc}
\newcounter{none}
\usepackage{xeCJK}
\usepackage{fontspec}
\usepackage{xcolor}
\usepackage{tikz}
\usetikzlibrary{positioning}
\usepackage[most]{tcolorbox}
\usepackage{booktabs,longtable,array,colortbl}
\usepackage{enumitem}
\usepackage{fancyhdr}
\usepackage{lastpage}
\usepackage{needspace}
\usepackage{adjustbox}
\usepackage[normalem]{ulem}
\usepackage{tocloft}
\tocloftpagestyle{fancy}
\usepackage[
  colorlinks=true,
  linkcolor=LinkColor,
  urlcolor=LinkColor,
  bookmarks=true,
  bookmarksopen=true,
  bookmarksopenlevel=2,
  bookmarksnumbered=false,
  pdfauthor={},
  pdfcreator={},
  pdfproducer={},
  pdftitle={泛函分析（第5-8章）期末闭卷冲刺与作业详解全集 v2}
]{hyperref}

% ------------------------- 配色系统 -------------------------
\definecolor{PaperBg}{HTML}{FFFFFF}
\definecolor{InkColor}{HTML}{1F2329}
\definecolor{MutedInk}{HTML}{656D76}
\definecolor{BorderColor}{HTML}{DFE4EA}
\definecolor{CodeBg}{HTML}{F3F5F8}
\definecolor{LinkColor}{HTML}{1E5AA8}

% 五篇专属学术主题色（第 0 篇深墨蓝 + 第 5-8 章分章色）
\definecolor{Ch0Main}{HTML}{273C56}
\definecolor{Ch0Soft}{HTML}{EFF3F8}
\definecolor{Ch5Main}{HTML}{264B7F}
\definecolor{Ch5Soft}{HTML}{EDF3FA}
\definecolor{Ch6Main}{HTML}{176B62}
\definecolor{Ch6Soft}{HTML}{EAF6F4}
\definecolor{Ch7Main}{HTML}{943A3B}
\definecolor{Ch7Soft}{HTML}{FAEEEE}
\definecolor{Ch8Main}{HTML}{55408A}
\definecolor{Ch8Soft}{HTML}{F1EEF9}

% 辅助功能色（思路/警告/减负）
\definecolor{AmberMain}{HTML}{B36B1D}
\definecolor{AmberSoft}{HTML}{FDF4E7}
\definecolor{CoralMain}{HTML}{B8433E}
\definecolor{CoralSoft}{HTML}{FCEEEE}
\definecolor{LightMain}{HTML}{5A6872}
\definecolor{LightSoft}{HTML}{F1F4F7}

\colorlet{AccentMain}{Ch0Main}
\colorlet{AccentSoft}{Ch0Soft}

\newcommand{\SwitchChapterTheme}[1]{%
  \ifnum#1=0\colorlet{AccentMain}{Ch0Main}\colorlet{AccentSoft}{Ch0Soft}\fi
  \ifnum#1=5\colorlet{AccentMain}{Ch5Main}\colorlet{AccentSoft}{Ch5Soft}\fi
  \ifnum#1=6\colorlet{AccentMain}{Ch6Main}\colorlet{AccentSoft}{Ch6Soft}\fi
  \ifnum#1=7\colorlet{AccentMain}{Ch7Main}\colorlet{AccentSoft}{Ch7Soft}\fi
  \ifnum#1=8\colorlet{AccentMain}{Ch8Main}\colorlet{AccentSoft}{Ch8Soft}\fi
}

% ------------------------- 字体与阅读舒适度 -------------------------
\setmainfont{Avenir Next}[
  BoldFont=Avenir Next Demi Bold,
  ItalicFont=Avenir Next Italic,
  BoldItalicFont=Avenir Next Demi Bold Italic
]
\setmonofont{Menlo}[Scale=0.84]
\setCJKmainfont{PingFang SC}[
  BoldFont=PingFang SC Semibold,
  ItalicFont=Kaiti SC,
  BoldItalicFont=PingFang SC Semibold
]
\setCJKmonofont{PingFang SC}[Scale=0.90]
\xeCJKsetup{
  CJKecglue={\hskip 0.18em plus 0.06em minus 0.04em},
  PunctStyle=kaiming
}

\setcounter{secnumdepth}{0}
\setcounter{tocdepth}{3}
\color{InkColor}
\linespread{1.40}
\setlength{\parindent}{0pt}
\setlength{\parskip}{5.2pt}
\sloppy
\emergencystretch=1.8em
\widowpenalty=1000
\clubpenalty=1000
\allowdisplaybreaks[3]

% ------------------------- 目录美化 -------------------------
\renewcommand{\contentsname}{\textcolor{Ch0Main}{\Large\bfseries 全书交互目录（点击任意条目直达）}}
\renewcommand{\cftsecleader}{\cftdotfill{\cftdotsep}}
\renewcommand{\cftsecfont}{\bfseries\color{InkColor}}
\renewcommand{\cftsecpagefont}{\bfseries\color{Ch0Main}}
\renewcommand{\cftsubsecfont}{\small\color{InkColor}}
\renewcommand{\cftsubsecpagefont}{\small\color{MutedInk}}
\renewcommand{\cftsubsubsecfont}{\footnotesize\color{MutedInk}}
\renewcommand{\cftsubsubsecpagefont}{\footnotesize\color{MutedInk}}
\setlength{\cftbeforesecskip}{7.5pt}
\setlength{\cftbeforesubsecskip}{2.6pt}
\setlength{\cftbeforesubsubsecskip}{1.2pt}
\setlength{\cftsubsecindent}{12pt}
\setlength{\cftsubsubsecindent}{26pt}

% ------------------------- 行内代码与列表 -------------------------
\let\oldtexttt\texttt
\renewcommand{\texttt}[1]{%
  \colorbox{CodeBg}{\oldtexttt{\color{AccentMain}#1}}%
}
\setlist[itemize]{leftmargin=16pt, itemsep=2.6pt, topsep=3.5pt, parsep=1.2pt, label=\textcolor{AccentMain}{\scriptsize$\bullet$}}
\setlist[enumerate]{leftmargin=19pt, itemsep=2.6pt, topsep=3.5pt, parsep=1.2pt, label=\textcolor{AccentMain}{\bfseries\arabic*.}}
\providecommand{\tightlist}{\setlength{\itemsep}{2pt}\setlength{\parskip}{0pt}}

% 任务清单勾选框兼容（Pandoc task_lists）
\providecommand{\checkbox}{$\square$}

% ------------------------- 动态页眉与可点击返回页脚 -------------------------
\newcommand{\CurrentChapterHeader}{泛函分析期末闭卷冲刺全集 v2}
\pagestyle{fancy}
\fancyhf{}
\renewcommand{\headrulewidth}{0pt}
\renewcommand{\footrulewidth}{0pt}

\fancyhead[L]{%
  \begin{tikzpicture}[remember picture, overlay]
    \fill[AccentMain] (current page.north west) rectangle ([yshift=-2.8mm]current page.north east);
    \draw[BorderColor, line width=0.5pt]
      ([xshift=21mm, yshift=-17.5mm]current page.north west) -- ([xshift=-21mm, yshift=-17.5mm]current page.north east);
  \end{tikzpicture}%
  {\footnotesize\color{MutedInk}\CurrentChapterHeader}%
}
\fancyhead[R]{%
  {\footnotesize\color{MutedInk}\hyperlink{book:portal}{\textcolor{AccentMain}{$\star$ 27 题传送门}}\;\textbar\;\hyperlink{book:toc}{\textcolor{MutedInk}{$\uparrow$ 目录}}}%
}
\fancyfoot[L]{%
  \begin{tikzpicture}[remember picture, overlay]
    \draw[BorderColor, line width=0.5pt]
      ([xshift=21mm, yshift=16mm]current page.south west) -- ([xshift=-21mm, yshift=16mm]current page.south east);
  \end{tikzpicture}%
  {\footnotesize\color{MutedInk}《泛函分析》（科学出版社）闭卷冲刺全集 v2}%
}
\fancyfoot[R]{%
  {\footnotesize\color{MutedInk}第 }\textcolor{AccentMain}{\footnotesize\bfseries\thepage}{\footnotesize\color{MutedInk} 页 \,/\, 共 \pageref*{LastPage} 页}%
}

% ------------------------- 矢量徽章与跳转按钮 -------------------------
\newcommand{\WarnBadge}{%
  \tcbox[on line, colback=CoralSoft, colframe=CoralMain!45!BorderColor, boxrule=0.35pt,
    arc=3pt, left=3.5pt, right=3.5pt, top=0.5pt, bottom=0.5pt, boxsep=0pt]{\scriptsize\bfseries\color{CoralMain}注意}\,%
}
\newcommand{\HotBadge}{%
  \tcbox[on line, colback=AmberSoft, colframe=AmberMain!45!BorderColor, boxrule=0.35pt,
    arc=3pt, left=3.5pt, right=3.5pt, top=0.5pt, bottom=0.5pt, boxsep=0pt]{\scriptsize\bfseries\color{AmberMain}高频}\,%
}
\newcommand{\FastBadge}{%
  \tcbox[on line, colback=AccentSoft, colframe=AccentMain!45!BorderColor, boxrule=0.35pt,
    arc=3pt, left=3.5pt, right=3.5pt, top=0.5pt, bottom=0.5pt, boxsep=0pt]{\scriptsize\bfseries\color{AccentMain}速通}\,%
}
\newcommand{\StarBadge}{%
  \tcbox[on line, colback=AmberSoft, colframe=AmberMain!45!BorderColor, boxrule=0.35pt,
    arc=3pt, left=3.5pt, right=3.5pt, top=0.5pt, bottom=0.5pt, boxsep=0pt]{\scriptsize\bfseries\color{AmberMain}核心}\,%
}
\newcommand{\CircNum}[1]{%
  \tikz[baseline=(char.base)]{\node[shape=circle, fill=AccentMain, text=white, inner sep=1pt, font=\tiny\bfseries] (char) {#1};}\,%
}

% 轻重缓急重要度标签（帮零基础一眼分清必背 vs 减负！）
\newcommand{\TagMust}[1]{%
  \;\tcbox[on line, colback=CoralSoft, colframe=CoralMain!45!BorderColor, boxrule=0.35pt,
    arc=3pt, left=4pt, right=4pt, top=1pt, bottom=1pt, boxsep=0pt]{\scriptsize\bfseries\color{CoralMain}#1}%
}
\newcommand{\TagHigh}[1]{%
  \;\tcbox[on line, colback=AmberSoft, colframe=AmberMain!45!BorderColor, boxrule=0.35pt,
    arc=3pt, left=4pt, right=4pt, top=1pt, bottom=1pt, boxsep=0pt]{\scriptsize\bfseries\color{AmberMain}#1}%
}
\newcommand{\TagLight}[1]{%
  \;\tcbox[on line, colback=LightSoft, colframe=LightMain!35!BorderColor, boxrule=0.35pt,
    arc=3pt, left=4pt, right=4pt, top=1pt, bottom=1pt, boxsep=0pt]{\scriptsize\color{LightMain}#1}%
}

% 可点击跳转胶囊与传送门按钮
\newcommand{\JumpPill}[2]{%
  \hyperlink{#1}{\tcbox[on line, colback=AccentSoft, colframe=AccentMain!40!BorderColor, boxrule=0.35pt,
    arc=3pt, left=4.5pt, right=4.5pt, top=1pt, bottom=1pt, boxsep=0pt]{\scriptsize\bfseries\color{AccentMain}#2}}%
}
\newcommand{\PortalBtn}[3]{%
  \hyperlink{#1}{\tcbox[on line, colback=#2, colframe=#2, boxrule=0pt,
    arc=3.5pt, left=6.5pt, right=6.5pt, top=2.5pt, bottom=2.5pt, boxsep=0pt]{\scriptsize\bfseries\color{white}#3}}%
}
\newcommand{\ProbChip}[3]{%
  \hyperlink{#1}{\tcbox[on line, colback=#2!8!white, colframe=#2!38!BorderColor, boxrule=0.4pt,
    arc=3.5pt, left=5pt, right=5pt, top=2pt, bottom=2pt, boxsep=0pt]{\scriptsize\color{InkColor}#3}}%
}

\newcommand{\SoftDivider}{%
  \par\vspace{5pt}\noindent\makebox[\linewidth][c]{\textcolor{BorderColor}{\rule{0.38\linewidth}{0.6pt}}}\par\vspace{5pt}%
}

% ------------------------- 各级标题组件 -------------------------
% 注：\hypertarget 均置于 \clearpage / \Needspace 之后，确保跨页跳转永不偏移到上一页
\newcommand{\ChapterBanner}[4]{%
  \clearpage
  \gdef\CurrentChapterHeader{#3}%
  \phantomsection
  \hypertarget{#2}{}%
  \addcontentsline{toc}{section}{#3}%
  \begin{tcolorbox}[
    enhanced,
    colback=AccentMain,
    colframe=AccentMain,
    boxrule=0pt,
    arc=5pt,
    left=12pt, right=12pt, top=10pt, bottom=10pt,
    before skip=0pt, after skip=10pt
  ]
    {\scriptsize\bfseries\color{AccentSoft}PART / CHAPTER #1}\hfill{\scriptsize\hyperlink{book:portal}{\textcolor{AccentSoft}{$\star$ 返回传送门}}}\\[3pt]
    {\Large\bfseries\color{white}#3}\\[4pt]
    {\footnotesize\color{AccentSoft}#4}
  \end{tcolorbox}%
}

\newcommand{\BookPart}[3]{%
  \Needspace{8\baselineskip}%
  \vspace{9pt}%
  \phantomsection
  \if\relax\detokenize{#1}\relax\else\hypertarget{#1}{}\fi
  \addcontentsline{toc}{subsection}{\texorpdfstring{#2}{#3}}%
  \begin{tcolorbox}[
    enhanced,
    colback=AccentSoft,
    colframe=AccentMain,
    boxrule=0pt,
    leftrule=3.5pt,
    arc=3pt,
    left=9pt, right=9pt, top=6pt, bottom=6pt,
    before skip=4pt, after skip=7pt
  ]
    {\large\bfseries\color{AccentMain}#2}
  \end{tcolorbox}%
}

\newcommand{\BookProblem}[4]{%
  \Needspace{10\baselineskip}%
  \vspace{10pt}%
  \phantomsection
  \hypertarget{#1}{}%
  \addcontentsline{toc}{subsubsection}{\texorpdfstring{#2}{#4}}%
  \begin{tcolorbox}[
    enhanced,
    colback=AccentMain,
    colframe=AccentMain,
    boxrule=0pt,
    arc=4pt,
    left=9.5pt, right=9.5pt, top=5.5pt, bottom=5.5pt,
    before skip=4pt, after skip=6.5pt
  ]
    \begin{minipage}[c]{0.68\linewidth}
      {\normalsize\bfseries\color{white}#2}
    \end{minipage}\hfill
    \begin{minipage}[c]{0.30\linewidth}
      \raggedleft{\scriptsize\color{AccentSoft}#3\quad\hyperlink{book:portal}{\textcolor{white}{[$\uparrow$ 传送门]}}}
    \end{minipage}
  \end{tcolorbox}%
}

\newcommand{\BookSubSection}[4]{%
  \Needspace{9.5\baselineskip}%
  \vspace{8pt}%
  \phantomsection
  \if\relax\detokenize{#1}\relax\else\hypertarget{#1}{}\fi
  \addcontentsline{toc}{subsubsection}{\texorpdfstring{#2}{#3}}%
  \noindent\textcolor{AccentMain}{\rule[-1.5pt]{2.8pt}{11pt}}\hspace{6pt}%
  {\normalsize\bfseries\color{InkColor}#2}#4\par\nopagebreak[4]%
  \vspace{3.5pt}%
}

\newcommand{\SubPillPrimary}[1]{%
  \Needspace{4.5\baselineskip}\vspace{4.5pt}\noindent
  \tcbox[on line, colback=AccentSoft, colframe=AccentMain!50!BorderColor, boxrule=0.45pt,
    arc=3.5pt, left=6pt, right=6pt, top=2pt, bottom=2pt]{\small\bfseries\color{AccentMain}#1}%
  \par\vspace{3pt}%
}
\newcommand{\SubPillAmber}[1]{%
  \Needspace{4.5\baselineskip}\vspace{5.5pt}\noindent
  \tcbox[on line, colback=AmberSoft, colframe=AmberMain!45!BorderColor, boxrule=0.45pt,
    arc=3.5pt, left=6pt, right=6pt, top=2pt, bottom=2pt]{\small\bfseries\color{AmberMain}#1}%
  \par\vspace{3pt}%
}
\newcommand{\SubPillMain}[1]{%
  \Needspace{4.5\baselineskip}\vspace{5.5pt}\noindent
  \tcbox[on line, colback=AccentMain, colframe=AccentMain, boxrule=0pt,
    arc=3.5pt, left=6pt, right=6pt, top=2pt, bottom=2pt]{\small\bfseries\color{white}#1}%
  \par\vspace{3pt}%
}
\newcommand{\SubPillCoral}[1]{%
  \Needspace{4.5\baselineskip}\vspace{5.5pt}\noindent
  \tcbox[on line, colback=CoralSoft, colframe=CoralMain!45!BorderColor, boxrule=0.45pt,
    arc=3.5pt, left=6pt, right=6pt, top=2pt, bottom=2pt]{\small\bfseries\color{CoralMain}#1}%
  \par\vspace{3pt}%
}
\newcommand{\SubSubHeading}[1]{%
  \Needspace{6.5\baselineskip}\vspace{5.5pt}\noindent
  \textcolor{AccentMain}{\scriptsize$\blacktriangleright$}\hspace{4pt}%
  {\small\bfseries\color{AccentMain}#1}\par\nopagebreak[4]\vspace{2.5pt}%
}
\newcommand{\StepHeading}[1]{%
  \Needspace{3.5\baselineskip}\vspace{4.5pt}\noindent
  {\small\bfseries\color{AccentMain}#1}\par\vspace{2pt}%
}

% ------------------------- 可跨页引用框与防溢边公式 -------------------------
\renewenvironment{quote}{%
  \begin{tcolorbox}[
    enhanced,
    breakable,
    colback=AccentSoft!48!white,
    colframe=AccentMain,
    boxrule=0pt,
    leftrule=2.8pt,
    arc=2.5pt,
    left=9pt, right=8pt, top=6pt, bottom=6pt,
    before skip=5.5pt, after skip=5.5pt
  ]\small
}{%
  \end{tcolorbox}%
}

\newcommand{\FitDisplayMath}[1]{%
  \par\vspace{4.5pt}\noindent
  \begin{adjustbox}{max width=\linewidth, center}
    $\displaystyle #1$
  \end{adjustbox}
  \par\vspace{4.5pt}%
}
\newcommand{\FitTaggedMath}[2]{%
  \par\vspace{4.5pt}\noindent
  \begin{minipage}{\linewidth}
    \centering
    \begin{adjustbox}{max width=0.88\linewidth}
      $\displaystyle #1$
    \end{adjustbox}\hfill{\small\color{MutedInk}(#2)}
  \end{minipage}
  \par\vspace{4.5pt}%
}

\setlength{\LTpre}{6pt}
\setlength{\LTpost}{6pt}
\setlength{\heavyrulewidth}{0.85pt}
\setlength{\lightrulewidth}{0.45pt}
\setlength{\cmidrulewidth}{0.45pt}
\arrayrulecolor{BorderColor}

\begin{document}

% ========================= 精美 A4 封面页（v2） =========================
\thispagestyle{empty}
\begin{tikzpicture}[remember picture, overlay]
  \fill[Ch0Main] (current page.north west) rectangle ([yshift=-8.5mm]current page.north east);
  \fill[Ch0Soft, opacity=0.85] ([xshift=-12mm, yshift=-26mm]current page.north east) circle (42mm);
  \draw[Ch0Main, opacity=0.14, line width=1pt] ([xshift=-20mm, yshift=-34mm]current page.north east) circle (54mm);
  \node[anchor=north east, font=\fontsize{68}{68}\selectfont\bfseries, text=Ch0Main, opacity=0.09]
    at ([xshift=-15mm, yshift=-23mm]current page.north east) {v2};
  \draw[BorderColor, line width=0.7pt]
    ([xshift=21mm, yshift=22mm]current page.south west) -- ([xshift=-21mm, yshift=22mm]current page.south east);
  \node[anchor=west, font=\small, text=MutedInk]
    at ([xshift=21mm, yshift=15mm]current page.south west) {对照教材：《泛函分析》（科学出版社“十四五”规划教材）};
  \node[anchor=east, font=\small\bfseries, text=Ch0Main]
    at ([xshift=-21mm, yshift=15mm]current page.south east) {全书共 \pageref*{LastPage} 页 · 含双向交互跳转};
\end{tikzpicture}

\vspace*{14mm}
\noindent
\tcbox[on line, colback=Ch0Soft, colframe=Ch0Main!35!BorderColor, boxrule=0.5pt,
  arc=10pt, left=9pt, right=9pt, top=3.5pt, bottom=3.5pt]{\small\bfseries\color{Ch0Main}FUNCTIONAL ANALYSIS \;\textbar\; 期末闭卷冲刺全集 v2.0}
\hspace{4pt}
\tcbox[on line, colback=CoralSoft, colframe=CoralMain!45!BorderColor, boxrule=0.5pt,
  arc=10pt, left=8pt, right=8pt, top=3.5pt, bottom=3.5pt]{\small\bfseries\color{CoralMain}闭卷考试 · 平时 50\% + 期末 50\%}

\vspace{9mm}
{\fontsize{27}{36}\selectfont\bfseries\color{InkColor}泛函分析（第 5–8 章）\par}
\vspace{2mm}
{\fontsize{17.5}{25}\selectfont\bfseries\color{Ch0Main}末课划重点总纲 · 零基础精讲 · 27 道作业详解\par}

\vspace{3.5mm}
{\normalsize\color{MutedInk}以老师末课录音框架为纲 · 补全高频卡壳点白话解读 · 全书目录/题号/自测双向一键跳转\par}

\vspace{5.5mm}
\noindent\textcolor{Ch0Main}{\rule{34mm}{2.6pt}}

\vspace{8mm}
\begin{tcolorbox}[
  enhanced,
  colback=Ch0Soft!45!white,
  colframe=Ch0Main!25!BorderColor,
  boxrule=0.6pt,
  arc=6pt,
  left=12pt, right=12pt, top=10pt, bottom=10pt
]
  {\normalsize\bfseries\color{Ch0Main}v2 精编版五大板块（点击书内页眉随时回目录/传送门）}\\[5pt]
  \begin{itemize}[leftmargin=14pt, itemsep=5pt, topsep=2pt]
    \item \textcolor{Ch0Main}{\bfseries 第 0 篇 · 末课老师划重点总纲与录音整理}：4 条防挂科真相、TikZ 全书沙漏主线图、闭卷必背定义/证明/关系、12 类考场变形题通法、25 问自测表（全部带跳转）及 13:49 完整录音稿。
    \item \textcolor{Ch5Main}{\bfseries 第 5 章 · 紧算子和 Fredholm 算子}：含第 1–4 章前置速通、零化子空间 $\mathcal{N}(M)$ 解码、Kato 闭值域 5 步证明；详解习题 5 第 \textbf{5、6、14、15、20、22} 题（6 题）。
    \item \textcolor{Ch6Main}{\bfseries 第 6 章 · 有界线性算子的谱理论}：新增「谱点 vs 特征值 4 例白话拆解」与「谱测度 $E$ 与算子 $A$ 关系双例」，6.4 节标明减负提示；详解习题 6 第 \textbf{2、5、13、23、24、31、37、39、45} 题（9 题）。
    \item \textcolor{Ch7Main}{\bfseries 第 7 章 · Hilbert 空间上的无界算子}：稠定与共轭 $A^*$、例 7.2 三个求导算子对比、闭算子与自伴判别准则、无界谱分解；详解习题 7 第 \textbf{3、4、6、7} 题（4 题）。
    \item \textcolor{Ch8Main}{\bfseries 第 8 章 · 广义函数与 Sobolev 空间}：新增例 8.5 $\operatorname{p.v.}\frac{1}{x}$ 白话推导、本书 $2\pi$ 约定 Fourier 变换与 $H^s(\mathbb{R}^n)$；详解习题 8 第 \textbf{7、9、11、14、17(2)(3)、20、28} 题及 \textbf{p.337 练习 8.8.2}（8 题）。
  \end{itemize}
\end{tcolorbox}

\vspace{6mm}
{\small\color{MutedInk}
\textbf{零基础防挂科三步读法}：\CircNum{1} 先看紧随目录后的\textbf{「1 页极速传送门」}与\textbf{第 0 篇}抓牢全局框架与闭卷必背定义；\CircNum{2} 按「\textcolor{CoralMain}{\textbf{闭卷必背}} $\to$ \textcolor{AmberMain}{\textbf{高频思路}} $\to$ \textcolor{LightMain}{大概了解}」的徽章分级分配时间；\CircNum{3} 用第 0 篇的 25 问自测表盖住答案默写，卡壳处直接点按钮跳到正文精讲或作业思路！
}

\vfill
\clearpage

% ========================= 全书可点击目录 =========================
\hypertarget{book:toc}{}
\gdef\CurrentChapterHeader{全书交互目录}
\tableofcontents
\clearpage

% ========================= 1 页极速传送门（考务红线 + 27 题一键跳转） =========================
\hypertarget{book:portal}{}
\gdef\CurrentChapterHeader{考前极速传送门}
\phantomsection
\addcontentsline{toc}{section}{考前极速传送门（4 条红线与 27 题一键跳转）}

\begin{tcolorbox}[
  enhanced,
  colback=CoralSoft,
  colframe=CoralMain,
  boxrule=0pt,
  leftrule=3.6pt,
  arc=4pt,
  left=10pt, right=10pt, top=7pt, bottom=7pt,
  before skip=0pt, after skip=8pt
]
  {\normalsize\bfseries\color{CoralMain}末课录音 4 条防挂科铁律（先读完再翻书！）}\\[4pt]
  \begin{enumerate}[leftmargin=16pt, itemsep=2.5pt, topsep=1pt, parsep=0pt]
    \small
    \item \textbf{闭卷考试}（录音 13:45 实锤）：核心定义、算子关系、谱半径公式、核心短证明必须能默写！
    \item \textbf{平时 50\% + 期末 50\%}（录音 12:44）：平时分是安全垫，但期末卷面仍需认真拿稳基本盘。
    \item \textbf{老师亲自辟谣：不考作业原题，考作业背后的框架与思路！}（录音 11:51–12:39）——切忌死背答案，每道题重点吃透「破题思路」。
    \item \textbf{减负策略}（录音 05:20 / 07:45 / 11:39）：上课略讲内容不死磕；6.4 谱测度/谱积分\textbf{大概了解}即可；6.6 末拓展应用小节可跳过。\hfill\JumpPill{part:audio}{看完整 13:49 录音稿 $\to$}
  \end{enumerate}
\end{tcolorbox}

\begin{tcolorbox}[
  enhanced,
  colback=Ch0Soft!45!white,
  colframe=Ch0Main!25!BorderColor,
  boxrule=0.55pt,
  arc=4.5pt,
  left=10pt, right=10pt, top=7pt, bottom=7pt,
  before skip=0pt, after skip=8pt
]
  \raggedright
  {\normalsize\bfseries\color{Ch0Main}全书五大板块与高频专题一键直达}\par\vspace{5pt}
  \PortalBtn{part:0}{Ch0Main}{第 0 篇 · 末课总纲}\hspace{3pt}%
  \PortalBtn{part:12types}{AmberMain}{12 类考场通法}\hspace{3pt}%
  \PortalBtn{part:25q}{CoralMain}{25 问闭卷自测}\hspace{3pt}%
  \PortalBtn{part:audio}{MutedInk}{13:49 录音逐字稿}\par\vspace{4.5pt}
  \PortalBtn{ch:5}{Ch5Main}{第 5 章 · 紧算子与 Fredholm}\hspace{3pt}%
  \PortalBtn{ch:6}{Ch6Main}{第 6 章 · 有界算子谱理论}\hspace{3pt}%
  \PortalBtn{ch:7}{Ch7Main}{第 7 章 · 无界算子}\hspace{3pt}%
  \PortalBtn{ch:8}{Ch8Main}{第 8 章 · 广义函数与 Sobolev}
\end{tcolorbox}

\begin{tcolorbox}[
  enhanced,
  colback=white,
  colframe=BorderColor,
  boxrule=0.6pt,
  arc=4.5pt,
  left=10pt, right=10pt, top=8pt, bottom=8pt,
  before skip=0pt, after skip=6pt
]
  \raggedright
  {\normalsize\bfseries\color{InkColor}全部 27 道指定作业一键传送门（点击题号瞬间跳转，每页右上角可一键返回）}\par\vspace{6pt}

  {\small\bfseries\color{Ch5Main}第 5 章 · 习题 5（共 6 题）}\par\vspace{3pt}
  \ProbChip{prob:5-5}{Ch5Main}{第 5 题（紧算子无有界逆）}\;
  \ProbChip{prob:5-6}{Ch5Main}{第 6 题（值域闭子空间·排雷）}\;
  \ProbChip{prob:5-14}{Ch5Main}{第 14 题（$C[0,1]$ 积分紧性）}\par\vspace{3pt}
  \ProbChip{prob:5-15}{Ch5Main}{第 15 题（有限秩 $\Rightarrow$ HS）}\;
  \ProbChip{prob:5-20}{Ch5Main}{第 20 题（Fredholm 紧扰动）}\;
  \ProbChip{prob:5-22}{Ch5Main}{第 22 题（$BA\in\mathcal{F}$ 推 $A,B\in\mathcal{F}$）}\par\vspace{6pt}

  {\small\bfseries\color{Ch6Main}第 6 章 · 习题 6（共 9 题）}\par\vspace{3pt}
  \ProbChip{prob:6-2}{Ch6Main}{第 2 题（乘法算子求谱）}\;
  \ProbChip{prob:6-5}{Ch6Main}{第 5 题（预解式连续性）}\;
  \ProbChip{prob:6-13}{Ch6Main}{第 13 题（预解式谱映射）}\par\vspace{3pt}
  \ProbChip{prob:6-23}{Ch6Main}{第 23 题（Cayley 变换证酉）}\;
  \ProbChip{prob:6-24}{Ch6Main}{第 24 题（Volterra 算子范数）}\;
  \ProbChip{prob:6-31}{Ch6Main}{第 31 题（HS 三条件等价）}\par\vspace{3pt}
  \ProbChip{prob:6-37}{Ch6Main}{第 37 题（$\lVert A\rVert_2$ 奇异值和）}\;
  \ProbChip{prob:6-39}{Ch6Main}{第 39 题（酉等价函数演算）}\;
  \ProbChip{prob:6-45}{Ch6Main}{第 45 题（解析幂级数演算）}\par\vspace{6pt}

  {\small\bfseries\color{Ch7Main}第 7 章 · 习题 7（共 4 题）}\par\vspace{3pt}
  \ProbChip{prob:7-3}{Ch7Main}{第 3 题（Legendre 算子特征值）}\;
  \ProbChip{prob:7-4}{Ch7Main}{第 4 题（乘法算子闭且无界）}\par\vspace{3pt}
  \ProbChip{prob:7-6}{Ch7Main}{第 6 题（闭对称/自伴与 $A^{**}$）}\;
  \ProbChip{prob:7-7}{Ch7Main}{第 7 题（非负对称与 $\mathcal{R}(A+I)$ 闭）}\par\vspace{6pt}

  {\small\bfseries\color{Ch8Main}第 8 章 · 习题 8 与正文练习（共 8 题）}\par\vspace{3pt}
  \ProbChip{prob:8-7}{Ch8Main}{第 7 题（$L^1$ 不能乘 $\delta$）}\;
  \ProbChip{prob:8-9}{Ch8Main}{第 9 题（Poisson/热核极限）}\;
  \ProbChip{prob:8-11}{Ch8Main}{第 11 题（4 个广义导数·例 8.5/8.6）}\par\vspace{3pt}
  \ProbChip{prob:8-14}{Ch8Main}{第 14 题（主值极限）}\;
  \ProbChip{prob:8-17}{Ch8Main}{第 17(2)(3) 题（Fourier 变换）}\;
  \ProbChip{prob:8-20}{Ch8Main}{第 20 题（半平面 Laplace）}\par\vspace{3pt}
  \ProbChip{prob:8-28}{Ch8Main}{第 28 题（$\lvert x\rvert^{-k}\in H^{-m}$）}\;
  \ProbChip{prob:8-882}{Ch8Main}{p.337 练习 8.8.2（椭圆先验估计）}
\end{tcolorbox}

\clearpage

\SwitchChapterTheme{0}
\ChapterBanner{00}{part:0}{第 0 篇 · 末课老师划重点：闭卷保过总纲与录音整理}{先抓老师强调的「框架 $\to$ 定义 $\to$ 基本性质 $\to$ 算子关系」，再对照 12 类变形通法与 25 问自测闭卷冲刺}

\begin{tcolorbox}[
  enhanced,
  colback=Ch0Soft!45!white,
  colframe=Ch0Main!25!BorderColor,
  boxrule=0.6pt,
  arc=5pt,
  left=8pt, right=8pt, top=8pt, bottom=8pt,
  before skip=6pt, after skip=8pt
]
  \centering
  {\small\bfseries\color{Ch0Main}全书核心逻辑一图打通：老师末课强调的「先收窄 $\to$ 再扩宽 $\to$ 推广到无界」主线}\\[6pt]
  \begin{tikzpicture}[
    node distance=5.5mm and 6mm,
    box/.style={rectangle, rounded corners=3.5pt, draw=#1, fill=#1!10!white, line width=0.7pt,
                text width=146mm, align=left, inner xsep=7pt, inner ysep=4.5pt, font=\footnotesize},
    arr/.style={->, >=stealth, line width=0.9pt, color=MutedInk}
  ]
    \node[box=Ch5Main] (n5) {\textcolor{Ch5Main}{\bfseries 第 5 章 · 四类算子全景}\quad
      $\text{有限秩} \subsetneq \text{Hilbert-Schmidt} \subsetneq \text{紧算子 } \mathcal{C}$ \;\textbar\; $\mathcal{C}$ 与 $\text{Fredholm } \mathcal{F}$ 靠 Atkinson 定理「模紧可逆」相连};
    \node[box=Ch6Main, below=of n5] (n6a) {\textcolor{Ch6Main}{\bfseries 第 6 章（起点）· 一般有界线性算子 $\mathcal{L}(X)$}\quad
      复平面拆为 $\rho(A)$ 与 $\sigma_p \cup \sigma_c \cup \sigma_r$；Gelfand 谱半径 $\Gamma(A)=\lim\lVert A^n\rVert^{1/n}$};
    \node[box=Ch6Main, below=of n6a] (n6b) {\textcolor{Ch6Main}{\bfseries 第 6 章（收窄 1 $\to$ 2）· 紧算子 $\to$ Hilbert 自伴紧算子}\quad
      Riesz-Schauder（非零谱全为有限重特征值）$\to$ $\Gamma(A)=\lVert A\rVert$，Hilbert-Schmidt 谱展开 $Ax=\sum\lambda_n(x,e_n)e_n$};
    \node[box=MutedInk, below=of n6b, dashed] (n6c) {\textcolor{MutedInk}{\bfseries 6.4 过渡桥梁（老师减负：大概了解）}\quad
      谱测度 $E(\Delta)$、谱系 $\{E_\lambda\}$ 与谱积分 $\int f\,\mathrm{d}E$（记牢 4 条代数运算法则即可！）};
    \node[box=Ch6Main, below=of n6c] (n6d) {\textcolor{Ch6Main}{\bfseries 第 6 章（扩宽 1 $\to$ 2）· 酉算子 $\xrightarrow{\text{Cayley 变换}}$ 一般有界自伴算子}\quad
      酉算子（保距满射，$\sigma\subset\{\lvert\lambda\rvert=1\}$）$\to$ 有界自伴算子（$\sigma\subset[-\lVert A\rVert,\lVert A\rVert]$，$A=\int_{\sigma(A)}\lambda\,\mathrm{d}E_\lambda$）};
    \node[box=Ch7Main, below=of n6d] (n7) {\textcolor{Ch7Main}{\bfseries 第 7 章（突破有界）· Hilbert 空间上的无界算子}\quad
      稠定 $\to$ 共轭 $A^*$（恒闭）$\to$ 对称 $(A\subset A^*)$ $\to$ 自伴 $(A=A^*)$ $\xrightarrow{\text{Cayley}}$ 谱在整条实轴 $\mathbb{R}$，$A=\int_{\mathbb{R}}\lambda\,\mathrm{d}E_\lambda$};
    \node[box=Ch8Main, below=of n7] (n8) {\textcolor{Ch8Main}{\bfseries 第 8 章（补充工具）· 广义函数与 Sobolev 空间}\quad
      $\mathcal{D}\subset\mathcal{S}\subset L^2\subset\mathcal{S}'\subset\mathcal{D}'$；广义导数甩负号；本书 $2\pi$ Fourier 变换；实指数 $H^s(\mathbb{R}^n)$ 与椭圆估计};

    \draw[arr] (n5) -- (n6a);
    \draw[arr] (n6a) -- (n6b);
    \draw[arr] (n6b) -- (n6c);
    \draw[arr] (n6c) -- (n6d);
    \draw[arr] (n6d) -- (n7);
    \draw[arr] (n7) -- (n8);
  \end{tikzpicture}
\end{tcolorbox}

\begin{quote}
\textbf{适用对象}：完全没学过或听课云里雾里、即将参加\textbf{闭卷期末考试}、目标稳稳过关不挂科的同学\\
\textbf{核心依据}：最后一节复习课录音（全长 \texttt{13:49}）+ 配套教材《泛函分析》（科学出版社）第 5--8 章正文与定理编号\\
\textbf{阅读方式}：本文是可以直接从头背到尾的\textbf{闭卷详细精讲手册}——不仅有老师完整录音整理，还把老师点名的每一个\textbf{定义、算子关系、完整可默写证明、具体算子实例、易混对比表}以及\textbf{12 类考场变形题通法}全部展开讲透。遇到需要看某道具体作业完整解答时，直接点击页眉「传送门」或目录跳至本书第 5–8 章作业详解即可。
\end{quote}

\BookPart{part:truth}{\WarnBadge{} 考前必读：录音里老师亲口确认的 4 条防挂科真相}{【注意】 考前必读：录音里老师亲口确认的 4 条防挂科真相}

\begin{enumerate}
\def\labelenumi{\arabic{enumi}.}
\tightlist
\item
  \textbf{考试形式实锤：闭卷！（录音 \texttt{13:43–13:47}）}

  \begin{itemize}
  \tightlist
  \item
    下课前同学追问``老师，考试开卷吗？''，老师干脆回答：\textbf{``闭卷！''}
  \item
    \textbf{应对铁律}：不能指望翻书！核心定义、算子包含/等价关系、核心定理陈述与关键短证明（尤其是老师专门提到``当时严格证明过''的 Gelfand 谱半径公式、自伴算子性质、闭算子/自伴判别准则）必须亲手默写过一遍。
  \end{itemize}
\item
  \textbf{成绩构成：平时 50\% + 期末 50\%（录音 \texttt{12:40–12:55}）}

  \begin{itemize}
  \tightlist
  \item
    老师明确说：``平时成绩占 50\%，已经比较高了；但就算平时拿 100 分，期末不好好考还是有点困难，期末考试还是要认真对待。''
  \item
    \textbf{含义}：平时分是你最大的安全垫，但期末卷面绝不能交白卷或胡写，必须拿下基本盘（概念定义 + 性质证明 + 常规计算题）。
  \end{itemize}
\item
  \textbf{老师亲自辟谣：不是考作业原题！（录音 \texttt{11:51–12:39}）}

  \begin{itemize}
  \tightlist
  \item
    老师专门纠偏：``上次有一句误导的话我要纠正一下——作业在复习时可以参考，但并不是说考试题就是作业原题。''
  \item
    她直接点名提醒：不管作业是自己做的，还是在豆包、DeepSeek、ChatGPT 等 AI 帮助下做出来的，\textbf{回去一定要看懂它到底是怎么做的、背后用了什么思路}！因为作业就是顺着课程\textbf{主干框架}布置的，考题会换壳，但考的还是同一套框架与思路。
  \end{itemize}
\item
  \textbf{老师明确划出的减负范围（录音 \texttt{05:20}、\texttt{07:45}、\texttt{11:39}）}

  \begin{itemize}
  \tightlist
  \item
    \textbf{上课讲得比较简单、一笔带过的内容，不用花精力死磕}；
  \item
    第 6 章 \textbf{6.4 节（谱系、谱测度、谱积分）} 老师明确说``\textbf{大概了解一下}''（掌握概念、与谱分解的关系以及谱积分的 4 条运算法则即可，不用背最底层构造性长证明）；
  \item
    第 6 章 \textbf{6.6 节最末尾的拓展应用小节}（循环元/生成元酉等价于乘法算子等）老师明说``有兴趣可以看，没兴趣可以不用管''。
  \end{itemize}
\end{enumerate}

\begin{quote}
\textbf{关于第 8 章的诚实说明}：本段 13 分 49 秒录音从第 5 章末句切入，完整串讲了 \textbf{第 5、6、7 章} 和考务要求，\textbf{未录到第 8 章的口述总结}。由于你此前收到的作业清单包含「第八章增」的 8 道题，本文第五部分专门补齐了\textbf{第 8 章闭卷详细保命精讲}，保证复习不留死角。
\end{quote}

\BookPart{}{零基础极速热身：看懂 5--8 章必会的 6 个底层概念}{零基础极速热身：看懂 5–8 章必会的 6 个底层概念}

如果你连第 1--4 章都没学过，先花 \textbf{5 分钟}把下面 6 个词记住，后面整份笔记就畅通无阻：

\begin{enumerate}
\def\labelenumi{\arabic{enumi}.}
\tightlist
\item
  \textbf{Banach 空间 \(X\)}：有长度（范数 \(\lVert x\rVert\)）且``没漏洞''（任何 Cauchy 列必收敛，即\textbf{完备}）的线性空间。常见例子：连续函数空间 \(C([a,b])\)（\(\lVert f\rVert_\infty = \max\lvert f\rvert\)）、数列空间 \(\ell^p\)、积分空间 \(L^p\)。
\item
  \textbf{Hilbert 空间 \(H\)}：带\textbf{内积} \((x, y)\)（能算夹角与垂直！）的 Banach 空间，由内积诱导范数 \(\lVert x\rVert = \sqrt{(x,x)}\)。

  \begin{itemize}
  \tightlist
  \item
    \textbf{本书内积约定（必记！）}：第一变元线性、第二变元共轭线性，即 \((\alpha x, \beta y) = \alpha\overline{\beta}(x, y)\)，\((y, x) = \overline{(x, y)}\)。
  \item
    \textbf{标准正交基 \(\{e_n\}\)}：两两垂直且长度为 1，任何 \(x \in H\) 可展开为 \(x = \sum_n (x, e_n)e_n\)，且满足 \textbf{Parseval 等式} \(\lVert x\rVert^2 = \sum_n \lvert(x, e_n)\rvert^2\)（无穷维勾股定理！）。
  \end{itemize}
\item
  \textbf{有界线性算子 \(\mathcal{L}(X, Y)\)}：把空间 \(X\) 的元素线性地映到 \(Y\)，且存在常数 \(M\) 使 \(\lVert Ax\rVert \le M\lVert x\rVert\)（对线性算子，``有界 \(\iff\) 连续''）。算子范数 \(\lVert A\rVert \triangleq \sup_{\lVert x\rVert = 1}\lVert Ax\rVert\)。若 \(X=Y\) 简记 \(\mathcal{L}(X)\)。
\item
  \textbf{三大支柱定理（闭卷随时要援引的三大护法）}：

  \begin{itemize}
  \tightlist
  \item
    \textbf{Banach 逆算子定理}：Banach 空间之间的\textbf{连续双射} \(A \in \mathcal{L}(X,Y)\)，其逆算子 \(A^{-1}\) \textbf{自动有界}（\(A^{-1} \in \mathcal{L}(Y,X)\)）。
  \item
    \textbf{开映射定理}：Banach 空间之间的\textbf{连续满射}把开集映成开集。
  \item
    \textbf{闭图像定理}：定义在整个 Banach 空间 \(X\) 上的线性算子 \(A: X \to Y\)，只要它是\textbf{闭算子}（\(x_n \to x, Ax_n \to y \Rightarrow Ax=y\)），就\textbf{自动有界}！
  \item
    \textbf{共鸣定理（一致有界原理）}：Banach 空间上一族有界算子/泛函若``逐点有界''，则必``范数一致有界''（Gelfand 定理证明会用到！）。
  \end{itemize}
\item
  \textbf{Riesz 引理（无穷维空间的``紧性杀手''）}：赋范空间 \(X\) 的\textbf{闭单位球列紧 \(\iff \dim X < +\infty\)}。换句话说：\textbf{在无穷维空间里，闭单位球一定不列紧，恒等算子 \(I_X\) 一定不是紧算子！}
\item
  \textbf{共轭算子 \(A^*\)（有界情形）}：在 Hilbert 空间 \(H\) 上，对 \(A \in \mathcal{L}(H)\)，唯一存在 \(A^* \in \mathcal{L}(H)\) 使 \((Ax, y) = (x, A^*y)\)（线性代数里就是矩阵的\textbf{共轭转置 \(A^H\)}）。满足 \(\lVert A^*\rVert = \lVert A\rVert\)、\(\lVert A^*A\rVert = \lVert AA^*\rVert = \lVert A\rVert^2\)、\((AB)^* = B^*A^*\)。若 \(A = A^*\) 称\textbf{自伴算子}（类比实对称/Hermite 矩阵）。
\end{enumerate}

\BookPart{}{第 5 章深度精讲 —— 四类算子的定义、性质、关系与必考证明}{第 5 章深度精讲 —— 四类算子的定义、性质、关系与必考证明}

\begin{quote}
\textbf{老师划重点回顾（\texttt{00:00–00:46}）}：紧算子、有限秩算子、Hilbert-Schmidt 算子、Fredholm 算子——掌握\textbf{定义、基本性质、彼此之间的关系}。
\end{quote}

\BookSubSection{}{2.1 四类算子的精准定义（闭卷逐字背诵）}{2.1 四类算子的精准定义（闭卷逐字背诵）}{}

\SubSubHeading{\CircNum{1} 紧算子 \(\mathcal{C}(X, Y)\)（定义 5.1，教材 p.182）}

\begin{itemize}
\tightlist
\item
  \textbf{定义}：设 \(X, Y\) 为赋范线性空间（常为 Banach 空间），\(A \in \mathcal{L}(X, Y)\)。如果 \(A\) 把 \(X\) 中\textbf{每个有界集}都映射为 \(Y\) 中的\textbf{列紧集}（相对紧集），则称 \(A\) 是紧算子，记作 \(A \in \mathcal{C}(X, Y)\)（当 \(X=Y\) 时简记 \(\mathcal{C}(X)\)）。
\item
  \textbf{最常用的等价序列定义（做题首选！）}： 
\FitDisplayMath{A \in \mathcal{C}(X, Y) \iff \text{对 } X \text{ 中任意有界序列 } \{x_n\}\text{，其像序列 } \{Ax_n\} \text{ 在 } Y \text{ 中必有收敛子列。}}

\item
  \textbf{人话直觉}：有限维里``有界集必列紧''，无穷维里``有界集不列紧''；\textbf{紧算子就像一台强力压缩机}，把无穷维里松散的有界球压扁成一个几乎像有限维一样紧凑的集合。
\end{itemize}

\SubSubHeading{\CircNum{2} 有限秩算子 \(\mathcal{F}_0(X, Y)\)（定义 5.2，教材 p.189）}

\begin{itemize}
\tightlist
\item
  \textbf{定义}：设 \(A \in \mathcal{L}(X, Y)\)，如果值域的维数有限，即 \(\dim\mathcal{R}(A) < +\infty\)，则称 \(A\) 是有限秩算子。
\item
  \textbf{人话直觉}：输出空间干脆就是一个有限维子空间（就像把无穷维向量投影到三维空间里），是``最极端的紧算子''。
\end{itemize}

\SubSubHeading{\CircNum{3} Hilbert-Schmidt 算子 \(\mathcal{L}^2(H, G)\)（定义 5.4，教材 p.196）}

\begin{itemize}
\tightlist
\item
  \textbf{定义}：设 \(H, G\) 为可分 Hilbert 空间，\(A \in \mathcal{L}(H, G)\)。若存在 \(H\) 的一个标准正交基 \(\{e_n\}_{n=1}^\infty\) 使得 
\FitDisplayMath{\lVert A\rVert_2 \triangleq \left(\sum_{n=1}^{\infty} \lVert Ae_n\rVert^2\right)^{\frac{1}{2}} < +\infty,}
 则称 \(A\) 是 \(H\) 到 \(G\) 的 \textbf{Hilbert-Schmidt 算子}，记作 \(A \in \mathcal{L}^2(H, G)\)，\(\lVert A\rVert_2\) 称为 \(A\) 的 Hilbert-Schmidt 范数。
\item
  \textbf{重要验证（为什么与基的选取无关？）}：若 \(\{e_m\}, \{\tilde{e}_m\}\) 是 \(H\) 的任意两组标准正交基，\(\{f_n\}\) 是 \(G\) 的标准正交基，由 Parseval 等式两次换基： 
\FitDisplayMath{\sum_{m=1}^\infty \lVert Ae_m\rVert^2 = \sum_{m,n} \lvert(Ae_m, f_n)\rvert^2 = \sum_{n}\sum_m \lvert(e_m, A^*f_n)\rvert^2 = \sum_n \lVert A^*f_n\rVert^2 = \sum_n \sum_m \lvert(\tilde{e}_m, A^*f_n)\rvert^2 = \sum_{m=1}^\infty \lVert A\tilde{e}_m\rVert^2.}
 （顺手还同时证明了 \textbf{\(\lVert A^*\rVert_2 = \lVert A\rVert_2\)}！闭卷考``证明 HS 范数与标准正交基选取无关''直接照抄这两行！）
\item
  \textbf{人话直觉}：对于 \(n\times n\) 矩阵 \(A=(a_{ij})\)，\(\lVert A\rVert_2 = \sqrt{\sum_{i,j}\lvert a_{ij}\rvert^2}\) 就是矩阵的 \textbf{Frobenius 范数}；对于 \(L^2(\Omega)\) 上的积分算子 \((Af)(x)=\int_\Omega K(x,y)f(y)\,\mathrm{d}y\)，它是 HS 算子当且仅当积分核平方可积 \(\iint \lvert K(x,y)\rvert^2\,\mathrm{d}x\mathrm{d}y < \infty\)。
\end{itemize}

\SubSubHeading{\CircNum{4} Fredholm 算子 \(\mathcal{F}(X, Y)\)（定义 5.5--5.6，教材 p.198--199）}

\begin{itemize}
\tightlist
\item
  \textbf{定义}：设 \(X, Y\) 为 Banach 空间，\(A \in \mathcal{L}(X, Y)\)。如果同时满足三条：

  \begin{enumerate}
  \def\labelenumi{\arabic{enumi}.}
  \tightlist
  \item
    \(\dim\operatorname{Ker}A < +\infty\)（核空间有限维）；
  \item
    \(\mathcal{R}(A)\) 是 \(Y\) 中的\textbf{闭}子空间；
  \item
    \(\operatorname{codim}\mathcal{R}(A) \triangleq \dim(Y / \mathcal{R}(A)) < +\infty\)（值域余维数有限）； 则称 \(A\) 是 \textbf{Fredholm 算子}，记作 \(A \in \mathcal{F}(X, Y)\)。
  \end{enumerate}
\item
  \textbf{指标（Index）定义}： 
\FitDisplayMath{\operatorname{ind} A \triangleq \dim\operatorname{Ker}A - \operatorname{codim}\mathcal{R}(A) \in \mathbb{Z}.}

\item
  \textbf{人话直觉}：如果一个算子完全可逆，那么 \(\dim\operatorname{Ker}A = 0\)（单射）且 \(\operatorname{codim}\mathcal{R}(A) = 0\)（满射），指标为 \(0\)。Fredholm 算子就是\textbf{``差一点点就可逆''的算子}——只丢了有限维的信息，也只少盖住有限维的目标空间！
\end{itemize}

\SoftDivider

\BookSubSection{}{2.2 四类算子之间的全部关系（老师重点强调！含完整证明与反例）}{2.2 四类算子之间的全部关系（老师重点强调！含完整证明与反例）}{}

\SubSubHeading{关系 1：严格真包含链（在无穷维可分 Hilbert 空间 \(H\) 上）}


\FitDisplayMath{\text{有限秩算子 } \mathcal{F}_0(H) \;\subsetneq\; \text{Hilbert-Schmidt 算子 } \mathcal{L}^2(H) \;\subsetneq\; \text{紧算子 } \mathcal{C}(H) \;\subsetneq\; \text{有界线性算子 } \mathcal{L}(H).}


下面把\textbf{每一环的包含证明 + 严格真包含反例}全部给齐（闭卷极易考其中任意一环的证明或举反例！）：

\begin{enumerate}
\def\labelenumi{\arabic{enumi}.}
\tightlist
\item
  \textbf{证明「有限秩 \(\subset\) HS」（即\hyperlink{prob:5-15}{\textcolor{LinkColor}{习题 5 第 15 题}}！）}：

  \begin{itemize}
  \tightlist
  \item
    因 \(A\) 有限秩，可证 \(\dim(\operatorname{Ker}A)^\perp = \dim\mathcal{R}(A) = m < \infty\)。取 \((\operatorname{Ker}A)^\perp\) 的标准正交基 \(\{e_1, \dots, e_m\}\)，并上 \(\operatorname{Ker}A\) 的标准正交基 \(\{z_k\}\)，合起来构成 \(H\) 的一组标准正交基。因 \(Az_k = 0\)，直接求和： 
\FitDisplayMath{\sum_{j=1}^m \lVert Ae_j\rVert^2 + \sum_k \lVert Az_k\rVert^2 = \sum_{j=1}^m \lVert Ae_j\rVert^2 < +\infty \implies A \in \mathcal{L}^2(H, G).}

  \item
    \textbf{真包含反例（HS 但非有限秩）}：在 \(\ell^2\) 上定义对角算子 \(Ae_n = \frac{1}{n}e_n\)。\(\sum_{n=1}^\infty \lVert Ae_n\rVert^2 = \sum_{n=1}^\infty \frac{1}{n^2} = \frac{\pi^2}{6} < \infty\)，故 \(A \in \mathcal{L}^2(\ell^2)\)；但每个 \(e_n = A(ne_n) \in \mathcal{R}(A)\)，值域无穷维，故不是有限秩算子。
  \end{itemize}
\item
  \textbf{证明「HS \(\subset\) 紧算子」（定理 5.12 核心！）}：

  \begin{itemize}
  \tightlist
  \item
    设 \(A \in \mathcal{L}^2(H, G)\)，\(\{e_n\}\) 是 \(H\) 的标准正交基。先验证 \(\lVert A\rVert \le \lVert A\rVert_2\)：对任意 \(x = \sum_{n=1}^\infty (x, e_n)e_n \in H\)，由 Cauchy-Schwarz 不等式 
\FitDisplayMath{\lVert Ax\rVert = \left\lVert\sum_{n=1}^\infty (x, e_n)Ae_n\right\rVert \le \sum_{n=1}^\infty \lvert(x, e_n)\rvert\lVert Ae_n\rVert \le \left(\sum_{n=1}^\infty \lvert(x, e_n)\rvert^2\right)^{\frac{1}{2}}\left(\sum_{n=1}^\infty \lVert Ae_n\rVert^2\right)^{\frac{1}{2}} = \lVert x\rVert\lVert A\rVert_2.}

  \item
    再构造有限秩截断算子 \(A_N x \triangleq \sum_{n=1}^N (x, e_n)Ae_n\)（显然秩 \(\le N\)，故 \(A_N\) 有限秩，从而紧）。余项满足 
\FitDisplayMath{\lVert A - A_N\rVert \le \lVert A - A_N\rVert_2 = \left(\sum_{n=N+1}^\infty \lVert Ae_n\rVert^2\right)^{\frac{1}{2}} \xrightarrow{N\to\infty} 0.}

  \item
    因为紧算子集按算子范数闭（定理 5.4），故 \(A \in \mathcal{C}(H, G)\)！
  \item
    \textbf{真包含反例（紧但非 HS）}：在 \(\ell^2\) 上定义 \(Be_n = \frac{1}{\sqrt{n}}e_n\)。截断 \(B_N\) 满足 \(\lVert B - B_N\rVert = \frac{1}{\sqrt{N+1}} \to 0\)，故 \(B\) 是紧算子；但 \(\sum_{n=1}^\infty \lVert Be_n\rVert^2 = \sum_{n=1}^\infty \frac{1}{n} = +\infty\)，故 \(B \notin \mathcal{L}^2(\ell^2)\)！
  \end{itemize}
\item
  \textbf{证明「有限秩 \(\subset\) 紧算子」（一般 Banach 空间也成立！）}：

  \begin{itemize}
  \tightlist
  \item
    若 \(\dim\mathcal{R}(A) < \infty\)，任意有界集 \(B \subset X\) 的像 \(A(B)\) 是有限维空间 \(\mathcal{R}(A)\) 中的有界集，由有限维 Bolzano-Weierstrass 定理必列紧，故 \(A \in \mathcal{C}(X, Y)\)。
  \item
    \textbf{反过来，紧算子能否被有限秩算子逼近？} 在可分 Hilbert 空间（定理 5.9）或有 Schauder 基的 Banach 空间（定理 5.11）上：\textbf{任何紧算子都是一列有限秩算子的算子范数极限}！
  \end{itemize}
\item
  \textbf{「紧算子 \(\subsetneq\) 有界算子」的反例}：无穷维 Banach 空间上的恒等算子 \(I_X \in \mathcal{L}(X)\)，但其闭单位球像还是自身，由 Riesz 引理不列紧，故 \(I_X \notin \mathcal{C}(X)\)。
\end{enumerate}

\SubSubHeading{关系 2：紧算子与 Fredholm 算子的``互补与模紧可逆''关系}

\begin{enumerate}
\def\labelenumi{\arabic{enumi}.}
\tightlist
\item
  \textbf{在无穷维空间中，紧算子与 Fredholm 算子完全互斥}（\(\mathcal{C}(X) \cap \mathcal{F}(X) = \emptyset\)）：

  \begin{itemize}
  \tightlist
  \item
    若 \(A \in \mathcal{C}(X)\)，则 \(\mathcal{R}(A)\) 不包含任何无穷维闭子空间（习题 5 第 6/17 题）；而 \(A \in \mathcal{F}(X)\) 要求 \(\mathcal{R}(A)\) 闭且余维有限，在无穷维空间里 \(\mathcal{R}(A)\) 必无穷维，直接矛盾！
  \end{itemize}
\item
  \textbf{恒等算子减紧算子必是指标为 \(0\) 的 Fredholm 算子}（Riesz-Fredholm 定理 5.8）： 
\FitDisplayMath{A \in \mathcal{C}(X) \implies I - A \in \mathcal{F}(X) \quad \text{且} \quad \operatorname{ind}(I - A) = \dim\operatorname{Ker}(I - A) - \operatorname{codim}\mathcal{R}(I - A) = 0.}

\item
  \textbf{Atkinson 定理（定理 5.13，教材 p.199，极其核心！）}： 
\FitDisplayMath{A \in \mathcal{F}(X, Y) \iff \exists\, S \in \mathcal{L}(Y, X),\; K_1 \in \mathcal{C}(X),\; K_2 \in \mathcal{C}(Y) \text{ 使得 } SA = I_X - K_1,\; AS = I_Y - K_2.}


  \begin{itemize}
  \tightlist
  \item
    \textbf{大白话翻译}：把所有紧算子当成``可以忽略的小零头''（商代数 \(\mathcal{L}(X)/\mathcal{C}(X)\)），\textbf{Fredholm 算子就是``模紧算子可逆''的算子}！
  \end{itemize}
\item
  \textbf{紧扰动不变性（\hyperlink{prob:5-20}{\textcolor{LinkColor}{习题 5 第 20 题}}）}：若 \(B \in \mathcal{F}(X, Y), A \in \mathcal{C}(X, Y)\)，则 \(B + A \in \mathcal{F}(X, Y)\) 且 \(\operatorname{ind}(B + A) = \operatorname{ind} B\)。
\end{enumerate}

\SoftDivider

\BookSubSection{}{2.3 第 5 章核心基本性质与闭卷必考短证明}{2.3 第 5 章核心基本性质与闭卷必考短证明}{\TagMust{闭卷必背}}

{\def\LTcaptype{none} % do not increment counter
{\small\renewcommand{\arraystretch}{1.38}\setlength{\tabcolsep}{5pt}
\begin{longtable}[]{@{}
  >{\raggedright\arraybackslash}p{(\linewidth - 4\tabcolsep) * \real{0.3333}}
  >{\raggedright\arraybackslash}p{(\linewidth - 4\tabcolsep) * \real{0.3333}}
  >{\raggedright\arraybackslash}p{(\linewidth - 4\tabcolsep) * \real{0.3333}}@{}}
\toprule\noalign{}
\rowcolor{AccentSoft}
\begin{minipage}[b]{\linewidth}\raggedright
性质/定理
\end{minipage} & \begin{minipage}[b]{\linewidth}\raggedright
精确结论
\end{minipage} & \begin{minipage}[b]{\linewidth}\raggedright
闭卷必背短证明 / 核心机制
\end{minipage} \\
\midrule\noalign{}
\endhead
\bottomrule\noalign{}
\endlastfoot
\textbf{定理 5.1（变弱为强）} & \(A \in \mathcal{C}(X, Y)\) 且 \(x_n \rightharpoonup x_0 \implies Ax_n \to Ax_0\)（强收敛） & 反证：若 \(\lVert Ax_n - Ax_0\rVert \ge \varepsilon_0 > 0\)，因弱收敛列必有界，由 \(A\) 紧抽子列 \(Ax_{n_k} \to y\)；又连续线性算子保弱收敛得 \(Ax_{n_k} \rightharpoonup Ax_0\)，弱极限唯一故 \(y = Ax_0\)，矛盾！（注记 5.1：\(X\) 自反时逆命题成立；非自反如 \(\ell^1\) 不成立） \\
\textbf{定理 5.3（双边理想）} & \(A_1, A_2 \in \mathcal{C}(X, Y) \Rightarrow \alpha_1 A_1 + \alpha_2 A_2 \in \mathcal{C}(X, Y)\)；\(A \in \mathcal{C}, B \in \mathcal{L} \Rightarrow AB, BA \in \mathcal{C}\) & \(BA\)：有界列 \(\xrightarrow{A\text{紧}}\) 收敛子列 \(\xrightarrow{B\text{连续}}\) 收敛子列；\(AB\)：有界列 \(\xrightarrow{B\text{有界}}\) 有界列 \(\xrightarrow{A\text{紧}}\) 收敛子列！ \\
\textbf{定理 5.4（范数闭性）} & \(A_n \in \mathcal{C}(X, Y), \lVert A_n - A\rVert \to 0 \implies A \in \mathcal{C}(X, Y)\)（\(Y\) 完备） & 对角线法！对有界列 \(\{x_k\}\) 依次抽子列使每个 \(A_n\) 像都收敛，取对角子列 \(\{x_{k_k}\}\)，再用 \(\lVert Ax_{k_i} - Ax_{k_j}\rVert \le \lVert A - A_n\rVert(\lVert x_{k_i}\rVert+\lVert x_{k_j}\rVert) + \lVert A_n x_{k_i} - A_n x_{k_j}\rVert\) 证 Cauchy 列！ \\
\textbf{定理 5.5（无穷维不可逆）} & \(\dim X = \infty, A \in \mathcal{C}(X, Y)\) 且单射 \(\implies \mathcal{R}(A) \neq Y\)（特别 \(A \in \mathcal{C}(X)\) 无有界逆） & 反证：若 \(\mathcal{R}(A) = Y\)，由 Banach 逆算子定理 \(A^{-1} \in \mathcal{L}(Y, X)\)，从而 \(I_X = A^{-1}A \in \mathcal{C}(X)\)，推出 \(X\) 闭单位球列紧，与 Riesz 引理矛盾！ \\
\textbf{定理 5.6（值域闭）} & \(A \in \mathcal{C}(X), \lambda \neq 0 \implies \mathcal{R}(\lambda I - A)\) 是闭子空间 & 化为 \(T = I - A\)，在商空间 \(X/\operatorname{Ker}T\) 上诱导单射 \(\tilde{T}\)，反证 \(\tilde{T}^{-1}\) 连续（取 \(\lVert\tilde{x}_n\rVert=1, \tilde{T}\tilde{x}_n\to 0\)，选代表元 \(\lVert x_n\rVert<2\)，用 \(A\) 紧抽子列得矛盾）。 \\
\textbf{定理 5.7（Schauder）} & \(A \in \mathcal{C}(X, Y) \implies A^* \in \mathcal{C}(Y', X')\)（实为充要条件） & 记结论即可，常配合紧算子谱论中处理 \(A^*\)。 \\
\textbf{定理 5.14（指标可加）} & \(A_1 \in \mathcal{F}(X, Y), A_2 \in \mathcal{F}(Y, Z) \implies A_2 A_1 \in \mathcal{F}(X, Z)\) 且 \(\operatorname{ind}(A_2 A_1) = \operatorname{ind}A_2 + \operatorname{ind}A_1\) & 记住与 Atkinson 定理联合使用：若 \(S\) 是 \(B\) 和 \(B+A\) 的公共拟逆，则 \(\operatorname{ind}S + \operatorname{ind}B = 0 = \operatorname{ind}S + \operatorname{ind}(B+A) \Rightarrow \operatorname{ind}(B+A)=\operatorname{ind}B\)！ \\
\textbf{定理 5.15（小扰动稳定）} & \(A \in \mathcal{F}(X, Y) \implies \exists\,\varepsilon > 0\)，当 \(\lVert S\rVert < \varepsilon\) 时 \(A+S \in \mathcal{F}(X, Y)\) 且 \(\operatorname{ind}(A+S)=\operatorname{ind}A\) & 直观：\(\mathcal{F}(X, Y)\) 是开集，指标是局部常值的整数连续函数！ \\
\end{longtable}
}
}

\BookPart{}{第 6 章深度精讲 —— 有界线性算子谱理论的``沙漏主线''}{第 6 章深度精讲 —— 有界线性算子谱理论的“沙漏主线”}

\begin{quote}
\textbf{老师划重点回顾（\texttt{00:46–08:21}，全场占时超过一半！）}： 整章是一条\textbf{``先收窄，再扩宽''}的经典逻辑链： \texttt{一般有界线性算子} \(\xrightarrow{\text{收窄1}}\) \texttt{紧算子（Riesz-Schauder）} \(\xrightarrow{\text{收窄2}}\) \texttt{Hilbert\ 空间自伴紧算子（Hilbert-Schmidt）} \(\xrightarrow{\text{桥梁}}\) \texttt{谱测度/谱系/谱积分（大概了解）} \(\xrightarrow{\text{扩宽1}}\) \texttt{酉算子谱分解} \(\xrightarrow{\text{扩宽2}}\) \texttt{一般有界自伴算子谱分解}。
\end{quote}

\BookSubSection{}{3.1 第一层：一般有界线性算子 \(A \in \mathcal{L}(X)\)（复 Banach 空间）}{3.1 第一层：一般有界线性算子 A ∈ L(X)（复 Banach 空间）}{\TagMust{闭卷必背}}

\SubSubHeading{1. 为什么无穷维里只有``特征值''不够？}

在线性代数里，矩阵 \(A\) 是 \(n\times n\) 方阵，\(\lambda I - A\) 只要是单射就自动是满射（\(\det(\lambda I - A) \neq 0\)），所以``不可逆 \(\iff\) 有特征向量''。但到了无穷维空间，一个算子完全可能\textbf{是单射（没有特征向量）却不是满射}，或者\textbf{逆算子存在却无界}！因此必须把复平面 \(\mathbb{C}\) 一分为二，再把``坏点''细分为三类：

\begin{itemize}
\item
  \textbf{预解集（正则点集）\(\rho(A)\)（定义 6.1）}： 
\FitDisplayMath{\rho(A) \triangleq \left\{\lambda \in \mathbb{C} : (\lambda I - A)^{-1} \text{ 存在，定义域为全空间 } X \text{ 且是有界线性算子}\right\}.}


  \begin{quote}
  \textbf{省力提醒（注记 6.1）}：在 Banach 空间上，由 Banach 逆算子定理，只要 \(\lambda I - A\) 是\textbf{单射且满射}（双射），逆算子 \((\lambda I - A)^{-1}\) 自动有界！预解式记作 \(R(\lambda, A) \triangleq (\lambda I - A)^{-1} \in \mathcal{L}(X)\)。
  \end{quote}
\item
  \textbf{谱集 \(\sigma(A)\)（定义 6.2）}：\(\sigma(A) \triangleq \mathbb{C} \setminus \rho(A)\)，互斥拆为三大类：

  \begin{enumerate}
  \def\labelenumi{\arabic{enumi}.}
  \tightlist
  \item
    \textbf{点谱（特征值集）\(\sigma_p(A)\)}：\(\lambda I - A\) \textbf{不是单射}（即存在 \(x \neq 0\) 使 \(Ax = \lambda x\)）；
  \item
    \textbf{连续谱 \(\sigma_c(A)\)}：\(\lambda I - A\) \textbf{是单射}，值域 \(\mathcal{R}(\lambda I - A)\) 在 \(X\) 中\textbf{稠密但不等于 \(X\)}（等价：逆存在且稠定但\textbf{无界}）；
  \item
    \textbf{剩余谱 \(\sigma_r(A)\)}：\(\lambda I - A\) \textbf{是单射}，且值域 \(\mathcal{R}(\lambda I - A)\) 在 \(X\) 中\textbf{不稠密}（\(\overline{\mathcal{R}(\lambda I - A)} \neq X\)）。
  \end{enumerate}
\end{itemize}

\SubSubHeading{2. 闭卷必会：四个典型算子的谱分类实战（彻底分清 \(\sigma_p, \sigma_c, \sigma_r\)！）}

{\def\LTcaptype{none} % do not increment counter
{\small\renewcommand{\arraystretch}{1.38}\setlength{\tabcolsep}{5pt}
\begin{longtable}[]{@{}
  >{\raggedright\arraybackslash}p{(\linewidth - 6\tabcolsep) * \real{0.2500}}
  >{\raggedright\arraybackslash}p{(\linewidth - 6\tabcolsep) * \real{0.2500}}
  >{\raggedright\arraybackslash}p{(\linewidth - 6\tabcolsep) * \real{0.2500}}
  >{\raggedright\arraybackslash}p{(\linewidth - 6\tabcolsep) * \real{0.2500}}@{}}
\toprule\noalign{}
\rowcolor{AccentSoft}
\begin{minipage}[b]{\linewidth}\raggedright
经典算子
\end{minipage} & \begin{minipage}[b]{\linewidth}\raggedright
空间与公式
\end{minipage} & \begin{minipage}[b]{\linewidth}\raggedright
谱集 \(\sigma(A)\)
\end{minipage} & \begin{minipage}[b]{\linewidth}\raggedright
精确三分类（闭卷杀手锏！）
\end{minipage} \\
\midrule\noalign{}
\endhead
\bottomrule\noalign{}
\endlastfoot
\textbf{左移算子 \(L\)} & \(\ell^2\)：\(L(x_1, x_2, \dots) = (x_2, x_3, \dots)\) & \(\{\lvert\lambda\rvert \le 1\}\) & \(\sigma_p = \{\lvert\lambda\rvert < 1\}\)（取 \(x=(1,\lambda,\lambda^2,\dots)\in\ell^2\)）；\(\sigma_c = \{\lvert\lambda\rvert = 1\}\)；\(\sigma_r = \emptyset\) \\
\textbf{右移算子 \(R\)} & \(\ell^2\)：\(R(x_1, x_2, \dots) = (0, x_1, x_2, \dots)\) & \(\{\lvert\lambda\rvert \le 1\}\) & \(\sigma_p = \emptyset\)（\textbf{没有任何特征值！}）；\(\sigma_r = \{\lvert\lambda\rvert < 1\}\)（如 \(\lambda=0\) 时值域缺第一位不稠！）；\(\sigma_c = \{\lvert\lambda\rvert = 1\}\) \\
\textbf{实连续乘法算子}（\hyperlink{prob:6-2}{\textcolor{LinkColor}{习题 6 第 2 题}}） & \(L^2(0,1)\)：\((Af)(x) = a(x)f(x)\)，\(a \in C([0,1];\mathbb{R})\) & 值域 \(a([0,1])\) & \(\lambda \in \sigma_p \iff a^{-1}(\{\lambda\})\) 有正测度；若 \(a\) 严格单调（如 \(a(x)=x\)），则 \(\sigma_p = \sigma_r = \emptyset\)，\textbf{全是连续谱} \(\sigma_c = [0,1]\)！ \\
\textbf{Volterra 积分算子}（\hyperlink{prob:6-24}{\textcolor{LinkColor}{习题 6 第 24 题}}） & \(L^2(0,1)\)：\((Af)(x) = \int_0^x f(t)\,\mathrm{d}t\) & \(\{0\}\) & \(\sigma_p = \emptyset, \sigma_r = \emptyset, \sigma_c = \{0\}\)；\(\Gamma(A) = 0\)，但 \(\lVert A\rVert = \frac{2}{\pi} \neq 0\)！ \\
\end{longtable}
}
}

\SubSubHeading{3. 预解集与谱的四大基本定理（含闭卷必会证明）}

\begin{enumerate}
\def\labelenumi{\arabic{enumi}.}
\tightlist
\item
  \textbf{定理 6.1（Neumann 级数展开，教材 p.209）}：记 \(r = \lim_{n\to\infty}\sqrt[n]{\lVert A^n\rVert}\)（引理 6.1 保证极限存在且 \(r \le \lVert A\rVert\)）。当 \(\lvert\lambda\rvert > r\)（特别当 \(\lvert\lambda\rvert > \lVert A\rVert\)）时，\(\lambda \in \rho(A)\)，且 
\FitDisplayMath{R(\lambda, A) = (\lambda I - A)^{-1} = \sum_{n=0}^{\infty} \frac{A^n}{\lambda^{n+1}} \quad \text{（按算子范数绝对收敛）}.}


  \begin{itemize}
  \tightlist
  \item
    \textbf{完整简明证明}：因 \(\lim_{n\to\infty}\left(\frac{\lVert A^n\rVert}{\lvert\lambda\rvert^n}\right)^{1/n} = \frac{r}{\lvert\lambda\rvert} < 1\)，由根值判别法知 \(\sum_{n=0}^\infty \frac{\lVert A^n\rVert}{\lvert\lambda\rvert^{n+1}} < \infty\)。因 \(\mathcal{L}(X)\) 完备，级数 \(S \triangleq \sum_{n=0}^\infty \frac{A^n}{\lambda^{n+1}} \in \mathcal{L}(X)\) 收敛。直接乘开（利用连续性逐项乘入）： 
\FitDisplayMath{(\lambda I - A)S = \sum_{n=0}^\infty \frac{A^n}{\lambda^n} - \sum_{n=0}^\infty \frac{A^{n+1}}{\lambda^{n+1}} = I, \qquad S(\lambda I - A) = I.}
 故 \(S = (\lambda I - A)^{-1} \in \mathcal{L}(X)\)，即 \(\lambda \in \rho(A)\)！
  \end{itemize}
\item
  \textbf{定理 6.2 \& 推论 6.1（\(\rho(A)\) 开集、\(\sigma(A)\) 有界闭集，教材 p.210--211）}：

  \begin{itemize}
  \tightlist
  \item
    若 \(\lambda_0 \in \rho(A)\) 且 \(\lvert\lambda - \lambda_0\rvert < \frac{1}{\lVert R(\lambda_0, A)\rVert}\)，则 \(\lambda \in \rho(A)\)，且 
\FitDisplayMath{R(\lambda, A) = \sum_{n=0}^{\infty} (\lambda_0 - \lambda)^n \left[R(\lambda_0, A)\right]^{n+1}.}

  \item
    \textbf{完整简明证明（也是\hyperlink{prob:6-5}{\textcolor{LinkColor}{习题 6 第 5 题}}的核心套路！）}：把 \(\lambda I - A\) 强行提出 \((\lambda_0 I - A)\)： 
\FitDisplayMath{\lambda I - A = (\lambda_0 I - A) - (\lambda_0 - \lambda)I = (\lambda_0 I - A)\left[I - (\lambda_0 - \lambda)R(\lambda_0, A)\right].}
 当 \(\lvert\lambda - \lambda_0\rvert\lVert R(\lambda_0, A)\rVert < 1\) 时，括号里的 \([I - B]^{-1} = \sum_{n=0}^\infty B^n\) 由 Neumann 级数存在且有界！故两边求逆即得 \(\lambda \in \rho(A)\) 及上式展开！由此立得 \(\rho(A)\) 是开集，\(\sigma(A)\) 是有界闭集（紧集）。
  \end{itemize}
\item
  \textbf{定理 6.4（谱非空定理，教材 p.213）}：在\textbf{复} Banach 空间 \(X \neq \{0\}\) 上，对任意 \(A \in \mathcal{L}(X)\)，必有 \textbf{\(\sigma(A) \neq \emptyset\)}！

  \begin{itemize}
  \tightlist
  \item
    \textbf{闭卷证明思路（反证 + Liouville 定理）}：假设 \(\sigma(A) = \emptyset\)，则 \(\rho(A) = \mathbb{C}\)。对任意 \(f \in \mathcal{L}(X)'\)，复函数 \(g(\lambda) = f(R(\lambda, A))\) 在整个复平面 \(\mathbb{C}\) 上解析（整函数）；又当 \(\lvert\lambda\rvert > \lVert A\rVert\) 时，\(\lvert g(\lambda)\rvert \le \lVert f\rVert \lVert R(\lambda, A)\rVert \le \frac{\lVert f\rVert}{\lvert\lambda\rvert - \lVert A\rVert} \to 0\)（\(\lvert\lambda\rvert \to \infty\)），故 \(g\) 是有界整函数。由复变函数 \textbf{Liouville 定理}，\(g(\lambda) \equiv 0\)。由 \(f \in \mathcal{L}(X)'\) 任意性（Hahn-Banach 推论），\(R(\lambda, A) \equiv 0\)，但 \((\lambda I - A)R(\lambda, A) = I \neq 0\)，矛盾！（注：实空间可以无谱，例如 \(\mathbb{R}^2\) 上 \(90^\circ\) 旋转矩阵！）
  \end{itemize}
\item
  \textbf{\WarnBadge{} 老师特别点名「当时严格证明过」的定理 6.5（Gelfand 谱半径公式，教材 p.214）}：

  \begin{itemize}
  \tightlist
  \item
    \textbf{陈述}：设 \(X\) 为复 Banach 空间，\(A \in \mathcal{L}(X)\)，记谱半径 \(\Gamma(A) \triangleq \sup_{\lambda \in \sigma(A)}\lvert\lambda\rvert\)，则 
\FitDisplayMath{\Gamma(A) = \lim_{n \to \infty} \sqrt[n]{\lVert A^n\rVert}.}

  \item
    \textbf{完整零跳步闭卷可默写证明（严格对齐教材 p.214！）}：

    \begin{itemize}
    \tightlist
    \item
      \textbf{先备引理 6.1（极限存在）}：记 \(a_n = \lVert A^n\rVert^{1/n}\)，\(r = \inf_{n\ge 1} a_n\)。对任意 \(\varepsilon > 0\)，取正整数 \(m\) 使 \(a_m < r + \varepsilon\)。任意 \(n > m\) 写带余除法 \(n = pm + q\)（\(0 \le q < m\)），则 \(\lVert A^n\rVert \le \lVert A^m\rVert^p \lVert A\rVert^q\)，开 \(n\) 次方得 \(a_n \le a_m^{pm/n} \lVert A\rVert^{q/n}\)。令 \(n \to \infty\)（此时 \(pm/n \to 1, q/n \to 0\)），得 \(\limsup_{n\to\infty} a_n \le a_m < r + \varepsilon\)。由 \(\varepsilon\) 任意性，\(\limsup a_n \le r \le \liminf a_n\)，故极限 \(\lim_{n\to\infty}\sqrt[n]{\lVert A^n\rVert} = r\) 存在！
    \item
      \textbf{第一步（证 \(\Gamma(A) \le r\)）}：由定理 6.1，只要 \(\lvert\lambda\rvert > r = \lim_{n\to\infty}\sqrt[n]{\lVert A^n\rVert}\)，就有 \(\lambda \in \rho(A)\)，因此一切 \(\lambda \in \sigma(A)\) 必满足 \(\lvert\lambda\rvert \le r\)，取上确界得 \(\Gamma(A) \le r\)。
    \item
      \textbf{第二步（证 \(r \le \Gamma(A)\)）}：对任意 \(\varepsilon > 0\)，圆环域 \(D = \{\lambda \in \mathbb{C} : \lvert\lambda\rvert > \Gamma(A)\}\) 全落在 \(\rho(A)\) 中。任取有界线性泛函 \(f \in \mathcal{L}(X)'\)，复函数 \(g(\lambda) = f\big[(\lambda I - A)^{-1}\big]\) 在 \(D\) 上解析（引理 6.2）。在无穷远点邻域 \(\lvert\lambda\rvert > \lVert A\rVert\) 上，由 Neumann 级数有 Laurent 展开 
\FitDisplayMath{g(\lambda) = \sum_{n=0}^{\infty} \frac{f(A^n)}{\lambda^{n+1}}.}
 由解析函数 Laurent 展开的唯一性，该级数在整个圆环域 \(\lvert\lambda\rvert > \Gamma(A)\) 内处处收敛！特别取 \(\lambda_0 = \Gamma(A) + \varepsilon \in D\)，级数 \(\sum_{n=0}^\infty \frac{f(A^n)}{\lambda_0^{n+1}}\) 收敛，故通项趋于 \(0\)，从而有界：存在常数 \(M_f > 0\) 使 
\FitDisplayMath{\sup_{n \ge 0} \left\lvert f\!\left(\frac{A^n}{\lambda_0^{n+1}}\right) \right\rvert \le M_f, \qquad \forall f \in \mathcal{L}(X)'.}
 将 \(\left\{\frac{A^n}{\lambda_0^{n+1}}\right\}\) 视作 \(\mathcal{L}(X)'\) 上的有界线性泛函族，由\textbf{共鸣定理（一致有界原理）}，存在与 \(n\) 无关的常数 \(M > 0\) 使 
\FitDisplayMath{\left\lVert \frac{A^n}{\lambda_0^n} \right\rVert \le M \implies \lVert A^n\rVert \le M \big(\Gamma(A) + \varepsilon\big)^n, \qquad \forall n \ge 1.}
 两边开 \(n\) 次方并令 \(n \to \infty\)，得 \(r = \lim_{n\to\infty}\sqrt[n]{\lVert A^n\rVert} \le \Gamma(A) + \varepsilon\)。由 \(\varepsilon > 0\) 的任意性，立得 \(r \le \Gamma(A)\)。综上 \(\Gamma(A) = \lim_{n\to\infty}\sqrt[n]{\lVert A^n\rVert}\)。\(\blacksquare\)
    \end{itemize}
  \end{itemize}
\end{enumerate}

\SoftDivider

\BookSubSection{}{3.2 第二层（收窄 1）：紧算子的谱理论 --- Riesz-Schauder 定理}{3.2 第二层（收窄 1）：紧算子的谱理论 — Riesz-Schauder 定理}{\TagMust{闭卷必背}}

\begin{quote}
\textbf{老师划重点（\texttt{02:11–03:19}）}：紧算子的谱有哪些类？有什么样的谱点？谱点之间什么关系？——核心就是 \textbf{Riesz-Schauder 定理（教材定理 6.6--6.9，p.215--219）}！
\end{quote}

设 \(X\) 为复 Banach 空间，\(A \in \mathcal{C}(X)\)，必背 \textbf{5 条核心结论}：

\begin{enumerate}
\def\labelenumi{\arabic{enumi}.}
\tightlist
\item
  \textbf{非零谱点全是点谱（特征值）}：\(\sigma(A) \setminus \{0\} \subset \sigma_p(A)\)（不可能有非零的连续谱或剩余谱！）；
\item
  \textbf{每个非零特征值的特征空间必有限维}：对任意 \(\lambda \neq 0\)，\(\dim\operatorname{Ker}(\lambda I - A) < +\infty\)；
\item
  \textbf{非零谱点孤立、至多可数，唯一可能的聚点是 \(0\)}：对任意 \(r > 0\)，\(\{\lambda \in \sigma(A) : \lvert\lambda\rvert \ge r\}\) 只有有限个点；
\item
  \textbf{无穷维空间必有 \(0 \in \sigma(A)\)}（因为定理 5.5：无穷维紧算子没有有界逆！但注意 \(0\) 可能是 \(\sigma_p, \sigma_c, \sigma_r\) 中的任意一种）；
\item
  \textbf{Fredholm 择一性}：对 \(\lambda \neq 0\)，齐次方程 \((\lambda I - A)x = 0\) 仅有零解 \(\iff\) 对任意 \(y \in X\) 非齐次方程 \((\lambda I - A)x = y\) 存在唯一解（此时 \(\lambda \in \rho(A)\)）。
\end{enumerate}

\SubSubHeading{闭卷高频必背：Riesz-Schauder 的两大极简短证明（Riesz 引理反证法！）}

\begin{itemize}
\tightlist
\item
  \textbf{短证明 A：为什么 \(\lambda \neq 0\) 时 \(\dim\operatorname{Ker}(\lambda I - A) < \infty\)？}

  \begin{itemize}
  \tightlist
  \item
    在闭子空间 \(M = \operatorname{Ker}(\lambda I - A)\) 上，对任意 \(x \in M\) 有 \(Ax = \lambda x \implies I_M x = \frac{1}{\lambda}Ax\)。因 \(A\) 紧且 \(\lambda \neq 0\)，恒等算子 \(I_M = \frac{1}{\lambda}A|_M\) 是 \(M\) 上的紧算子，故 \(M\) 闭单位球列紧；由 Riesz 引理，\(\dim M < +\infty\)！
  \end{itemize}
\item
  \textbf{短证明 B：为什么非零特征值不能有非零聚点（即对任意 \(r > 0\)，模长 \(\ge r\) 的特征值只有有限个）？}

  \begin{itemize}
  \tightlist
  \item
    反证：假设存在两两互异的特征值列 \(\{\lambda_n\}\) 满足 \(\lvert\lambda_n\rvert \ge r > 0\)，对应非零特征向量 \(x_n\)。线代基础知互异特征值的特征向量线性无关。令 \(M_n = \operatorname{span}\{x_1, \dots, x_n\}\)，则 \(M_1 \subsetneq M_2 \subsetneq \cdots\) 是严格递增的有限维闭子空间列，且 \(A(M_n) \subset M_n\)，\((A - \lambda_n I)(M_n) \subset M_{n-1}\)。
  \item
    由 \textbf{Riesz 引理}，对每个 \(n \ge 2\) 可取 \(y_n \in M_n\) 满足 \(\lVert y_n\rVert = 1\) 且 \(\operatorname{dist}(y_n, M_{n-1}) \ge \frac{1}{2}\)。
  \item
    对任意 \(m < n\)，展开写： 
\FitDisplayMath{\frac{Ay_n}{\lambda_n} - \frac{Ay_m}{\lambda_m} = y_n - \underbrace{\left(\frac{Ay_m}{\lambda_m} - \frac{(A - \lambda_n I)y_n}{\lambda_n}\right)}_{\in M_{n-1}}.}
 因此 \(\left\lVert\frac{Ay_n}{\lambda_n} - \frac{Ay_m}{\lambda_m}\right\rVert \ge \operatorname{dist}(y_n, M_{n-1}) \ge \frac{1}{2}\)！这说明序列 \(\left\{A\!\left(\frac{y_n}{\lambda_n}\right)\right\}\) 没有任何收敛子列；但 \(\left\lVert\frac{y_n}{\lambda_n}\right\rVert = \frac{1}{\lvert\lambda_n\rvert} \le \frac{1}{r}\) 是\textbf{有界序列}，与 \(A\) 是紧算子矛盾！\(\blacksquare\)
  \end{itemize}
\end{itemize}

\SoftDivider

\BookSubSection{}{3.3 第三层（收窄 2，性质最好）：Hilbert 空间上的自伴算子与自伴紧算子}{3.3 第三层（收窄 2，性质最好）：Hilbert 空间上的自伴算子与自伴紧算子}{\TagMust{闭卷必背}}

\begin{quote}
\textbf{老师划重点（\texttt{03:19–05:19}）}：为什么还要放到 Hilbert 空间并加``自伴''？因为有了内积几何结构，谱半径公式变简单了（与一般 Gelfand 公式有区别）、谱被锁在实轴 \([-\lVert A\rVert, \lVert A\rVert]\) 里、还能像对称矩阵一样正交对角展开（\textbf{Hilbert-Schmidt 定理}）！
\end{quote}

\SubSubHeading{1. 有界自伴算子 \(A = A^* \in \mathcal{L}(H)\) 的五大必考短证明（闭卷送分题！）}

\begin{enumerate}
\def\labelenumi{\arabic{enumi}.}
\tightlist
\item
  \textbf{\((Ax, x) \in \mathbb{R}\)，特征值全实，不同特征值的特征向量正交（定理 6.10）}：

  \begin{itemize}
  \tightlist
  \item
    \(\overline{(Ax, x)} = (x, Ax) = (Ax, x) \implies (Ax, x) \in \mathbb{R}\)。
  \item
    若 \(Ax = \lambda x\)（\(x \neq 0\)），则 \(\lambda\lVert x\rVert^2 = (Ax, x) \in \mathbb{R} \implies \lambda \in \mathbb{R}\)。
  \item
    若 \(Ax = \lambda x, Ay = \mu y\) 且 \(\lambda \neq \mu\)，则 \(\lambda(x,y) = (Ax,y) = (x,Ay) = \mu(x,y) \implies (\lambda-\mu)(x,y)=0 \implies (x,y)=0\)。
  \end{itemize}
\item
  \textbf{谱全在实轴上 \(\sigma(A) \subset \mathbb{R}\)，且剩余谱为空 \(\sigma_r(A) = \emptyset\)}：

  \begin{itemize}
  \tightlist
  \item
    \textbf{证 \(\sigma(A) \subset \mathbb{R}\)}：设 \(\lambda = \alpha + \mathrm{i}\beta\)（\(\beta \neq 0\)）。因 \(A - \alpha I\) 自伴，直接展开（或套内积虚部）： 
\FitDisplayMath{\lVert(A - \lambda I)x\rVert^2 = \lVert(A - \alpha I)x\rVert^2 + \beta^2\lVert x\rVert^2 \ge \beta^2\lVert x\rVert^2 \implies \lVert(\lambda I - A)x\rVert \ge \lvert\beta\rvert\lVert x\rVert.}
 由此立得 \(\lambda I - A\) 是单射且值域闭！又若 \(y \perp \mathcal{R}(\lambda I - A)\)，则对一切 \(x\) 有 \(((\lambda I - A)x, y) = 0 \Rightarrow (x, (\bar{\lambda}I - A)y) = 0 \Rightarrow (\bar{\lambda}I - A)y = 0\)；因 \(\bar{\lambda} \notin \mathbb{R}\) 不是特征值，故 \(y = 0\)，即值域稠密！既闭又稠密 \(\implies \mathcal{R}(\lambda I - A) = H\)，故 \(\lambda \in \rho(A)\)！
  \item
    \textbf{证 \(\sigma_r(A) = \emptyset\)}：若 \(\lambda \in \sigma_r(A)\)，则 \(\lambda \in \mathbb{R}\) 且 \(\overline{\mathcal{R}(\lambda I - A)} \neq H\)，存在非零 \(y \perp \mathcal{R}(\lambda I - A)\)，同理得 \((\lambda I - A)y = 0\)，即 \(\lambda \in \sigma_p(A)\)，与 \(\sigma_r \cap \sigma_p = \emptyset\) 矛盾！
  \end{itemize}
\item
  \textbf{二次型算子范数公式（定理 6.11，教材 p.221）}： 
\FitDisplayMath{\lVert A\rVert = \sup_{\lVert x\rVert = 1} \lvert(Ax, x)\rvert.}


  \begin{itemize}
  \tightlist
  \item
    （直观：自伴算子的最大拉伸倍数一定在某个``同方向投影 \((Ax, x)\)''上就能测出来，且谱被锁在 \([m, M] \subset [-\lVert A\rVert, \lVert A\rVert]\) 内，其中 \(m = \inf_{\lVert x\rVert=1}(Ax,x),\; M = \sup_{\lVert x\rVert=1}(Ax,x)\)。）
  \end{itemize}
\item
  \textbf{\WarnBadge{} 老师强调的谱半径公式区别：自伴算子必有 \(\Gamma(A) = \lVert A\rVert\)（推论 6.2）}：

  \begin{itemize}
  \tightlist
  \item
    \textbf{超短闭卷神证（4 行拿下！）}：

    \begin{itemize}
    \tightlist
    \item
      先证 \(\lVert A^2\rVert = \lVert A\rVert^2\)：一方面 \(\lVert A^2\rVert \le \lVert A\rVert^2\)；另一方面对任意 \(x\)，由自伴性有 \(\lVert Ax\rVert^2 = (Ax, Ax) = (A^2 x, x) \le \lVert A^2\rVert\lVert x\rVert^2\)，取上确界得 \(\lVert A\rVert^2 \le \lVert A^2\rVert\)，故 \(\lVert A^2\rVert = \lVert A\rVert^2\)！
    \item
      因为 \(A^2, A^4, \dots, A^{2^k}\) 均自伴，数学归纳立得 \(\lVert A^{2^k}\rVert = \lVert A\rVert^{2^k}\)！
    \item
      代入 Gelfand 谱半径公式（取子列 \(n = 2^k\)）： 
\FitDisplayMath{\Gamma(A) = \lim_{k\to\infty} \lVert A^{2^k}\rVert^{\frac{1}{2^k}} = \lim_{k\to\infty} \left(\lVert A\rVert^{2^k}\right)^{\frac{1}{2^k}} = \lVert A\rVert! \qquad \blacksquare}

    \end{itemize}
  \end{itemize}
\end{enumerate}

\SubSubHeading{2. 自伴紧算子的两大支柱定理（定理 6.12 \& 定理 6.13）}

\begin{itemize}
\tightlist
\item
  \textbf{定理 6.12（必达范数特征值，教材 p.221）}：若 \(A\) 是 Hilbert 空间上的非零自伴紧算子，则 \(+\lVert A\rVert\) 或 \(-\lVert A\rVert\) 至少有一个是 \(A\) 的特征值！

  \begin{itemize}
  \tightlist
  \item
    \textbf{证明思路（3 步记牢）}：

    \begin{enumerate}
    \def\labelenumi{\arabic{enumi}.}
    \tightlist
    \item
      由定理 6.11，取单位向量列 \(\{x_n\}\)（\(\lVert x_n\rVert = 1\)）使 \(\lvert(Ax_n, x_n)\rvert \to \lVert A\rVert\)。因 \((Ax_n, x_n) \in \mathbb{R}\)，可抽子列使 \((Ax_{n_k}, x_{n_k}) \to \lambda\)，其中 \(\lambda = +\lVert A\rVert\) 或 \(-\lVert A\rVert\)（\(\lambda \neq 0\)）；
    \item
      展开 \(\lVert Ax_{n_k} - \lambda x_{n_k}\rVert^2 = \lVert Ax_{n_k}\rVert^2 - 2\lambda(Ax_{n_k}, x_{n_k}) + \lambda^2\lVert x_{n_k}\rVert^2 \le \lVert A\rVert^2 - 2\lambda(Ax_{n_k}, x_{n_k}) + \lambda^2 \to 0\)，即 \textbf{\(Ax_{n_k} - \lambda x_{n_k} \to 0\)}；
    \item
      因为 \(A\) 紧、\(\{x_{n_k}\}\) 有界，再抽子列使 \(Ax_{n_{k_j}} \to y\)，结合上式立得 \(x_{n_{k_j}} \to \frac{y}{\lambda} \triangleq x_0\)（\(\lVert x_0\rVert = 1 \neq 0\)），再由 \(A\) 连续即得 \(Ax_0 = \lambda x_0\)！\(\blacksquare\)
    \end{enumerate}
  \end{itemize}
\item
  \textbf{定理 6.13（Hilbert-Schmidt 谱展开定理，教材 p.223）}：

  \begin{itemize}
  \tightlist
  \item
    \textbf{陈述}：设 \(A\) 是 Hilbert 空间 \(H\) 上的非零自伴紧算子，则存在有限或可数个非零实特征值 \(\{\lambda_n\}\)（按绝对值递减 \(\lvert\lambda_1\rvert \ge \lvert\lambda_2\rvert \ge \cdots > 0\)，无穷时 \(\lambda_n \to 0\)）及对应的标准正交特征向量系 \(\{e_n\}\)，使得对任意 \(x \in H\)，有唯一正交分解与按范数收敛的展开式： 
\FitDisplayMath{x = x_0 + \sum_{n} (x, e_n)e_n \quad (x_0 \in \operatorname{Ker}A), \qquad Ax = \sum_{n} \lambda_n (x, e_n)e_n.}

  \item
    \textbf{直觉构造（为什么一定能一步步扒皮扒干净？）}：

    \begin{itemize}
    \tightlist
    \item
      第 1 步：在 \(H_1 = H\) 上用定理 6.12 拿出模最大的特征值 \(\lambda_1\)（\(\lvert\lambda_1\rvert = \lVert A\rVert\)）和单位特征向量 \(e_1\)；
    \item
      第 2 步：限制到正交补 \(H_2 = \{e_1\}^\perp\)（因为 \(A\) 自伴，\(H_2\) 仍是 \(A\)-不变闭子空间！），再拿出次大特征值 \(\lambda_2\) 和 \(e_2\)；
    \item
      不断剥离下去，剩余空间上算子范数 \(\lVert A|_{H_k}\rVert = \lvert\lambda_k\rvert \to 0\)，最终剩下的只能全在 \(\operatorname{Ker}A\) 里！
    \end{itemize}
  \end{itemize}
\end{itemize}

\SoftDivider

\BookSubSection{}{3.4 第四层（过渡桥梁，老师要求``大概了解''）：谱测度、谱系与谱积分（6.4 节）}{3.4 第四层（过渡桥梁，老师要求“大概了解”）：谱测度、谱系与谱积分（6.4 节）}{\TagLight{老师减负·大概了解}}

\begin{quote}
\textbf{零基础 1 分钟顿悟}：

\begin{itemize}
\tightlist
\item
  在有限维对称矩阵或自伴紧算子里，谱展开是\textbf{离散求和}：\(A = \sum_n \lambda_n P_n\)，其中 \(P_n\) 是往第 \(n\) 个特征空间投的正交投影算子（满足 \(P_n^2 = P_n = P_n^*\)，且不同特征空间互相垂直 \(P_m P_n = 0\)）。
\item
  可对于一般的\textbf{非紧}有界自伴算子（比如乘法算子 \((Af)(x) = x f(x)\) 在 \(L^2(0,1)\) 上），它\textbf{根本没有特征值}，谱是一整段连续区间 \([0, 1]\)！
\item
  怎么办？把``按孤立点求和 \(\sum \lambda_n P_n\)''升级成``按区间微元做积分 \(\int \lambda\,\mathrm{d}E_\lambda\)''——每给一个区间 \(\Delta = (a, b]\)，就对应一个\textbf{正交投影算子} \(E(\Delta)\)（例如对乘法算子就是乘特征函数 \(\chi_\Delta(x)\)）！
\end{itemize}
\end{quote}

只需掌握以下 \textbf{3 个概念 + 4 条运算法则}（足够应付闭卷概念题和\hyperlink{prob:6-39}{\textcolor{LinkColor}{习题 6 第 39 题}}、\hyperlink{prob:6-45}{\textcolor{LinkColor}{第 45 题}}！）：

\begin{enumerate}
\def\labelenumi{\arabic{enumi}.}
\tightlist
\item
  \textbf{谱测度 \(E\)（定义 6.7，教材 p.237）}：定义在可测空间 \((\Omega, \Sigma)\) 上、取值于 \(H\) 上\textbf{正交投影算子}的映射 \(E: \Sigma \to \mathcal{L}(H)\)，满足：

  \begin{itemize}
  \tightlist
  \item
    对任意 \(\Delta \in \Sigma\)，\(E(\Delta)\) 是正交投影（\(E(\Delta)^2 = E(\Delta) = E(\Delta)^*\)）；
  \item
    \(E(\emptyset) = 0,\; E(\Omega) = I\)；
  \item
    \(E(\Delta_1 \cap \Delta_2) = E(\Delta_1)E(\Delta_2)\)（特别互斥时正交 \(E(\Delta_1)E(\Delta_2) = 0\)）；
  \item
    可数互斥并满足\textbf{强可加性}：对任意 \(x \in H\)，\(E\!\left(\bigcup_{n=1}^\infty \Delta_n\right)x = \sum_{n=1}^\infty E(\Delta_n)x\)。
  \end{itemize}
\item
  \textbf{谱系 \(\{E_\lambda\}_{\lambda\in\mathbb{R}}\)（定义 6.10，教材 p.243）}：随 \(\lambda \in \mathbb{R}\) 变化的单调正交投影族，满足：

  \begin{itemize}
  \tightlist
  \item
    单调性：\(E_\lambda E_\mu = E_{\min(\lambda, \mu)}\)；
  \item
    强右连续：\(\lim_{\mu \to \lambda^+} E_\mu x = E_\lambda x\)；
  \item
    两端极限：\(\lim_{\lambda \to -\infty} E_\lambda x = 0,\; \lim_{\lambda \to +\infty} E_\lambda x = x\)。
  \item
    \textbf{与谱测度的对应}：\(E((a, b]) = E_b - E_a\)（类比概率论中``分布函数 \(F(\lambda)\) 决定概率测度 \(\mu((a,b]) = F(b)-F(a)\)''！）。
  \end{itemize}
\item
  \textbf{谱积分 \(\int_\Omega f(\omega)\,\mathrm{d}E(\omega)\) 的四大黄金运算法则（定理 6.15--6.16、推论 6.5--6.8）}： 对任意有界可测函数 \(f, g \in B(\Omega, \Sigma)\)，记 \(T_f = \int_\Omega f(\omega)\,\mathrm{d}E(\omega) \in \mathcal{L}(H)\)：

  \begin{itemize}
  \tightlist
  \item
    \textbf{法则 1（线性性）}：\(T_{\alpha f + \beta g} = \alpha T_f + \beta T_g\)；
  \item
    \textbf{法则 2（乘法变复合——谱积分最神奇的性质！）}：\(T_f T_g = T_{fg}\)，特别对多项式有 \(\int \lambda^n\,\mathrm{d}E_\lambda = \left(\int \lambda\,\mathrm{d}E_\lambda\right)^n\)；
  \item
    \textbf{法则 3（复共轭变共轭算子）}：\((T_f)^* = T_{\bar{f}}\)（故 \(f\) 为实函数时 \(T_f\) 自伴；\(\lvert f\rvert \equiv 1\) 时 \(T_f\) 为酉算子！）；
  \item
    \textbf{法则 4（范数估计与弱内积定义）}：\(\lVert T_f\rVert \le \sup_{\omega\in\Omega}\lvert f(\omega)\rvert\)，且 \((T_f x, y) = \int_\Omega f(\omega)\,\mathrm{d}(E(\omega)x, y)\)。
  \end{itemize}
\end{enumerate}

\SoftDivider

\BookSubSection{}{3.5 第五层（扩宽 1）：酉算子的定义、等价刻画与谱分解（6.5 节）}{3.5 第五层（扩宽 1）：酉算子的定义、等价刻画与谱分解（6.5 节）}{\TagMust{闭卷必背}}

\begin{quote}
\textbf{老师划重点（\texttt{05:32–06:06}）}：特别点名了\textbf{酉算子的定义（保距且满射）}、\textbf{保距/酉的等价定义方式}、以及\textbf{酉算子的谱分解定理}！
\end{quote}

\begin{enumerate}
\def\labelenumi{\arabic{enumi}.}
\tightlist
\item
  \textbf{严格定义（定义 6.11--6.12，教材 p.247）}：

  \begin{itemize}
  \tightlist
  \item
    \textbf{保距算子（等距算子）}：\(U \in \mathcal{L}(H_1, H_2)\) 满足 \(\lVert Ux\rVert = \lVert x\rVert\)（\(\forall x \in H_1\)）。（注：右移算子 \(R\) 是保距算子，但\textbf{不是满射}，所以不是酉算子！）
  \item
    \textbf{酉算子}：\(U\) 是\textbf{保距算子并且是满射}（刚体旋转/镜像反射：不改变长度、不改变夹角，且覆盖整个目标空间！）。
  \end{itemize}
\item
  \textbf{等价刻画（引理 6.6，闭卷必会证明！）}： 设 \(U \in \mathcal{L}(H)\)，以下四条等价：

  \begin{itemize}
  \tightlist
  \item
    \begin{enumerate}
    \def\labelenumii{(\roman{enumii})}
    \tightlist
    \item
      \(U\) 是酉算子（保距且满射）；
    \end{enumerate}
  \item
    (ii) \(U\) 保内积且满射：\((Ux, Uy) = (x, y)\)（\(\forall x, y \in H\)）且 \(\mathcal{R}(U) = H\)；
  \item
    (iii) \(U^*U = UU^* = I\)；
  \item
    (iv) \(U\) 可逆且 \(U^{-1} = U^*\)。
  \item
    \textbf{关键过渡（极化恒等式，保距 \(\Rightarrow\) 保内积）}：对复内积空间， 
\FitDisplayMath{(x, y) = \frac{1}{4}\left(\lVert x+y\rVert^2 - \lVert x-y\rVert^2 + \mathrm{i}\lVert x+\mathrm{i}y\rVert^2 - \mathrm{i}\lVert x-\mathrm{i}y\rVert^2\right).}
 因 \(U\) 线性且保范，把每一项里的向量换成 \(Ux, Uy\) 范数全不变，立刻推出 \((Ux, Uy) = (x, y)\)！从而 \((U^*Ux, y) = (Ux, Uy) = (x, y) \Rightarrow U^*U = I\)。再结合 \(U\) 是双射（保距自动单射 + 满射），右乘 \(U^{-1}\) 立得 \(U^* = U^{-1}\)，进而 \(UU^* = I\)！
  \end{itemize}
\item
  \textbf{酉算子谱的性质（推论 6.10，闭卷高频短证明！）}：

  \begin{itemize}
  \tightlist
  \item
    \textbf{结论}：酉算子的谱必落在复平面\textbf{单位圆周}上：\(\sigma(U) \subset \{\lambda \in \mathbb{C} : \lvert\lambda\rvert = 1\}\)，且 \(\Gamma(U) = \lVert U\rVert = 1\)！
  \item
    \textbf{4 行完整证明}：

    \begin{itemize}
    \tightlist
    \item
      因 \(\lVert U\rVert = 1\)，当 \(\lvert\lambda\rvert > 1\) 时由定理 6.1 立得 \(\lambda \in \rho(U)\)；
    \item
      当 \(\lambda = 0\) 时，\(0I - U = -U\) 可逆（逆为 \(-U^*\)），故 \(0 \in \rho(U)\)；
    \item
      当 \(0 < \lvert\lambda\rvert < 1\) 时，改写 \(\lambda I - U = -\lambda U\!\left(\frac{1}{\lambda}I - U^{-1}\right) = -\lambda U\!\left(\frac{1}{\lambda}I - U^*\right)\)。因为 \(\lVert U^*\rVert = \lVert U\rVert = 1 < \frac{1}{\lvert\lambda\rvert}\)，括号内可逆，且 \(-\lambda U\) 可逆，故 \(\lambda I - U\) 可逆，即 \(\lambda \in \rho(U)\)！
    \item
      综上，所有谱点只能满足 \(\lvert\lambda\rvert = 1\)！\(\blacksquare\)
    \end{itemize}
  \end{itemize}
\item
  \textbf{酉算子谱分解定理（定理 6.20--6.21，教材 p.250、256）}： 设 \(U\) 是复 Hilbert 空间 \(H\) 上的酉算子，则存在 \([0, 2\pi]\) 上唯一谱系 \(\{E_\theta\}\)（或单位圆周上的唯一谱测度 \(E\)），使得 
\FitDisplayMath{U = \int_0^{2\pi} \mathrm{e}^{\mathrm{i}\theta}\,\mathrm{d}E_\theta = \int_{\sigma(U)} \lambda\,\mathrm{d}E(\lambda).}

\end{enumerate}

\SoftDivider

\BookSubSection{}{3.6 第六层（扩宽 2）：有界自伴算子的谱分解（6.6 节）}{3.6 第六层（扩宽 2）：有界自伴算子的谱分解（6.6 节）}{\TagMust{闭卷必背}}

\begin{quote}
\textbf{老师划重点（\texttt{06:06–07:43}）}：把酉算子的谱分解通过 \textbf{Cayley 变换} 拉回实轴，就完成了``先收窄 \(\to\) 再扩宽''的大闭环，得到了一般有界自伴算子的谱分解！
\end{quote}

\begin{enumerate}
\def\labelenumi{\arabic{enumi}.}
\tightlist
\item
  \textbf{核心桥梁——Cayley 变换（定义 6.15，对应\hyperlink{prob:6-23}{\textcolor{LinkColor}{习题 6 第 23 题}}！）}： 设 \(A \in \mathcal{L}(H)\) 是有界自伴算子，定义 
\FitDisplayMath{U \triangleq (A - \mathrm{i}I)(A + \mathrm{i}I)^{-1}.}


  \begin{itemize}
  \tightlist
  \item
    \textbf{几何直觉}：复变函数中的分式线性变换 \(\varphi(t) = \frac{t - \mathrm{i}}{t + \mathrm{i}}\) 把整条实轴 \(\mathbb{R}\) 一一映到单位圆周除去点 \(1\) 的部分 \(S^1 \setminus \{1\}\)！
  \item
    \textbf{完整验证 \(U\) 是酉算子且 \(1 \in \rho(U)\)（闭卷高频题！）}：

    \begin{enumerate}
    \def\labelenumii{\arabic{enumii}.}
    \tightlist
    \item
      因 \(A\) 自伴故 \(\sigma(A) \subset \mathbb{R}\)，于是 \(\pm\mathrm{i} \in \rho(A)\)，\((A + \mathrm{i}I)^{-1}\) 和 \((A - \mathrm{i}I)^{-1}\) 均存在有界；
    \item
      算共轭：\(U^* = \big[(A + \mathrm{i}I)^{-1}\big]^*(A - \mathrm{i}I)^* = (A - \mathrm{i}I)^{-1}(A + \mathrm{i}I)\)；
    \item
      因 \(A - \mathrm{i}I\) 与 \(A + \mathrm{i}I\) 可交换（从而与各自的逆都可交换），直接相乘： 
\FitDisplayMath{U^*U = (A - \mathrm{i}I)^{-1}(A + \mathrm{i}I)(A - \mathrm{i}I)(A + \mathrm{i}I)^{-1} = I, \qquad UU^* = I \implies U \text{ 是酉算子！}}

    \item
      算 \(I - U\)：\(I - U = \big[(A+\mathrm{i}I) - (A-\mathrm{i}I)\big](A+\mathrm{i}I)^{-1} = 2\mathrm{i}(A+\mathrm{i}I)^{-1}\)，其逆为 \(\frac{1}{2\mathrm{i}}(A+\mathrm{i}I) \in \mathcal{L}(H)\)，故 \(1 \in \rho(U)\)，而且反解出 \(A = \mathrm{i}(I + U)(I - U)^{-1}\)！
    \end{enumerate}
  \end{itemize}
\item
  \textbf{有界自伴算子谱分解定理（定理 6.24--6.25，教材 p.259、262）}： 设 \(A\) 是复 Hilbert 空间 \(H\) 上的有界自伴算子，则存在唯一的谱测度 \(E\)（支集为紧集 \(\sigma(A) \subset [m, M] \subset [-\lVert A\rVert, \lVert A\rVert]\)）及对应的唯一谱系 \(\{E_\lambda\}\)（当 \(\lambda < m\) 时 \(E_\lambda = 0\)，当 \(\lambda \ge M\) 时 \(E_\lambda = I\)），使得 
\FitDisplayMath{A = \int_{\sigma(A)} \lambda\,\mathrm{d}E(\lambda) = \int_{m-0}^{M} \lambda\,\mathrm{d}E_\lambda.}

\item
  \textbf{Borel 函数演算（推论 6.15--6.17，教材 p.264）}： 对任意有界 Borel 函数 \(f \in B(\sigma(A))\)，定义算子 \(f(A) \triangleq \int_{\sigma(A)} f(\lambda)\,\mathrm{d}E(\lambda)\)。它保持一切代数运算：

  \begin{itemize}
  \tightlist
  \item
    \((f+g)(A) = f(A)+g(A)\)，\((fg)(A) = f(A)g(A)\)，\(\bar{f}(A) = f(A)^*\)，\(\lVert f(A)\rVert \le \sup_{\lambda\in\sigma(A)}\lvert f(\lambda)\rvert\)；
  \item
    若 \(A \ge 0\)（即 \((Ax,x)\ge 0\)，等价于 \(\sigma(A)\subset[0,\infty)\)），则存在唯一正自伴平方根 \(A^{1/2} = \int_{\sigma(A)} \sqrt{\lambda}\,\mathrm{d}E(\lambda)\) 满足 \((A^{1/2})^2 = A\)；
  \item
    若 \(f(\lambda) = \sum_{n=1}^\infty \alpha_n \lambda^n\) 在 \(\lvert\lambda\rvert < \Gamma(A) + \varepsilon\) 解析，则 \(f(A) = \sum_{n=1}^\infty \alpha_n A^n\)（\hyperlink{prob:6-45}{\textcolor{LinkColor}{习题 6 第 45 题}}）；
  \item
    若 \(\tilde{A} = U^{-1}AU\)（酉等价），则 \(f(\tilde{A}) = U^{-1}f(A)U\)（\hyperlink{prob:6-39}{\textcolor{LinkColor}{习题 6 第 39 题}}）。
  \end{itemize}
\end{enumerate}

\SubSubHeading{闭卷终极对照表：第 6 章 + 第 7 章六类算子谱与分解全景}

{\def\LTcaptype{none} % do not increment counter
{\small\renewcommand{\arraystretch}{1.38}\setlength{\tabcolsep}{5pt}
\begin{longtable}[]{@{}
  >{\raggedright\arraybackslash}p{(\linewidth - 10\tabcolsep) * \real{0.1389}}
  >{\centering\arraybackslash}p{(\linewidth - 10\tabcolsep) * \real{0.0972}}
  >{\raggedright\arraybackslash}p{(\linewidth - 10\tabcolsep) * \real{0.2778}}
  >{\centering\arraybackslash}p{(\linewidth - 10\tabcolsep) * \real{0.0972}}
  >{\centering\arraybackslash}p{(\linewidth - 10\tabcolsep) * \real{0.1528}}
  >{\raggedright\arraybackslash}p{(\linewidth - 10\tabcolsep) * \real{0.2361}}@{}}
\toprule\noalign{}
\rowcolor{AccentSoft}
\begin{minipage}[b]{\linewidth}\raggedright
算子类别
\end{minipage} & \begin{minipage}[b]{\linewidth}\centering
空间
\end{minipage} & \begin{minipage}[b]{\linewidth}\raggedright
谱集 \(\sigma(A)\) 位置与特征
\end{minipage} & \begin{minipage}[b]{\linewidth}\centering
剩余谱 \(\sigma_r\)
\end{minipage} & \begin{minipage}[b]{\linewidth}\centering
谱半径 \(\Gamma(A)\)
\end{minipage} & \begin{minipage}[b]{\linewidth}\raggedright
谱分解 / 展开形式
\end{minipage} \\
\midrule\noalign{}
\endhead
\bottomrule\noalign{}
\endlastfoot
\textbf{一般有界算子} & 复 Banach & 非空紧集，\(\subset\{\lvert\lambda\rvert\le\lVert A\rVert\}\) & 可非空 & \(\lim\lVert A^n\rVert^{1/n}\) & 无一般分解 \\
\textbf{紧算子} & 复 Banach & 至多可数，\(\sigma\setminus\{0\}\subset\sigma_p\)，仅 \(0\) 可聚 & 仅可能含 \(0\) & \(\sup_{\lambda\in\sigma_p}\lvert\lambda\rvert\) & Riesz-Schauder 结构 \\
\textbf{自伴紧算子} & 复 Hilbert & 实轴 \([-\lVert A\rVert, \lVert A\rVert]\) 上至多可数实特征值，仅 \(0\) 可聚 & \(\emptyset\) & \(\lVert A\rVert = \max\lvert\lambda_n\rvert\) & \(Ax = \sum_n \lambda_n(x,e_n)e_n\) \\
\textbf{酉算子} & 复 Hilbert & 单位圆周 \(\{\lvert\lambda\rvert = 1\}\) 闭子集 & \(\emptyset\) & \(1\) & \(U = \int_0^{2\pi}\mathrm{e}^{\mathrm{i}\theta}\,\mathrm{d}E_\theta\) \\
\textbf{一般有界自伴算子} & 复 Hilbert & 实轴有界闭区间 \([m, M]\subset[-\lVert A\rVert, \lVert A\rVert]\)（可连续） & \(\emptyset\) & \(\lVert A\rVert\) & \(A = \int_{\sigma(A)}\lambda\,\mathrm{d}E_\lambda\) \\
\textbf{无界自伴算子（第7章）} & 复 Hilbert & 实轴 \(\mathbb{R}\) 的闭子集（\textbf{可无界}！） & \(\emptyset\) & 可为 \(+\infty\) & \(A = \int_{\mathbb{R}}\lambda\,\mathrm{d}E_\lambda\) \\
\end{longtable}
}
}

\BookPart{}{第 7 章深度精讲 —— Hilbert 空间上的无界算子}{第 7 章深度精讲 —— Hilbert 空间上的无界算子}

\begin{quote}
\textbf{老师划重点（\texttt{08:22–10:30}）}：实际问题里最基本的算子（如\textbf{求导算子}）都是无界的！因此必须先把无界算子放到 Hilbert 空间里，通过\textbf{共轭、对称、自伴}做层层限定，再通过 Cayley 变换建立\textbf{无界自伴算子的谱分解}。
\end{quote}

\BookSubSection{}{4.1 为什么微分算子必然无界？定义域为什么决定算子的生死？}{4.1 为什么微分算子必然无界？定义域为什么决定算子的生死？}{}

\begin{enumerate}
\def\labelenumi{\arabic{enumi}.}
\item
  \textbf{直观例子（求导算子无界）}：在 \(L^2(0, 2\pi)\) 上取 \(u_n(x) = \frac{1}{\sqrt{\pi}}\sin(nx)\)，则 \(\lVert u_n\rVert_2 = 1\)，但求导后 \(u_n'(x) = \frac{n}{\sqrt{\pi}}\cos(nx)\)，范数 \(\lVert u_n'\rVert_2 = n \to +\infty\)！不存在任何有限常数 \(M\) 能压住所有 \(n\)。
\item
  \textbf{Hellinger-Toeplitz 定理（定理 7.3(i)，教材 p.280）}：

  \begin{quote}
  如果 \(A\) 是 \(H\) 上的对称算子且\textbf{定义在整个空间 \(\mathcal{D}(A) = H\) 上}，则 \(A\) 必是自伴算子，而且\textbf{必是有界算子} \(A \in \mathcal{L}(H)\)！
  \end{quote}

  \begin{itemize}
  \tightlist
  \item
    \textbf{极简 2 行证明}：因 \(A\) 对称且 \(\mathcal{D}(A) = H\)，有 \(H = \mathcal{D}(A) \subset \mathcal{D}(A^*) \subset H \Rightarrow A = A^*\)（自伴）；由定理 7.1 自伴算子必为闭算子；再由\textbf{闭图像定理}，定义在全空间 \(H\) 上的闭算子必有界！\(\blacksquare\)
  \item
    \textbf{惊人推论}：\textbf{只要想研究无界对称/自伴算子，它的定义域 \(\mathcal{D}(A)\) 绝不可能是全空间 \(H\)}，只能是某个真稠密子空间——因此描述一个无界算子，必须同时写清\textbf{「定义域 \(\mathcal{D}(A)\)」+「作用法则 \(Ax\)」}！
  \end{itemize}
\end{enumerate}

\BookSubSection{}{4.2 稠定算子、共轭算子、对称算子与自伴算子（7.1.1--7.1.2）}{4.2 稠定算子、共轭算子、对称算子与自伴算子（7.1.1–7.1.2）}{}

\begin{enumerate}
\def\labelenumi{\arabic{enumi}.}
\tightlist
\item
  \textbf{扩张（定义 7.1）}：记 \(A \subset B\)，意指 \(\mathcal{D}(A) \subset \mathcal{D}(B)\) 且对一切 \(x \in \mathcal{D}(A)\) 有 \(Ax = Bx\)。
\item
  \textbf{稠定算子与共轭算子 \(A^*\)（定义 7.2--7.3，教材 p.273）}：

  \begin{itemize}
  \tightlist
  \item
    设 \(\overline{\mathcal{D}(A)} = H\)（称 \(A\) \textbf{稠定}）。定义共轭算子 \(A^*\) 的定义域为： 
\FitDisplayMath{\mathcal{D}(A^*) \triangleq \left\{u \in H : \exists\, v \in H \text{ 使得 } (Ax, u) = (x, v),\; \forall x \in \mathcal{D}(A)\right\},}
 并规定 \(A^*u \triangleq v\)。
  \item
    \textbf{为什么必须要求 \(\mathcal{D}(A)\) 稠密？} 若存在两个 \(v_1, v_2\) 都满足等式，则 \((x, v_1 - v_2) = 0\) 对一切 \(x \in \mathcal{D}(A)\) 成立；因为 \(\mathcal{D}(A)\) 稠密，必有 \(v_1 - v_2 = 0\)，从而 \(v\) 唯一，保证 \(A^*\) 是单值良定义映射！
  \item
    \textbf{共轭的基本性质（引理 7.2，必背！）}：

    \begin{itemize}
    \tightlist
    \item
      \textbf{扩张反序律}：若 \(A, B\) 稠定且 \(A \subset B\)，则 \textbf{\(B^* \subset A^*\)}！（直觉：对 \(B^*\) 的考核要在更大的集合 \(\mathcal{D}(B)\) 上通过，条件更严，能过关的 \(u\) 反而更少！）
    \item
      \((\alpha A)^* = \bar{\alpha}A^*\)（\(\alpha \neq 0\)）；\(A^* + B^* \subset (A+B)^*\)（若 \(B \in \mathcal{L}(H)\) 则取等号）；\(B^*A^* \subset (AB)^*\)（若 \(A \in \mathcal{L}(H)\) 则取等号）。
    \end{itemize}
  \end{itemize}
\item
  \textbf{对称算子 vs 自伴算子（定义 7.4，教材 p.275，全章头号易混点！）}：

  \begin{itemize}
  \tightlist
  \item
    \textbf{对称算子}：\(A\) 稠定且 \textbf{\(A \subset A^*\)}（等价于：\(\forall x, y \in \mathcal{D}(A)\)，\((Ax, y) = (x, Ay)\)，但允许 \(\mathcal{D}(A) \subsetneq \mathcal{D}(A^*)\)！）；
  \item
    \textbf{自伴算子}：\(A\) 稠定且 \textbf{\(A = A^*\)}（即不仅在 \(\mathcal{D}(A)\) 上 \((Ax, y) = (x, Ay)\)，而且\textbf{定义域也必须精确相等 \(\mathcal{D}(A) = \mathcal{D}(A^*)\)}！）。
  \end{itemize}
\item
  \textbf{\StarBadge{} 全章灵魂活例：教材 p.275 例 7.2（三个求导算子彻底讲透对称与自伴的区别！）}： 在 \(H = L^2(0, 1)\) 上，定义\textbf{完全相同的微分公式} \(A_k f = \mathrm{i} f'\)（\(k = 1, 2, 3\)），仅改变定义域的边界条件：

  \begin{itemize}
  \tightlist
  \item
    \(\mathcal{D}(A_1) = \{f \in AC[0,1] : f' \in L^2(0,1)\}\)（\textbf{无任何边界条件}，定义域最大）；
  \item
    \(\mathcal{D}(A_2) = \{f \in \mathcal{D}(A_1) : f(0) = f(1)\}\)（\textbf{周期边界条件}，定义域居中）；
  \item
    \(\mathcal{D}(A_3) = \{f \in \mathcal{D}(A_1) : f(0) = f(1) = 0\}\)（\textbf{两端为零条件}，定义域最小）。 显然 \(A_3 \subsetneq A_2 \subsetneq A_1\)。分部积分得 \((A_k f, g) - (f, A_k g) = \mathrm{i}\big[f(1)\overline{g(1)} - f(0)\overline{g(0)}\big]\)，据此精确算出三者的共轭：
  \item
    \textbf{\(A_1^* = A_3 \subsetneq A_1\)}：定义域太大，连对称都不是！
  \item
    \textbf{\(A_3^* = A_1 \supsetneq A_3\)}：定义域太小，\(A_3 \subsetneq A_3^*\)，只是\textbf{对称算子，但不是自伴算子}！
  \item
    \textbf{\(A_2^* = A_2\)}：定义域不多不少刚刚好，是真正的\textbf{自伴算子}！（且验证了反序律 \(A_3 \subset A_2 \subset A_1 \implies A_1^* \subset A_2^* \subset A_3^*\)！）
  \item
    （顺带一提，习题 7 第 5 题就考了这三个算子的谱：\(\sigma_p(A_1) = \mathbb{C}\)；\(\sigma(A_2) = \{2n\pi : n\in\mathbb{Z}\} \subset \mathbb{R}\)；\(\sigma(A_3) = \mathbb{C}\)！可见\textbf{只有真正自伴的 \(A_2\)，谱才乖乖落在实轴上}，仅对称的 \(A_3\) 谱甚至铺满整个复平面！）
  \end{itemize}
\end{enumerate}

\SoftDivider

\BookSubSection{}{4.3 算子图像、闭算子与 \(A^{**}\)（7.1.4，必考短证明！）}{4.3 算子图像、闭算子与 A**（7.1.4，必考短证明！）}{}

\begin{enumerate}
\def\labelenumi{\arabic{enumi}.}
\tightlist
\item
  \textbf{算子图像与闭算子}：

  \begin{itemize}
  \tightlist
  \item
    图像 \(\mathcal{G}(A) \triangleq \{(x, Ax) : x \in \mathcal{D}(A)\} \subset H \times H\)；
  \item
    \(A\) 称为\textbf{闭算子} \(\iff \mathcal{G}(A)\) 是 \(H \times H\) 中的闭子空间 \(\iff\) 对任意 \(\{x_n\} \subset \mathcal{D}(A)\)，若 \(x_n \to x\) 且 \(Ax_n \to y\)，则必有 \(x \in \mathcal{D}(A)\) 且 \(Ax = y\)。
  \item
    \textbf{直观}：无界算子不再``连续''（\(x_n \to x\) 推不出 \(Ax_n\) 收敛），退而求其次，只要 \(x_n\) 和 \(Ax_n\) \textbf{双双收敛}，极限还能稳稳落在图像上不跑偏，这就是闭算子！
  \end{itemize}
\item
  \textbf{\WarnBadge{} 闭卷高频必考短证明：定理 7.1（稠定算子的共轭 \(A^*\) 恒为闭算子！）}：

  \begin{itemize}
  \tightlist
  \item
    \textbf{极简序列证法（4 行直接拿满分！）}： 设 \(\{u_n\} \subset \mathcal{D}(A^*)\) 满足 \(u_n \to u\) 且 \(A^*u_n \to v\)（于 \(H\) 中）。任取 \(x \in \mathcal{D}(A)\)，由共轭算子定义有 
\FitDisplayMath{(Ax, u_n) = (x, A^*u_n), \qquad \forall n \ge 1.}
 令 \(n \to \infty\)，由内积的连续性立得 
\FitDisplayMath{(Ax, u) = (x, v), \qquad \forall x \in \mathcal{D}(A).}
 对照 \(\mathcal{D}(A^*)\) 的定义，立刻知 \(u \in \mathcal{D}(A^*)\) 且 \(A^*u = v\)！故 \(A^*\) 是闭算子。特别地，若 \(A\) 自伴（\(A = A^*\)），则 \(A\) 必为闭算子！\(\blacksquare\)
  \item
    \textbf{教材几何证法（引理 7.6 \& 推论 7.1，p.279--280）}：在 \(H \times H\) 上定义酉翻转算子 \(V(x, y) = (-y, x)\)，则 \(\mathcal{G}(A^*) = \big[V\mathcal{G}(A)\big]^\perp\)。因为正交补永远是闭子空间，故 \(\mathcal{G}(A^*)\) 必闭！
  \end{itemize}
\item
  \textbf{定理 7.2（双共轭定理，教材 p.280）}： 若 \(A\) 是 \(H\) 上的\textbf{稠定闭算子}，则 \(\mathcal{D}(A^*)\) 必在 \(H\) 中稠密，且 \textbf{\(A^{**} = A\)}。
\item
  \textbf{对称算子的四级阶梯（直接秒杀\hyperlink{prob:7-6}{\textcolor{LinkColor}{习题 7 第 6 题}}！）}： 对任意对称算子 \(A\)，因 \(\mathcal{D}(A) \subset \mathcal{D}(A^*)\)，故 \(\mathcal{D}(A^*)\) 自动稠密，从而 \(A^{**}\) 必然存在（且等于 \(A\) 的最小闭扩张 \(\overline{A}\)），满足黄金链条： 
\FitDisplayMath{A \;\subset\; \overline{A} = A^{**} \;\subset\; A^*.}


  \begin{itemize}
  \tightlist
  \item
    \textbf{一般对称算子}：\(A \subset A^{**} \subset A^*\)；
  \item
    \textbf{闭对称算子}：\(A = A^{**} \subset A^*\)；
  \item
    \textbf{本质自伴算子}：\(\overline{A} = A^{**} = A^*\)；
  \item
    \textbf{自伴算子}：\(A = A^{**} = A^*\)（且引理 7.4：自伴算子是极大对称算子，若 \(A\) 自伴、\(B\) 对称且 \(A \subset B\)，则必有 \(A = B\)！）。
  \end{itemize}
\end{enumerate}

\SoftDivider

\BookSubSection{}{4.4 对称算子成为自伴算子的核心判别准则（7.1.5，必背证明！）}{4.4 对称算子成为自伴算子的核心判别准则（7.1.5，必背证明！）}{\TagMust{闭卷必背}}

\begin{enumerate}
\def\labelenumi{\arabic{enumi}.}
\tightlist
\item
  \textbf{万能能量恒等式（式 7.9，教材 p.281）}：设 \(A\) 是对称算子，对任意 \(x \in \mathcal{D}(A)\)，直接展开： 
\FitDisplayMath{\lVert(A \pm \mathrm{i}I)x\rVert^2 = (Ax \pm \mathrm{i}x, Ax \pm \mathrm{i}x) = \lVert Ax\rVert^2 \mp \mathrm{i}(Ax, x) \pm \mathrm{i}(x, Ax) + \lVert x\rVert^2 = \lVert Ax\rVert^2 + \lVert x\rVert^2.}
 （因为 \(A\) 对称，\((Ax, x) = (x, Ax)\)，中间两项一正一负完美抵消！）
\item
  \textbf{由式 (7.9) 秒推两大引理（引理 7.7--7.8）}：

  \begin{itemize}
  \tightlist
  \item
    \textbf{引理 7.7}：\(\lVert(A \pm \mathrm{i}I)x\rVert \ge \lVert x\rVert \implies \operatorname{Ker}(A \pm \mathrm{i}I) = \{0\}\)（\(A \pm \mathrm{i}I\) 单射且逆有界 \(\le 1\)！）；
  \item
    \textbf{引理 7.8}：若 \(A\) 是\textbf{闭对称算子}，则 \textbf{\(\mathcal{R}(A \pm \mathrm{i}I)\) 是闭子空间}！（证：若 \((A+\mathrm{i}I)x_n \to y\)，由式 7.9 知 \(\lVert x_m - x_n\rVert \le \lVert(A+\mathrm{i}I)(x_m-x_n)\rVert \to 0\)，故 \(x_n \to x\)；从而 \(Ax_n = (A+\mathrm{i}I)x_n - \mathrm{i}x_n \to y - \mathrm{i}x\)；由 \(A\) 闭得 \(x \in \mathcal{D}(A)\) 且 \(Ax = y - \mathrm{i}x\)，即 \((A+\mathrm{i}I)x = y \in \mathcal{R}(A+\mathrm{i}I)\)！——对比\hyperlink{prob:7-7}{\textcolor{LinkColor}{习题 7 第 7 题}}的非负对称算子 \(A+I\)，套路一模一样！）
  \end{itemize}
\item
  \textbf{引理 7.9（值域正交补公式）}：对任意稠定算子 \(A\)，有 \(\operatorname{Ker}A^* = (\mathcal{R}(A))^\perp\)，\(\overline{\mathcal{R}(A)} = (\operatorname{Ker}A^*)^\perp\)。
\item
  \textbf{\StarBadge{} 定理 7.4（自伴判别基本准则，教材 p.281--282，闭卷绝对必背！）}： 设 \(A\) 是复 Hilbert 空间 \(H\) 上的对称算子，则以下三条\textbf{等价}：

  \begin{itemize}
  \tightlist
  \item
    \begin{enumerate}
    \def\labelenumii{(\roman{enumii})}
    \tightlist
    \item
      \(A\) 是\textbf{自伴算子}（\(A = A^*\)）；
    \end{enumerate}
  \item
    (ii) \(A\) 是\textbf{闭算子}，且 \(\operatorname{Ker}(A^* + \mathrm{i}I) = \operatorname{Ker}(A^* - \mathrm{i}I) = \{0\}\)；
  \item
    (iii) \textbf{\(\mathcal{R}(A + \mathrm{i}I) = \mathcal{R}(A - \mathrm{i}I) = H\)}。
  \end{itemize}

  \begin{quote}
  \textbf{串联记忆}：(ii) \(\Leftrightarrow\) (iii) 就是用引理 7.8（闭性）+ 引理 7.9（正交补 \(= \{0\} \Leftrightarrow\) 稠密）秒出！而 (iii) \(\Rightarrow\) (i) 只需任取 \(u \in \mathcal{D}(A^*)\)，由满射存在 \(x \in \mathcal{D}(A)\) 使 \((A+\mathrm{i}I)x = (A^*+\mathrm{i}I)u\)，因 \(A \subset A^*\) 即 \((A^*+\mathrm{i}I)(u-x)=0\)，而 \(u-x \in \operatorname{Ker}(A^*+\mathrm{i}I) = \mathcal{R}(A-\mathrm{i}I)^\perp = H^\perp = \{0\}\)，故 \(u = x \in \mathcal{D}(A)\)，即 \(A^* = A\)！
  \end{quote}
\end{enumerate}

\SoftDivider

\BookSubSection{}{4.5 无界自伴算子的谱、Cayley 变换与谱分解（7.2 节）}{4.5 无界自伴算子的谱、Cayley 变换与谱分解（7.2 节）}{\TagMust{闭卷必背}}

\begin{enumerate}
\def\labelenumi{\arabic{enumi}.}
\tightlist
\item
  \textbf{无界闭算子的预解集与谱（定义 7.6--7.7，教材 p.282）}：

  \begin{itemize}
  \tightlist
  \item
    设 \(A\) 为闭算子。\(\lambda \in \rho(A) \iff \lambda I - A: \mathcal{D}(A) \to H\) 是双射且 \((\lambda I - A)^{-1} \in \mathcal{L}(H)\)（由闭图像定理，只要是双射即自动有界！）；
  \item
    \(\sigma(A) = \mathbb{C} \setminus \rho(A)\)，同样分 \(\sigma_p, \sigma_c, \sigma_r\)。
  \end{itemize}
\item
  \textbf{无界自伴算子的谱必在实轴上（定理 7.5 \& 推论 7.2，教材 p.283）}：

  \begin{itemize}
  \tightlist
  \item
    \textbf{定理 7.5}：对自伴算子 \(A\)，\(\lambda \in \rho(A) \iff \exists\, C > 0\) 使 \(\lVert(\lambda I - A)x\rVert \ge C\lVert x\rVert\)（\(\forall x \in \mathcal{D}(A)\)）；
  \item
    \textbf{推论 7.2（证 \(\sigma(A) \subset \mathbb{R}\)）}：对任意非实数 \(\lambda = \alpha + \mathrm{i}\beta\)（\(\beta \neq 0\)），将式 (7.9) 用于对称算子 \(\frac{1}{\beta}(A - \alpha I)\) 得 \(\lVert(\lambda I - A)x\rVert^2 = \lVert(A-\alpha I)x\rVert^2 + \beta^2\lVert x\rVert^2 \ge \beta^2\lVert x\rVert^2\)，由定理 7.5 立得 \(\lambda \in \rho(A)\)！故 \(\sigma(A) \subset \mathbb{R}\)。
  \end{itemize}
\item
  \textbf{无界算子的 Cayley 变换（定理 7.6--7.7，教材 p.285--286）}：

  \begin{itemize}
  \tightlist
  \item
    对对称算子 \(A\)，因 \(\lVert(A-\mathrm{i}I)x\rVert = \lVert(A+\mathrm{i}I)x\rVert\)，定义在 \(\mathcal{R}(A+\mathrm{i}I)\) 上的映射 
\FitDisplayMath{U \triangleq (A - \mathrm{i}I)(A + \mathrm{i}I)^{-1} : \mathcal{R}(A + \mathrm{i}I) \to \mathcal{R}(A - \mathrm{i}I)}
 是保距算子（且 \(1\) 不是 \(U\) 的特征值）。
  \item
    \textbf{核心等价}：\(A\) 是\textbf{自伴算子} \(\iff \mathcal{R}(A+\mathrm{i}I) = \mathcal{R}(A-\mathrm{i}I) = H \iff\) \textbf{\(U\) 是整个 \(H\) 上的酉算子}！
  \item
    对比第 6 章有界情形：有界自伴算子对应 \(1 \in \rho(U)\)（\(1\) 不在谱里）；无界自伴算子对应 \(1 \notin \sigma_p(U)\)（但 \(1 \in \sigma_c(U)\)，因为 \(A = \mathrm{i}(I+U)(I-U)^{-1}\) 无界！）。
  \end{itemize}
\item
  \textbf{定理 7.8（无界自伴算子谱分解定理，教材 p.290）}： 设 \(A\) 是 \(H\) 上的自伴算子，则存在 \(\mathbb{R}\) 上唯一的谱系 \(\{E_\lambda\}_{\lambda\in\mathbb{R}}\)，使得： 
\FitDisplayMath{\mathcal{D}(A) = \left\{x \in H : \int_{\mathbb{R}} \lambda^2\,\mathrm{d}\lVert E_\lambda x\rVert^2 < +\infty\right\}, \qquad (Ax, y) = \int_{\mathbb{R}} \lambda\,\mathrm{d}(E_\lambda x, y) \quad (\forall x\in\mathcal{D}(A), y\in H),}
 简记为 \(A = \int_{\mathbb{R}} \lambda\,\mathrm{d}E_\lambda\)。

  \begin{itemize}
  \tightlist
  \item
    \textbf{典型实例（命题 7.1--7.2，与\hyperlink{prob:7-4}{\textcolor{LinkColor}{习题 7 第 4 题}}对应！）}：对 \(L^2(\mathbb{R})\) 上的位置乘法算子 \((Ax)(t) = tx(t)\)，其谱为整条实轴 \(\sigma(A) = \mathbb{R}\)，谱系就是截断投影 \((E_\lambda x)(t) = \chi_{(-\infty, \lambda]}(t) x(t)\)！
  \end{itemize}
\end{enumerate}

\BookPart{}{第 8 章详细补充精讲（对应「第八章增」8 道作业）}{第 8 章详细补充精讲（对应「第八章增」8 道作业）}

\begin{quote}
虽然本段录音未录到第 8 章，但考虑到布置了第 8 章新增作业，这里按同样的「框架 \(\to\) 定义 \(\to\) 性质 \(\to\) 思路」把第 8 章闭卷必拿分的核心知识全部拆清！
\end{quote}

\BookSubSection{}{5.1 全章空间包含链与本书两大核心约定（防踩坑必背！）}{5.1 全章空间包含链与本书两大核心约定（防踩坑必背！）}{}

\begin{enumerate}
\def\labelenumi{\arabic{enumi}.}
\tightlist
\item
  \textbf{五大函数/分布空间包含链}： 
\FitDisplayMath{\mathcal{D}(\mathbb{R}^n) = C_c^\infty(\mathbb{R}^n) \;\subset\; \mathcal{S}(\mathbb{R}^n) \;\subset\; L^2(\mathbb{R}^n) \;\subset\; \mathcal{S}'(\mathbb{R}^n) \;\subset\; \mathcal{D}'(\mathbb{R}^n).}


  \begin{itemize}
  \tightlist
  \item
    试验函数空间 \(\mathcal{D}(\Omega) = C_c^\infty(\Omega)\)（无穷次可导 + 紧支集，边界外全为 \(0\)，分部积分零边界项！）；
  \item
    速降函数空间 \(\mathcal{S}(\mathbb{R}^n)\)（Schwartz 类：任意阶导数乘任意多项式有界，\(\lim_{\lvert x\rvert\to\infty} x^\alpha\partial^\beta f = 0\)）；
  \item
    缓增分布 \(\mathcal{S}'(\mathbb{R}^n)\)（\(\mathcal{S}\) 的连续对偶，允许做 Fourier 变换！）；
  \item
    广义函数 \(\mathcal{D}'(\Omega)\)（\(\mathcal{D}\) 的连续对偶，最宽泛的分布空间）。
  \end{itemize}
\item
  \textbf{\WarnBadge{} 本书两大核心约定（切勿与其他教材/网课混淆！）}：

  \begin{itemize}
  \tightlist
  \item
    \textbf{约定 A（常规广义函数配对，p.304 例 8.2）}：\(g \in L^1_{\mathrm{loc}}(\Omega)\) 对应的广义函数定义为 \(\lambda_g(\varphi) \triangleq \int_\Omega \varphi(x)\overline{g(x)}\,\mathrm{d}x\)（光滑函数乘法 \((h\lambda)(\varphi) = \lambda(\bar{h}\varphi)\)；对实值函数 \(g\) 直接退化为 \(\int g(x)\varphi(x)\,\mathrm{d}x\)）。
  \item
    \textbf{约定 B（Fourier 变换指数带 \(2\pi\)，p.312 式 8.9）}： 
\FitDisplayMath{\hat{f}(\xi) \triangleq \int_{\mathbb{R}^n} f(x)\mathrm{e}^{-2\pi\mathrm{i}\xi\cdot x}\,\mathrm{d}x, \qquad f(x) = \int_{\mathbb{R}^n} \hat{f}(\xi)\mathrm{e}^{2\pi\mathrm{i}\xi\cdot x}\,\mathrm{d}\xi.}
 在此约定下：

    \begin{itemize}
    \tightlist
    \item
      \textbf{求导 \(\leftrightarrow\) 乘多项式}：\(\widehat{\partial_j f}(\xi) = 2\pi\mathrm{i}\xi_j \hat{f}(\xi)\)，\(\partial_{\xi_j}\hat{f}(\xi) = -2\pi\mathrm{i}\widehat{x_j f}(\xi)\)；
    \item
      \textbf{卷积 \(\leftrightarrow\) 乘积}：\(\widehat{f * g}(\xi) = \hat{f}(\xi)\hat{g}(\xi)\)；
    \item
      \textbf{Plancherel 等距（无额外常数！）}：\(\int_{\mathbb{R}^n}\lvert f(x)\rvert^2\,\mathrm{d}x = \int_{\mathbb{R}^n}\lvert\hat{f}(\xi)\rvert^2\,\mathrm{d}\xi\)。
    \end{itemize}
  \end{itemize}
\end{enumerate}

\BookSubSection{}{5.2 广义函数求导、极限与乘法三大必会模块}{5.2 广义函数求导、极限与乘法三大必会模块}{}

\begin{enumerate}
\def\labelenumi{\arabic{enumi}.}
\tightlist
\item
  \textbf{广义导数定义}：\(\partial^\alpha \lambda(\varphi) \triangleq (-1)^{\lvert\alpha\rvert}\lambda(\partial^\alpha \varphi)\)（口诀：``求导甩给试验函数，甩一次加一个负号''）。

  \begin{itemize}
  \tightlist
  \item
    \textbf{四大必背 1 维经典广义导数（\hyperlink{prob:8-11}{\textcolor{LinkColor}{习题 8 第 11 题}}）}：

    \begin{itemize}
    \tightlist
    \item
      \((\operatorname{sgn} x)' = 2\delta\)；\((H(x))' = \delta\)（Heaviside 阶跃函数）；
    \item
      \((\lvert x\rvert)' = \operatorname{sgn} x\)（弱导数存在，因 \(\operatorname{sgn} x \in L^1_{\mathrm{loc}}\)；但 \(\operatorname{sgn} x\) 没有弱导数，因 \(2\delta \notin L^1_{\mathrm{loc}}\)！）；
    \item
      \((\ln\lvert x\rvert)' = \operatorname{p.v.}\frac{1}{x}\)，其中主值 \(\left(\operatorname{p.v.}\frac{1}{x}\right)(\varphi) = \lim_{\varepsilon\to 0^+}\int_{\lvert x\rvert > \varepsilon}\frac{\varphi(x)}{x}\,\mathrm{d}x\)；
    \item
      \((\ln\lvert x\rvert\operatorname{sgn} x)' = \lambda_1\)（教材 p.305 例 8.6：\(\lambda_1(\varphi) = \int_{\lvert x\rvert<1}\frac{\varphi(x)-\varphi(0)}{\lvert x\rvert}\,\mathrm{d}x + \int_{\lvert x\rvert\ge 1}\frac{\varphi(x)}{\lvert x\rvert}\,\mathrm{d}x\)）。
    \end{itemize}
  \end{itemize}
\item
  \textbf{广义函数列收敛（\(\lambda_k \to \lambda\) 于 \(\mathcal{D}' \iff \forall\varphi\in\mathcal{D}, \lambda_k(\varphi)\to\lambda(\varphi)\)）}：

  \begin{itemize}
  \tightlist
  \item
    \textbf{核心性质}：\(L^1_{\mathrm{loc}}\) 收敛 \(\implies \mathcal{D}'\) 收敛；且 \(\mathcal{D}'\) 收敛可与任意阶广义求导\textbf{无条件交换}（\(\lambda_k \to \lambda \implies \partial^\alpha\lambda_k \to \partial^\alpha\lambda\)）！
  \item
    \textbf{必背三大极限（\hyperlink{prob:8-9}{\textcolor{LinkColor}{习题 8 第 9 题}}、\hyperlink{prob:8-14}{\textcolor{LinkColor}{第 14 题}}）}：

    \begin{itemize}
    \tightlist
    \item
      \(\lim_{\varepsilon\to 0^+}\frac{\varepsilon}{\varepsilon^2 + x^2} = \pi\delta\)（\(\varepsilon\to 0^-\) 为 \(-\pi\delta\)）；
    \item
      \(\lim_{\varepsilon\to 0^+}\frac{1}{\sqrt{\pi\varepsilon}}\mathrm{e}^{-x^2/\varepsilon} = \delta\)；
    \item
      \(\lim_{\varepsilon\to 0}\left[(x+\mathrm{i}\varepsilon)^{-1} + (x-\mathrm{i}\varepsilon)^{-1}\right] = \lim_{\varepsilon\to 0}\frac{2x}{x^2+\varepsilon^2} = \lim_{\varepsilon\to 0}\big(\ln(x^2+\varepsilon^2)\big)' = (2\ln\lvert x\rvert)' = 2\operatorname{p.v.}\frac{1}{x}\)！
    \end{itemize}
  \end{itemize}
\item
  \textbf{广义函数乘法红线（\hyperlink{prob:8-7}{\textcolor{LinkColor}{习题 8 第 7 题}}）}：只能用 \(C^\infty\) 光滑函数乘广义函数；\(f \in L^1(\mathbb{R})\) \textbf{不能}与 \(\delta\) 相乘（\(L^1\) 中单点 \(x=0\) 无良定义，且不同光滑 \(L^1\)-逼近列乘 \(\delta\) 极限不唯一，如 \(0\) 与 \(\mathrm{e}^{-k^2 x^2}\) 均 \(L^1\to 0\)，乘 \(\delta\) 分别得 \(0\) 和 \(\delta\)）。
\end{enumerate}

\BookSubSection{}{5.3 Fourier 变换必背对、Sobolev 空间 \(H^s(\mathbb{R}^n)\) 与椭圆估计}{5.3 Fourier 变换必背对、Sobolev 空间 H\textasciicircum{}s(R\textasciicircum{}n) 与椭圆估计}{}

\begin{enumerate}
\def\labelenumi{\arabic{enumi}.}
\tightlist
\item
  \textbf{必背两对 Fourier 变换（本书 \(2\pi\) 约定，\hyperlink{prob:8-17}{\textcolor{LinkColor}{习题 8 第 17(2)(3) 题}}、\hyperlink{prob:8-20}{\textcolor{LinkColor}{第 20 题}}）}：

  \begin{itemize}
  \tightlist
  \item
    \textbf{Lorentz / Poisson 核}：\(\widehat{\frac{1}{a^2 + x^2}}(\xi) = \frac{\pi}{\lvert a\rvert}\mathrm{e}^{-2\pi\lvert a\rvert\lvert\xi\rvert}\)（\(a \neq 0\)）；逆过来 \(\mathrm{e}^{-2\pi y\lvert\xi\rvert}\) 的逆变换就是上半平面 Poisson 核 \(P_y(x) = \frac{1}{\pi}\frac{y}{x^2+y^2}\)，直接解出上半平面 Laplace 方程 \(u(x,y) = (P_y * g)(x) = \frac{y}{\pi}\int_{-\infty}^\infty \frac{g(t)}{(x-t)^2+y^2}\,\mathrm{d}t\)！
  \item
    \textbf{Gauss 核}：\(\widehat{\mathrm{e}^{-ax^2}}(\xi) = \sqrt{\frac{\pi}{a}}\mathrm{e}^{-\frac{\pi^2\xi^2}{a}}\)（\(a > 0\)；特例 \(a=\pi\) 时自反演 \(\widehat{\mathrm{e}^{-\pi x^2}}(\xi)=\mathrm{e}^{-\pi\xi^2}\)）。
  \end{itemize}
\item
  \textbf{整数阶与实指数 Sobolev 空间}：

  \begin{itemize}
  \tightlist
  \item
    \textbf{整数阶 \(W^{m,p}(\Omega)\)（定义 8.18）}：\(\{f \in L^p(\Omega) : \forall\lvert\alpha\rvert\le m,\; \partial^\alpha f \in L^p(\Omega)\}\)，\(H^m(\Omega) \triangleq W^{m,2}(\Omega)\)；
  \item
    \textbf{实指数 \(H^s(\mathbb{R}^n)\)（p.336，对任意 \(s \in \mathbb{R}\) 统一定义！）}：记 \(\langle\xi\rangle \triangleq (1 + \lvert\xi\rvert^2)^{1/2}\)， 
\FitDisplayMath{H^s(\mathbb{R}^n) \triangleq \left\{f \in \mathcal{S}'(\mathbb{R}^n) : \langle\xi\rangle^s \hat{f}(\xi) \in L^2(\mathbb{R}^n)\right\}, \qquad \lVert f\rVert_{H^s} \triangleq \lVert\langle\xi\rangle^s \hat{f}\rVert_{L^2}.}

  \item
    \textbf{负指标判定 \(\lvert x\rvert^{-k} \in H^{-m}(\mathbb{R}^n)\)（\hyperlink{prob:8-28}{\textcolor{LinkColor}{习题 8 第 28 题}}）}：由 \(\widehat{\lvert x\rvert^{-k}}(\xi) = c_{n,k}\lvert\xi\rvert^{k-n}\)（\(1\le k<n\)），检查球坐标积分 \(\int_0^\infty (1+r^2)^{-m} r^{2k-n-1}\,\mathrm{d}r < \infty\) 的\textbf{两端}：

    \begin{itemize}
    \tightlist
    \item
      低频 \(r \to 0^+\) 要求 \(k > \frac{n}{2}\)（故 \(k \le \frac{n}{2}\) 时\textbf{无解}！）；
    \item
      高频 \(r \to +\infty\) 要求 \(m > k - \frac{n}{2}\)（故 \(\frac{n}{2} < k < n\) 时最小自然数 \(m = \lfloor k - \frac{n}{2}\rfloor + 1\)！）。
    \end{itemize}
  \item
    \textbf{椭圆算子先验估计（p.337 \hyperlink{prob:8-882}{\textcolor{LinkColor}{练习 8.8.2}}）}：若 \(m\) 阶常系数齐次算子 \(L = \sum_{\lvert\alpha\rvert=m} c_\alpha\partial^\alpha\) 象征满足 \(\lvert l(\xi)\rvert \ge c\lvert\xi\rvert^m\)，由 \(\widehat{Lf} = (2\pi\mathrm{i})^m l(\xi)\hat{f}\) 与 \(\langle\xi\rangle^m \le 2^{m/2}(1+\lvert\xi\rvert^m)\)，两边同乘 \(\langle\xi\rangle^s\lvert\hat{f}\rvert\) 取 \(L^2\) 范数立得 \(\lVert f\rVert_{H^{s+m}} \le C(\lVert f\rVert_{H^s} + \lVert Lf\rVert_{H^s})\)！
  \end{itemize}
\end{enumerate}

\BookPart{part:12types}{闭卷考场 12 类高频变形题通法模板}{闭卷考场 12 类高频变形题通法模板}

\begin{quote}
老师亲口说过：\textbf{不是考作业原题，但考作业背后体现的框架与思路！} 考场拿到一道陌生证明/计算题时，直接对照下表匹配思路：
\end{quote}

{\def\LTcaptype{none} % do not increment counter
{\small\renewcommand{\arraystretch}{1.38}\setlength{\tabcolsep}{5pt}
\begin{longtable}[]{@{}
  >{\centering\arraybackslash}p{(\linewidth - 6\tabcolsep) * \real{0.1270}}
  >{\raggedright\arraybackslash}p{(\linewidth - 6\tabcolsep) * \real{0.2381}}
  >{\raggedright\arraybackslash}p{(\linewidth - 6\tabcolsep) * \real{0.4286}}
  >{\raggedright\arraybackslash}p{(\linewidth - 6\tabcolsep) * \real{0.2063}}@{}}
\toprule\noalign{}
\rowcolor{AccentSoft}
\begin{minipage}[b]{\linewidth}\centering
题型编号
\end{minipage} & \begin{minipage}[b]{\linewidth}\raggedright
常见关键词 / 问法
\end{minipage} & \begin{minipage}[b]{\linewidth}\raggedright
第一反应与三步通法
\end{minipage} & \begin{minipage}[b]{\linewidth}\raggedright
对应精讲/作业
\end{minipage} \\
\midrule\noalign{}
\endhead
\bottomrule\noalign{}
\endlastfoot
\textbf{题型 1} & 证 \(C([a,b])\) 上某积分算子是紧算子 & \textbf{Arzelà-Ascoli 定理}：\CircNum{1} 验证线性有界（映入 \(C([a,b])\)）；\CircNum{2} 对有界集 \(\lVert f\rVert_\infty\le M\) 证一致有界；\CircNum{3} 估计 \(\lvert Af(x_1)-Af(x_2)\rvert\) 证等度连续（强调 \(\delta\) 与 \(f\) 无关）。 & \hyperlink{prob:5-14}{\textcolor{LinkColor}{习题 5 第 14 题}} \\
\textbf{题型 2} & 证 \(\ell^p / L^2\) 上对角算子或积分算子是紧算子 & 两条路二选一：\CircNum{1} 构造\textbf{有限秩截断算子 \(A_N\)}，算 \(\lVert A - A_N\rVert \to 0\) 套定理 5.4；\CircNum{2} 在 \(L^2\)/可分 Hilbert 上直接算 \textbf{Hilbert-Schmidt 范数} \(\lVert A\rVert_2 < \infty\) 套定理 5.12！ & 定理 5.4 / 5.12 \\
\textbf{题型 3} & 证无穷维上某紧算子无有界逆 / 某算子不紧 / 值域不含无穷维闭子空间 & \textbf{Riesz 引理反证大法}：凑出无穷维空间上的恒等算子 \(I\)（或其某闭子空间 \(M\) 上的 \(I_M\)）是紧算子 \(\implies\) 闭单位球列紧 \(\implies\) 矛盾于无穷维！（值域子空间题先收缩到 \(X_0 = A^{-1}(M)\) 用开映射定理！） & \hyperlink{prob:5-5}{\textcolor{LinkColor}{习题 5 第 5 题}}、\hyperlink{prob:5-6}{\textcolor{LinkColor}{第 6 题}}（及 17 题） \\
\textbf{题型 4} & 证算子是 Fredholm 算子 / 证指标相等 & 优先 \textbf{Atkinson 定理（定理 5.13）}凑左右拟逆 \(SA=I-K_1, AS=I-K_2\)；证指标相等用\textbf{公共拟逆 + 乘积指标公式（定理 5.14）+ \(\operatorname{ind}(I-K)=0\)}；若按定义硬证，\textbf{千万别漏证值域闭}！ & \hyperlink{prob:5-20}{\textcolor{LinkColor}{习题 5 第 20 题}}、\hyperlink{prob:5-22}{\textcolor{LinkColor}{第 22 题}} \\
\textbf{题型 5} & 求乘法算子或对角算子的谱集 \(\sigma(A)\) & \textbf{双向夹逼}：\CircNum{1} \(\lambda \notin \text{目标集}\) 时直接写出显式逆算子并验证有界（\(\lambda\in\rho(A)\)）；\CircNum{2} \(\lambda \in \text{目标集}\) 时，有特征向量直接标点谱，没有就构造单位向量列使 \(\lVert(\lambda I-A)f_n\rVert\to 0\)（Weyl 序列）证无有界逆！ & \hyperlink{prob:6-2}{\textcolor{LinkColor}{习题 6 第 2 题}} \\
\textbf{题型 6} & 证明预解集开性、预解式连续性或预解式恒等式 & 遇扰动 \(A_n\to A\) 或 \(\lambda\to\lambda_0\)：\textbf{强行提取 \((\lambda_0 I-A)\)} 化成 \((\lambda_0 I-A)[I - B]\)，对 \(\lVert B\rVert < 1\) 上 \textbf{Neumann 级数}；遇 \(B=R(\alpha,A)\) 先把 \(\mu I - B\) 因式分解成 \(\mu(\alpha I-A)^{-1}(\lambda I - A)\)！ & 定理 6.2、\hyperlink{prob:6-5}{\textcolor{LinkColor}{习题 6 第 5 题}}、\hyperlink{prob:6-13}{\textcolor{LinkColor}{第 13 题}} \\
\textbf{题型 7} & 证明某算子是酉算子（特别 Cayley 变换类） & \CircNum{1} 先证分母 \((A+\bar{\lambda}I)\) 可逆（自伴谱在实轴或用下界 \(\lVert(A+\bar{\lambda}I)x\rVert\ge\lvert\operatorname{Im}\lambda\rvert\lVert x\rVert\)）；\CircNum{2} 利用 \(A+\lambda I\) 与 \((A+\bar{\lambda}I)^{-1}\) 可交换，直接算 \(U^*U = UU^* = I\)！ & \hyperlink{prob:6-23}{\textcolor{LinkColor}{习题 6 第 23 题}} \\
\textbf{题型 8} & 求非自伴紧积分算子（如 Volterra）的算子范数 \(\lVert A\rVert\) & \textbf{绝不直接求 \(\Gamma(A)\)}！固定三步：\CircNum{1} Fubini 求 \(A^*\)；\CircNum{2} 算正自伴紧算子 \(AA^*\)（或 \(A^*A\)）的特征方程，逐次求导化二阶 ODE 边值问题求最大正特征值 \(\lambda_1 = \Gamma(AA^*)\)；\CircNum{3} \(\lVert A\rVert = \sqrt{\lambda_1}\)！ & \hyperlink{prob:6-24}{\textcolor{LinkColor}{习题 6 第 24 题}} \\
\textbf{题型 9} & 涉及 \((A^*A)^{1/2}\)、HS 算子奇异值分解或 \(\lVert A\rVert_2\) & 一切围绕恒等式 \textbf{\(\lVert Ax\rVert^2 = (A^*Ax, x) = \lVert(A^*A)^{1/2}x\rVert^2\)}，对紧正自伴算子 \(T=(A^*A)^{1/2}\) 用 Hilbert-Schmidt 谱展开取特征基 \(\{x_n\}\)，令 \(y_n = Ax_n/\lambda_n\)！ & \hyperlink{prob:6-31}{\textcolor{LinkColor}{习题 6 第 31 题}}、\hyperlink{prob:6-37}{\textcolor{LinkColor}{第 37 题}} \\
\textbf{题型 10} & 涉及 Borel 函数演算 \(f(A)\)、酉等价或解析级数展开 & 酉等价题证 \(\tilde{E}(\Delta)=U^{-1}E(\Delta)U\) 是谱测度 + \textbf{唯一性}；解析函数题利用紧集 \(\sigma(A)\) 上一致收敛 + 谱积分范数界 \(\lVert\int(f-S_N)\,\mathrm{d}E\rVert \le \sup_{\sigma(A)}\lvert f-S_N\rvert \to 0\)！ & \hyperlink{prob:6-39}{\textcolor{LinkColor}{习题 6 第 39 题}}、\hyperlink{prob:6-45}{\textcolor{LinkColor}{第 45 题}} \\
\textbf{题型 11} & 证明无界算子闭/无界/对称/自伴/值域闭 & 无界用滑动截断函数 \(\chi_{[n,n+1]}\) 或高频 \(\sin(nx)\)；闭性用序列定义（\(L^2\) 上\textbf{嵌套抽两次几乎处处收敛子列}）；对称/非负算子闭性必先展开\textbf{能量不等式} \(\lVert(A+\mathrm{i}I)x\rVert^2 = \lVert Ax\rVert^2 + \lVert x\rVert^2\) 或 \(\lVert(A+I)x\rVert^2 \ge \lVert Ax\rVert^2 + \lVert x\rVert^2\)！ & \hyperlink{prob:7-4}{\textcolor{LinkColor}{习题 7 第 4 题}}、\hyperlink{prob:7-6}{\textcolor{LinkColor}{第 6 题}}、\hyperlink{prob:7-7}{\textcolor{LinkColor}{第 7 题}} \\
\textbf{题型 12} & 广义导数、分布极限、Fourier 变换与 \(H^s\) 估计 & 求导``甩给 \(\varphi\) 加负号 + 挖奇点分部积分''；极限优先``换元+控制收敛''或``原函数 \(L^1_{\mathrm{loc}}\) 收敛再求导''；Fourier 变换牢记\textbf{本书 \(2\pi\) 约定}；\(H^s\) 题立刻切频域检查 \(\langle\xi\rangle^s\hat{f}\in L^2\)（\textbf{低频和高频两端都要查}！）。 & \hyperlink{ch:8}{\textcolor{LinkColor}{直达第 8 章全 8 题 $\to$}} \\
\end{longtable}
}
}

\BookPart{part:25q}{考前 25 问闭卷自测默写清单（全答对 = 稳过！）}{考前 25 问闭卷自测默写清单（全答对 = 稳过！）}

把答案盖住，逐条在草稿纸上默写或口答，打勾通过为止：

\BookSubSection{}{第 5 章（6 问）}{第 5 章（6 问）}{\TagMust{闭卷必背}}

\begin{itemize}
\tightlist
\item[\textcolor{AccentMain}{$\square$}]
  \textbf{Q1}：写出紧算子、有限秩算子、Hilbert-Schmidt 算子、Fredholm 算子及其指标 \(\operatorname{ind}A\) 的严格定义。 \JumpPill{sec:5-1}{第5章精讲 $\to$}
\item[\textcolor{AccentMain}{$\square$}]
  \textbf{Q2}：写出无穷维可分 Hilbert 空间上四类算子的包含链，并各举一个反例证明「有限秩 \(\subsetneq\) HS \(\subsetneq\) 紧 \(\subsetneq\) 有界」。 \JumpPill{sec:5-2}{包含链与反例 $\to$}
\item[\textcolor{AccentMain}{$\square$}]
  \textbf{Q3}：证明 Hilbert-Schmidt 范数 \(\lVert A\rVert_2\) 与标准正交基选取无关，且 \(\lVert A\rVert \le \lVert A\rVert_2\)。 \JumpPill{sec:5-2}{5.2 HS算子 $\to$}
\item[\textcolor{AccentMain}{$\square$}]
  \textbf{Q4}：证明无穷维 Banach 空间上的紧算子没有有界逆（定理 5.5 / \hyperlink{prob:5-5}{\textcolor{LinkColor}{习题 5 第 5 题}}）。 \JumpPill{prob:5-5}{习题5第5题 $\to$}
\item[\textcolor{AccentMain}{$\square$}]
  \textbf{Q5}：紧算子与 Fredholm 算子是什么关系？写出 Atkinson 定理（定理 5.13）和乘积指标公式（定理 5.14）。 \JumpPill{sec:5-3}{5.3 Fredholm $\to$}
\item[\textcolor{AccentMain}{$\square$}]
  \textbf{Q6}：如何用 Arzelà-Ascoli 定理证明 \(C([0,1])\) 上的积分算子是紧算子？需验证哪两个条件？ \JumpPill{prob:5-14}{习题5第14题 $\to$}
\end{itemize}

\BookSubSection{}{第 6 章（10 问——绝对大头！）}{第 6 章（10 问——绝对大头！）}{\TagMust{闭卷必背}}

\begin{itemize}
\tightlist
\item[\textcolor{AccentMain}{$\square$}]
  \textbf{Q7}：默写预解集 \(\rho(A)\)、预解式 \(R(\lambda, A)\) 以及点谱 \(\sigma_p\)、连续谱 \(\sigma_c\)、剩余谱 \(\sigma_r\) 的精确定义。 \JumpPill{sec:6-1}{6.1 谱定义 $\to$}
\item[\textcolor{AccentMain}{$\square$}]
  \textbf{Q8}：举出至少两个例子：一个算子只有连续谱没有特征值；一个算子 \(\Gamma(A) = 0\) 但 \(\lVert A\rVert > 0\)。 \JumpPill{sec:6-1}{6.1 四大算子 $\to$}
\item[\textcolor{AccentMain}{$\square$}]
  \textbf{Q9}：默写 \textbf{Gelfand 谱半径公式}（定理 6.5），并完整默写教材 p.214 的两步证明（Neumann 级数 + Laurent 展开 + 共鸣定理）。 \JumpPill{sec:6-1}{6.1 Gelfand $\to$}
\item[\textcolor{AccentMain}{$\square$}]
  \textbf{Q10}：默写紧算子 \textbf{Riesz-Schauder 定理}的 5 条核心结论，并能用 Riesz 引理证明``\(\lambda\neq 0\) 时特征空间有限维''。 \JumpPill{sec:6-2}{6.2 紧算子谱 $\to$}
\item[\textcolor{AccentMain}{$\square$}]
  \textbf{Q11}：证明有界自伴算子的特征值全实、异特征向量正交、\(\sigma(A) \subset \mathbb{R}\)、\(\sigma_r(A) = \emptyset\)。 \JumpPill{sec:6-3}{6.3 自伴算子 $\to$}
\item[\textcolor{AccentMain}{$\square$}]
  \textbf{Q12}：自伴算子的谱半径公式与一般 Gelfand 公式有什么区别？用 \(\lVert A^2\rVert = \lVert A\rVert^2\) 写出 \(\Gamma(A) = \lVert A\rVert\) 的 4 行证明！ \JumpPill{sec:6-3}{6.3 谱半径 $\to$}
\item[\textcolor{AccentMain}{$\square$}]
  \textbf{Q13}：默写 \textbf{Hilbert-Schmidt 谱展开定理}（定理 6.13）的结论，并说出第 6 章``先收窄、再扩宽''的完整逻辑链。 \JumpPill{sec:6-3}{6.3 谱展开 $\to$}
\item[\textcolor{AccentMain}{$\square$}]
  \textbf{Q14}：写出谱积分 \(\int f\,\mathrm{d}E\) 的 4 条基本运算法则（线性、乘法、共轭、范数界）。 \JumpPill{sec:6-4}{6.4 谱积分 $\to$}
\item[\textcolor{AccentMain}{$\square$}]
  \textbf{Q15}：默写\textbf{酉算子}的定义（保距且满射）及 3 条等价刻画，证明 \(\sigma(U) \subset \{\lvert\lambda\rvert = 1\}\)，并写出酉算子谱分解公式。 \JumpPill{sec:6-5}{6.5 酉算子 $\to$}
\item[\textcolor{AccentMain}{$\square$}]
  \textbf{Q16}：写出有界自伴算子的 \textbf{Cayley 变换} \(U = (A-\mathrm{i}I)(A+\mathrm{i}I)^{-1}\)，证明 \(U\) 是酉算子且 \(1\in\rho(U)\)，并写出有界自伴算子的谱分解与 Borel 函数演算。 \JumpPill{prob:6-23}{Cayley/第23题 $\to$}
\end{itemize}

\BookSubSection{}{第 7 章（6 问）}{第 7 章（6 问）}{\TagMust{闭卷必背}}

\begin{itemize}
\tightlist
\item[\textcolor{AccentMain}{$\square$}]
  \textbf{Q17}：举出一个无界算子的例子并证明它无界；为什么无界对称算子不可能定义在整个 Hilbert 空间 \(H\) 上？ \JumpPill{prob:7-4}{第7章/第4题 $\to$}
\item[\textcolor{AccentMain}{$\square$}]
  \textbf{Q18}：默写稠定算子共轭算子 \(A^*\) 的定义域 \(\mathcal{D}(A^*)\) 与作用法则；为什么必须要求稠定？为什么 \(A \subset B \Rightarrow B^* \subset A^*\)？ \JumpPill{sec:7-1-1}{7.1.1 共轭 $\to$}
\item[\textcolor{AccentMain}{$\square$}]
  \textbf{Q19}：\textbf{对称算子与自伴算子的本质区别是什么？} 能讲出教材例 7.2 三个求导算子 \(A_1, A_2, A_3\) 各自的边界条件与对称/自伴结论吗？ \JumpPill{sec:7-1-2}{7.1.2 例7.2 $\to$}
\item[\textcolor{AccentMain}{$\square$}]
  \textbf{Q20}：默写闭算子定义，并 4 行证明 \textbf{定理 7.1（稠定算子的共轭 \(A^*\) 必闭）}；写出对称算子的四级阶梯 \(A \subset \overline{A} = A^{**} \subset A^*\)。 \JumpPill{sec:7-1-4}{7.1.4 闭算子 $\to$}
\item[\textcolor{AccentMain}{$\square$}]
  \textbf{Q21}：证明对称算子恒等式 \(\lVert(A\pm\mathrm{i}I)x\rVert^2 = \lVert Ax\rVert^2 + \lVert x\rVert^2\)，并默写\textbf{定理 7.4（对称算子成为自伴算子的充要条件）}。 \JumpPill{sec:7-1-5}{7.1.5 判别准则 $\to$}
\item[\textcolor{AccentMain}{$\square$}]
  \textbf{Q22}：对比\textbf{有界自伴算子}与\textbf{无界自伴算子}的谱集位置、Cayley 变换（点 \(1\) 的归属）与谱分解表达式的异同！ \JumpPill{sec:7-2}{7.2 无界谱分解 $\to$}
\end{itemize}

\BookSubSection{}{第 8 章补充自测（3 问）}{第 8 章补充自测（3 问）}{\TagHigh{作业核心}}

\begin{itemize}
\tightlist
\item[\textcolor{AccentMain}{$\square$}]
  \textbf{Q23}：默写广义导数定义与 4 个 1 维经典广义导数：\((\operatorname{sgn} x)'\)、\((\lvert x\rvert)'\)、\((\ln\lvert x\rvert)'\)、\((\ln\lvert x\rvert\operatorname{sgn} x)'\)。 \JumpPill{prob:8-11}{习题8第11题 $\to$}
\item[\textcolor{AccentMain}{$\square$}]
  \textbf{Q24}：默写本书 \textbf{\(2\pi\) 约定}下的 Fourier 变换、求导法则，以及 \(\frac{1}{a^2+x^2}\) 和 \(\mathrm{e}^{-ax^2}\) 的变换结果。 \JumpPill{prob:8-17}{8.4/第17题 $\to$}
\item[\textcolor{AccentMain}{$\square$}]
  \textbf{Q25}：默写实指数 Sobolev 空间 \(H^s(\mathbb{R}^n)\) 的频域定义，并能口述判断 \(\frac{1}{\lvert x\rvert^k} \in H^{-m}(\mathbb{R}^n)\)（查两端！）和证椭圆先验估计（\hyperlink{prob:8-882}{\textcolor{LinkColor}{练习 8.8.2}}）的核心步骤。 \JumpPill{prob:8-28}{8.8/第28题 $\to$}
\end{itemize}

\BookPart{part:audio}{附：末课录音完整逐字整理（全长 13:49，带精确时间戳）}{附：末课录音完整逐字整理（全长 13:49，带精确时间戳）}

\begin{quote}
说明：为避免零基础同学被口语谐音误导，下文中已将专业名词还原为准确教材术语（例如录音 \texttt{05:35} 处口语读作``有算子、保距并且满射''即教材 6.5 节的\textbf{酉算子}；``里斯-绍德尔''即 \textbf{Riesz-Schauder}；``盖尔范德''即 \textbf{Gelfand}；``弗雷德霍姆''即 \textbf{Fredholm}）。现场环境音遮挡导致的极个别模糊字句已如实括号标注，绝不编造。
\end{quote}

\BookSubSection{}{\texttt{{[}00:00–00:46{]}} 第五章总结 \& 全书复习方法论}{[00:00–00:46] 第五章总结 \& 全书复习方法论}{}

\begin{quote}
（录音开头紧接第五章的紧算子、）有限秩算子、Hilbert-Schmidt（希尔伯特-施密特）算子、Fredholm（弗雷德霍姆）算子，这四类算子我们都学过。那么这些算子的\textbf{定义}、\textbf{它们之间的关系}、\textbf{它们的基本性质}（是复习的核心）。

大家在复习的时候要有一个\textbf{框架}。因为内容太多的话，如果一开始陷入某一个细节，可能就把握不住全局，学完之后好像什么也没有得到。所以这也就是数学理论学习的一般顺序：

\begin{enumerate}
\def\labelenumi{\arabic{enumi}.}
\tightlist
\item
  先掌握它的\textbf{定义}；
\item
  再掌握它的\textbf{基本性质}；
\item
  再弄清这几类算子\textbf{之间的关系}。
\end{enumerate}

这些最基本的一定要知道，然后再去填里面的细节。第五章讲的就是这些了。
\end{quote}

\BookSubSection{}{\texttt{{[}00:46–02:11{]}} 第六章（上）：一般有界线性算子的谱理论与 Gelfand 谱半径公式}{[00:46–02:11] 第六章（上）：一般有界线性算子的谱理论与 Gelfand 谱半径公式}{}

\begin{quote}
第六章讲的是\textbf{算子的谱理论}。算子的谱理论分成几块：

最基本的是\textbf{有界线性算子的谱理论}：把复平面上的点分成\textbf{正则点}（预解集中的点）以及\textbf{谱点}；而谱又分成\textbf{三种类型}（点谱、连续谱、剩余谱）。

有了基本性质之后，对于有界线性算子，我们也定义了它的\textbf{谱半径}。谱半径公式是由 \textbf{Gelfand（盖尔范德）定理} 给出来的——这个当时都是\textbf{严格证明}了的，我估计你们现在都已经交还给我了。对于什么样的算子有什么样的谱、谱半径公式，这些都是讲过的。

讲的有界线性算子的谱实际上是一个一般的情形。因为我们这学期整个课本里前面大部分讲的都是有界线性算子（两个线性赋范空间之间的有界线性算子）的性质。
\end{quote}

\BookSubSection{}{\texttt{{[}02:11–03:19{]}} 第六章（中）：收窄到紧算子的谱理论（Riesz-Schauder 定理）}{[02:11–03:19] 第六章（中）：收窄到紧算子的谱理论（Riesz-Schauder 定理）}{}

\begin{quote}
有了有界算子的基本性质之后，我们自然会考虑上一章专门讲过的\textbf{紧算子}。主要原因在于，我们实际生活中碰到的很多算子都是紧算子，这也是第五章为什么把紧算子单独拿出来讲的原因。我们最终处理的很多问题就是 \(Ax = y\) 求解 \(x\) 的问题（以及对应的特征值/谱问题）。

所以紧算子的谱理论，包括：

\begin{itemize}
\tightlist
\item
  紧算子的谱有哪些类；
\item
  紧算子到底有什么样的谱点、这些谱点之间有什么关系——也就是紧算子的 \textbf{Riesz-Schauder（里斯-绍德尔）定理}。
\end{itemize}

这就是紧算子谱理论的主要内容。
\end{quote}

\BookSubSection{}{\texttt{{[}03:19–05:19{]}} 第六章（中）：继续收窄到 Hilbert 空间上的自伴紧算子（Hilbert-Schmidt 定理）}{[03:19–05:19] 第六章（中）：继续收窄到 Hilbert 空间上的自伴紧算子（Hilbert-Schmidt 定理）}{}

\begin{quote}
只限定到紧算子还不够，为什么呢？因为我们在一般 Banach 空间里，空间本身没有比较好的几何结构。在这个不太好的空间里去研究好的算子，性质上还是差了一点。好在我们实际问题中考虑的很多都是在 \textbf{Hilbert（希尔伯特）空间}里。（\emph{老师插叙：十年前泛函分析是数学专业和大面积工科一起上的，工科同学反映 Hilbert 空间以及其中的紧算子理论在科研中非常核心。}）

所以把紧算子放到 \textbf{Hilbert 空间} 上，而且再要求它是\textbf{自伴的紧算子}，它的谱理论、谱的性质就更好了！因为 Hilbert 空间里有内积几何结构：

\begin{itemize}
\tightlist
\item
  它对应的\textbf{谱半径公式}跟前面一般 Gelfand 定理给出的公式是有区别的；
\item
  谱半径公式还告诉我们，谱的范围是在 \textbf{\([-\lVert A\rVert, \lVert A\rVert]\)} 之间（且落在实轴上）；
\item
  有了这些之后，就明确了谱展开/谱积分的表达形式在什么范围，对于 Hilbert 空间上的自伴紧算子，核心定理就是 \textbf{Hilbert-Schmidt（希尔伯特-施密特）定理}。
\end{itemize}
\end{quote}

\BookSubSection{}{\texttt{{[}05:19–06:06{]}} 第六章（下）：谱系/谱测度/谱积分 \& 酉算子谱分解}{[05:19–06:06] 第六章（下）：谱系/谱测度/谱积分 \& 酉算子谱分解}{}

\begin{quote}
后面的\textbf{谱系、谱测度、谱积分}，大家需要\textbf{大概了解一下}（当时也给大家演示过具体怎么用）。

对紧算子研究完之后，接着研究\textbf{酉算子}（\emph{录音口语读音近``有算子''}），酉算子的性质非常特殊：

\begin{itemize}
\tightlist
\item
  首先对于一个酉算子，它的\textbf{定义}——\textbf{保距并且满射}，以及上面给出的几种等价定义方式；
\item
  其次是\textbf{酉算子的谱分解}、\textbf{酉算子的谱分解定理}。
\end{itemize}
\end{quote}

\BookSubSection{}{\texttt{{[}06:06–08:21{]}} 第六章全局脉络：先缩小、再延拓 \(\to\) 一般有界自伴算子的谱分解}{[06:06–08:21] 第六章全局脉络：先缩小、再延拓 → 一般有界自伴算子的谱分解}{}

\begin{quote}
大家看我们第六章整个过程：

\begin{enumerate}
\def\labelenumi{\arabic{enumi}.}
\tightlist
\item
  先从一般的\textbf{有界线性算子}，缩小到\textbf{紧算子}；
\item
  再缩小到 \textbf{Hilbert 空间上的自伴紧算子}，在这个最窄的地方做了比较深入的研究（引入了谱系、谱测度和谱积分）；
\item
  有了这组非常好的结论之后，自然的想法就是把它\textbf{延拓}到更一般的算子上去——先开延到\textbf{酉算子的谱分解}，再延拓到一般的\textbf{有界自伴算子的谱分解}！
\end{enumerate}

所以有界自伴算子也有谱积分表达式，但注意它与前面和后面的区别：

\begin{itemize}
\tightlist
\item
  \textbf{自伴紧算子}：谱局限在实轴 \([-\lVert A\rVert, \lVert A\rVert]\) 上且是离散可数特征值（仅 \(0\) 可为聚点）；
\item
  \textbf{一般有界自伴算子}：谱仍在实轴有界闭区间 \([-\lVert A\rVert, \lVert A\rVert]\) 上，但不限于离散特征值；
\item
  \textbf{无界自伴算子}（过渡到第七章）：谱同样在实轴 \(\mathbb{R}\) 上，但不再局限在有限区间 \([-\lVert A\rVert, \lVert A\rVert]\) 里，而是分布在整条实轴 \(\mathbb{R}\) 上。
\end{itemize}

最后面那一小节（\emph{\texttt{{[}07:49{]}} 此处录音略模糊，结合教材位置对应 6.6 节末尾生成元/循环元酉等价于乘法算子等偏应用拓展内容}），相对来说比较偏应用，时间关系没有展开来讲，有兴趣的同学可以看一下，没兴趣的同学也可以不用管。这本书几乎每一节、每一章展开来都能写成一本很深的专著。
\end{quote}

\BookSubSection{}{\texttt{{[}08:22–10:30{]}} 第七章：Hilbert 空间上的无界算子 \& 课程定位}{[08:22–10:30] 第七章：Hilbert 空间上的无界算子 \& 课程定位}{}

\begin{quote}
第七章就是 \textbf{Hilbert 空间上的无界算子}。为什么要研究无界算子？主要因为很多实际中的算子都是无界的，最基本的例子就是\textbf{求导算子}。前面讲的很多理论都建立在有界线性算子上，因此这一章专门研究无界算子，包括无界算子的\textbf{谱积分}怎么做、\textbf{谱分解定理}怎么建立。

研究无界算子肯定也要基于比较特殊的空间（Hilbert 空间）和特殊的算子类：

\begin{itemize}
\tightlist
\item
  通过\textbf{共轭}、\textbf{对称}、\textbf{自伴}对无界算子做限定；
\item
  限定之后再去研究\textbf{无界自伴算子的谱分解}。
\end{itemize}

结合第六章和第七章来看，我们讲的主要是几类最重要的算子性质和谱分解。我们学院目前没有专门做纯泛函分析理论研究的老师，大家主要是把泛函分析作为科研的基本工具和手段来用，重点关注怎么用；但作为一门正式课程，它的基本原理必须清楚。
\end{quote}

\BookSubSection{}{\texttt{{[}10:30–13:49{]}} 复习方法、作业辟谣、成绩比例与闭卷确认}{[10:30–13:49] 复习方法、作业辟谣、成绩比例与闭卷确认}{}

\begin{quote}
所以大家每次复习时，先看\textbf{基本结构框架}，抓住三件事：\textbf{定义、基本性质、它们之间的关联}。如果已经在上讨论班或做科研，看到这些符号记号可以想想怎么联系起来用。现在虽然 AI 很发达，但作为研究生的基本考核任务还是必须完成的。

时间还是比较紧张，需要花时间复习：

\begin{itemize}
\tightlist
\item
  \textbf{减负原则}：如果记了笔记，\textbf{上课讲得比较简单的内容就可以不用去复习}。
\item
  \textbf{纠正误导（关键！）}：上次有一句误导的话我要纠正一下——\textbf{作业在复习的时候可以参考，但是并不是说考试题就是作业原题}，我没有这么说过啊。也就是说，如果作业是自己做的，或者做了一部分，或者在``某些老师''的帮助之下（比如豆包、DeepSeek、ChatGPT 这些``老师''）把作业做出来了，回去一定要看一看，\textbf{要清楚它到底是怎么去做的，这些思路至少要知道}。因为每次布置作业，也是根据我们讲过的框架去布置的。
\item
  \textbf{成绩与考试形式}：\textbf{平时成绩占 50\%}，占比已经比较高了；但平时哪怕拿 100 分，期末如果不好好考还是有点困难，期末考试还是要认真对待。这门课内容相对抽象，听课确实辛苦，研究生课当堂听懂 30\% 就不错了，课下必须花时间消化。
\item
  （学生问：``老师，考试开卷吗？''）老师答：\textbf{闭卷！}
\end{itemize}
\end{quote}

\SoftDivider


% =========================================================================
% 第 5 章 · 紧算子和 Fredholm 算子
% =========================================================================
\SwitchChapterTheme{5}
\ChapterBanner{05}{ch:5}{第 5 章 · 紧算子和 Fredholm 算子}{第 1–4 章前置速通 · 四类算子定义/性质/关系 · 零化子空间与 Kato 闭值域 · 习题 5（6 题）}
\begin{quote}
\textbf{适用对象}：完全没学过泛函分析、需要快速打通核心逻辑并稳拿期末与作业分的本科生\\
\textbf{对照教材}：《泛函分析》（科学出版社``十四五''普通高等教育本科规划教材，第 5 章，教材第 182--205 页）\\
\textbf{覆盖作业}：习题 5 第 \textbf{5、6、14、15、20、22} 题（共 6 题，附教材笔误勘误与严密排雷）
\end{quote}

\BookPart{}{零、零基础前置急救包（第 1--4 章核心概念 5 分钟速通）}{零、零基础前置急救包（第 1–4 章核心概念 5 分钟速通）}

\textbf{什么是泛函分析？} 一句话总结：\textbf{把线性代数推广到无穷维，再加上数学分析里的``极限与收敛''。}\\
在线性代数里，向量是有限维数组 \((x_1, \dots, x_n)\)，线性变换是 \(n \times n\) 矩阵；在泛函分析里，``向量''可以是一个无穷数列或一个连续函数，``矩阵''就升级成了作用在函数/数列空间上的\textbf{线性算子}。

下面按``第 5 章真的会用到什么''提炼 4 组必背前置知识，每条都标明了它在第 5 章的直接用处。

\BookSubSection{sec:5-01}{0.1 度量空间、完备性、列紧性与 Arzelà-Ascoli 定理}{0.1 度量空间、完备性、列紧性与 Arzelà-Ascoli 定理}{\TagLight{前置速通}}

\begin{enumerate}
\def\labelenumi{\arabic{enumi}.}
\tightlist
\item
  \textbf{度量空间 \((X, d)\)}：就是定义了``两点之间距离 \(d(x, y)\)''的集合，只要距离满足非负性、对称性和三角不等式，就能讨论序列收敛 \(x_n \to x\)（即 \(d(x_n, x) \to 0\)）。
\item
  \textbf{完备性（Cauchy 列必收敛）}：如果空间里任何``越挤越近''的 Cauchy 序列（\(\lim_{m,n\to\infty} d(x_m, x_n) = 0\)）都能在空间内部找到极限点（空间没有``针孔漏洞''），就称这个空间是\textbf{完备的}。

  \begin{itemize}
  \tightlist
  \item
    反例：有理数 \(\mathbb{Q}\) 不完备（\(1, 1.4, 1.41, \dots\) 趋向于 \(\sqrt{2} \notin \mathbb{Q}\)）；实数 \(\mathbb{R}\) 完备。
  \end{itemize}
\item
  \textbf{列紧集（相对紧集）}：子集 \(E \subset X\) 称为\textbf{列紧的}，如果 \(E\) 中\textbf{任意一个序列} \(\{x_n\} \subset E\) 都能抽出一个在 \(X\) 中收敛的子列 \(\{x_{n_k}\}\)。（注意：子列的极限允许落在 \(E\) 的闭包 \(\overline{E}\) 中，不要求一定在 \(E\) 内。）

  \begin{itemize}
  \tightlist
  \item
    \textbf{高数类比}：在有限维空间 \(\mathbb{R}^n\) 里，Bolzano-Weierstrass 定理告诉我们``有界集 \(\iff\) 列紧集''。\textbf{但在无穷维空间里，有界集通常不再列紧！} 这正是引入``紧算子''的全部动机。
  \end{itemize}
\item
  \textbf{Arzelà-Ascoli 定理（判定连续函数集列紧的头号神器 \(\to\) \hyperlink{prob:5-14}{\textcolor{LinkColor}{习题 5 第 14 题}}直接要用！）}： 设 \(C([a, b])\) 是闭区间 \([a, b]\) 上全体连续函数组成的空间，配最大值范数 \(\|f\|_\infty = \max_{x \in [a, b]} |f(x)|\)。子集 \(\mathcal{F} \subset C([a, b])\) 是\textbf{列紧集}，当且仅当它同时满足：

  \begin{itemize}
  \tightlist
  \item
    \textbf{一致有界}：存在统一常数 \(M > 0\)，使对一切 \(f \in \mathcal{F}\) 和一切 \(x \in [a, b]\)，都有 \(|f(x)| \le M\)；
  \item
    \textbf{等度连续}：对任意 \(\varepsilon > 0\)，存在统一的 \(\delta > 0\)（只依赖于 \(\varepsilon\)，\textbf{与 \(f \in \mathcal{F}\) 及点 \(x_1, x_2\) 均无关}），使得只要 \(|x_1 - x_2| < \delta\)，就对一切 \(f \in \mathcal{F}\) 同时有 \(|f(x_1) - f(x_2)| < \varepsilon\)。
  \end{itemize}
\end{enumerate}

\BookSubSection{sec:5-02}{0.2 赋范线性空间、Banach 空间、四大基础定理与 Riesz 引理}{0.2 赋范线性空间、Banach 空间、四大基础定理与 Riesz 引理}{\TagLight{前置速通}}

\begin{enumerate}
\def\labelenumi{\arabic{enumi}.}
\tightlist
\item
  \textbf{赋范线性空间与 Banach 空间}：

  \begin{itemize}
  \tightlist
  \item
    \textbf{范数 \(\|x\|\)} 就是``向量的长度''，由它自然诱导距离 \(d(x, y) = \|x - y\|\)。
  \item
    \textbf{Banach 空间} = \textbf{完备的赋范线性空间}。期末最常见的三类 Banach 空间：
  \end{itemize}
\end{enumerate}

{\def\LTcaptype{none} % do not increment counter
{\small\renewcommand{\arraystretch}{1.38}\setlength{\tabcolsep}{5pt}
\begin{longtable}[]{@{}
  >{\raggedright\arraybackslash}p{(\linewidth - 6\tabcolsep) * \real{0.2500}}
  >{\raggedright\arraybackslash}p{(\linewidth - 6\tabcolsep) * \real{0.2500}}
  >{\raggedright\arraybackslash}p{(\linewidth - 6\tabcolsep) * \real{0.2500}}
  >{\raggedright\arraybackslash}p{(\linewidth - 6\tabcolsep) * \real{0.2500}}@{}}
\toprule\noalign{}
\rowcolor{AccentSoft}
\begin{minipage}[b]{\linewidth}\raggedright
空间记号
\end{minipage} & \begin{minipage}[b]{\linewidth}\raggedright
元素是什么
\end{minipage} & \begin{minipage}[b]{\linewidth}\raggedright
范数定义
\end{minipage} & \begin{minipage}[b]{\linewidth}\raggedright
备注
\end{minipage} \\
\midrule\noalign{}
\endhead
\bottomrule\noalign{}
\endlastfoot
\(C([0,1])\) & \([0,1]\) 上的连续函数 \(f\) & \(\lVert f\rVert_\infty = \max_{x \in [0,1]} \lvert f(x)\rvert\) & 收敛即``一致收敛''，非自反 \\
\(\ell^p\) (\(1 \le p < \infty\)) & 满足 \(\sum_{n=1}^\infty \lvert x_n\rvert^p < \infty\) 的数列 & \(\lVert x\rVert_p = \left(\sum_{n=1}^\infty \lvert x_n\rvert^p\right)^{1/p}\) & \(p=2\) 是 Hilbert 空间；\(p=1\) 非自反 \\
\(L^p(\Omega)\) (\(1 \le p < \infty\)) & \(\int_\Omega \lvert f\rvert^p < \infty\) 的可测函数（几乎处处相等视为同一元素） & \(\lVert f\rVert_p = \left(\int_\Omega \lvert f(x)\rvert^p \mathrm{d}x\right)^{1/p}\) & \(p=2\) 是 Hilbert 空间 \\
\end{longtable}
}
}

\begin{enumerate}
\def\labelenumi{\arabic{enumi}.}
\setcounter{enumi}{1}
\tightlist
\item
  \textbf{有界线性算子 \(\mathcal{L}(X, Y)\) 与算子范数}：

  \begin{itemize}
  \tightlist
  \item
    设 \(X, Y\) 是赋范线性空间，线性映射 \(A: X \to Y\) 若满足存在常数 \(M \ge 0\) 使 \(\|Ax\| \le M\|x\|\)（\(\forall x \in X\)），则称 \(A\) 为\textbf{有界线性算子}。对线性算子而言，\textbf{有界 \(\iff\) 连续}。
  \item
    全体有界线性算子记为 \(\mathcal{L}(X, Y)\)（\(X = Y\) 时简记为 \(\mathcal{L}(X)\)），其\textbf{算子范数}定义为： 
\FitDisplayMath{\|A\| = \sup_{\|x\| \le 1} \|Ax\| = \sup_{x \neq 0} \frac{\|Ax\|}{\|x\|}.}
 （直觉：\(A\) 对单位球最大的``拉伸倍数''。）若 \(Y\) 是 Banach 空间，则 \(\mathcal{L}(X, Y)\) 也是 Banach 空间。
  \item
    常用记号：核空间 \(\operatorname{Ker} A = \{x \in X : Ax = 0\}\)（总为闭子空间），值域 \(\mathcal{R}(A) = \{Ax : x \in X\}\)（\textbf{不一定闭！}）。
  \end{itemize}
\item
  \textbf{Banach 空间三大核心定理（第 5 章多处用到！）}：

  \begin{itemize}
  \tightlist
  \item
    \textbf{开映射定理}：若 \(X, Y\) 为 Banach 空间，\(A \in \mathcal{L}(X, Y)\) 是\textbf{满射}，则 \(A\) 把开集映成开集（特别地，\(A\) 把 \(X\) 的单位开球映成的集合必然包含 \(Y\) 中某个以原点为心的开球 \(B_Y(0, r)\)）。\(\to\) \textbf{第 6 题 4 行秒杀证明的核心！}
  \item
    \textbf{Banach 逆算子定理}：若 \(X, Y\) 为 Banach 空间，\(A \in \mathcal{L}(X, Y)\) 是\textbf{双射}（既单又满），则逆算子 \(A^{-1}\) 自动有界，即 \(A^{-1} \in \mathcal{L}(Y, X)\)。\(\to\) \textbf{定理 5.5、定理 5.13、第 5/6 题核心！}
  \item
    \textbf{闭图像定理}：若 \(X, Y\) 为 Banach 空间，线性算子 \(A: X \to Y\) 的图像 \(\{(x, Ax) : x \in X\}\) 在 \(X \times Y\) 中是闭集，则 \(A\) 自动有界。
  \end{itemize}
\item
  \textbf{Riesz 引理与无穷维的根本特征（第 5、6 题与引理 5.1 的核心！）}：

  \begin{itemize}
  \tightlist
  \item
    \textbf{Riesz 引理}：设 \(M\) 是赋范线性空间 \(X\) 的\textbf{真闭子空间}（\(M \subsetneq X\)）。对任意 \(0 < \theta < 1\)（例如 \(\theta = \frac{1}{2}\)），存在单位向量 \(x_\theta \in X\)（\(\|x_\theta\| = 1\)）使得它到 \(M\) 的距离至少为 \(\theta\)： 
\FitDisplayMath{\operatorname{dist}(x_\theta, M) = \inf_{y \in M} \|x_\theta - y\| \ge \theta.}

  \item
    \textbf{最关键推论（闭单位球列紧 \(\iff\) 有限维）}：赋范线性空间 \(X\) 是\textbf{有限维的}，当且仅当其闭单位球 \(B_X = \{x \in X : \|x\| \le 1\}\) 是\textbf{列紧集}。

    \begin{itemize}
    \tightlist
    \item
      证明直觉：若 \(\dim X = \infty\)，任取单位向量 \(x_1\)；令 \(M_1 = \operatorname{span}\{x_1\}\)（有限维必真闭），由 Riesz 引理取单位向量 \(x_2\) 使 \(\|x_2 - x_1\| \ge \frac{1}{2}\)；再令 \(M_2 = \operatorname{span}\{x_1, x_2\}\) 取 \(x_3\) 使到 \(M_2\) 距离 \(\ge \frac{1}{2}\)\ldots\ldots 如此得单位向量列 \(\{x_n\}\)，任意两项距离 \(\|x_m - x_n\| \ge \frac{1}{2}\) (\(m \neq n\))，绝不可能有收敛子列！
    \end{itemize}
  \end{itemize}
\end{enumerate}

\BookSubSection{sec:5-03}{0.3 Hilbert 空间、标准正交基与共轭算子}{0.3 Hilbert 空间、标准正交基与共轭算子}{\TagLight{前置速通}}

\begin{enumerate}
\def\labelenumi{\arabic{enumi}.}
\tightlist
\item
  \textbf{内积与 Hilbert 空间}：

  \begin{itemize}
  \tightlist
  \item
    \textbf{Hilbert 空间} \(H\) = 带有内积 \((\cdot, \cdot)\)、且由诱导范数 \(\|x\| = \sqrt{(x, x)}\) 构成\textbf{完备空间}的线性空间（即``带角度与垂直概念的无穷维欧氏空间''）。
  \item
    \textbf{本教材内积约定}（务必看准）：\textbf{第一变元线性、第二变元共轭线性}，即对任意 \(\alpha, \beta \in \mathbb{C}\)： 
\FitDisplayMath{(\alpha x + \beta y, z) = \alpha(x, z) + \beta(y, z), \qquad (x, \alpha y + \beta z) = \bar{\alpha}(x, y) + \bar{\beta}(x, z), \qquad (y, x) = \overline{(x, y)}.}

  \item
    \textbf{Cauchy-Schwarz 不等式}：\(|(x, y)| \le \|x\|\|y\|\)。
  \end{itemize}
\item
  \textbf{正交投影定理（第 15 题核心！）}：

  \begin{itemize}
  \tightlist
  \item
    设 \(M\) 是 Hilbert 空间 \(H\) 的\textbf{闭}线性子空间，记其正交补 \(M^\perp = \{y \in H : (y, x) = 0, \, \forall x \in M\}\)。则 \(M^\perp\) 也是闭子空间，且 \(H\) 具有唯一的正交直和分解： 
\FitDisplayMath{H = M \oplus M^\perp,}
 即任意 \(x \in H\) 都可唯一写成 \(x = x_1 + x_2\)（\(x_1 \in M, x_2 \in M^\perp\)），且满足勾股定理 \(\|x\|^2 = \|x_1\|^2 + \|x_2\|^2\)。
  \end{itemize}
\item
  \textbf{标准正交基与 Parseval 等式（5.2 节与第 15 题核心！）}：

  \begin{itemize}
  \tightlist
  \item
    若 \(\{e_n\}_{n=1}^\infty\) 是可分 Hilbert 空间 \(H\) 的一组\textbf{标准正交基}（两两正交、单位长度、且极大/张成稠密子空间），则对任意 \(x \in H\) 都有：

    \begin{itemize}
    \tightlist
    \item
      \textbf{Fourier 展开}：\(x = \sum_{n=1}^\infty (x, e_n) e_n\)；
    \item
      \textbf{Parseval 等式（无穷维勾股定理）}： 
\FitDisplayMath{\|x\|^2 = \sum_{n=1}^\infty |(x, e_n)|^2.}

    \item
      若仅满足标准正交（未必极大），则把等号换成 \(\le\)，称为 \textbf{Bessel 不等式}。
    \end{itemize}
  \end{itemize}
\item
  \textbf{Hilbert 空间上的共轭算子 \(A^*\)}：

  \begin{itemize}
  \tightlist
  \item
    对 \(A \in \mathcal{L}(H, G)\)，唯一存在有界线性算子 \(A^* \in \mathcal{L}(G, H)\) 使得对一切 \(x \in H, y \in G\) 都有： 
\FitDisplayMath{(Ax, y)_G = (x, A^* y)_H.}

  \item
    性质：\(\|A^*\| = \|A\|\)，\((A^*)^* = A\)，\((AB)^* = B^* A^*\)。（矩阵类比：复矩阵的共轭转置 \(A^H = \overline{A}^T\)。）
  \end{itemize}
\end{enumerate}

\BookSubSection{sec:5-04}{0.4 对偶空间、弱收敛、商空间与余维数}{0.4 对偶空间、弱收敛、商空间与余维数}{\TagLight{前置速通}}

\begin{enumerate}
\def\labelenumi{\arabic{enumi}.}
\tightlist
\item
  \textbf{对偶空间 \(X'\) 与 Banach 空间共轭算子}：

  \begin{itemize}
  \tightlist
  \item
    设 \(X\) 为赋范线性空间，所有连续线性泛函（即值域在数域 \(\mathbb{C}\) 或 \(\mathbb{R}\) 中的有界线性算子）构成的空间 \(X' = \mathcal{L}(X, \mathbb{C})\) 称为 \(X\) 的\textbf{对偶空间}（或共轭空间）。
  \item
    对 \(A \in \mathcal{L}(X, Y)\)，其\textbf{Banach 空间共轭算子} \(A^* \in \mathcal{L}(Y', X')\) 定义为：对任意 \(y' \in Y'\)，\((A^* y')(x) = y'(Ax)\)（\(\forall x \in X\)）。同样有 \(\|A^*\| = \|A\|\)。
  \end{itemize}
\item
  \textbf{强收敛 vs. 弱收敛（定理 5.1 与第 6 题反例用到！）}：

  \begin{itemize}
  \tightlist
  \item
    \textbf{强收敛}（按范数收敛）：\(\lim_{n\to\infty} \|x_n - x_0\| = 0\)，记作 \(x_n \to x_0\)。
  \item
    \textbf{弱收敛}：对\textbf{每一个}有界线性泛函 \(f \in X'\)，都有数列的极限 \(\lim_{n\to\infty} f(x_n) = f(x_0)\)（在 Hilbert 空间中由 Riesz 表示定理等价于：对一切 \(y \in H\)，\((x_n, y) \to (x_0, y)\)），常记作 \(x_n \rightharpoonup x_0\)。
  \item
    \textbf{直觉对比}：在 \(\ell^2\) 中，标准正交基 \(\{e_n\}\) 满足对任意 \(y = (y_1, y_2, \dots) \in \ell^2\) 都有 \((e_n, y) = \bar{y}_n \to 0\)，所以 \(e_n \rightharpoonup 0\)（弱收敛到 \(0\)）；但 \(\|e_n - 0\| = 1 \not\to 0\)，所以\textbf{不强收敛}！
  \item
    \textbf{重要特例}：强收敛必推弱收敛；反之一般不成立。但 \(\ell^1\) 空间具有特殊的 \textbf{Schur 性质}：在 \(\ell^1\) 中，\textbf{弱收敛与强收敛等价}（教材第 183 页注记 5.1）。
  \item
    \textbf{自反空间}：若 \(X\) 到其二次对偶 \(X''\) 的自然嵌入是满射，则称 \(X\) 为\textbf{自反 Banach 空间}。所有 Hilbert 空间和 \(\ell^p, L^p\)（\(1 < p < \infty\)）都自反；\(C([0,1])\) 和 \(\ell^1\) 不自反。自反空间的有界闭球弱列紧（任意有界列必有弱收敛子列）。
  \end{itemize}
\item
  \textbf{商空间 \(X / M\) 与余维数 \(\operatorname{codim} M\)（5.3 节 Fredholm 算子与第 6、22 题必考！）}：

  \begin{itemize}
  \tightlist
  \item
    设 \(M\) 是线性空间 \(X\) 的线性子空间。把相差一个 \(M\) 中元素的向量视为``同一类''，即定义等价类（陪集）\(\tilde{x} = x + M = \{x + m : m \in M\}\)。所有等价类组成的线性空间称为\textbf{商空间}，记作 \(X / M\)。
  \item
    \textbf{几何直觉}：把 \(\mathbb{R}^3\) 模掉 \(xy\)-平面 \(M\)，就是把所有水平高度相同的点捏成一个点，剩下的空间 \(\mathbb{R}^3 / M\) 本质上就是 1 维的 \(z\)-轴！
  \item
    \textbf{商范数}：当 \(M\) 是 Banach 空间 \(X\) 的\textbf{闭}子空间时，在 \(X / M\) 上定义范数（几何上即点 \(x\) 到闭子空间 \(M\) 的距离）： 
\FitDisplayMath{\|\tilde{x}\| = \inf_{m \in M} \|x + m\| = \inf_{z \in \tilde{x}} \|z\|,}
 此时 \(X / M\) 成为一个 \textbf{Banach 空间}，且自然商映射 \(q: X \to X/M\)（\(q(x) = \tilde{x}\)）是连续线性满射（\(\|q\| \le 1\)）。
  \item
    \textbf{余维数（codimension）}：定义子空间 \(M\) 在 \(X\) 中的余维数为商空间的维数： 
\FitDisplayMath{\operatorname{codim} M \triangleq \dim(X / M).}
 \textbf{一句话理解}：\(\operatorname{codim} M\) 就是``为了把 \(M\) 补全成整个空间 \(X\)，还需要额外添加多少个线性无关的向量''。例如对算子 \(A: X \to Y\)，\(\operatorname{codim} \mathcal{R}(A)\) 就是值域 \(\mathcal{R}(A)\) 距离铺满 \(Y\) 还``差了几个维度''。
  \end{itemize}
\end{enumerate}

\BookPart{}{一、第 5 章核心知识零基础精讲}{一、第 5 章核心知识零基础精讲}

本章恒设 \(X, Y\) 为 Banach 空间（教材第 182 页约定）。整章围绕两条主线展开：

\begin{itemize}
\tightlist
\item
  \textbf{紧算子}：无穷维空间里``脾气最像有限维矩阵''的算子（把庞大松散的无穷维单位球``压扁、拧紧''成列紧集）；
\item
  \textbf{Fredholm 算子}：无穷维空间里``几乎可逆''的算子（只在有限维上丢失信息、只在有限维上没覆盖到）。
\end{itemize}

\BookSubSection{sec:5-1}{5.1 紧算子 \(\mathcal{C}(X, Y)\)}{5.1 紧算子 C(X, Y)}{\TagMust{闭卷必背}}

\SubSubHeading{1. 定义与直觉（定义 5.1，教材第 182 页）}

\begin{quote}
\textbf{定义 5.1}：设 \(A \in \mathcal{L}(X, Y)\)。如果 \(A\) 把 \(X\) 中每个有界集都映射为 \(Y\) 中的列紧集，则称 \(A\) 为\textbf{紧算子}（或\textbf{全连续算子}）。\(X\) 到 \(Y\) 的全体紧算子记为 \(\mathcal{C}(X, Y)\)；当 \(X = Y\) 时简记为 \(\mathcal{C}(X)\)。
\end{quote}

\begin{itemize}
\tightlist
\item
  \textbf{等价序列刻画（做题最常用！）}：\(A \in \mathcal{L}(X, Y)\) 是紧算子 \(\iff\) 对 \(X\) 中\textbf{任意有界序列} \(\{x_n\}\)，其像序列 \(\{Ax_n\}\) 在 \(Y\) 中\textbf{必有收敛子列}。
\item
  \textbf{反面典型（例 5.2，教材第 183 页）}：在无穷维空间（如 \(\ell^2\)）上，\textbf{恒等算子 \(I\) 不是紧算子}！因为标准正交基 \(\{e_n\}\) 有界（\(\|e_n\|=1\)），但 \(\|Ie_j - Ie_k\| = \|e_j - e_k\| = \sqrt{2}\)（\(j \neq k\)），根本抽不出收敛子列。
\item
  \textbf{正面典型（例 5.1，教材第 182--183 页）}：若 \(K(s, t)\) 在 \([0,1] \times [0,1]\) 上连续，则 \(C([0,1])\) 上的 Fredholm 积分算子 
\FitDisplayMath{(Ax)(s) = \int_0^1 K(s, t) x(t) \mathrm{d}t}
 是紧算子（证明靠 Arzelà-Ascoli 定理：先证一致有界，再利用 \(K\) 的一致连续性证等度连续）。\textbf{直觉}：积分有``平滑/抚平毛刺''的作用，所以积分算子常常是紧算子；微分算子会放大震荡，通常是无界且不紧的。
\end{itemize}

\SubSubHeading{2. 紧算子的核心性质（定理 5.1--5.8，教材第 183--187 页）}

{\def\LTcaptype{none} % do not increment counter
{\small\renewcommand{\arraystretch}{1.38}\setlength{\tabcolsep}{5pt}
\begin{longtable}[]{@{}
  >{\raggedright\arraybackslash}p{(\linewidth - 4\tabcolsep) * \real{0.1951}}
  >{\raggedright\arraybackslash}p{(\linewidth - 4\tabcolsep) * \real{0.4390}}
  >{\raggedright\arraybackslash}p{(\linewidth - 4\tabcolsep) * \real{0.3659}}@{}}
\toprule\noalign{}
\rowcolor{AccentSoft}
\begin{minipage}[b]{\linewidth}\raggedright
教材编号
\end{minipage} & \begin{minipage}[b]{\linewidth}\raggedright
精准结论
\end{minipage} & \begin{minipage}[b]{\linewidth}\raggedright
零基础直觉与考试用途
\end{minipage} \\
\midrule\noalign{}
\endhead
\bottomrule\noalign{}
\endlastfoot
\textbf{定理 5.1} & \(A \in \mathcal{C}(X, Y)\) 且 \(x_n \rightharpoonup x_0\)（弱收敛） \(\implies Ax_n \to Ax_0\)（强收敛）。（注记 5.1：\(X\) 自反时逆命题成立；非自反时如 \(\ell^1\) 不成立。） & 紧算子有``变弱为强''的魔力；常用于处理弱收敛极限（如习题 13、16）。 \\
\textbf{定理 5.2} & \(A \in \mathcal{C}(X, Y) \implies \mathcal{R}(A)\) 是可分的（有可数稠密子集）。 & 说明紧算子的像集``体量不会太大''。 \\
\textbf{定理 5.3} & \textbf{(i)} \(\alpha_1 A_1 + \alpha_2 A_2 \in \mathcal{C}(X, Y)\)；\textbf{(ii)} \(B_1 \in \mathcal{L}(Y, Z), B_2 \in \mathcal{L}(Z, X) \implies A_1 B_2 \in \mathcal{C}(Z, Y), \; B_1 A_1 \in \mathcal{C}(X, Z)\)。 & \textbf{双边理想性质（超高频！）}：紧算子左乘或右乘\textbf{任何}有界算子，结果还是紧算子！第 5、20、22 题全靠它。 \\
\textbf{定理 5.4} & 若 \(A_n \in \mathcal{C}(X, Y)\) 且 \(\lVert A_n - A\rVert \to 0\)，则 \(A \in \mathcal{C}(X, Y)\)。 & \textbf{按算子范数极限封闭}！证明算子紧的核武器：找一列有限秩算子 \(A_n\) 逼近 \(A\)。 \\
\textbf{推论 5.1} & \(\mathcal{C}(X, Y)\) 是 \(\mathcal{L}(X, Y)\) 的闭线性子空间，本身是 Banach 空间。 & 由定理 5.3 + 定理 5.4 立得。 \\
\textbf{定理 5.5} & 若 \(X\) 无穷维、\(A \in \mathcal{C}(X, Y)\) 且 \(A\) 单射，则 \(\mathcal{R}(A) \neq Y\)。 & \textbf{无穷维紧算子不可能有界可逆！}（若可逆则 \(I_X = A^{-1}A\) 紧，矛盾于 Riesz 引理。）\(\to\) 对应\hyperlink{prob:5-5}{\textcolor{LinkColor}{习题 5 第 5 题}}。 \\
\textbf{定理 5.6} & \(A \in \mathcal{C}(X), \, \lambda \neq 0 \implies \mathcal{R}(\lambda I - A)\) 是\textbf{闭子空间}。 & 单个紧算子 \(A\) 的值域通常不闭，但加上 \(\lambda I\) 后值域一定闭！ \\
\textbf{定理 5.7}\emph{(Schauder)} & \(A \in \mathcal{C}(X, Y) \implies A^* \in \mathcal{C}(Y', X')\)。（事实上互为充要。） & 紧性在取共轭时保持不变。 \\
\textbf{引理 5.1 \&}\textbf{定理 5.8}\emph{(Riesz-Fredholm)} & 设 \(A \in \mathcal{C}(X), \, \lambda \neq 0\)，则：(1) \(\mathcal{R}(\lambda I - A) = \mathcal{N}'(\operatorname{Ker}(\lambda I - A^*))\)；(2) \(\mathcal{R}(\lambda I - A^*) = \mathcal{N}(\operatorname{Ker}(\lambda I - A))\)；(3) \(\dim \operatorname{Ker}(\lambda I - A) = \dim \operatorname{Ker}(\lambda I - A^*)\)；(4) \(\operatorname{codim} \mathcal{R}(\lambda I - A) = \dim \operatorname{Ker}(\lambda I - A) < \infty\)。 & \textbf{超级核心结论}：\CircNum{1} \textbf{Fredholm 择一性}：对 \(\lambda \neq 0\)，\(\lambda I - A\) 单射 \(\iff\) 满射 \(\iff\) 有界可逆（像极了有限维方阵！）；\CircNum{2} 令 \(\lambda = 1\) 得 \(I - A \in \mathcal{F}(X)\) 且 \textbf{\(\operatorname{ind}(I - A) = 0\)}（第 20 题关键基石！）。 \\
\end{longtable}
}
}

\begin{quote}
\textbf{教材勘误提示}：教材第 187 页定理 5.8(2) 末尾多印了一个``\(< \infty\)''，正确表述为 \(\mathcal{R}(\lambda I - A^*) = \mathcal{N}(\operatorname{Ker}(\lambda I - A))\)。
\end{quote}

\begin{quote}
\textbf{符号解码：\(\mathcal{N}(M)\) 与 \(\mathcal{N}'(N)\) 是什么意思？（教材 4.4 节零化子空间）}

\begin{enumerate}
\def\labelenumi{\arabic{enumi}.}
\tightlist
\item
  \textbf{严格定义}：设 \(X\) 是 Banach 空间、\(X'\) 是其对偶空间，

  \begin{itemize}
  \tightlist
  \item
    对子集 \(M \subset X\)，其\textbf{零化子空间（Annihilator）}定义为 
\FitDisplayMath{\mathcal{N}(M) \triangleq \{f \in X' : f(x) = 0,\; \forall x \in M\} \subset X',}
 大白话就是：\(X'\) 中所有把 \(M\) 里的向量全都``吃成 0''的有界线性泛函 \(f\)。
  \item
    对子集 \(N \subset X'\)，其\textbf{前零化子空间（Pre-annihilator）}定义为 
\FitDisplayMath{\mathcal{N}'(N) \triangleq \{x \in X : f(x) = 0,\; \forall f \in N\} \subset X,}
 大白话就是：\(X\) 中所有被 \(N\) 里的泛函全都``吃成 0''的向量 \(x\)。
  \end{itemize}
\item
  \textbf{零基础直观理解（Banach 空间版的正交补 \(M^\perp\)）}： 在 Hilbert 空间或线性代数里，有内积可以写 \((x, y) = 0\) 表示``垂直''；一般 Banach 空间没有内积，就用对偶配对 \(f(x) = 0\) 来代替``垂直''——因此 \(\mathcal{N}(M)\) 在许多教材里直接记作 \(M^\perp\)！
\item
  \textbf{翻成大白话看懂定理 5.8(1)(2)}：

  \begin{itemize}
  \tightlist
  \item
    \textbf{(1) \(\mathcal{R}(\lambda I - A) = \mathcal{N}'(\operatorname{Ker}(\lambda I - A^*))\)}：原方程 \((\lambda I - A)x = y\) 有解 \(\iff\) 右端项 \(y\) 被共轭齐次方程 \((\lambda I - A^*)f = 0\) 的所有解 \(f\) 都作用成 \(0\)（即 \(f(y) = 0\)）。
  \item
    \textbf{(2) \(\mathcal{R}(\lambda I - A^*) = \mathcal{N}(\operatorname{Ker}(\lambda I - A))\)}：共轭方程 \((\lambda I - A^*)f = g\) 有解 \(\iff\) 右端项 \(g\) 把原齐次方程 \((\lambda I - A)x = 0\) 的所有解 \(x\) 都作用成 \(0\)（即 \(g(x) = 0\)）。这正是大一线代里 \(\operatorname{Im}(B^T) = (\operatorname{Ker} B)^\perp\) 的无穷维推广！
  \end{itemize}
\end{enumerate}
\end{quote}

\SubSubHeading{3. 有限秩算子与逼近（定义 5.2、定理 5.9，教材第 189 页）}

\begin{itemize}
\tightlist
\item
  \textbf{定义 5.2}：设 \(A \in \mathcal{L}(X, Y)\)，如果 \(\dim \mathcal{R}(A) < \infty\)，则称 \(A\) 为\textbf{有限秩算子}。
\item
  \textbf{结构与紧性（例 5.3）}：\(A\) 是有限秩算子 \(\iff\) 存在 \(f_1, \dots, f_k \in X'\) 与 \(y_1, \dots, y_k \in Y\) 使 
\FitDisplayMath{Ax = \sum_{j=1}^k f_j(x) y_j.}
 由于有限维空间中的有界集必列紧，\textbf{任何有限秩算子都是紧算子}。
\item
  \textbf{定理 5.9（可分 Hilbert 空间上的有限秩逼近）}：设 \(A\) 是可分 Hilbert 空间 \(H\) 上的紧算子，则存在一列有限秩算子 \(\{A_n\}_{n=1}^\infty\) 使得 \(\lim_{n\to\infty} \|A_n - A\| = 0\)。（具体构造就是截断正交展开 \(A_n x = \sum_{j=1}^n (x, e_j) A e_j\)。）

  \begin{itemize}
  \tightlist
  \item
    \textbf{重要推论}：结合定理 5.4，在 Hilbert 空间上，\textbf{紧算子 \(\iff\) 有限秩算子按算子范数的极限}！
  \item
    \emph{延伸说明}：对一般 Banach 空间，Enflo (1973) 证明存在不具逼近性质的 Banach 空间，但本科常见的 \(C([0,1]), \ell^p, L^p\) 等空间都具备有限秩逼近性质。
  \end{itemize}
\end{itemize}

\SubPillMain{4. 期末实操：证明/否定``\(A\) 是紧算子''的四大高频套路}

{\def\LTcaptype{none} % do not increment counter
{\small\renewcommand{\arraystretch}{1.38}\setlength{\tabcolsep}{5pt}
\begin{longtable}[]{@{}
  >{\raggedright\arraybackslash}p{(\linewidth - 4\tabcolsep) * \real{0.3333}}
  >{\raggedright\arraybackslash}p{(\linewidth - 4\tabcolsep) * \real{0.3333}}
  >{\raggedright\arraybackslash}p{(\linewidth - 4\tabcolsep) * \real{0.3333}}@{}}
\toprule\noalign{}
\rowcolor{AccentSoft}
\begin{minipage}[b]{\linewidth}\raggedright
题目类型
\end{minipage} & \begin{minipage}[b]{\linewidth}\raggedright
优先武器
\end{minipage} & \begin{minipage}[b]{\linewidth}\raggedright
怎么操作（3 步套路）
\end{minipage} \\
\midrule\noalign{}
\endhead
\bottomrule\noalign{}
\endlastfoot
\textbf{套路 1}：\(C([a,b])\) 上的积分算子 & \textbf{Arzelà-Ascoli 定理} & \CircNum{1} 先证 \(A\) 线性有界；\CircNum{2} 取有界集 \(\lVert f\rVert_\infty \le M\)，证像集\textbf{一致有界}；\CircNum{3} 估计 \(\lvert Af(x_1) - Af(x_2)\rvert\) 证\textbf{等度连续}（如第 14 题）。 \\
\textbf{套路 2}：\(\ell^p\) 上的对角/移位/矩阵算子 & \textbf{有限秩逼近（定理 5.4）} & \CircNum{1} 构造前 \(N\) 项截断算子 \(A_N\)（显为有限秩 \(\implies\) 紧）；\CircNum{2} 估计余项算子范数 \(\lVert A - A_N\rVert \to 0\)；\CircNum{3} 由定理 5.4 立即得 \(A\) 紧。 \\
\textbf{套路 3}：Hilbert 空间（尤其是 \(L^2, \ell^2\)）上的积分/矩阵算子 & \textbf{Hilbert-Schmidt 判定（定理 5.12(5)）} & 验证 \(\sum_{n=1}^\infty \lVert Ae_n\rVert^2 < \infty\)（对 \(L^2\) 积分算子即验证 \(\iint \lvert K(x,y)\rvert^2 \mathrm{d}x\mathrm{d}y < \infty\)），立即推出紧性！ \\
\textbf{套路 4}：证明算子\textbf{不是}紧算子 / 紧算子不可能具备某性质 & \textbf{反证法 + Riesz 引理（或直接找反例有界列）} & \CircNum{1} 反设紧性，利用理想性质凑出 \(I_X \in \mathcal{C}(X)\)（第 5、6 题）；\CircNum{2} 或直接构造有界列 \(\{x_n\}\) 使 \(\lVert Ax_m - Ax_n\rVert \ge c > 0\)（例 5.2）。 \\
\end{longtable}
}
}

\SoftDivider

\BookSubSection{sec:5-2}{5.2 Hilbert-Schmidt 算子 \(\mathcal{L}^2(H, G)\)}{5.2 Hilbert-Schmidt 算子 L²(H, G)}{\TagMust{闭卷必背}}

\textbf{零基础直觉}：在线性代数中，衡量矩阵 \(A = (a_{ij})\) 大小的``平方和开根号''叫 \textbf{Frobenius 范数} \(\|A\|_F = \left(\sum_{i,j} |a_{ij}|^2\right)^{1/2}\)。把矩阵的列向量 \(Ae_n\) 的长度平方加起来再开方，推广到无穷维 Hilbert 空间，就是 \textbf{Hilbert-Schmidt 范数}！

本节默认 \(H, G\) 为可分 Hilbert 空间（定义对一般 Hilbert 空间亦可自然推广）。

\SubSubHeading{1. 定义与基无关性（定义 5.4、式 (5.8)，教材第 196 页）}

\begin{quote}
\textbf{定义 5.4}：设 \(A\) 是 \(H\) 到 \(G\) 的线性算子。如果存在 \(H\) 上的一组标准正交基 \(\{e_n\}_{n=1}^\infty\) 使得 
\FitDisplayMath{\|A\|_2 \triangleq \left(\sum_{n=1}^\infty \|Ae_n\|^2\right)^{\frac{1}{2}} < \infty,}
 则称 \(A\) 是 \(H\) 到 \(G\) 的 \textbf{Hilbert-Schmidt 算子}，全体记为 \(\mathcal{L}^2(H, G)\)；当 \(H = G\) 时简记为 \(\mathcal{L}^2(H)\)。
\end{quote}

\begin{itemize}
\tightlist
\item
  \textbf{为什么与标准正交基的选取无关？（式 (5.8)）}：设 \(\{e_m\}, \{\tilde{e}_m\}\) 为 \(H\) 的两组标准正交基，\(\{f_n\}\) 为 \(G\) 的标准正交基，由 Parseval 等式与 \(A^*\) 的定义： 
\FitDisplayMath{\sum_{m=1}^\infty \|Ae_m\|^2 = \sum_{m,n=1}^\infty |(Ae_m, f_n)|^2 = \sum_{m,n=1}^\infty |(e_m, A^* f_n)|^2 = \sum_{n=1}^\infty \|A^* f_n\|^2 = \sum_{m,n=1}^\infty |(\tilde{e}_m, A^* f_n)|^2 = \sum_{m=1}^\infty \|A\tilde{e}_m\|^2.}
 这个等式一石二鸟：同时证明了 \textbf{\(\|A\|_2\) 与基的选取无关} 和 \textbf{\(\|A\|_2 = \|A^*\|_2\)}！ \emph{(注：教材第 196 页式 (5.8) 中将复内积的模平方 \(|(Ae_m, f_n)|^2\) 漏写了绝对值符号，做题时请务必写绝对值。)}
\end{itemize}

\SubSubHeading{2. Hilbert-Schmidt 算子的六大性质（定理 5.12，教材第 197--198 页）}

设 \(A \in \mathcal{L}^2(H, G)\)，则：

\begin{enumerate}
\def\labelenumi{\arabic{enumi}.}
\tightlist
\item
  \textbf{共轭等距}：\(\|A\|_2 = \|A^*\|_2\)，故 \(A^* \in \mathcal{L}^2(G, H)\)；
\item
  \textbf{控制算子范数}：\(\|A\| \le \|A\|_2\)（特别地，HS 算子自动是有界算子）；
\item
  \textbf{双边理想与范数不等式}：若 \(B_1 \in \mathcal{L}(H), B_2 \in \mathcal{L}(G)\)，则 \(AB_1, B_2 A \in \mathcal{L}^2(H, G)\)，且 
\FitDisplayMath{\|B_2 A\|_2 \le \|B_2\|\|A\|_2, \qquad \|AB_1\|_2 \le \|B_1\|\|A\|_2;}

\item
  \textbf{自身构成 Hilbert 空间}：在定义于 \(\mathcal{L}^2(H, G)\) 的内积 
\FitDisplayMath{(A_1, A_2)_2 \triangleq \sum_{n=1}^\infty (A_1 e_n, A_2 e_n)}
 及诱导范数 \(\|\cdot\|_2\) 下，\(\mathcal{L}^2(H, G)\) 是 Hilbert 空间；
\item
  \textbf{必为紧算子}：每个 Hilbert-Schmidt 算子都是\textbf{紧算子}，且有限秩算子按 \(\|\cdot\|_2\) 范数在 \(\mathcal{L}^2(H, G)\) 中稠密；
\item
  \textbf{可分性}：\(\mathcal{L}^2(H, G)\) 是可分 Hilbert 空间。
\end{enumerate}

\begin{itemize}
\tightlist
\item
  \textbf{必背算子包含链（在无穷维可分 Hilbert 空间上均为严格真包含）}： 
\FitDisplayMath{\{\text{有限秩算子}\} \;\subsetneq\; \mathcal{L}^2(H, G) \;\subsetneq\; \mathcal{C}(H, G) \;\subsetneq\; \mathcal{L}(H, G).}
 \emph{(对角算子记忆法：\(A e_n = a_n e_n\)。有限秩 \(\iff\) 仅有限个 \(a_n \neq 0\)；HS \(\iff \sum |a_n|^2 < \infty\)；紧 \(\iff a_n \to 0\)；有界 \(\iff \{a_n\}\) 有界。例如 \(a_n = 1/\sqrt{n}\) 紧但非 HS，\(a_n = 1/n\) 是 HS 但非有限秩！)}
\end{itemize}

\SoftDivider

\BookSubSection{sec:5-3}{5.3 Fredholm 算子 \(\mathcal{F}(X, Y)\)}{5.3 Fredholm 算子 F(X, Y)}{\TagMust{闭卷必背}}

\textbf{零基础直觉}：一个有界线性算子 \(A: X \to Y\) 可逆，意味着 \(\dim \operatorname{Ker} A = 0\)（不压缩任何维度）且 \(\operatorname{codim} \mathcal{R}(A) = 0\)（完全铺满 \(Y\)）。Fredholm 算子把``等于 0''放宽到``\textbf{有限}''——也就是\textbf{``几乎可逆''的算子}！

\SubSubHeading{1. 定义与指标（定义 5.5、定义 5.6，教材第 198--199 页）}

\begin{quote}
\textbf{定义 5.5}：设 \(X, Y\) 均为 Banach 空间，\(A \in \mathcal{L}(X, Y)\)。如果同时满足：

\begin{enumerate}
\def\labelenumi{\arabic{enumi}.}
\tightlist
\item
  \(\mathcal{R}(A)\) 是 \(Y\) 中的\textbf{闭}子空间；
\item
  \(\dim \operatorname{Ker}(A) < +\infty\)（核空间有限维）；
\item
  \(\operatorname{codim} \mathcal{R}(A) = \dim(Y / \mathcal{R}(A)) < +\infty\)（值域余维有限），
\end{enumerate}

则称 \(A\) 为 \textbf{Fredholm 算子}，全体记为 \(\mathcal{F}(X, Y)\)；当 \(X = Y\) 时简记为 \(\mathcal{F}(X)\)。
\end{quote}

\begin{quote}
\textbf{重要注记（深度理解：为什么条件 (3) \(\operatorname{codim}\mathcal{R}(A) < +\infty\) 能自动推出条件 (1) \(\mathcal{R}(A)\) 闭？）}：

\begin{itemize}
\tightlist
\item
  \textbf{先避开一个直觉误区}：一个普通的有限余维线性子空间\textbf{并不一定是闭的}！（反例：在无穷维 Banach 空间中，不连续线性泛函 \(f\) 的核 \(\operatorname{Ker} f\) 余维数为 \(1\)，却在空间中处处稠密而不闭。）
\item
  \textbf{这里为何能自动变闭？} 全靠两大前提联手：\textbf{\(X, Y\) 完备（Banach 空间）} \(+\) \textbf{\(\mathcal{R}(A)\) 是有界线性算子 \(A\) 的值域}！既然 \(\mathcal{R}(A)\) 距铺满 \(Y\) 只差一个有限维小缺口 \(W\)，我们把 \(W\) 补上去就能造出一个 Banach 空间之间的连续双射，再用开映射定理（Banach 逆算子定理）直接逼出 \(\mathcal{R}(A)\) 必闭！
\item
  \textbf{5 步完整证明（Kato 引理）}：设 \(\operatorname{codim}\mathcal{R}(A) = \dim(Y/\mathcal{R}(A)) = k < +\infty\)。

  \begin{enumerate}
  \def\labelenumi{\arabic{enumi}.}
  \tightlist
  \item
    \textbf{消核变单射}：因 \(A\) 连续，\(\operatorname{Ker}A\) 闭，故商空间 \(\tilde{X} \triangleq X/\operatorname{Ker}A\) 是 Banach 空间，诱导算子 \(A_0([x]) \triangleq Ax\) 为有界单射，且 \(\mathcal{R}(A_0) = \mathcal{R}(A)\)。
  \item
    \textbf{找 \(k\) 维补丁 \(W\)}：取 \(Y\) 的 \(k\) 维子空间 \(W\) 使代数直和 \(Y = \mathcal{R}(A) \oplus W\) 成立。
  \item
    \textbf{拼成直和 Banach 空间}：令 \(Z \triangleq \tilde{X} \oplus W\)，配范数 \(\lVert([x], w)\rVert_Z \triangleq \lVert[x]\rVert_{\tilde{X}} + \lVert w\rVert_Y\)。因 \(\tilde{X}\) 完备、\(W\) 有限维必完备，故 \(Z\) 是 Banach 空间。
  \item
    \textbf{构造连续双射 \(T: Z \to Y\)}：定义 \(T([x], w) \triangleq A_0[x] + w = Ax + w\)。显然 \(T\) 有界（\(\lVert T([x],w)\rVert_Y \le \max(\lVert A_0\rVert, 1)\lVert([x],w)\rVert_Z\)）；由 \(Y = \mathcal{R}(A) \oplus W\) 与 \(A_0\) 单射立得 \(T\) 既是单射又是满射。
  \item
    \textbf{开映射定理收尾}：由 Banach 逆算子定理（开映射定理推论），\(T^{-1}: Y \to Z\) 连续，即 \(T\) 是拓扑同构（把闭集映成闭集）。因 \(\tilde{X} \oplus \{0\}\) 显然是 \(Z\) 中的闭子空间，故其像 \(\mathcal{R}(A) = T(\tilde{X} \oplus \{0\})\) 必为 \(Y\) 中的\textbf{闭子空间}！\(\blacksquare\)
  \end{enumerate}
\item
  \textbf{考场提醒}：虽然条件 (3) 能推出条件 (1)，但教材定义 5.5 明确列出了全部 3 条，\textbf{闭卷写定义与做常规证明题时仍请按定义 5.5 把这 3 条写全（尤其别漏证``\(\mathcal{R}(A)\) 闭''）！}
\end{itemize}
\end{quote}

\begin{quote}
\textbf{定义 5.6（Fredholm 指标）}：设 \(A \in \mathcal{F}(X, Y)\)，定义 \(A\) 的\textbf{指标（index）}为整数： 
\FitDisplayMath{\operatorname{ind} A = \dim \operatorname{Ker}(A) - \operatorname{codim} \mathcal{R}(A).}

\end{quote}

\begin{itemize}
\tightlist
\item
  \textbf{三个必记标杆例子（例 5.8--5.9，教材第 199 页）}：

  \begin{enumerate}
  \def\labelenumi{\arabic{enumi}.}
  \tightlist
  \item
    \textbf{可逆算子}（注记 5.3）：\(\dim \operatorname{Ker} A = 0, \, \operatorname{codim} \mathcal{R}(A) = 0 \implies A \in \mathcal{F}(X, Y)\) 且 \(\operatorname{ind} A = 0\)。
  \item
    \textbf{紧算子的恒等扰动 \(I - K\)}（例 5.8）：对任意 \(K \in \mathcal{C}(X)\)，由定理 5.6 与定理 5.8(4) 知 \(I - K \in \mathcal{F}(X)\) 且 \textbf{\(\operatorname{ind}(I - K) = 0\)}！（第 20 题必用！）
  \item
    \textbf{\(\ell^2\) 上的左/右平移算子}（例 5.9，最能体现``无穷维与有限维不同''的例子！）：

    \begin{itemize}
    \tightlist
    \item
      左平移 \(L(x_1, x_2, x_3, \dots) = (x_2, x_3, \dots)\)：丢掉了第 1 个分量，故 \(\operatorname{Ker} L = \operatorname{span}\{e_1\}\)（1 维），但它是满射 \(\mathcal{R}(L) = \ell^2\)（闭，余维 0），所以 \(L \in \mathcal{F}(\ell^2)\) 且 \textbf{\(\operatorname{ind} L = 1 - 0 = 1\)}！
    \item
      右平移 \(R(x_1, x_2, \dots) = (0, x_1, x_2, \dots) = L^*\)：单射（核 0 维），但值域是第 1 个分量为 0 的闭超平面（余维 1），所以 \(R \in \mathcal{F}(\ell^2)\) 且 \textbf{\(\operatorname{ind} R = 0 - 1 = -1\)}！
    \item
      进一步：\(\operatorname{ind}(L^n) = n\)，\(\operatorname{ind}(R^n) = -n\)。\emph{(注：有限维方阵的秩-零度定理强制要求 \(\dim \operatorname{Ker} A - \operatorname{codim} \mathcal{R}(A) = 0\)，只有无穷维才会出现非零指标！)}
    \end{itemize}
  \end{enumerate}
\end{itemize}

\SubSubHeading{2. Fredholm 算子的三大主定理（定理 5.13--5.15，教材第 199--203 页）}

\begin{enumerate}
\def\labelenumi{\arabic{enumi}.}
\tightlist
\item
  \textbf{定理 5.13（Atkinson 定理 --- 判定 Fredholm 算子的超级神器！）}：

  \begin{itemize}
  \tightlist
  \item
    \textbf{(1) 必要性}：若 \(A \in \mathcal{F}(X, Y)\)，则存在 \(S \in \mathcal{L}(Y, X)\)（称为 \(A\) 的正则化子/拟逆）及紧算子 \(A_1 \in \mathcal{C}(X), \, A_2 \in \mathcal{C}(Y)\)，使得 
\FitDisplayMath{SA = I_X - A_1, \qquad AS = I_Y - A_2.}
 \emph{(注记 5.4：此处的 \(S\) 自动也满足 \(S \in \mathcal{F}(Y, X)\)！)}
  \item
    \textbf{(2) 充分性}：若 \(A \in \mathcal{L}(X, Y)\)，且存在 \(S_1, S_2 \in \mathcal{L}(Y, X)\) 及紧算子 \(A_1 \in \mathcal{C}(X), \, A_2 \in \mathcal{C}(Y)\)，使得 
\FitDisplayMath{S_1 A = I_X - A_1, \qquad A S_2 = I_Y - A_2,}
 则 \(A \in \mathcal{F}(X, Y)\)。
  \item
    \textbf{一句话核心思想（注记 5.5）}：\textbf{Fredholm 算子 \(\iff\) ``模掉紧算子之后可逆''！}（即在商代数 \(\mathcal{L}(X)/\mathcal{C}(X)\) 中可逆。）只要能凑出``乘完等于 \(I\) 减紧算子''，就可以完全绕开繁琐的``核维数 + 余维数 + 值域闭''逐项检验！\emph{(教材勘误：第 200 页定理 5.13(2) 末尾把``则 \(A \in \mathcal{F}(X, Y)\)''误印成了 \(T\)。)}
  \end{itemize}
\item
  \textbf{定理 5.14（乘积指标公式 --- 对数法则）}： 设 \(X, Y, Z\) 为 Banach 空间，\(A_1 \in \mathcal{F}(X, Y), \, A_2 \in \mathcal{F}(Y, Z)\)，则 \(A_2 A_1 \in \mathcal{F}(X, Z)\) 且 
\FitDisplayMath{\operatorname{ind}(A_2 A_1) = \operatorname{ind} A_2 + \operatorname{ind} A_1.}

\item
  \textbf{定理 5.15（小扰动稳定性）}： 设 \(A \in \mathcal{F}(X, Y)\)，则存在 \(\varepsilon > 0\)，使得对一切满足 \(\|S\| < \varepsilon\) 的 \(S \in \mathcal{L}(X, Y)\)，均有 \(A + S \in \mathcal{F}(X, Y)\) 且 
\FitDisplayMath{\operatorname{ind}(A + S) = \operatorname{ind} A.}
 \emph{(直觉：Fredholm 算子集在 \(\mathcal{L}(X, Y)\) 中是开集；指标是整数值连续函数，微小扰动掰不动它！)}
\item
  \textbf{紧扰动不变性（即\hyperlink{prob:5-20}{\textcolor{LinkColor}{习题 5 第 20 题}}）}： 设 \(B \in \mathcal{F}(X, Y)\)，\(A \in \mathcal{C}(X, Y)\)，则 \(B + A \in \mathcal{F}(X, Y)\) 且 \(\operatorname{ind}(B + A) = \operatorname{ind} B\)。\emph{(不管紧扰动 \(A\) 的范数有多大，都不改变 Fredholm 性与指标！)}
\end{enumerate}

\BookPart{}{二、第 5 章课后作业逐题详解（第 5、6、14、15、20、22 题）}{二、第 5 章课后作业逐题详解（第 5、6、14、15、20、22 题）}

\SoftDivider

\BookProblem{prob:5-5}{习题 5 第 5 题}{教材第 204 页}{习题 5 第 5 题}

\SubPillPrimary{【原题】}

\begin{quote}
\begin{enumerate}
\def\labelenumi{\arabic{enumi}.}
\setcounter{enumi}{4}
\tightlist
\item
  设 \(X\) 为无穷维 Banach 空间，\(A \in \mathcal{C}(X)\)，证明：\(A\) 没有有界逆。
\end{enumerate}
\end{quote}

\SubPillAmber{【零基础破题思路】}

\begin{itemize}
\tightlist
\item
  \textbf{第一反应}：题目出现两个黄金关键词——\textbf{``无穷维''} + \textbf{``紧算子 \(A\)''} + \textbf{``没有有界逆''}。只要在无穷维空间里证明一个紧算子做不到某件事，\textbf{99\% 是用反证法导出恒等算子 \(I_X\) 是紧算子，再用 Riesz 引理矛盾！}
\item
  \textbf{逻辑链条}：如果 \(A\) 有定义在全空间上的有界逆 \(A^{-1} \in \mathcal{L}(X)\)，那么 \(I_X = A^{-1} A\)。由紧算子的双边理想性质（定理 5.3(ii)），有界算子乘以紧算子还是紧算子，所以 \(I_X \in \mathcal{C}(X)\)。但 \(I_X\) 紧意味着 \(X\) 的闭单位球列紧，这与``\(X\) 是无穷维空间''（Riesz 引理）直接冲突！
\item
  \textbf{更严谨的覆盖}：教材和试卷中``\(A\) 没有有界逆''通常指不存在全空间上的 \(A^{-1} \in \mathcal{L}(X)\)；但有时也指更强的情形——即便 \(A\) 不是满射，仅要求 \(A\) 单射且在其值域 \(\mathcal{R}(A)\) 上有界可逆（即 \(A^{-1}: \mathcal{R}(A) \to X\) 有界，等价于存在下界 \(\|Ax\| \ge c\|x\|\)）。下面我们先给出标准简洁证法，再把``值域上有界逆也不存在''的加强版证法一并补齐，做到滴水不漏！
\end{itemize}

\SubPillMain{【规范解答】}

\textbf{证明（标准版：不存在全空间上的有界逆 \(A^{-1} \in \mathcal{L}(X)\)）：}

用反证法。假设 \(A\) 存在有界逆 \(A^{-1} \in \mathcal{L}(X)\)，则有 
\FitDisplayMath{I_X = A^{-1} A,}
 其中 \(I_X\) 为 \(X\) 上的恒等算子。

因为 \(A \in \mathcal{C}(X)\) 且 \(A^{-1} \in \mathcal{L}(X)\)，由紧算子的理想性质（定理 5.3(ii)）可知 
\FitDisplayMath{I_X = A^{-1} A \in \mathcal{C}(X).}


令 \(B_X = \{x \in X : \|x\| \le 1\}\) 为 \(X\) 中的闭单位球。由于 \(B_X\) 是有界集且 \(I_X\) 是紧算子，根据紧算子定义（定义 5.1），像集 \(I_X(B_X) = B_X\) 是 \(X\) 中的列紧集。

然而，\(X\) 是无穷维赋范线性空间，由 Riesz 引理可知 \(X\) 的闭单位球 \(B_X\) 不是列紧集（否则可利用 Riesz 引理递推构造出单位向量列 \(\{x_n\}_{n=1}^\infty \subset B_X\) 满足 \(\|x_m - x_n\| \ge \frac{1}{2}\) 对一切 \(m \neq n\)，它没有任何收敛子列），矛盾！

因此，\(A\) 没有有界逆。 \(\blacksquare\)

\SoftDivider

\textbf{补充证明（加强版：即便不要求 \(A\) 满射，也不存在定义在 \(\mathcal{R}(A)\) 上的有界逆 \(A^{-1}: \mathcal{R}(A) \to X\)）：}

假设 \(A\) 单射且 \(A^{-1}: \mathcal{R}(A) \to X\) 是有界线性算子，记 \(M = \|A^{-1}\| > 0\)，则对一切 \(x \in X\) 均有 
\FitDisplayMath{\|x\| = \|A^{-1}(Ax)\| \le M \|Ax\|.}


任取 \(X\) 中的有界序列 \(\{x_n\}_{n=1}^\infty\)。因为 \(A \in \mathcal{C}(X)\)，所以像序列 \(\{Ax_n\}_{n=1}^\infty\) 在 \(X\) 中必有收敛子列 \(\{Ax_{n_k}\}_{k=1}^\infty\)。特别地，\(\{Ax_{n_k}\}\) 是 Cauchy 列。于是对任意 \(j, k \in \mathbb{N}\) 有 
\FitDisplayMath{\|x_{n_k} - x_{n_j}\| \le M \|A(x_{n_k} - x_{n_j})\| = M \|Ax_{n_k} - Ax_{n_j}\| \to 0 \qquad (j, k \to \infty).}


由此可见 \(\{x_{n_k}\}_{k=1}^\infty\) 是 \(X\) 中的 Cauchy 列。由于 \(X\) 是 Banach 空间（完备），\(\{x_{n_k}\}\) 在 \(X\) 中收敛。这表明 \(X\) 中任意有界序列都有收敛子列，即闭单位球 \(B_X\) 是列紧集，与 \(\dim X = \infty\)（Riesz 引理）矛盾！故即使仅限定在 \(\mathcal{R}(A)\) 上，\(A\) 也没有有界逆。 \(\blacksquare\)

\SubPillCoral{【考点复盘与避坑提示】}

\begin{enumerate}
\def\labelenumi{\arabic{enumi}.}
\tightlist
\item
  \textbf{与定理 5.5 的关系}：定理 5.5 证明了``单射紧算子不可能是满射''（实际上同样是通过反设满射 \(\Rightarrow\) Banach 逆算子定理得 \(A^{-1} \in \mathcal{L}(Y, X)\) \(\Rightarrow\) \(I_X = A^{-1}A\) 紧，得出矛盾）。本题直接从定义 + 定理 5.3 + Riesz 引理出发，自包含且最稳。
\item
  \textbf{加强版的隐藏考点}：在证明加强版时，\textbf{必须用到 \(X\) 的完备性}（从 Cauchy 列推出收敛子列）；如果只证标准版，只需 \(X\) 是无穷维赋范线性空间即可。
\item
  \textbf{常见误区}：注意题干给的是``无穷维''！有限维空间中每个线性算子都是紧算子，且可逆矩阵显然有有界逆，所以千万别漏写``与 \(\dim X = \infty\)（Riesz 引理）矛盾''这一步！
\end{enumerate}

\SoftDivider

\BookProblem{prob:5-6}{习题 5 第 6 题【重点排雷题！】}{教材第 204 页}{习题 5 第 6 题【重点排雷题！】}

\SubPillPrimary{【原题】}

\begin{quote}
\begin{enumerate}
\def\labelenumi{\arabic{enumi}.}
\setcounter{enumi}{5}
\tightlist
\item
  设 \(X\) 为无穷维 Banach 空间. 证明：\(A \in \mathcal{C}(X)\) 当且仅当 \(\mathcal{R}(A)\) 不包含无穷维的子空间.
\end{enumerate}
\end{quote}

\begin{quote}
\textbf{\WarnBadge{} 极其重要的严谨性排雷（必看！）}： 原题表述存在 \textbf{两处实质性的数学瑕疵}，在理解和作答时必须心中有数：

\begin{enumerate}
\def\labelenumi{\arabic{enumi}.}
\tightlist
\item
  \textbf{漏写了关键限定词``闭''}：如果不要求子空间是\textbf{闭}的，那么连必要性 \((\Rightarrow)\) 都是\textbf{假命题}！

  \begin{itemize}
  \tightlist
  \item
    \textbf{反例}：取 \(X = \ell^2\)，定义紧算子 \(A(x_1, x_2, \dots) = \left(x_1, \frac{x_2}{2}, \dots, \frac{x_n}{n}, \dots\right)\)（显然是 Hilbert-Schmidt 算子，故 \(A \in \mathcal{C}(\ell^2)\)）。对每个标准基向量 \(e_n\)，都有 \(e_n = A(ne_n) \in \mathcal{R}(A)\)，于是所有有限非零项序列构成的子空间 \(c_{00} = \operatorname{span}\{e_1, e_2, \dots\} \subset \mathcal{R}(A)\)，它显然是\textbf{无穷维}线性子空间（只是在 \(\ell^2\) 中不闭）。
  \item
    \textbf{教材自身佐证}：同书下一页（第 205 页）习题 17(ii) 就正确地表述为``如果\textbf{闭}线性子空间 \(M \subset \mathcal{R}(A)\)，则 \(\dim M < +\infty\)''。
  \end{itemize}
\item
  \textbf{补上``闭''字之后，充分性 \((\Leftarrow)\) 仅在 Hilbert 空间上成立，在一般 Banach 空间上仍不成立！}

  \begin{itemize}
  \tightlist
  \item
    \textbf{必要性 \((\Rightarrow)\)}（若 \(A \in \mathcal{C}(X)\)，则 \(\mathcal{R}(A)\) 不包含无穷维\textbf{闭}子空间）：对所有 Banach 空间\textbf{完全成立}！
  \item
    \textbf{充分性 \((\Leftarrow)\)}（若 \(A \in \mathcal{L}(X)\) 且 \(\mathcal{R}(A)\) 不包含无穷维闭子空间，则 \(A \in \mathcal{C}(X)\)）：

    \begin{itemize}
    \tightlist
    \item
      在 \textbf{Hilbert 空间} 上\textbf{成立}；
    \item
      在 \textbf{一般 Banach 空间} 上\textbf{不成立}！（满足``值域不含无穷维闭子空间''的算子在泛函分析中称为``严格奇异算子''，范围严格大于紧算子；下文给出使用教材第 183 页注记 5.1 即可完全看懂的标准反例！）
    \end{itemize}
  \end{itemize}
\end{enumerate}
\end{quote}

\SubPillAmber{【零基础破题思路】}

\begin{enumerate}
\def\labelenumi{\arabic{enumi}.}
\tightlist
\item
  \textbf{怎么证必要性 \((\Rightarrow)\)？（核心必考，对应习题 17(ii)！）}

  \begin{itemize}
  \tightlist
  \item
    假设存在一个\textbf{闭}子空间 \(M \subset \mathcal{R}(A)\)，我们的目标是证 \(\dim M < \infty\)，也就是证 \textbf{\(M\) 的某个闭球是列紧集}（再靠 Riesz 引理秒杀！）。
  \item
    难点在于：\(A\) 未必是单射，所以不能直接写 \(A^{-1}\)！怎么把 \(A\) 变成从某个 Banach 空间到 \(M\) 的满射/双射？
  \item
    \textbf{核心桥梁}：把 \(A\) 的定义域收缩到原像 \(X_0 = A^{-1}(M) = \{x \in X : Ax \in M\}\) 上！因为 \(A\) 连续、\(M\) 闭，所以 \(X_0\) 也是闭子空间（即 Banach 空间），且 \(A|_{X_0}: X_0 \to M\) 是\textbf{连续满射}！
  \item
    接下来有两条康庄大道：

    \begin{itemize}
    \tightlist
    \item
      \textbf{路子 A（开映射定理，4 行秒杀）}：\(A|_{X_0}: X_0 \to M\) 是 Banach 空间间的满射 \(\implies\) 它把 \(X_0\) 的单位开球映成的集合包含 \(M\) 里的某个开球 \(B_M(0, r)\) \(\implies\) \(B_M(0, r)\) 被包在一个列紧集 \(A(B_{X_0}(0, 1))\) 里 \(\implies\) \(M\) 的闭球列紧 \(\implies\) \(\dim M < \infty\)！
    \item
      \textbf{路子 B（商空间 + Banach 逆算子定理，与教材风格一致）}：模掉 \(\operatorname{Ker} A\)，得到双射 \(\tilde{A}: X_0 / \operatorname{Ker} A \to M\)，由逆算子定理得 \(\tilde{A}^{-1}\) 有界，再对 \(M\) 的单位球里的任意序列利用商范数拉回 \(X_0\) 中找有界代表元，借 \(A\) 的紧性抽出收敛子列 \(\implies\) \(\dim M < \infty\)！
    \end{itemize}
  \end{itemize}
\item
  \textbf{怎么理解充分性 \((\Leftarrow)\) 的一般 Banach 空间反例？}

  \begin{itemize}
  \tightlist
  \item
    我们需要一个``不是紧算子、但值域里装不下任何无穷维闭子空间''的算子。利用两个``几何性质格格不入''的空间之间的恒等嵌入 \(J: \ell^1 \hookrightarrow \ell^2\)：

    \begin{itemize}
    \tightlist
    \item
      \(J\) 不紧（\(Je_n = e_n\) 间距恒为 \(\sqrt{2}\)）；
    \item
      如果 \(J(\ell^1)\) 里藏着一个无穷维闭子空间 \(M_2 \subset \ell^2\)，那 \(M_2\) 就会通过 \(J^{-1}\) 拓扑同构于 \(\ell^1\) 的一个闭子空间 \(N\)——可 \(M_2\) 里有``弱收敛于 \(0\) 但不强收敛''的正交列，而 \(\ell^1\) 里弱收敛等价于强收敛（Schur 性质，教材第 183 页注记 5.1），直接爆出矛盾！
    \end{itemize}
  \item
    为凑成 \(X \to X\) 的算子，只需把两个空间拼成直和 \(X = \ell^1 \oplus \ell^2\)，令 \(A(x, y) = (0, Jx)\) 即可！
  \end{itemize}
\end{enumerate}

\SubPillMain{【规范解答】}

\textbf{一、修正原题为``不含无穷维闭子空间''后的必要性 \((\Rightarrow)\) 严格证明}

\begin{quote}
\textbf{待证命题}：设 \(X\) 为 Banach 空间，\(A \in \mathcal{C}(X)\)。若 \(M\) 是 \(X\) 的\textbf{闭}线性子空间且 \(M \subset \mathcal{R}(A)\)，则 \(\dim M < +\infty\)。
\end{quote}

\textbf{证法一（商空间 + Banach 逆算子定理，最贴合教材 5.1/5.3 节风格）：}

\begin{enumerate}
\def\labelenumi{\arabic{enumi}.}
\tightlist
\item
  \textbf{构造原像子空间}：令 
\FitDisplayMath{X_0 = A^{-1}(M) = \{x \in X : Ax \in M\}.}
 因为 \(A \in \mathcal{L}(X)\) 连续且 \(M\) 是闭子空间，所以 \(X_0\) 是 \(X\) 的闭线性子空间，从而 \(X_0\) 与 \(M\) 均为 Banach 空间。
\item
  \textbf{构造商空间上的双射}：显然 \(\operatorname{Ker} A \subset X_0\) 且 \(\operatorname{Ker} A\) 是 \(X_0\) 的闭线性子空间。作商空间 \(\tilde{X} = X_0 / \operatorname{Ker} A\)，赋予商范数 
\FitDisplayMath{\|\tilde{x}\| = \inf_{z \in \operatorname{Ker} A} \|x + z\|, \qquad \tilde{x} = x + \operatorname{Ker} A \in \tilde{X},}
 则 \(\tilde{X}\) 是 Banach 空间。定义映射 \(\tilde{A}: \tilde{X} \to M\) 为 
\FitDisplayMath{\tilde{A}(\tilde{x}) = Ax \qquad (x \in \tilde{x}).}


  \begin{itemize}
  \tightlist
  \item
    \textbf{良定性}：若 \(x_1 - x_2 \in \operatorname{Ker} A\)，则 \(Ax_1 = Ax_2\)，故 \(\tilde{A}\) 与代表元选取无关。
  \item
    \textbf{线性与有界性}：\(\tilde{A}\) 显然线性；对任意 \(x \in \tilde{x}\) 与 \(z \in \operatorname{Ker} A\)，有 \(\|\tilde{A}\tilde{x}\| = \|A(x + z)\| \le \|A\|\|x + z\|\)，对所有 \(z \in \operatorname{Ker} A\) 取下确界得 \(\|\tilde{A}\tilde{x}\| \le \|A\|\|\tilde{x}\|\)，故 \(\tilde{A} \in \mathcal{L}(\tilde{X}, M)\)。
  \item
    \textbf{单射与满射}：若 \(\tilde{A}\tilde{x} = 0\)，则 \(Ax = 0 \implies x \in \operatorname{Ker} A \implies \tilde{x} = \tilde{0}\)，故单射；对任意 \(y \in M \subset \mathcal{R}(A)\)，存在 \(x \in X\) 使 \(Ax = y \in M\)，于是 \(x \in X_0\) 且 \(\tilde{A}\tilde{x} = y\)，故满射。
  \end{itemize}
\item
  \textbf{利用逆算子定理与紧性}：\(\tilde{A}\) 是两个 Banach 空间之间的有界线性双射，由 \textbf{Banach 逆算子定理}，其逆算子 \(\tilde{A}^{-1} \in \mathcal{L}(M, \tilde{X})\) 有界。
\item
  \textbf{证明 \(M\) 的闭单位球列紧}：任取 \(M\) 的闭单位球中的序列 \(\{y_n\}_{n=1}^\infty \subset M\)（\(\|y_n\| \le 1\)）。令 \(\tilde{x}_n = \tilde{A}^{-1} y_n \in \tilde{X}\)，则 
\FitDisplayMath{\|\tilde{x}_n\| \le \|\tilde{A}^{-1}\|\|y_n\| \le \|\tilde{A}^{-1}\|.}
 根据商范数的下确界定义，对每个 \(n \in \mathbb{N}\)，可选取代表元 \(x_n \in \tilde{x}_n \subset X_0\) 使得 
\FitDisplayMath{\|x_n\| \le \|\tilde{x}_n\| + 1 \le \|\tilde{A}^{-1}\| + 1, \qquad \text{且 } Ax_n = \tilde{A}\tilde{x}_n = y_n.}
 由此可见 \(\{x_n\}_{n=1}^\infty\) 是 \(X\) 中的有界序列。因为 \(A \in \mathcal{C}(X)\) 是紧算子，序列 \(\{Ax_n\} = \{y_n\}\) 在 \(X\) 中必有收敛子列 \(\{y_{n_k}\}\)；又因 \(M\) 是闭子空间，该子列的极限仍落在 \(M\) 中。
\item
  \textbf{结论}：\(M\) 的闭单位球是列紧集。由 \textbf{Riesz 引理}，必有 \(\dim M < +\infty\)。 \(\blacksquare\)
\end{enumerate}

\SoftDivider

\textbf{证法二（开映射定理，4 行极速版）：}

令 \(X_0 = A^{-1}(M)\)。因为 \(A\) 连续且 \(M\) 闭，故 \(X_0\) 与 \(M\) 均为 Banach 空间，且限制算子 \(A_0 = A|_{X_0}: X_0 \to M\) 是有界线性满射。记 \(U = \{x \in X_0 : \|x\| < 1\}\) 为 \(X_0\) 的单位开球，由\textbf{开映射定理}，\(A_0(U)\) 是 \(M\) 中包含 \(0\) 的开集，故存在 \(r > 0\) 使 \(M\) 中的闭球 \(B_M(0, r) = \{y \in M : \|y\| \le r\} \subset A_0(U) = A(U)\)。因为 \(U\) 有界且 \(A\) 为紧算子，\(A(U)\) 列紧，其子集 \(B_M(0, r)\) 亦列紧（且因 \(M\) 闭，子列极限必在 \(M\) 中），从而 \(M\) 的闭单位球列紧。由 \textbf{Riesz 引理}，\(\dim M < +\infty\)。 \(\blacksquare\)

\SoftDivider

\textbf{二、充分性 \((\Leftarrow)\) 的深度分析：一般 Banach 空间反例与 Hilbert 空间说明}

\begin{enumerate}
\def\labelenumi{\arabic{enumi}.}
\tightlist
\item
  \textbf{一般 Banach 空间中充分性 \((\Leftarrow)\) 不成立（标准反例）}： 取 Banach 空间 \(X = \ell^1 \oplus \ell^2\)，赋予范数 \(\|(x, y)\| = \|x\|_1 + \|y\|_2\)。记 \(J: \ell^1 \to \ell^2\) 为自然恒等嵌入 \(Jx = x\)（因为 \(\|x\|_2 \le \|x\|_1\)，故 \(J\) 线性有界）。定义 \(A \in \mathcal{L}(X)\) 为 
\FitDisplayMath{A(x, y) = (0, Jx), \qquad (x, y) \in \ell^1 \oplus \ell^2.}


  \begin{itemize}
  \tightlist
  \item
    \textbf{\(A\) 不是紧算子}：取标准单位向量 \(u_n = (e_n, 0) \in X\)，有 \(\|u_n\| = \|e_n\|_1 = 1\)，但 \(Au_n = (0, e_n)\) 满足 \(\|Au_m - Au_n\| = \|e_m - e_n\|_2 = \sqrt{2}\)（\(m \neq n\)），无收敛子列，故 \(A \notin \mathcal{C}(X)\)。
  \item
    \textbf{\(\mathcal{R}(A)\) 不包含无穷维闭子空间}：反设 \(M \subset \mathcal{R}(A) = \{0\} \oplus J(\ell^1)\) 是无穷维闭子空间，将其投影到第二分量得到 \(\ell^2\) 的无穷维闭子空间 \(M_2 \subset J(\ell^1)\)。令 \(N = J^{-1}(M_2) \subset \ell^1\)，则 \(N\) 是 \(\ell^1\) 的闭子空间，且 \(J|_N: N \to M_2\) 为 Banach 空间之间的连续双射，由 Banach 逆算子定理，\(N\) 与 \(M_2\) 拓扑同构。 作为无穷维 Hilbert 空间 \(\ell^2\) 的无穷维闭子空间，\(M_2\) 中存在标准正交列 \(\{g_n\}_{n=1}^\infty\)：由 Bessel 不等式，对任意 \(h \in M_2\) 有 \((g_n, h) \to 0\)，即 \(g_n \rightharpoonup 0\)（弱收敛于 \(0\)），但 \(\|g_n\|_2 = 1\) 不强收敛于 \(0\)。令 \(x_n = (J|_N)^{-1} g_n \in N \subset \ell^1\)，因拓扑同构保持弱收敛与强收敛，故 \(\{x_n\}\) 在 \(\ell^1\) 中弱收敛于 \(0\) 但不强收敛于 \(0\)，这与 \textbf{\(\ell^1\) 的 Schur 性质}（教材第 183 页注记 5.1：\(\ell^1\) 中弱收敛等价于强收敛）矛盾！因此 \(\mathcal{R}(A)\) 不包含无穷维闭子空间。 \(\blacksquare\)
  \end{itemize}
\item
  \textbf{当 \(X = H\) 为 Hilbert 空间时，充分性 \((\Leftarrow)\) 成立（思路简述）}： 设 \(A \in \mathcal{L}(H)\) 且 \(A\) 不是紧算子，则正自伴算子 \(|A| = (A^*A)^{1/2}\) 也不是紧算子（因 \(\|Ax\| = \||A|x\|\)）。利用自伴算子的谱分解（第 6 章），存在 \(\delta > 0\) 使谱投影对应的闭子空间 \(N = E([\delta, \infty))H\) 为无穷维闭子空间，且对一切 \(x \in N\) 满足 \(\|Ax\| = \||A|x\| \ge \delta \|x\|\)。这种有正下界的有界算子把闭子空间 \(N\) 映成闭子空间 \(M = A(N) \subset \mathcal{R}(A)\)，且与 \(N\) 同构，故 \(M\) 是无穷维闭子空间。逆否即得充分性！
\end{enumerate}

\SubPillCoral{【考点复盘与避坑提示】}

\begin{enumerate}
\def\labelenumi{\arabic{enumi}.}
\tightlist
\item
  \textbf{期末考场实战策略（怎么写最稳拿满分？）}：

  \begin{itemize}
  \tightlist
  \item
    如果考卷直接照搬原题，先简明写一句：``注：当 \(A \in \mathcal{C}(X)\) 时，若不要求子空间闭，取 \(\ell^2\) 上 \(Ae_n = \frac{1}{n}e_n\) 其值域含无穷维非闭子空间 \(c_{00}\)，故题意应理解为不含无穷维\textbf{闭}子空间。''
  \item
    然后把上面的 \textbf{必要性 \((\Rightarrow)\) 完整证明（证法一或证法二）} 漂漂亮亮写出来——这正是出题人核心想考的知识点（也是习题 17(ii)）。
  \item
    最后用 3--4 行注明：对于充分性 \((\Leftarrow)\)，在 Hilbert 空间情形下成立，而在一般 Banach 空间下不成立（简述 \(X = \ell^1 \oplus \ell^2, \, A(x, y) = (0, Jx)\) 反例）。这样无论阅卷老师按哪套标准批改，都能稳拿满分！
  \end{itemize}
\item
  \textbf{高频副产品}：证法一（商空间 + 逆算子定理）和证法二（开映射定理）直接秒杀同页的 \textbf{习题 17(ii)}（只需把 \(A \in \mathcal{C}(X)\) 换成 \(A \in \mathcal{C}(X, Y)\)，字都不用改！）。
\end{enumerate}

\SoftDivider

\BookProblem{prob:5-14}{习题 5 第 14 题}{教材第 205 页}{习题 5 第 14 题}

\SubPillPrimary{【原题】}

\begin{quote}
\begin{enumerate}
\def\labelenumi{\arabic{enumi}.}
\setcounter{enumi}{13}
\tightlist
\item
  定义 \(C([0,1])\) 上的算子 \(A\) 为 
\FitDisplayMath{Af(x) = \int_0^{1-x} f(t)\mathrm{d}t, \qquad x \in [0,1], \quad f \in C([0,1]).}
 证明：\(A\) 是紧算子。
\end{enumerate}
\end{quote}

\SubPillAmber{【零基础破题思路】}

\begin{itemize}
\tightlist
\item
  \textbf{第一反应}：空间是 \(C([0,1])\)，算子是积分算子 \(\implies\) 立刻想到 \textbf{Arzelà-Ascoli 定理}！
\item
  \textbf{千万别踩坑}：能不能直接说``由教材例 5.1 知积分算子是紧算子''？\textbf{绝对不行！} 因为如果把积分上限 \(1-x\) 强行改写成 \(\int_0^1 K(x, t) f(t)\mathrm{d}t\)，核函数 \(K(x, t) = \begin{cases} 1, & 0 \le t \le 1-x \\ 0, & 1-x < t \le 1 \end{cases}\) 在对角线 \(x + t = 1\) 上是\textbf{间断的}，不满足例 5.1 要求的``\(K\) 连续''！所以必须亲自老老实实用 Arzelà-Ascoli 定理检验。
\item
  \textbf{三步走战略}：

  \begin{enumerate}
  \def\labelenumi{\arabic{enumi}.}
  \tightlist
  \item
    确认 \(A \in \mathcal{L}(C([0,1]))\)（把连续函数映成连续函数、线性、有界）；
  \item
    对任意有界集 \(B = \{f \in C([0,1]) : \|f\|_\infty \le M\}\)，证 \(A(B)\) \textbf{一致有界}；
  \item
    估计 \(|Af(x_1) - Af(x_2)|\)，得到 Lipschitz 估计 \(\le M|x_1 - x_2|\)，立刻推出 \textbf{等度连续}！
  \end{enumerate}
\end{itemize}

\SubPillMain{【规范解答】}

记 \(C([0,1])\) 上的范数为一致范数 \(\|f\|_\infty = \max_{x \in [0,1]} |f(x)|\)。

\textbf{第一步：验证 \(A \in \mathcal{L}(C([0,1]))\)。}

\begin{itemize}
\tightlist
\item
  \textbf{良定性（值域落在 \(C([0,1])\) 中）}：对任意 \(f \in C([0,1])\)，变上限积分 \(g(u) = \int_0^u f(t)\mathrm{d}t\) 在 \([0,1]\) 上连续（事实上连续可微），复合连续函数 \(u = 1 - x\) 后，\(Af(x) = g(1 - x) \in C([0,1])\)。
\item
  \textbf{线性性}：对任意 \(f, g \in C([0,1])\) 及常数 \(\alpha, \beta\)，由定积分的线性性有 
\FitDisplayMath{A(\alpha f + \beta g)(x) = \int_0^{1-x} (\alpha f(t) + \beta g(t))\mathrm{d}t = \alpha Af(x) + \beta Ag(x).}

\item
  \textbf{有界性}：对任意 \(f \in C([0,1])\) 和任意 \(x \in [0,1]\)，有 
\FitDisplayMath{|Af(x)| = \left|\int_0^{1-x} f(t)\mathrm{d}t\right| \le \int_0^{1-x} |f(t)|\mathrm{d}t \le (1 - x)\|f\|_\infty \le \|f\|_\infty.}
 对 \(x \in [0,1]\) 取最大值，得 \(\|Af\|_\infty \le \|f\|_\infty\)。故 \(A\) 是 \(C([0,1])\) 上的有界线性算子（且 \(\|A\| \le 1\)）。
\end{itemize}

\textbf{第二步：利用 Arzelà-Ascoli 定理证明 \(A\) 将有界集映为列紧集。}

任取 \(C([0,1])\) 中的有界子集 \(B\)，则存在常数 \(M > 0\)，使得对一切 \(f \in B\) 均有 \(\|f\|_\infty \le M\)。下面验证像集 \(A(B) = \{Af : f \in B\}\) 在 \([0,1]\) 上一致有界且等度连续：

\begin{enumerate}
\def\labelenumi{\arabic{enumi}.}
\tightlist
\item
  \textbf{一致有界性}：对任意 \(f \in B\) 及任意 \(x \in [0,1]\)，由第一步的不等式立得 
\FitDisplayMath{|Af(x)| \le \|f\|_\infty \le M.}
 故 \(A(B)\) 在 \([0,1]\) 上一致有界。
\item
  \textbf{等度连续性}：任取 \(f \in B\) 及 \(x_1, x_2 \in [0,1]\)。不妨设 \(x_1 \le x_2\)（若 \(x_1 > x_2\)，交换二者角色即可），则 \(0 \le 1 - x_2 \le 1 - x_1 \le 1\)，从而 
\FitDisplayMath{|Af(x_1) - Af(x_2)| = \left|\int_0^{1-x_1} f(t)\mathrm{d}t - \int_0^{1-x_2} f(t)\mathrm{d}t\right| = \left|\int_{1-x_2}^{1-x_1} f(t)\mathrm{d}t\right| \le \int_{1-x_2}^{1-x_1} |f(t)|\mathrm{d}t \le M \bigl((1 - x_1) - (1 - x_2)\bigr) = M|x_1 - x_2|.}
 于是对任意 \(\varepsilon > 0\)，取 \(\delta = \frac{\varepsilon}{M} > 0\)（注意 \(\delta\) 只依赖于 \(\varepsilon\) 和 \(M\)，与 \(f \in B\) 以及点 \(x_1, x_2\) 均无关）。当 \(|x_1 - x_2| < \delta\) 时，对一切 \(f \in B\) 统一有 
\FitDisplayMath{|Af(x_1) - Af(x_2)| \le M|x_1 - x_2| < M \cdot \frac{\varepsilon}{M} = \varepsilon.}
 故 \(A(B)\) 在 \([0,1]\) 上等度连续。
\end{enumerate}

由 \textbf{Arzelà-Ascoli 定理}，像集 \(A(B)\) 是 \(C([0,1])\) 中的列紧集。根据定义 5.1，\(A\) 是 \(C([0,1])\) 上的紧算子。 \(\blacksquare\)

\SubPillCoral{【考点复盘与避坑提示】}

\begin{enumerate}
\def\labelenumi{\arabic{enumi}.}
\tightlist
\item
  \textbf{等度连续的``一致''二字怎么写才不扣分}：一定要明确点出``\(\delta = \varepsilon / M\) \textbf{与 \(f \in B\) 及 \(x_1, x_2 \in [0,1]\) 无关}''！如果 \(\delta\) 依赖于具体的函数 \(f\)，那只证明了单个函数一致连续，没证明一族函数等度连续。
\item
  \textbf{同类题通杀}：无论是 Volterra 算子 \((Vf)(x) = \int_0^x f(t)\mathrm{d}t\)，还是带连续核的变上限算子 \(\int_0^{\varphi(x)} K(x,t)f(t)\mathrm{d}t\)（\(\varphi\) 连续），统统照搬这个``一致有界 + 等度连续 + Arzelà-Ascoli''模板！
\item
  \textbf{细节规范}：证明 \(A\) 是紧算子时，别忘了先花两行交代 \(A\) 是 \(C([0,1])\) 上的有界线性算子（良定性 + 线性 + 有界）。
\end{enumerate}

\SoftDivider

\BookProblem{prob:5-15}{习题 5 第 15 题}{教材第 205 页}{习题 5 第 15 题}

\SubPillPrimary{【原题】}

\begin{quote}
\begin{enumerate}
\def\labelenumi{\arabic{enumi}.}
\setcounter{enumi}{14}
\tightlist
\item
  设 \(H, G\) 为 Hilbert 空间，\(A \in \mathcal{L}(H, G)\). 证明：如果 \(A\) 是有限秩算子，则 \(A\) 必是 Hilbert-Schmidt 算子.
\end{enumerate}
\end{quote}

\SubPillAmber{【零基础破题思路】}

\begin{itemize}
\tightlist
\item
  \textbf{回忆定义 5.4}：要证 \(A\) 是 Hilbert-Schmidt 算子，只需找一组 \(H\) 的标准正交基 \(\{e_\alpha\}\)（可分情形即 \(\{e_n\}_{n=1}^\infty\)），验证 \(\sum_\alpha \|Ae_\alpha\|^2 < +\infty\) 即可！（且由教材式 (5.8)，只要对\textbf{某一组}标准正交基成立，对所有标准正交基都成立且和相等！）
\item
  \textbf{灵光一现（证法一：挑一组最聪明的基！）}：

  \begin{itemize}
  \tightlist
  \item
    \(A\) 是有限秩算子，意思是它的像 \(\mathcal{R}(A)\) 只有有限维（设为 \(m\) 维）——这意味着空间 \(H\) 里绝大部分方向都被 \(A\) 拍成了 \(0\)（落进了 \(\operatorname{Ker} A\)）！
  \item
    用 Hilbert 空间的正交分解：\(H = \operatorname{Ker} A \oplus (\operatorname{Ker} A)^\perp\)。
  \item
    在 \((\operatorname{Ker} A)^\perp\) 上 \(A\) 是单射且映满 \(\mathcal{R}(A)\)，所以 \((\operatorname{Ker} A)^\perp\) 的维数也是有限的 \(m\) 维！
  \item
    于是我们把 \((\operatorname{Ker} A)^\perp\) 的 \(m\) 个标准正交基向量，与 \(\operatorname{Ker} A\) 的标准正交基拼成 \(H\) 的标准正交基——算 \(\sum \|Ae_\alpha\|^2\) 时，落在 \(\operatorname{Ker} A\) 里的基向量全被 \(A\) 映成 \(0\)，无穷级数瞬间坍缩成只有 \(m\) 项的有限和，当然有限！
  \end{itemize}
\item
  \textbf{另一条路（证法二：对任意标准正交基直接算，用到 \(A^*\) 和 Parseval！）}：

  \begin{itemize}
  \tightlist
  \item
    在有限维空间 \(\mathcal{R}(A)\) 里取标准正交基 \(\{f_1, \dots, f_m\}\)，把 \(Ax\) 展开成 \(\sum_{j=1}^m (Ax, f_j) f_j = \sum_{j=1}^m (x, A^* f_j) f_j\)，再用 Parseval 等式把对基向量 \(e_\alpha\) 的无穷和一键收成 \(\sum_{j=1}^m \|A^* f_j\|^2 < \infty\)！
  \end{itemize}
\end{itemize}

\SubPillMain{【规范解答】}

\textbf{证法一（正交分解 + 构造标准正交基，最直观简洁）：}

因为 \(A \in \mathcal{L}(H, G)\) 是有限秩算子，故 \(m \triangleq \dim \mathcal{R}(A) < +\infty\)。若 \(m = 0\)，则 \(A = 0\)，显然 \(\|A\|_2 = 0 < +\infty\)，结论成立。下设 \(m \ge 1\)。

\begin{enumerate}
\def\labelenumi{\arabic{enumi}.}
\item
  \textbf{正交分解与维数确定}：由于 \(A\) 连续，其核空间 \(\operatorname{Ker} A\) 是 \(H\) 的闭线性子空间。由 Hilbert 空间的正交分解定理，有 
\FitDisplayMath{H = \operatorname{Ker} A \oplus (\operatorname{Ker} A)^\perp.}
 考察限制映射 \(A_1 = A|_{(\operatorname{Ker} A)^\perp}: (\operatorname{Ker} A)^\perp \to \mathcal{R}(A)\)：

  \begin{itemize}
  \tightlist
  \item
    \textbf{单射}：若 \(x \in (\operatorname{Ker} A)^\perp\) 且 \(Ax = 0\)，则 \(x \in \operatorname{Ker} A \cap (\operatorname{Ker} A)^\perp = \{0\}\)；
  \item
    \textbf{满射}：任取 \(y = Ax \in \mathcal{R}(A)\)，将 \(x\) 正交分解为 \(x = x_1 + x_2\)（\(x_1 \in \operatorname{Ker} A, \, x_2 \in (\operatorname{Ker} A)^\perp\)），则 \(A_1 x_2 = A x_2 = A(x_1 + x_2) = Ax = y\)。
  \end{itemize}

  因此 \(A_1\) 是线性同构，从而 
\FitDisplayMath{\dim (\operatorname{Ker} A)^\perp = \dim \mathcal{R}(A) = m < +\infty.}

\item
  \textbf{构造 \(H\) 的标准正交基}：在有限维子空间 \((\operatorname{Ker} A)^\perp\) 中取一组标准正交基 \(\{e_1, e_2, \dots, e_m\}\)；在闭子空间 \(\operatorname{Ker} A\) 中任取一组标准正交基 \(\{u_\beta\}_{\beta \in \Lambda}\)（按教材 5.2 节可分假设，可编号为 \(\{e_{m+1}, e_{m+2}, \dots\}\)，有限或可数）。 因为 \((\operatorname{Ker} A)^\perp \perp \operatorname{Ker} A\) 且 \(H = \operatorname{Ker} A \oplus (\operatorname{Ker} A)^\perp\)，二者合并 \(\{e_1, \dots, e_m\} \cup \{u_\beta\}_{\beta \in \Lambda}\) 构成 \(H\) 的一组标准正交基。
\item
  \textbf{计算 Hilbert-Schmidt 范数}：由于每个 \(u_\beta \in \operatorname{Ker} A\) 均满足 \(A u_\beta = 0\)，按定义 5.4 有 
\FitDisplayMath{\|A\|_2^2 = \sum_{j=1}^m \|Ae_j\|^2 + \sum_{\beta \in \Lambda} \|Au_\beta\|^2 = \sum_{j=1}^m \|Ae_j\|^2 \le \sum_{j=1}^m \|A\|^2 \|e_j\|^2 = m \|A\|^2 < +\infty.}
 因此 \(\|A\|_2 \le \sqrt{m}\,\|A\| < +\infty\)，即 \(A\) 是 Hilbert-Schmidt 算子。 \(\blacksquare\)
\end{enumerate}

\SoftDivider

\textbf{证法二（任意标准正交基 + 共轭算子与 Parseval 等式）：}

设 \(m = \dim \mathcal{R}(A) < +\infty\)（\(m = 0\) 平凡，下设 \(m \ge 1\)），在有限维子空间 \(\mathcal{R}(A) \subset G\) 中取标准正交基 \(\{f_1, f_2, \dots, f_m\}\)。对任意 \(x \in H\)，因为 \(Ax \in \mathcal{R}(A)\)，由有限维正交展开与共轭算子 \(A^* \in \mathcal{L}(G, H)\) 的定义得 
\FitDisplayMath{Ax = \sum_{j=1}^m (Ax, f_j)_G f_j = \sum_{j=1}^m (x, A^* f_j)_H f_j.}
 由勾股定理，对一切 \(x \in H\) 有 
\FitDisplayMath{\|Ax\|^2 = \sum_{j=1}^m |(x, A^* f_j)_H|^2.}
 任取 \(H\) 的一组标准正交基 \(\{e_\alpha\}_{\alpha \in \Gamma}\)（可分情形即为 \(\{e_n\}_{n=1}^\infty\)），对所有基向量求和并交换有限和与级数顺序，再利用 \textbf{Parseval 等式}： 
\FitDisplayMath{\sum_{\alpha \in \Gamma} \|Ae_\alpha\|^2 = \sum_{\alpha \in \Gamma} \sum_{j=1}^m |(e_\alpha, A^* f_j)_H|^2 = \sum_{j=1}^m \sum_{\alpha \in \Gamma} |(e_\alpha, A^* f_j)_H|^2 = \sum_{j=1}^m \|A^* f_j\|^2 \le m\|A^*\|^2 = m\|A\|^2 < +\infty.}
 由定义 5.4，故 \(A\) 是 Hilbert-Schmidt 算子。 \(\blacksquare\)

\SubPillCoral{【考点复盘与避坑提示】}

\begin{enumerate}
\def\labelenumi{\arabic{enumi}.}
\tightlist
\item
  \textbf{副产品不等式}：两种证法都顺带给出了非常实用的估计式 \(\|A\|_2 \le \sqrt{\operatorname{rank}(A)}\,\|A\|\)（和有限维矩阵的 \(\|A\|_F \le \sqrt{\operatorname{rank}(A)}\,\|A\|_2\) 完全一致！）。
\item
  \textbf{为何不会受``可分/不可分''限制}：题目没明写``可分''，上面两种证法使用指标集 \(\Lambda / \Gamma\)（顺带括号注明可分情形就是 \(\{e_n\}_{n=1}^\infty\)），无论是按教材可分设定还是对一般 Hilbert 空间，都完全无懈可击！
\item
  \textbf{常见误区}：别只说``把 \(\mathcal{R}(A)\) 的基扩充一下''，切记定义 5.4 求和的是\textbf{定义域 \(H\)} 的标准正交基，而不是值域 \(G\) 的标准正交基！
\end{enumerate}

\SoftDivider

\BookProblem{prob:5-20}{习题 5 第 20 题}{教材第 205 页}{习题 5 第 20 题}

\SubPillPrimary{【原题】}

\begin{quote}
\begin{enumerate}
\def\labelenumi{\arabic{enumi}.}
\setcounter{enumi}{19}
\tightlist
\item
  设 \(X, Y\) 为 Banach 空间，\(B \in \mathcal{F}(X, Y)\)，\(A \in \mathcal{C}(X, Y)\). 证明：
\end{enumerate}

\begin{enumerate}
\def\labelenumi{(\roman{enumi})}
\tightlist
\item
  \(B + A \in \mathcal{F}(X, Y)\)；
\item
  \(\operatorname{ind}(T + A) = \operatorname{ind} T\).
\end{enumerate}
\end{quote}

\begin{quote}
\textbf{\WarnBadge{} 教材印刷笔误勘误}：原题第 (ii) 问中的 \(T\) 系原书印刷笔误（把题干里的 \(B\) 误印成了 \(T\)），应为证明 \textbf{\(\operatorname{ind}(B + A) = \operatorname{ind} B\)}。
\end{quote}

\SubPillAmber{【零基础破题思路】}

\begin{itemize}
\tightlist
\item
  \textbf{第 (i) 问（紧扰动不破坏 Fredholm 性）}：

  \begin{itemize}
  \tightlist
  \item
    千万不要傻傻地去按定义 5.5 算 \(\operatorname{Ker}(B + A)\) 和 \(\mathcal{R}(B + A)\)——加法之后核和值域都会彻底乱掉！
  \item
    记住金句：\textbf{``看到证明某算子属于 \(\mathcal{F}(X, Y)\)，首选 Atkinson 定理（定理 5.13）！''}
  \item
    \(B\) 是 Fredholm 算子 \(\implies\) 存在拟逆 \(S\) 使 \(SB = I_X - K_1, \, BS = I_Y - K_2\)。
  \item
    把同一个 \(S\) 乘到 \(B + A\) 上：\(S(B + A) = I_X - (K_1 - SA)\)，\((B + A)S = I_Y - (K_2 - AS)\)。
  \item
    因为紧算子是双边理想（定理 5.3），\(SA\) 和 \(AS\) 还是紧的，紧减紧还是紧！于是一步秒杀 \(B + A \in \mathcal{F}(X, Y)\)！
  \end{itemize}
\item
  \textbf{第 (ii) 问（紧扰动不改变指标）}：

  \begin{itemize}
  \tightlist
  \item
    \textbf{神仙 3 行代数法（最推荐！）}：观察刚刚算出的两条式子——\textbf{同一个 \(S\) 同时充当了 \(B\) 和 \(B + A\) 的左拟逆！}

    \begin{itemize}
    \tightlist
    \item
      由注记 5.4，\(S\) 自己也是 Fredholm 算子；
    \item
      由乘积指标公式（定理 5.14）与恒等减紧算子指标为 \(0\)（例 5.8 / 定理 5.8(4)）： 
\FitDisplayMath{\operatorname{ind} S + \operatorname{ind} B = \operatorname{ind}(SB) = \operatorname{ind}(I_X - K_1) = 0,}
 
\FitDisplayMath{\operatorname{ind} S + \operatorname{ind}(B + A) = \operatorname{ind}(S(B + A)) = \operatorname{ind}(I_X - (K_1 - SA)) = 0.}

    \item
      两式相减，\(\operatorname{ind} S\) 直接抵消，立得 \(\operatorname{ind}(B + A) = \operatorname{ind} B\)！（和教材第 203 页证明定理 5.15 的手法一模一样！）
    \end{itemize}
  \item
    \textbf{经典同伦法（附注备用）}：也可以构造连续路径 \(F(t) = B + tA\)（\(t \in [0, 1]\)），由 (i) 知整条路径都在 \(\mathcal{F}(X, Y)\) 里，由定理 5.15 知整数值函数 \(t \mapsto \operatorname{ind}(B + tA)\) 在连通区间 \([0, 1]\) 上局部常值（连续），故必为常数，从而 \(t=1\) 与 \(t=0\) 处相等！
  \end{itemize}
\end{itemize}

\SubPillMain{【规范解答】}

\emph{(注：第 (ii) 问中 \(T\) 应为题干中的 \(B\)，以下证 \(\operatorname{ind}(B + A) = \operatorname{ind} B\)。)}

\textbf{(i) 证明 \(B + A \in \mathcal{F}(X, Y)\)：}

因为 \(B \in \mathcal{F}(X, Y)\)，由 \textbf{Atkinson 定理（定理 5.13(1)）}，存在算子 \(S \in \mathcal{L}(Y, X)\) 以及紧算子 \(K_1 \in \mathcal{C}(X), \, K_2 \in \mathcal{C}(Y)\)，使得 
\FitDisplayMath{SB = I_X - K_1, \qquad BS = I_Y - K_2.}


将 \(S\) 分别左乘、右乘 \(B + A\)，得 
\FitDisplayMath{S(B + A) = SB + SA = I_X - (K_1 - SA),}
 
\FitDisplayMath{(B + A)S = BS + AS = I_Y - (K_2 - AS).}


由于 \(A \in \mathcal{C}(X, Y)\)，\(S \in \mathcal{L}(Y, X)\)，由紧算子的双边理想性质及线性性（定理 5.3），有 \(SA \in \mathcal{C}(X)\)、\(AS \in \mathcal{C}(Y)\)，从而 
\FitDisplayMath{\tilde{K}_1 \triangleq K_1 - SA \in \mathcal{C}(X), \qquad \tilde{K}_2 \triangleq K_2 - AS \in \mathcal{C}(Y).}


再次应用 \textbf{Atkinson 定理（定理 5.13(2)）}，立即得到 \(B + A \in \mathcal{F}(X, Y)\)。

\textbf{(ii) 证明 \(\operatorname{ind}(B + A) = \operatorname{ind} B\)：}

\textbf{证法一（Atkinson 定理 + 乘积指标公式，最简洁推荐）：}

承接 (i) 中的算子 \(S \in \mathcal{L}(Y, X)\)。由于 \(BS = I_Y - K_2\) 且 \(SB = I_X - K_1\)，将 \(B\) 视为 \(S\) 的正则化子，由定理 5.13(2)（即注记 5.4）可知 \(S \in \mathcal{F}(Y, X)\)。

另一方面，由定理 5.6 与定理 5.8(4)（即教材例 5.8），对任意紧算子 \(K \in \mathcal{C}(X)\)，均有 \(I_X - K \in \mathcal{F}(X)\) 且 \(\operatorname{ind}(I_X - K) = 0\)。

分别对乘积 \(SB\) 与 \(S(B + A)\) 应用 \textbf{Fredholm 算子乘积指标公式（定理 5.14）}，可得：、 
\FitDisplayMath{\operatorname{ind} S + \operatorname{ind} B = \operatorname{ind}(SB) = \operatorname{ind}(I_X - K_1) = 0,}
 
\FitDisplayMath{\operatorname{ind} S + \operatorname{ind}(B + A) = \operatorname{ind}\bigl(S(B + A)\bigr) = \operatorname{ind}(I_X - \tilde{K}_1) = 0.}


两式相减（或直接由两式均给出 \(-\operatorname{ind} S\)），即得 
\FitDisplayMath{\operatorname{ind}(B + A) = \operatorname{ind} B. \qquad \blacksquare}


\SoftDivider

\textbf{证法二（同伦路径 + 小扰动稳定性，几何直觉法）：}

对每个 \(t \in [0, 1]\)，令 \(F(t) = B + tA\)。因为 \(tA \in \mathcal{C}(X, Y)\)，由 (i) 知 \(F(t) \in \mathcal{F}(X, Y)\)。定义整数值函数 
\FitDisplayMath{\varphi(t) = \operatorname{ind}(B + tA), \qquad t \in [0, 1].}


任取 \(t_0 \in [0, 1]\)。由 \textbf{定理 5.15（小扰动指标不变性）}，存在 \(\varepsilon > 0\)，使得对任意满足 \(\|S\| < \varepsilon\) 的 \(S \in \mathcal{L}(X, Y)\)，都有 \(\operatorname{ind}(F(t_0) + S) = \operatorname{ind} F(t_0)\)。取 \(\delta > 0\) 满足 \(\delta \|A\| < \varepsilon\)（若 \(A = 0\) 则结论平凡，下设 \(\|A\| > 0\)，取 \(\delta = \frac{\varepsilon}{\|A\|}\)）。于是当 \(t \in [0, 1]\) 且 \(|t - t_0| < \delta\) 时，有 
\FitDisplayMath{\|F(t) - F(t_0)\| = |t - t_0| \cdot \|A\| < \varepsilon,}
 从而 \(\varphi(t) = \operatorname{ind}\bigl(F(t_0) + (F(t) - F(t_0))\bigr) = \operatorname{ind} F(t_0) = \varphi(t_0)\)。

这表明 \(\varphi: [0, 1] \to \mathbb{Z}\) 是连通闭区间 \([0, 1]\) 上的连续（局部常值）整数值函数，因此 \(\varphi\) 在 \([0, 1]\) 上恒为常数。特别地， 
\FitDisplayMath{\operatorname{ind}(B + A) = \varphi(1) = \varphi(0) = \operatorname{ind} B. \qquad \blacksquare}


\SubPillCoral{【考点复盘与避坑提示】}

\begin{enumerate}
\def\labelenumi{\arabic{enumi}.}
\tightlist
\item
  \textbf{两个超级高频定理的黄金组合}：证明指标相等时，\textbf{``找公共正则化子 \(S\) + 定理 5.14 + \(\operatorname{ind}(I - K) = 0\)''} 几乎是最强套路！不仅秒杀本题，也是教材证明定理 5.15 的核心。
\item
  \textbf{别忘交代 \(S \in \mathcal{F}(Y, X)\)}：在证法一中，使用定理 5.14（\(\operatorname{ind}(SB) = \operatorname{ind} S + \operatorname{ind} B\)）的前提是参与相乘的两个算子\textbf{都}必须是 Fredholm 算子！所以务必明确写上一句``由注记 5.4 / 定理 5.13(2) 知 \(S \in \mathcal{F}(Y, X)\)''，避免逻辑跳步被扣细节分。
\item
  \textbf{直觉升华}：定理 5.15 说``范数很小的任意扰动不改指标''，本题说``任意大的紧扰动不改指标''——二者合起来表明，Fredholm 指标是一个极其稳健的\textbf{拓扑不变量}！
\end{enumerate}

\SoftDivider

\BookProblem{prob:5-22}{习题 5 第 22 题}{教材第 205 页}{习题 5 第 22 题}

\SubPillPrimary{【原题】}

\begin{quote}
\begin{enumerate}
\def\labelenumi{\arabic{enumi}.}
\setcounter{enumi}{21}
\tightlist
\item
  设 \(X, Y, Z\) 均为 Banach 空间，\(A \in \mathcal{L}(X, Y)\)，\(B \in \mathcal{L}(Y, Z)\)，\(BA \in \mathcal{F}(X, Z)\) 且 \(\dim \operatorname{Ker} B < +\infty\). 证明：\(A \in \mathcal{F}(X, Y)\) 且 \(B \in \mathcal{F}(Y, Z)\).
\end{enumerate}
\end{quote}

\SubPillAmber{【零基础破题思路】}

\begin{itemize}
\tightlist
\item
  \textbf{先看条件给谁铺好了路}：

  \begin{itemize}
  \tightlist
  \item
    关于 \(B\)：题目直接给了 \(\dim \operatorname{Ker} B < \infty\)（条件 2 到手！）；而 \(\mathcal{R}(BA) \subset \mathcal{R}(B) \subset Z\)，既然小的 \(\mathcal{R}(BA)\) 在 \(Z\) 中余维有限，那么包着它的更大的 \(\mathcal{R}(B)\) 余维只会更小，必有限（条件 3 到手！）。
  \item
    \textbf{关键卡点（绝对不能漏！）}：还要证明 \textbf{\(\mathcal{R}(B)\) 是闭子空间}（条件 1）！怎么证？因为 \(\mathcal{R}(BA)\) 是闭子空间且余维有限，所以 \(\mathcal{R}(B)\) 只比 \(\mathcal{R}(BA)\) 多了一个\textbf{有限维子空间 \(W\)}，即 \(\mathcal{R}(B) = \mathcal{R}(BA) + W\)！利用经典小引理``\textbf{闭子空间 + 有限维子空间必闭}''（用商映射 2 行就能证完！），立刻补齐最后一块拼图，证得 \textbf{\(B \in \mathcal{F}(Y, Z)\)}！
  \end{itemize}
\item
  \textbf{再攻克 \(A \in \mathcal{F}(X, Y)\)}：

  \begin{itemize}
  \tightlist
  \item
    手里现在握着两个 Fredholm 算子：\(BA \in \mathcal{F}(X, Z)\) 和刚证出的 \(B \in \mathcal{F}(Y, Z)\)！
  \item
    想证 \(A \in \mathcal{F}(X, Y)\)，果断请出 \textbf{Atkinson 定理（定理 5.13）}：

    \begin{itemize}
    \tightlist
    \item
      \(BA\) 有双边拟逆 \(S_1\)：\(S_1 BA = I_X - K_1\)，\(BA S_1 = I_Z - K_2\)；
    \item
      \(B\) 有左拟逆 \(S_2\)：\(S_2 B = I_Y - K_3\)。
    \end{itemize}
  \item
    怎么凑出 \(A\) 的拟逆？看第一个式子 \((S_1 B) A = I_X - K_1\)，很自然猜 \(R = S_1 B\) 就是 \(A\) 的拟逆！左乘已经搞定，那右乘 \(AR = A S_1 B\) 怎么办？
  \item
    \textbf{神奇恒等式（乘法结合律）}：给 \(AR\) 左边乘上 \(S_2 B = I_Y - K_3\)：

    \begin{itemize}
    \tightlist
    \item
      从左边结合：\((S_2 B) AR = (I_Y - K_3) AR = AR - K_3 AR\)；
    \item
      从中间结合：\(S_2 (B A S_1) B = S_2 (I_Z - K_2) B = S_2 B - S_2 K_2 B = I_Y - K_3 - S_2 K_2 B\)！
    \item
      两边一相等，移项立马出现 \(AR = I_Y - (\text{一堆紧算子})\)！由 Atkinson 定理秒杀 \(A \in \mathcal{F}(X, Y)\)！
    \end{itemize}
  \end{itemize}
\end{itemize}

\SubPillMain{【规范解答】}

\textbf{预备小引理（闭子空间与有限维子空间之和必闭）：}

\begin{quote}
设 \(U\) 是 Banach 空间 \(Z\) 的闭线性子空间，\(W \subset Z\) 是有限维线性子空间，则 \(U + W\) 是 \(Z\) 的闭线性子空间。

\emph{简证}：考虑连续线性商映射 \(q: Z \to Z/U\)（\(q(z) = z + U\)）。因 \(U\) 闭，商空间 \(Z/U\) 赋予商范数后是赋范线性空间。像集 \(q(W)\) 是 \(Z/U\) 的有限维线性子空间，而赋范线性空间中的有限维子空间必闭，故 \(q(W)\) 在 \(Z/U\) 中闭。注意到 \(U + W = q^{-1}(q(W))\)，由 \(q\) 的连续性可知 \(U + W\) 是 \(Z\) 的闭线性子空间。
\end{quote}

\SoftDivider

\textbf{第一步：证明 \(B \in \mathcal{F}(Y, Z)\)。}

按照 Fredholm 算子的定义（定义 5.5），逐一验证以下三条：

\begin{enumerate}
\def\labelenumi{\arabic{enumi}.}
\tightlist
\item
  \textbf{\(\dim \operatorname{Ker} B < +\infty\)}：题设直接给出。
\item
  \textbf{\(\operatorname{codim} \mathcal{R}(B) < +\infty\)}： 对任意 \(x \in X\)，\(BAx = B(Ax) \in \mathcal{R}(B)\)，故 \(\mathcal{R}(BA) \subset \mathcal{R}(B) \subset Z\)。定义商空间之间的自然线性映射 
\FitDisplayMath{\pi: Z / \mathcal{R}(BA) \to Z / \mathcal{R}(B), \qquad \pi\bigl(z + \mathcal{R}(BA)\bigr) = z + \mathcal{R}(B).}
 若 \(z_1 - z_2 \in \mathcal{R}(BA)\)，则 \(z_1 - z_2 \in \mathcal{R}(B)\)，故 \(\pi\) 良定，且显然是线性满射。因为 \(BA \in \mathcal{F}(X, Z)\)，有 \(\dim\bigl(Z / \mathcal{R}(BA)\bigr) < +\infty\)，从而 
\FitDisplayMath{\operatorname{codim} \mathcal{R}(B) = \dim\bigl(Z / \mathcal{R}(B)\bigr) \le \dim\bigl(Z / \mathcal{R}(BA)\bigr) = \operatorname{codim} \mathcal{R}(BA) < +\infty.}

\item
  \textbf{\(\mathcal{R}(B)\) 是 \(Z\) 的闭子空间}： 因为 \(\mathcal{R}(BA) \subset \mathcal{R}(B) \subset Z\)，商空间 \(\mathcal{R}(B) / \mathcal{R}(BA)\) 是有限维空间 \(Z / \mathcal{R}(BA)\) 的线性子空间，故其维数 \(k \triangleq \dim\bigl(\mathcal{R}(B) / \mathcal{R}(BA)\bigr) < +\infty\)。 选取 \(w_1, \dots, w_k \in \mathcal{R}(B)\) 使它们的陪集构成 \(\mathcal{R}(B) / \mathcal{R}(BA)\) 的一组基，并记有限维子空间 \(W = \operatorname{span}\{w_1, \dots, w_k\} \subset \mathcal{R}(B)\)（若 \(k = 0\) 则取 \(W = \{0\}\)），于是有 
\FitDisplayMath{\mathcal{R}(B) = \mathcal{R}(BA) + W.}
 由 \(BA \in \mathcal{F}(X, Z)\) 知 \(\mathcal{R}(BA)\) 是 \(Z\) 的闭子空间。应用上述预备小引理，立即得出 \(\mathcal{R}(B)\) 是 \(Z\) 的闭子空间。
\end{enumerate}

综合 (1)(2)(3)，按定义 5.5 已证得 \textbf{\(B \in \mathcal{F}(Y, Z)\)}。

\SoftDivider

\textbf{第二步：证明 \(A \in \mathcal{F}(X, Y)\)。}

因为 \(BA \in \mathcal{F}(X, Z)\) 且第一步已证 \(B \in \mathcal{F}(Y, Z)\)，分别应用 \textbf{Atkinson 定理（定理 5.13(1)）}：

\begin{itemize}
\tightlist
\item
  存在 \(S_1 \in \mathcal{L}(Z, X)\) 及紧算子 \(K_1 \in \mathcal{C}(X), \, K_2 \in \mathcal{C}(Z)\)，使得 
\FitDisplayMath{S_1 (BA) = I_X - K_1, \qquad (BA) S_1 = I_Z - K_2;}

\item
  存在 \(S_2 \in \mathcal{L}(Z, Y)\) 及紧算子 \(K_3 \in \mathcal{C}(Y)\)，使得 
\FitDisplayMath{S_2 B = I_Y - K_3.}

\end{itemize}

令 \(R \triangleq S_1 B \in \mathcal{L}(Y, X)\)。下面验证 \(R\) 满足定理 5.13(2) 的条件：

\begin{enumerate}
\def\labelenumi{\arabic{enumi}.}
\tightlist
\item
  \textbf{左乘检验}： 
\FitDisplayMath{RA = (S_1 B) A = S_1 (BA) = I_X - K_1, \qquad K_1 \in \mathcal{C}(X).}

\item
  \textbf{右乘检验}：用两种结合方式计算乘积 \(S_2 B A S_1 B \in \mathcal{L}(Y)\)：

  \begin{itemize}
  \tightlist
  \item
    一方面，结合中间的 \(BA S_1 = I_Z - K_2\)，得 
\FitDisplayMath{S_2 (BA S_1) B = S_2 (I_Z - K_2) B = S_2 B - S_2 K_2 B = I_Y - K_3 - S_2 K_2 B;}

  \item
    另一方面，结合左端的 \(S_2 B = I_Y - K_3\) 以及 \(S_1 B = R\)，得 
\FitDisplayMath{(S_2 B) A (S_1 B) = (I_Y - K_3) AR = AR - K_3 AR.}

  \end{itemize}

  令两式相等并移项整理，得 
\FitDisplayMath{AR = I_Y - \bigl(K_3 + S_2 K_2 B - K_3 AR\bigr).}
 由紧算子的双边理想性质与线性性（定理 5.3）：\(K_3 \in \mathcal{C}(Y)\)，\(S_2 K_2 B \in \mathcal{C}(Y)\)，\(K_3 AR \in \mathcal{C}(Y)\)，故 
\FitDisplayMath{\tilde{K} \triangleq K_3 + S_2 K_2 B - K_3 AR \in \mathcal{C}(Y),}
 即有 \(AR = I_Y - \tilde{K}\)。
\end{enumerate}

综上，\(RA = I_X - K_1\) 且 \(AR = I_Y - \tilde{K}\)（其中 \(K_1 \in \mathcal{C}(X), \, \tilde{K} \in \mathcal{C}(Y)\)）。由 \textbf{Atkinson 定理（定理 5.13(2)）}，立即得出 \textbf{\(A \in \mathcal{F}(X, Y)\)}。 \(\blacksquare\)

\SubPillCoral{【考点复盘与避坑提示】}

\begin{enumerate}
\def\labelenumi{\arabic{enumi}.}
\tightlist
\item
  \textbf{最大的扣分黑洞}：第一步证 \(B \in \mathcal{F}(Y, Z)\) 时，很多同学只写了``\(\dim \operatorname{Ker} B < \infty\) 且 \(\operatorname{codim} \mathcal{R}(B) \le \operatorname{codim} \mathcal{R}(BA) < \infty\)''就宣布完毕，\textbf{漏证 \(\mathcal{R}(B)\) 是闭子空间}！请一定把``\(\mathcal{R}(B) = \mathcal{R}(BA) + W\) + 预备小引理''完整写上（教材第 201 页定理 5.13(2) 证明中正使用了完全相同的步骤）。
\item
  \textbf{为什么条件 \(\dim \operatorname{Ker} B < +\infty\) 缺一不可？（反例直觉）}： 如果去掉这个条件，结论不成立！反例：取 \(X = Z = \ell^2, \, Y = \ell^2 \oplus \ell^2\)，令 \(Ax = (x, 0)\)，\(B(x, y) = x\)。此时 \(BA = I_{\ell^2} \in \mathcal{F}(\ell^2)\)，但 \(\operatorname{Ker} B = \{0\} \oplus \ell^2\) 是无穷维的（\(B \notin \mathcal{F}\)），同时 \(\mathcal{R}(A) = \ell^2 \oplus \{0\}\) 在 \(Y\) 中余维无穷（\(A \notin \mathcal{F}\)）。
\item
  \textbf{延伸考点}：本题条件可对称替换为``\(BA \in \mathcal{F}(X, Z)\) 且 \(\operatorname{codim}\mathcal{R}(A) < +\infty\)''，同样能推出 \(A \in \mathcal{F}(X, Y)\) 且 \(B \in \mathcal{F}(Y, Z)\)（证明只需先用 \(\operatorname{Ker} A \subset \operatorname{Ker}(BA)\) 与 \(\operatorname{codim}\mathcal{R}(A)<\infty\) 证 \(A \in \mathcal{F}\)，再如法炮制用 Atkinson 定理证 \(B \in \mathcal{F}\)）。
\end{enumerate}

\BookPart{}{三、考前 3 分钟速背清单（全章定理与题型总览）}{三、考前 3 分钟速背清单（全章定理与题型总览）}

{\def\LTcaptype{none} % do not increment counter
{\small\renewcommand{\arraystretch}{1.38}\setlength{\tabcolsep}{5pt}
\begin{longtable}[]{@{}
  >{\raggedright\arraybackslash}p{(\linewidth - 4\tabcolsep) * \real{0.3333}}
  >{\raggedright\arraybackslash}p{(\linewidth - 4\tabcolsep) * \real{0.3333}}
  >{\raggedright\arraybackslash}p{(\linewidth - 4\tabcolsep) * \real{0.3333}}@{}}
\toprule\noalign{}
\rowcolor{AccentSoft}
\begin{minipage}[b]{\linewidth}\raggedright
考场遇到什么问题
\end{minipage} & \begin{minipage}[b]{\linewidth}\raggedright
立刻条件反射的工具
\end{minipage} & \begin{minipage}[b]{\linewidth}\raggedright
对应教材与作业
\end{minipage} \\
\midrule\noalign{}
\endhead
\bottomrule\noalign{}
\endlastfoot
证无穷维紧算子\textbf{不可逆} / 某算子\textbf{不紧} & 反证法凑出 \(I_X \in \mathcal{C}(X)\) + \textbf{Riesz 引理} & 定理 5.5、例 5.2、\textbf{第 5 题} \\
证紧算子值域\textbf{不含无穷维闭子空间} & 收缩到 \(X_0 = A^{-1}(M)\) + \textbf{开映射定理}（或商空间 + 逆算子定理）+ Riesz 引理 & \textbf{第 6 题}、习题 17(ii) \\
证 \(C([0,1])\) 上积分算子是\textbf{紧算子} & \textbf{Arzelà-Ascoli 定理}（一致有界 + 等度连续，强调 \(\delta\) 与 \(f\) 无关） & 例 5.1、\textbf{第 14 题} \\
证 \(\ell^p / L^2\) 算子是\textbf{紧算子} & 有限秩截断逼近（定理 5.4）或 \textbf{Hilbert-Schmidt 判定}（定理 5.12(5)） & 定理 5.4、定理 5.9、定理 5.12 \\
证有限秩算子是 \textbf{Hilbert-Schmidt 算子} & \(H = \operatorname{Ker} A \oplus (\operatorname{Ker} A)^\perp\) 凑基，或 Parseval 等式 + \(A^*\) & 定义 5.4、\textbf{第 15 题} \\
证算子是 \textbf{Fredholm 算子} & 优先 \textbf{Atkinson 定理}（定理 5.13：凑 \(SA = I - K_1, \, AS = I - K_2\)）；按定义证时\textbf{切勿漏证值域闭}！ & 定义 5.5、定理 5.13、\textbf{第 20/22 题} \\
算/证 \textbf{Fredholm 指标}相等 & \textbf{乘积指标公式}（定理 5.14）+ \textbf{\(\operatorname{ind}(I - K) = 0\)}（定理 5.8(4)）+ 公共拟逆 \(S\) & 定理 5.14--5.15、\textbf{第 20 题} \\
\end{longtable}
}
}


% =========================================================================
% 第 6 章 · 有界线性算子的谱理论
% =========================================================================
\SwitchChapterTheme{6}
\ChapterBanner{06}{ch:6}{第 6 章 · 有界线性算子的谱理论}{谱三分类与四例 · Gelfand 谱半径 · Riesz-Schauder · Hilbert-Schmidt · 谱分解 · 习题 6（9 题）}
\begin{quote}
\textbf{使用说明}：本资料面向\textbf{零泛函分析基础}的本科生期末急救与刷题，内容、记号与定理编号严格对齐教材第 6 章（p. 206--271）及第 5 章第 5.2 节（Hilbert-Schmidt 算子，p. 196--197）。所有作业解答均自包含、步骤严谨，可直接作为期末考试规范作答模板。
\end{quote}

\BookPart{}{零基础导读：3 分钟看懂第 6 章到底在学什么}{零基础导读：3 分钟看懂第 6 章到底在学什么}

在\textbf{线性代数}里，研究一个 \(n \times n\) 矩阵 \(A\) 最强大的武器是\textbf{特征值与对角化}：

\begin{itemize}
\tightlist
\item
  找特征值 \(\lambda\)：解 \(\det(\lambda I - A) = 0\)，也就是让矩阵 \(\lambda I - A\) \textbf{不可逆}的复数 \(\lambda\)；
\item
  如果是实对称矩阵（或复 Hermite 矩阵 \(A^* = A\)），一定可以正交对角化：存在标准正交特征向量 \(e_1, \dots, e_n\) 和实特征值 \(\lambda_1, \dots, \lambda_n\)，使得 
\FitDisplayMath{Ax = \sum_{k=1}^{n} \lambda_k (x, e_k) e_k, \qquad A = \sum_{k=1}^{n} \lambda_k P_k,}
 其中 \(P_k x = (x, e_k)e_k\) 是向第 \(k\) 个特征方向的正交投影。
\item
  一旦写成了 \(A = \sum \lambda_k P_k\)，求 \(A^m\)、\(A^{-1}\)、\(A^{1/2}\)、\(\mathrm{e}^{tA}\) 就全变成了对标量做运算：\(f(A) = \sum f(\lambda_k) P_k\)！
\end{itemize}

到了\textbf{泛函分析的无穷维空间}（如 \(L^2(0,1)\)、\(\ell^2\)、\(C([0,1])\)），发生了两件天翻地覆的事：

\begin{enumerate}
\def\labelenumi{\arabic{enumi}.}
\tightlist
\item
  \textbf{``\(\lambda I - A\) 不可逆''不再等价于``\(\lambda\) 是特征值''！} 有限维里，方阵单射 \(\iff\) 满射 \(\iff\) 可逆。但在无穷维里，\(\lambda I - A\) 可以是单射（即 \(\lambda\) 不是特征值），却\textbf{不是满射}，从而依然不可逆！最著名的例子是 \(L^2(0,1)\) 上的乘法算子 \((Af)(x) = x f(x)\)：它\textbf{一个特征值都没有}，但对所有 \(\lambda \in [0,1]\)，\(\lambda I - A\) 都不可逆。因此必须把``特征值集合''推广为\textbf{谱集 \(\sigma(A)\)}（所有让 \(\lambda I - A\) 不可逆的 \(\lambda\)）。
\item
  \textbf{一般自伴算子的谱可能是连续的区间（例如 \([0,1]\)），不再只有离散的点！} 如果算子是\textbf{自伴紧算子}（最像矩阵的无穷维算子），谱还是离散可数的，我们依然有级数展开 \(A = \sum_{n=1}^\infty \lambda_n P_n\)（定理 6.13）；但对于\textbf{一般的有界自伴算子}，离散求和``不够用了''，必须把求和号 \(\sum\) 升级为积分号 \(\int\)，把离散投影 \(P_k\) 升级为\textbf{谱测度（投影值测度）} \(\mathrm{d}E_\lambda\)，写成 
\FitDisplayMath{A = \int_{\sigma(A)} \lambda \,\mathrm{d}E_\lambda, \qquad f(A) = \int_{\sigma(A)} f(\lambda)\,\mathrm{d}E_\lambda.}

\end{enumerate}

\BookSubSection{}{有限维线性代数 vs 第 6 章泛函分析对照表}{有限维线性代数 vs 第 6 章泛函分析对照表}{}

{\def\LTcaptype{none} % do not increment counter
{\small\renewcommand{\arraystretch}{1.38}\setlength{\tabcolsep}{5pt}
\begin{longtable}[]{@{}
  >{\raggedright\arraybackslash}p{(\linewidth - 4\tabcolsep) * \real{0.3333}}
  >{\raggedright\arraybackslash}p{(\linewidth - 4\tabcolsep) * \real{0.3333}}
  >{\raggedright\arraybackslash}p{(\linewidth - 4\tabcolsep) * \real{0.3333}}@{}}
\toprule\noalign{}
\rowcolor{AccentSoft}
\begin{minipage}[b]{\linewidth}\raggedright
线性代数（有限维 \(\mathbb{C}^n\)）
\end{minipage} & \begin{minipage}[b]{\linewidth}\raggedright
泛函分析（无穷维 Banach / Hilbert 空间）
\end{minipage} & \begin{minipage}[b]{\linewidth}\raggedright
对应教材小节
\end{minipage} \\
\midrule\noalign{}
\endhead
\bottomrule\noalign{}
\endlastfoot
特征值集合（\(\det(\lambda I - A)=0\)） & 谱集 \(\sigma(A) = \sigma_p(A) \cup \sigma_c(A) \cup \sigma_r(A)\) & 6.1 节 \\
矩阵逆 \((\lambda I - A)^{-1}\) & 预解式 \(R(\lambda, A) = (\lambda I - A)^{-1}\) 与 Neumann 级数 & 6.1 节 \\
最大特征值模 \(\max \lvert \lambda_i \rvert\) & 谱半径 \(\Gamma(A) = \sup_{\lambda \in \sigma(A)} \lvert \lambda \rvert = \lim_{n\to\infty}\sqrt[n]{\lVert A^n \rVert}\) & 6.1.4 节 \\
有限维线性变换 & 紧算子（非零谱点必为特征值、至多可数、只可能聚于 \(0\)） & 6.2 节 \\
Hermite 矩阵正交对角化 \(A=\sum \lambda_k P_k\) & 自伴紧算子 Hilbert-Schmidt 谱展开 \(Ax = \sum \lambda_n (x,e_n)e_n\) & 6.3 节 \\
矩阵 Frobenius 范数与奇异值分解 & Hilbert-Schmidt 算子、\(\lvert A \rvert = (A^*A)^{1/2}\) 与迹类算子 & 5.2 / 6.3.2 节 \\
正交投影族 \(\{P_k\}\) 与离散求和 & 谱测度 \(E(\cdot)\)、谱系 \(\{E_\lambda\}\) 与谱积分 \(\int f \,\mathrm{d}E\) & 6.4 节 \\
酉矩阵 \(U^*U = UU^* = I\) 与 Cayley 变换 & 酉算子谱分解（单位圆上）\(\to\) 有界自伴算子谱分解（实轴上） & 6.5--6.6 节 \\
矩阵函数 \(f(A) = \sum f(\lambda_k) P_k\) & 有界 Borel 函数演算 \(f(A) = \int_{\sigma(A)} f(\lambda)\,\mathrm{d}E_\lambda\) & 6.6 节 \\
\end{longtable}
}
}

\BookSubSection{}{全章 9 道课后作业题考点快速索引}{全章 9 道课后作业题考点快速索引}{}

{\def\LTcaptype{none} % do not increment counter
{\small\renewcommand{\arraystretch}{1.38}\setlength{\tabcolsep}{5pt}
\begin{longtable}[]{@{}
  >{\centering\arraybackslash}p{(\linewidth - 6\tabcolsep) * \real{0.2778}}
  >{\centering\arraybackslash}p{(\linewidth - 6\tabcolsep) * \real{0.2778}}
  >{\raggedright\arraybackslash}p{(\linewidth - 6\tabcolsep) * \real{0.2222}}
  >{\raggedright\arraybackslash}p{(\linewidth - 6\tabcolsep) * \real{0.2222}}@{}}
\toprule\noalign{}
\rowcolor{AccentSoft}
\begin{minipage}[b]{\linewidth}\centering
题号
\end{minipage} & \begin{minipage}[b]{\linewidth}\centering
页码
\end{minipage} & \begin{minipage}[b]{\linewidth}\raggedright
核心考点
\end{minipage} & \begin{minipage}[b]{\linewidth}\raggedright
所用关键定理 / 技巧
\end{minipage} \\
\midrule\noalign{}
\endhead
\bottomrule\noalign{}
\endlastfoot
\hyperlink{prob:6-2}{\textbf{第 2 题}} & p. 266 & \(L^2(0,1)\) 乘法算子的自伴性与谱集 & 紧区间连续函数值域闭性 + 局部特征函数构造近似特征向量 \\
\hyperlink{prob:6-5}{\textbf{第 5 题}} & p. 266 & 预解集开性与预解式关于算子的连续性 & 算子因式分解 + Neumann 级数（定理 6.1） \\
\hyperlink{prob:6-13}{\textbf{第 13 题}} & p. 267 & 预解算子 \(B = R(\alpha, A)\) 的谱与预解式 & 恒等式 \(\mu I - B = \mu(\alpha I - A)^{-1}(\lambda I - A)\) \\
\hyperlink{prob:6-23}{\textbf{第 23 题}} & p. 268 & 广义 Cayley 变换 \((A+\lambda I)(A+\bar{\lambda}I)^{-1}\) 的酉性 & 下有界性证可逆 + 伴随运算验证 \(U^*U = UU^* = I\)（定义 6.12） \\
\hyperlink{prob:6-24}{\textbf{第 24 题}} & p. 268 & Volterra 积分算子伴随、\(AA^*\) 谱半径与算子范数 & Fubini 定理 + 紧自伴算子 \(\Gamma(AA^*)=\lVert AA^*\rVert\) + ODE 边值问题 \\
\hyperlink{prob:6-31}{\textbf{第 31 题}} & p. 269 & Hilbert-Schmidt 算子、\((A^*A)^{1/2}\) 与奇异值展开等价性 & 定义 5.4、定理 5.12、定理 6.13、推论 6.17 + 严谨基配对注记 \\
\hyperlink{prob:6-37}{\textbf{第 37 题}} & p. 270 & Hilbert-Schmidt 范数等于奇异值平方和开方 & \(\lVert Ax \rVert = \lVert (A^*A)^{1/2}x \rVert\) + 定理 6.13 特征基补全 \\
\hyperlink{prob:6-39}{\textbf{第 39 题}} & p. 270 & 酉相似算子 \(\tilde{A}=U^{-1}AU\) 的 Borel 函数演算 & 谱测度共轭变换 \(\tilde{E}=U^{-1}EU\) + 谱测度唯一性（定理 6.21 / 6.25） \\
\hyperlink{prob:6-45}{\textbf{第 45 题}} & p. 271 & 解析函数（幂级数）演算与算子幂级数的一致性 & 收敛圆内一致收敛 + 推论 6.7（乘法同态）+ 推论 6.8 / 6.16（范数界） \\
\end{longtable}
}
}

\BookPart{}{第一部分：第 6 章核心知识零基础精讲}{第一部分：第 6 章核心知识零基础精讲}

\begin{quote}
\textbf{全章基础约定}：如无特别说明，本章中 \(X\) 表示\textbf{复数域 \(\mathbb{C}\) 上的 Banach 空间}（完备赋范线性空间），\(H\) 表示\textbf{复数域 \(\mathbb{C}\) 上的 Hilbert 空间}（完备内积空间）；\(\mathcal{L}(X)\) 表示 \(X\) 上全体有界线性算子组成的 Banach 空间，\(\mathcal{C}(X)\) 表示 \(X\) 上全体紧算子；复内积 \((\cdot, \cdot)\) 对第一变元线性、对第二变元共轭线性。之所以必须在\textbf{复数域}上讨论谱理论，是因为像复平面上的 Liouville 定理、代数基本定理等核心工具在实数域上不成立（实矩阵如 \(90^\circ\) 旋转矩阵就没有实特征值）。
\end{quote}

\SoftDivider

\BookSubSection{sec:6-1}{6.1 有界线性算子谱的定义和基本性质（p. 206--215）}{6.1 有界线性算子谱的定义和基本性质（p. 206–215）}{\TagMust{闭卷必背}}

\SubSubHeading{1. 预解集、预解式与谱集（定义 6.1、注记 6.1，p. 206）}

设 \(A \in \mathcal{L}(X)\)，\(I\) 是 \(X\) 上的恒等算子。

\begin{itemize}
\item
  \textbf{预解集（正则点集）}： 
\FitDisplayMath{\rho(A) \triangleq \{\lambda \in \mathbb{C} \mid \lambda I - A \text{ 是 } X \to X \text{ 的双射}\}.}
 \(\rho(A)\) 中的点称为 \(A\) 的\textbf{正则点}。
\item
  \textbf{预解式（预解算子）}：当 \(\lambda \in \rho(A)\) 时，逆算子 
\FitDisplayMath{R(\lambda, A) \triangleq (\lambda I - A)^{-1}}
 称为 \(A\) 在 \(\lambda\) 处的\textbf{预解式}。

  \begin{quote}
  \textbf{注记 6.1（为什么不用额外要求逆算子有界？）}：因为 \(X\) 是完备的 Banach 空间且 \(\lambda I - A \in \mathcal{L}(X)\)，只要 \(\lambda I - A\) 是双射，由 \textbf{Banach 逆算子定理}自动保证 \((\lambda I - A)^{-1} \in \mathcal{L}(X)\) 有界！
  \end{quote}
\item
  \textbf{谱集}：预解集在复平面上的补集 
\FitDisplayMath{\sigma(A) \triangleq \mathbb{C} \setminus \rho(A)}
 称为 \(A\) 的\textbf{谱集}，\(\sigma(A)\) 中的点称为\textbf{谱点}。
\end{itemize}

\SubSubHeading{2. 谱的三大分类（定义 6.2，p. 206）}

当 \(\lambda \in \sigma(A)\)（即 \(\lambda I - A\) 不是双射）时，失败的方式只有三种互不相交的情形：

{\def\LTcaptype{none} % do not increment counter
{\small\renewcommand{\arraystretch}{1.38}\setlength{\tabcolsep}{5pt}
\begin{longtable}[]{@{}
  >{\raggedright\arraybackslash}p{(\linewidth - 6\tabcolsep) * \real{0.2353}}
  >{\centering\arraybackslash}p{(\linewidth - 6\tabcolsep) * \real{0.2941}}
  >{\raggedright\arraybackslash}p{(\linewidth - 6\tabcolsep) * \real{0.2353}}
  >{\raggedright\arraybackslash}p{(\linewidth - 6\tabcolsep) * \real{0.2353}}@{}}
\toprule\noalign{}
\rowcolor{AccentSoft}
\begin{minipage}[b]{\linewidth}\raggedright
谱的类型
\end{minipage} & \begin{minipage}[b]{\linewidth}\centering
记号
\end{minipage} & \begin{minipage}[b]{\linewidth}\raggedright
严格判定条件
\end{minipage} & \begin{minipage}[b]{\linewidth}\raggedright
大白话理解
\end{minipage} \\
\midrule\noalign{}
\endhead
\bottomrule\noalign{}
\endlastfoot
\textbf{点谱（特征值）} & \(\sigma_p(A)\) & \(\lambda I - A\) \textbf{不是单射}（\(\exists x \neq 0\) 使 \(Ax = \lambda x\)） & 和线代完全一样，真有非零特征向量 \\
\textbf{连续谱} & \(\sigma_c(A)\) & \(\lambda I - A\) 是单射，值域 \(\mathcal{R}(\lambda I - A) \neq X\)，但 \(\overline{\mathcal{R}(\lambda I - A)} = X\) & 没打满整个空间，但``几乎打满''（值域稠密），此时逆算子定义在稠密子空间上但无界 \\
\textbf{剩余谱} & \(\sigma_r(A)\) & \(\lambda I - A\) 是单射，且 \(\overline{\mathcal{R}(\lambda I - A)} \neq X\) & 没打满且``缺了一大块''（值域不稠密，连闭包都盖不住全空间） \\
\end{longtable}
}
}

\begin{quote}
\textbf{注}：教材定义 6.2(ii)--(iii) 文字中默许在``(i) 不成立（即 \(\lambda I - A\) 为单射）''的前提下按值域是否稠密继续细分，因此复平面分解为四个\textbf{两两不相交}集合的并： 
\FitDisplayMath{\mathbb{C} = \rho(A) \cup \sigma(A) = \rho(A) \cup \sigma_p(A) \cup \sigma_c(A) \cup \sigma_r(A).}

\end{quote}

\begin{quote}
\textbf{零基础防懵：谱点和特征值到底是什么关系？}

\begin{itemize}
\tightlist
\item
  \textbf{一句话}：特征值只是谱点里的``其中一类''（即点谱 \(\sigma_p(A)\)），满足 \(\sigma_p(A) \subset \sigma(A) = \sigma_p(A) \cup \sigma_c(A) \cup \sigma_r(A)\)；因此\textbf{特征值一定是谱点，但谱点不一定是特征值！}
\item
  \textbf{为什么会多出不是特征值的谱点？} 有限维方阵里``单射 \(\iff\) 满射''，只要没有非零特征向量就一定可逆，所以不可能出现 \(\sigma_c\) 和 \(\sigma_r\)；而无穷维空间像``无限房间旅馆''，完全可以\textbf{一对一不撞车（单射、无特征值），却空出房间没住满（不满射）}——按``值域稠不稠密''就长出了连续谱 \(\sigma_c\) 和剩余谱 \(\sigma_r\)！
\end{itemize}
\end{quote}

\textbf{零基础必背的 4 个标志性算子例子（例 6.1--6.4、例 6.7，p. 206--208, 215）——逐个看懂``为什么坏在这里''}：

\begin{enumerate}
\def\labelenumi{\arabic{enumi}.}
\tightlist
\item
  \textbf{有限维方阵（例 6.1）}：\(\sigma(A) = \sigma_p(A)\)，\(\sigma_c(A) = \sigma_r(A) = \emptyset\)（有限维方阵单射 \(\iff\) 满射，坏点只能是不单射的特征值）。
\item
  \textbf{Volterra 积分算子（例 6.2、例 6.5）}：\((Ax)(t) = \int_a^t x(\tau)\,\mathrm{d}\tau\) 在 \(C([a,b])\) 上，\textbf{无特征值}（\(\lambda = 0\) 时积分恒为 \(0\) 求导得 \(x\equiv 0\)；\(\lambda \neq 0\) 时化为一阶 ODE \(x' = \frac{1}{\lambda}x, x(a)=0\) 解唯一为 \(x\equiv 0\)）。对任意 \(\lambda \neq 0\) 可显式解出唯一连续解，故 \(\lambda \in \rho(A)\)；而 \(\lambda = 0\) 时值域里每个函数必满足 \(y(a)=0\)，在 \(C([a,b])\) 一致范数下与常函数 \(h \equiv 1\) 距离恒 \(\ge 1\)，值域\textbf{不稠密}，故 \(0 \in \sigma_r(A)\)、\(\sigma(A) = \{0\}\)。
\item
  \textbf{乘法算子 \((Af)(t) = t f(t)\)（同一个算子，换一把距离尺子就从 \(\sigma_r\) 变 \(\sigma_c\)！）}：

  \begin{itemize}
  \tightlist
  \item
    先看无特征值：\((\lambda - t)f(t) = 0\) 仅可能在单点 \(t=\lambda\) 非零，在 \(C([a,b])\) 由连续性、在 \(L^2(0,1)\) 由单点零测，都迫使 \(f = 0\)，故两空间均单射（\(\sigma_p = \emptyset\)）；区间外除以 \(\lambda - t\) 即有界逆。
  \item
    \textbf{在 \(C([a,b])\) 上（例 6.3）}：\(\sigma(A) = \sigma_r(A) = [a,b]\)。因为值域函数 \(g(t)=(\lambda-t)f(t)\) 在 \(t=\lambda\) 处必为 \(0\)，而 \(C([a,b])\) 看逐点最大误差，与常函数 \(h \equiv 1\) 的最大差恒 \(\ge 1\)，逼近不了，值域\textbf{不稠密} \(\implies\) 剩余谱！
  \item
    \textbf{在 \(L^2(0,1)\) 上（例 6.4）}：\(\sigma(A) = \sigma_c(A) = [0,1]\)。因为 \(L^2\) 看的是积分面积，对任意 \(h \in L^2\) 只需把单点 \(\lambda\) 附近极窄缝 \((\lambda-\varepsilon, \lambda+\varepsilon)\) 挖掉（外面令 \(f_\varepsilon(t) = \frac{h(t)}{\lambda-t}\chi_{|t-\lambda|>\varepsilon} \in L^2\)），面积误差随 \(\varepsilon \to 0\) 趋于 \(0\)，值域\textbf{稠密}（但如 \(h\equiv 1\) 无原像故不满）\(\implies\) 连续谱！
  \end{itemize}
\item
  \textbf{\(\ell^2\) 上的右移算子（例 6.7）}：\(U(x_1, x_2, \dots) = (0, x_1, x_2, \dots)\)。

  \begin{itemize}
  \tightlist
  \item
    \textbf{无特征值（\(\sigma_p(U)=\emptyset\)）}：\(Ux=\lambda x\) 逐项比对必推出 \(x_1 = x_2 = \cdots = 0\)。
  \item
    \textbf{\(\sigma_r(U) = \{\lambda : \lvert\lambda\rvert < 1\}\)}：值域不稠密 \(\iff\) 存在非零 \(y \in \ell^2\) 垂直于值域 \(\iff U^*y = \bar{\lambda}y\)（\(U^*\) 为左移算子），解得 \(y = (1, \bar{\lambda}, \bar{\lambda}^2, \dots)\)；当 \(\lvert\lambda\rvert < 1\) 时级数平方和收敛，\(y \in \ell^2\) 合法存在，故值域不稠密！
  \item
    \textbf{\(\sigma_c(U) = \{\lambda : \lvert\lambda\rvert = 1\}\)}：当 \(\lvert\lambda\rvert = 1\) 时上述 \(y\) 掉出 \(\ell^2\)，值域变得稠密；又因谱集 \(\sigma(U)\) 必闭，开圆盘边界 \(\lvert\lambda\rvert = 1\) 仍属谱集，故只能是连续谱！
  \end{itemize}
\end{enumerate}

\SubSubHeading{3. Neumann 级数、预解集的开性与谱集的紧性（引理 6.1、定理 6.1--6.2、推论 6.1，p. 208--211）}

把标量几何级数 \(\frac{1}{\lambda - a} = \sum_{n=0}^\infty \frac{a^n}{\lambda^{n+1}}\) 推广到算子，就得到本章最有力的计算武器：

\begin{itemize}
\item
  \textbf{引理 6.1（极限存在性，p. 208）}：对任意 \(A \in \mathcal{L}(X)\)，极限 \(\lim_{n\to\infty} \sqrt[n]{\|A^n\|}\) 一定存在，且等于 \(\inf_{n \in \mathbb{N}} \sqrt[n]{\|A^n\|} \le \|A\|\)。
\item
  \textbf{定理 6.1（Neumann 级数展开，p. 209）}：记 \(\tau = \lim_{n\to\infty}\sqrt[n]{\|A^n\|}\)。当 \(|\lambda| > \tau\)（特别地，当 \(|\lambda| > \|A\|\)）时，必有 \(\lambda \in \rho(A)\)，且预解式按算子范数收敛于级数 
\FitDisplayMath{(\lambda I - A)^{-1} = \sum_{n=0}^\infty \frac{A^n}{\lambda^{n+1}}.}


  \begin{quote}
  \textbf{极高频特例（取 \(\lambda = 1\)）}：只要 \(\|B\| < 1\)，则 \(I - B\) 必可逆，且 
\FitDisplayMath{(I - B)^{-1} = \sum_{k=0}^\infty B^k, \qquad \|(I - B)^{-1}\| \le \frac{1}{1 - \|B\|}.}
 （这是证明预解集开性、第 5 题及各类算子扰动问题的``万能钥匙''！）
  \end{quote}
\item
  \textbf{定理 6.2（预解式的局部幂级数展开，p. 210）}：若 \(\lambda_0 \in \rho(A)\)，且 \(|\lambda - \lambda_0| < \frac{1}{\|(\lambda_0 I - A)^{-1}\|}\)，则 \(\lambda \in \rho(A)\)，且 
\FitDisplayMath{(\lambda I - A)^{-1} = \sum_{n=0}^\infty (-1)^n (\lambda_0 I - A)^{-(n+1)} (\lambda - \lambda_0)^n.}

\item
  \textbf{推论 6.1（p. 211）}：对任意 \(A \in \mathcal{L}(X)\)，\(\rho(A)\) 是复平面上的\textbf{开集}，\(\sigma(A)\) 是\textbf{有界闭集（紧集）}（有界性来自定理 6.1：\(\sigma(A) \subset \{\lambda : |\lambda| \le \|A\|\}\)）。
\item
  \textbf{第一预解式恒等式（式 (6.10)，p. 213）}：对任意 \(\lambda, \mu \in \rho(A)\)， 
\FitDisplayMath{R(\lambda, A) - R(\mu, A) = (\mu - \lambda) R(\lambda, A) R(\mu, A).}
 记忆窍门：完全就是小学分式通分 \(\frac{1}{\lambda - A} - \frac{1}{\mu - A} = \frac{\mu - \lambda}{(\lambda - A)(\mu - A)}\)！
\end{itemize}

\SubSubHeading{4. 谱集非空性与 Gelfand 谱半径公式（定理 6.4--6.5，p. 213--214）}

\begin{itemize}
\item
  \textbf{定理 6.4（谱非空定理，p. 213）}：对复 Banach 空间 \(X \neq \{0\}\) 上的任意 \(A \in \mathcal{L}(X)\)，恒有 \(\sigma(A) \neq \emptyset\)。

  \begin{itemize}
  \tightlist
  \item
    \textbf{证明思路一句话}：反设 \(\sigma(A) = \emptyset\)，则 \(R(\lambda, A)\) 在全复平面解析，且当 \(|\lambda| \to \infty\) 时 \(\|R(\lambda, A)\| \to 0\)；对任意 \(f \in X', x \in X\)，复函数 \(f(R(\lambda, A)x)\) 是有界整函数，由复变 Liouville 定理恒为 \(0\)，推得 \(R(\lambda, A) \equiv 0\)，与 \(R(\lambda, A)(\lambda I - A) = I\) 矛盾！
  \end{itemize}
\item
  \textbf{定理 6.5（Gelfand 谱半径公式，p. 214）}：定义 \(A\) 的\textbf{谱半径}为谱集包含在原点为心的最小闭圆盘半径： 
\FitDisplayMath{\Gamma(A) \triangleq \sup_{\lambda \in \sigma(A)} |\lambda|.}
 则有精确计算公式： 
\FitDisplayMath{\Gamma(A) = \lim_{n\to\infty} \sqrt[n]{\|A^n\|} = \inf_{n \in \mathbb{N}} \sqrt[n]{\|A^n\|} \le \|A\|.}


  \begin{quote}
  \textbf{记号统一提示}：教材在定理 6.5（p. 214）与习题 6（第 6 题、第 45 题）中统一将谱半径记作大写 \(\Gamma(A)\)，但在推论 6.2（p. 222）中亦写作小写 \(\gamma(A)\)，二者完全同义。本资料统一使用 \(\Gamma(A)\)。
  \end{quote}
\end{itemize}

\SoftDivider

\BookSubSection{sec:6-2}{6.2 紧算子的谱理论：Riesz-Schauder 理论（p. 215--220）}{6.2 紧算子的谱理论：Riesz-Schauder 理论（p. 215–220）}{\TagMust{闭卷必背}}

\textbf{什么是紧算子？（第 5 章回顾）} 把任何有界点列映成``存在收敛子列''的点列的线性算子 \(A \in \mathcal{C}(X)\)。直觉上，紧算子把无穷维空间里的单位球``压扁''成了近似有限维的列紧集，因此它的\textbf{非零谱行为几乎和有限维矩阵一模一样}！

\SubSubHeading{Riesz-Schauder 理论核心结论（6.2 节，定理 6.6--6.9，p. 215--220）}

\begin{enumerate}
\def\labelenumi{\arabic{enumi}.}
\tightlist
\item
  \textbf{定理 6.6（满射即正则，p. 215）}：设 \(A \in \mathcal{C}(X)\)，\(\lambda \neq 0\)。若 \((\lambda I - A)X = X\)，则 \(\lambda \in \rho(A)\)。
\item
  \textbf{非零谱点的完全刻画（Riesz-Schauder 理论，6.2 节）}：设 \(A \in \mathcal{C}(X)\)，则：

  \begin{itemize}
  \tightlist
  \item
    \textbf{无穷维时 \(0\) 必在谱中}：若 \(\dim X = \infty\)，则 \(0 \in \sigma(A)\)（因为如果 \(0 \in \rho(A)\)，则恒等算子 \(I = A \cdot A^{-1}\) 为紧算子，迫使单位球紧，仅在有限维可能）。
  \item
    \textbf{非零谱点全是特征值}：若 \(\lambda \in \sigma(A)\) 且 \(\lambda \neq 0\)，则必有 \(\lambda \in \sigma_p(A)\)（即紧算子的连续谱和剩余谱最多只能包含 \(0\)！）。
  \item
    \textbf{特征空间有限维}：对任意非零特征值 \(\lambda\)，特征子空间 \(\operatorname{Ker}(\lambda I - A)\) 必为\textbf{有限维}，且值域 \(\mathcal{R}(\lambda I - A)\) 必为\textbf{闭子空间}。
  \item
    \textbf{至多可数且唯一可能聚点是 \(0\)}：\(\sigma(A)\) 至多是可数集；若存在无穷多个互异特征值 \(\{\lambda_n\}\)，必有 \(\lim_{n\to\infty} \lambda_n = 0\)（对任意 \(\varepsilon > 0\)，圆盘 \(|\lambda| \ge \varepsilon\) 外只有有限个谱点）。
  \end{itemize}
\item
  \textbf{定理 6.9（紧算子与共轭算子、Fredholm 择一性，p. 219）}：设 \(A \in \mathcal{C}(X)\)，则：

  \begin{itemize}
  \tightlist
  \item
    \textbf{(i)} \(\sigma(A) = \sigma(A^*)\)（这里 \(A^* : X' \to X'\) 为 Banach 空间共轭算子；若在 Hilbert 空间用 Riesz 表示定义的伴随算子 \(A^*\)，则 \(\sigma(A^*) = \overline{\sigma(A)}\)）；
  \item
    \textbf{(ii)} 若 \(\lambda, \mu \in \sigma(A)\) 且 \(\lambda \neq \mu\)，则 \(A\) 对应 \(\lambda\) 的特征向量 \(x\) 与 \(A^*\) 对应 \(\mu\) 的特征泛函 \(f\) 正交（\(f(x) = 0\)）；
  \item
    \textbf{(iii)--(iv) Fredholm 择一性（方程可解判据）}：对非零 \(\lambda \in \sigma(A)\)，方程 \((\lambda I - A)x = y\) 有解 \(\iff y\) 与 \(A^*\) 对应于 \(\lambda\) 的一切特征泛函正交；共轭方程 \((\lambda I - A^*)\varphi = f\) 有解 \(\iff f\) 与 \(A\) 对应于 \(\lambda\) 的一切特征向量正交。
  \end{itemize}
\end{enumerate}

\SoftDivider

\BookSubSection{sec:6-3}{6.3 Hilbert 空间上自伴算子与自伴紧算子的谱理论（p. 220--231）}{6.3 Hilbert 空间上自伴算子与自伴紧算子的谱理论（p. 220–231）}{\TagMust{闭卷必背}}

\SubSubHeading{1. 自伴算子与正算子的基础（结合第 3 章与 6.3 / 6.6 节）}

设 \(A \in \mathcal{L}(H)\)：

\begin{itemize}
\tightlist
\item
  \textbf{自伴（自共轭）算子}：\(A^* = A\)，即对一切 \(x, y \in H\) 有 \((Ax, y) = (x, Ay)\)。等价地，对一切 \(x \in H\) 都有 \((Ax, x) \in \mathbb{R}\)。
\item
  \textbf{正算子（记作 \(A \ge 0\)）}：\(A\) 是自伴算子且对一切 \(x \in H\) 有 \((Ax, x) \ge 0\)。（注：复 Hilbert 空间上只要 \((Ax, x) \ge 0\) 对一切 \(x\) 成立，由极化恒等式自动推出 \(A\) 自伴。对任意 \(T \in \mathcal{L}(H, G)\)，\(T^*T \ge 0\) 且 \(TT^* \ge 0\) 恒成立！）
\end{itemize}

\SubSubHeading{2. 自伴算子的核心谱性质（定理 6.10--6.11、引理 6.10，p. 220--221, 263）}

\begin{itemize}
\tightlist
\item
  \textbf{定理 6.10（p. 220）}：设 \(A \in \mathcal{L}(H)\) 自伴，则：

  \begin{itemize}
  \tightlist
  \item
    \begin{enumerate}
    \def\labelenumi{(\roman{enumi})}
    \tightlist
    \item
      \(A\) 的特征值都是实数（进一步，整个谱集都在实轴上：\(\sigma(A) \subset \mathbb{R}\)，且 \(\sigma_r(A) = \emptyset\)）；
    \end{enumerate}
  \item
    (ii) \(A\) 对应于不同特征值的特征向量互相正交。
  \end{itemize}
\item
  \textbf{定理 6.11（二次型范数公式，p. 221）}：设 \(A \in \mathcal{L}(H)\) 自伴，则 
\FitDisplayMath{\|A\| = \sup_{\|x\|=1} |(Ax, x)|.}

\item
  \textbf{自伴算子的谱半径等于范数（推论 6.2 / 引理 6.10）}：

  \begin{itemize}
  \tightlist
  \item
    对一般有界自伴算子 \(A\)，由 \(\|A^2\| = \|A^*A\| = \|A\|^2\) 归纳得 \(\|A^{2^k}\| = \|A\|^{2^k}\)，代入 Gelfand 公式立得： 
\FitDisplayMath{\Gamma(A) = \lim_{k\to\infty} \|A^{2^k}\|^{1/2^k} = \|A\|.}

  \item
    更精细地，\textbf{引理 6.10（p. 263）} 给出谱的上下界恰好是二次型的上下确界： 
\FitDisplayMath{\sup_{\lambda \in \sigma(A)} \lambda = \sup_{\|x\|=1} (Ax, x), \qquad \inf_{\lambda \in \sigma(A)} \lambda = \inf_{\|x\|=1} (Ax, x).}
 由此立得极常用的充要判据：\(A \ge 0 \iff \sigma(A) \subset [0, \infty)\)。
  \end{itemize}
\end{itemize}

\SubSubHeading{3. 自伴紧算子的谱展开定理（定理 6.12--6.13、推论 6.2--6.4，p. 221--224）}

当算子同时满足\textbf{自伴}和\textbf{紧}两个好性质时，它就拥有完美的``无穷维正交对角化''：

\begin{itemize}
\tightlist
\item
  \textbf{定理 6.12（最大特征值存在性，p. 221）}：设 \(A \in \mathcal{C}(H)\) 是非零自伴紧算子，则 \(\|A\|\) 或 \(-\|A\|\) 必至少有一个是 \(A\) 的特征值。
\item
  \textbf{推论 6.2（p. 222）}：自伴紧算子 \(A\) 的谱半径 \(\Gamma(A) = \|A\|\)。
\item
  \textbf{推论 6.3（瑞利商极值，p. 223）}：非零自伴紧算子存在单位向量 \(w\)（\(\|w\|=1\)）达到极大值 \(|(Aw, w)| = \sup_{\|x\|\le 1}|(Ax, x)| = \|A\|\)，且该 \(w\) 必为对应于 \(\pm\|A\|\) 的特征向量。
\item
  \textbf{定理 6.13（Hilbert-Schmidt 谱展开定理，p. 223）}：设 \(A\) 是无穷维 Hilbert 空间 \(H\)（教材原文写``实''，但对复 Hilbert 空间同样完全成立）上的自伴紧算子，则 \(A\) 的非零谱点都是实特征值且至多可数。按绝对值从大到小排列（\(m\) 重特征值重复写 \(m\) 次），若有无穷多个则 \(\lim_{n\to\infty} |\lambda_n| = 0\)。进一步，存在对应于 \(\lambda_n\) 的单位正交特征向量列 \(\{e_n\}\) 构成 \(\overline{\mathcal{R}(A)}\) 的标准正交基（若补上 \(\operatorname{Ker} A\) 的标准正交基 \(\{v_j\}\)，则并集是整个 \(H\) 的标准正交基），且对任意 \(x \in H\) 有展开 
\FitDisplayMath{Ax = \sum_{n=1}^{\infty} \lambda_n (x, e_n) e_n.}

\item
  \textbf{推论 6.4（投影级数形式，p. 224）}：记 \(P_n\) 为 \(H\) 到 \(\operatorname{span}\{e_n\}\) 的正交投影，则 \(A = \sum_{n=1}^\infty \lambda_n P_n\)（不仅强收敛，事实上由于 \(|\lambda_n| \to 0\)，它按算子范数也一致收敛！）。
\end{itemize}

\SubSubHeading{4. Hilbert-Schmidt 算子与绝对值算子 \(|A| = (A^*A)^{1/2}\)（第 5.2 节 p. 196--197；第 6.3.2 节 p. 230--231）}

\begin{quote}
\textbf{作业必备衔接}：习题 6 第 24、31、37 题直接用到本段知识！
\end{quote}

\begin{itemize}
\tightlist
\item
  \textbf{定义 5.4（Hilbert-Schmidt 算子，p. 196）}：设 \(H, G\) 为可分 Hilbert 空间，\(A : H \to G\) 是线性算子。若存在 \(H\) 的标准正交基 \(\{e_n\}_{n=1}^\infty\) 使得 
\FitDisplayMath{\|A\|_2 \triangleq \left(\sum_{n=1}^\infty \|Ae_n\|^2\right)^{\frac{1}{2}} < \infty,}
 则称 \(A\) 是 \(H\) 到 \(G\) 的 \textbf{Hilbert-Schmidt 算子}，全体记作 \(\mathcal{L}^2(H, G)\)（当 \(H=G\) 时简记 \(\mathcal{L}^2(H)\)）。由 Parseval 等式（式 (5.8)），\(\|A\|_2\) \textbf{与标准正交基的选取无关}。
\item
  \textbf{定理 5.12（Hilbert-Schmidt 算子的核心性质，p. 197）}：若 \(A \in \mathcal{L}^2(H, G)\)，则：

  \begin{enumerate}
  \def\labelenumi{\arabic{enumi}.}
  \tightlist
  \item
    \(\|A\|_2 = \|A^*\|_2\)；
  \item
    \(\|A\| \le \|A\|_2\)（从而自动有界）；
  \item
    理想性质：对 \(B_1 \in \mathcal{L}(H), B_2 \in \mathcal{L}(G)\)，有 \(\|B_2 A\|_2 \le \|B_2\|\|A\|_2\)，\(\|AB_1\|_2 \le \|B_1\|\|A\|_2\)；
  \item
    \textbf{核心结论 (5)}：Hilbert-Schmidt 算子\textbf{必是紧算子}（\(A \in \mathcal{C}(H, G)\)）！
  \end{enumerate}

  \begin{itemize}
  \tightlist
  \item
    \textbf{积分算子快捷判据（例 5.4 / 例 5.7）}：若 \((Af)(x) = \int_a^b K(x, t)f(t)\mathrm{d}t\) 在 \(L^2(a, b)\) 上满足 \(\int_a^b\int_a^b |K(x, t)|^2 \mathrm{d}t\,\mathrm{d}x < \infty\)，则 \(A\) 是 Hilbert-Schmidt 算子，从而必是紧算子，且 \(\|A\|_2 = \|K\|_{L^2}\)！
  \end{itemize}
\item
  \textbf{绝对值算子 \(|A| = (A^*A)^{\frac{1}{2}}\) 与奇异值（6.3.2 节，p. 230--231）}：

  \begin{itemize}
  \tightlist
  \item
    对任意紧正算子 \(T \ge 0\)（由定理 6.13 有 \(Tu = \sum \lambda_n (u, e_n)e_n\)，\(\lambda_n \ge 0\)），直接定义 \(T^{\frac{1}{2}} u = \sum \lambda_n^{\frac{1}{2}}(u, e_n)e_n\)（对一般有界正算子，由 6.6 节推论 6.17 用谱积分定义 \(T^{\frac{1}{2}} = \int_{\sigma(T)} \sqrt{\lambda}\,\mathrm{d}E_\lambda\)，存在且唯一）。
  \item
    对任意 \(A \in \mathcal{L}(H, G)\)，记 \(|A| \triangleq (A^*A)^{\frac{1}{2}} \in \mathcal{L}(H)\)，则 \(|A| \ge 0\)，\(|A|^2 = A^*A\)，且对一切 \(x \in H\) 满足最核心的等距恒等式： 
\FitDisplayMath{\lVert\,|A|x\,\rVert^2 = (|A|^2 x, x) = (A^*Ax, x) = \lVert Ax\rVert^2 \implies \lVert\,|A|x\,\rVert = \lVert Ax\rVert.}

  \item
    当 \(A\) 是紧算子时，\(|A|\) 的特征值 \(\lambda_n\) 就是 \(A\) 的\textbf{奇异值}；\(\sum \lambda_n < \infty\) 对应\textbf{迹类算子} \(\mathcal{L}^1(H)\)（定义 6.5、定理 6.14），\(\sum \lambda_n^2 < \infty\) 对应 \textbf{Hilbert-Schmidt 算子} \(\mathcal{L}^2(H, G)\)（习题 31、37）。
  \end{itemize}
\end{itemize}

\SoftDivider

\BookSubSection{sec:6-4}{6.4 谱测度、谱系与谱积分（p. 237--247）}{6.4 谱测度、谱系与谱积分（p. 237–247）}{\TagLight{老师减负·大概了解}}

\SubSubHeading{零基础全景导读：正交投影、谱测度，以及第 6 章后半部分到底在干什么？}

\begin{quote}
\textbf{先看整部电影的剧情线（6.4 \(\to\) 6.5 \(\to\) 6.6 三步``迂回包抄战''）}： 第 6 章的终极梦想只有一件事——把线性代数里``实对称矩阵正交对角化 \(A = \sum \lambda_k P_k\)''推广到\textbf{任意有界自伴算子}。前半章（6.3 节）靠``紧算子''这根拐杖做到了离散求和 \(A = \sum_{n=1}^\infty \lambda_n P_n\)；但一扔掉紧性就卡死了——例如 \(L^2(0,1)\) 上的乘法算子 \((Ax)(t)=tx(t)\) \textbf{一个特征值都没有}，单点投影全是 \(0\)，根本没东西可加！于是后半章分三步破局：

\begin{enumerate}
\def\labelenumi{\arabic{enumi}.}
\tightlist
\item
  \textbf{6.4 节（先造工具箱，老师减负：大概了解）}：故意先不提具体算子 \(A\)，像造锤子一样先把``谱测度 \(E\)、谱系 \(E_\lambda\)、谱积分 \(\int f\,\mathrm{d}E\)''和 \textbf{4 条代数运算法则}造好备用；
\item
  \textbf{6.5 节（挑软柿子先攻酉算子 \(U\)）}：酉算子的谱落在单位圆周 \(\{\mathrm{e}^{\mathrm{i}\theta}\}\) 上，天然对接 Fourier 级数 / 三角多项式，容易先帮每个酉算子造出谱分解 \(U = \int_0^{2\pi} \mathrm{e}^{\mathrm{i}\theta}\,\mathrm{d}E_\theta\)；
\item
  \textbf{6.6 节（借 Cayley 变换搬回有界自伴算子 \(A\) + 附赠函数演算）}：用 \(U = (A - \mathrm{i}I)(A + \mathrm{i}I)^{-1}\) 把实轴上的 \(A\) 变成单位圆上的 \(U\)，再把 6.5 节造好的谱测度搬回实轴，立刻拿下任意有界自伴算子的谱分解 \(A = \int_{\sigma(A)} \lambda\,\mathrm{d}E(\lambda)\) 与 Borel 函数演算 \(f(A) = \int_{\sigma(A)} f(\lambda)\,\mathrm{d}E(\lambda)\)！
\end{enumerate}
\end{quote}

\StepHeading{先补齐地基：什么是``正交投影'' \(P^2 = P = P^*\)？}

一句话直觉：\textbf{正交投影 \(P\) 就是``拿手电筒垂直照下去，只保留向量在某个方向/子空间上的影子，把其他方向全归零''}。拿具体向量 \(x = (3, 4, 7)\) 和 \(3\times 3\) 对角矩阵 \(A = \operatorname{diag}(2, 5, 9)\) 看一眼就彻底通透：

\begin{itemize}
\tightlist
\item
  三个坐标轴正交投影就是最朴素的 \(0\)-\(1\) 对角矩阵： 
\FitDisplayMath{P_1 = \begin{pmatrix} 1 & 0 & 0 \\ 0 & 0 & 0 \\ 0 & 0 & 0 \end{pmatrix}, \qquad P_2 = \begin{pmatrix} 0 & 0 & 0 \\ 0 & 1 & 0 \\ 0 & 0 & 0 \end{pmatrix}, \qquad P_3 = \begin{pmatrix} 0 & 0 & 0 \\ 0 & 0 & 0 \\ 0 & 0 & 1 \end{pmatrix}.}
 于是 \(P_1(3,4,7) = (3,0,0)\)（在 \(x_1\) 轴上的垂直影子），扔掉的误差 \(x - P_1 x = (0,4,7)\) 与影子内积为 \(0\)（夹角 \(90^\circ\)，故称\textbf{正交}）。
\item
  矩阵 \(A\) 其实就是把三个影子分别拉长 \(2, 5, 9\) 倍再拼回去： 
\FitDisplayMath{A = \begin{pmatrix} 2 & 0 & 0 \\ 0 & 5 & 0 \\ 0 & 0 & 9 \end{pmatrix} = 2P_1 + 5P_2 + 9P_3.}

\item
  \textbf{解码教材方程 \(P^2 = P = P^*\)}：

  \begin{enumerate}
  \def\labelenumi{\arabic{enumi}.}
  \tightlist
  \item
    \textbf{\(P^2 = P\)（拍两次 = 拍一次，决定它是``投影''）}：\((3,4,7)\) 拍到 \(x_1\) 轴得到 \((3,0,0)\)，再拍一次它已经在轴上了，纹丝不动；
  \item
    \textbf{\(P = P^*\)（自伴/对称，保证是``垂直''拍而非斜着拍）}：反例 \(Q = \begin{pmatrix} 1 & 1 \\ 0 & 0 \end{pmatrix}\) 同样满足 \(Q^2 = Q\)，但它把 \((0,1)\) \textbf{斜着}拍到 \((1,0)\) 上，影子与误差不垂直，所以 \(Q^T \neq Q\)。只有垂直拍下去才满足 \(P = P^*\)！
  \end{enumerate}
\end{itemize}

\StepHeading{用``概率论''一分钟秒懂谱测度 \(E(\Delta)\) 与算子 \(A\) 的关系}

\begin{itemize}
\tightlist
\item
  \textbf{掷骰子（离散，类比紧算子）}：每个点数 \(k\) 都有非零概率 \(p_k\)，期望直接求和 \(\sum k p_k\)。
\item
  \textbf{测身高（连续，类比一般自伴算子）}：身高\textbf{恰好}等于 \(175.000000\dots\) cm 的概率是 \textbf{\(0\)}！单点问不出信息，只能改问``落在一段区间 \(\Delta = [170, 175]\) 里的概率 \(\mu(\Delta)\)''，期望随之升级为积分 \(\int x\,\mathrm{d}\mu(x)\)。
\item
  \textbf{谱测度 \(E(\Delta)\) 完全同理，只是把``输出概率数字''换成了``输出正交投影算子''}（生活比喻：把算子 \(A\) 想成合成音乐，\(E(\Delta)\) 就是\textbf{频段滤波器}，只留 \(\Delta\) 频段的分量、其余清零）：

  \begin{itemize}
  \tightlist
  \item
    \textbf{谱测度 \(E(\Delta)\)}：输入谱区间 \(\Delta \subset \sigma(A)\)，输出该谱段对应的正交投影算子（替代线代里的单点投影 \(P_k\)）；
  \item
    \textbf{谱系 \(E_\lambda \triangleq E((-\infty, \lambda])\)}：就是概率论里的\textbf{累积分布函数 \(F(x) = \mu((-\infty, x])\)}，与谱测度一一对应；
  \item
    \textbf{谱积分 \(\int_{\sigma(A)} \lambda\,\mathrm{d}E(\lambda)\)}：把连续谱区间切成无数极细小段 \(\Delta_k\)，各取影子 \(E(\Delta_k)x\) 乘倍率 \(\lambda_k\) 拼回去取极限，即把离散求和 \(\sum \lambda_k P_k\) 升级为连续积分！
  \end{itemize}
\end{itemize}

再对照两个具体例子，彻底打通``有限维到无穷维''：

\begin{itemize}
\tightlist
\item
  \textbf{例子 A（有限维对角矩阵，有特征值）}：设 \(A = \operatorname{diag}(2, 5, 9)\)，三个坐标轴正交投影为 \(P_1, P_2, P_3\)。对任意实数集 \(\Delta \subset \mathbb{R}\)，定义 \(E(\Delta) \triangleq \sum_{\lambda_k \in \Delta} P_k\)（例如 \(E([1,3])=P_1\)，\(E([1,6])=P_1+P_2\)，\(E([6,8])=0\)，\(E(\mathbb{R})=P_1+P_2+P_3=I\)）。此时离散谱积分就是求和：\(\int \lambda\,\mathrm{d}E(\lambda) = 2P_1 + 5P_2 + 9P_3 = A\)，而任意函数演算 \(\int f(\lambda)\,\mathrm{d}E(\lambda) = f(2)P_1 + f(5)P_2 + f(9)P_3 = f(A)\)！
\item
  \textbf{例子 B（无穷维连续谱乘法算子，无特征值）}：在 \(L^2(0,1)\) 上 \((Ax)(t) = t x(t)\)，它一个特征值都没有（单点投影 \(E(\{0.5\})=0\)），谱是连续区间 \([0,1]\)。它专属的谱测度就是\textbf{``区间截断投影''}：\((E(\Delta)x)(t) = \chi_\Delta(t)x(t)\)（只保留 \(t\in\Delta\) 上的图像，其余清零）。把 \([0,1]\) 切成无数极小段 \(\Delta_k\)、各乘倍率 \(\lambda_k\) 拼回去，极限正好是 \(\left(\int_0^1 \lambda\,\mathrm{d}E(\lambda)\,x\right)(t) = t x(t) = (Ax)(t)\)！
\end{itemize}

\SubSubHeading{1. 谱测度的严格定义与性质（定义 6.7、引理 6.5，p. 237）}

设 \(\Omega\) 为集合，\(\Sigma\) 为 \(\Omega\) 上子集构成的 \(\sigma\)-代数，\(\mathcal{P}\) 为 \(H\) 上正交投影算子全体。映射 \(E : \Sigma \to \mathcal{P}\) 称为 \((\Omega, \Sigma)\) 上的\textbf{谱测度}（三元组 \((\Omega, \Sigma, E)\) 称为\textbf{谱测度空间}），若满足：

\begin{enumerate}
\def\labelenumi{\arabic{enumi}.}
\tightlist
\item
  \textbf{规范性}：\(E(\Omega) = I\)；
\item
  \textbf{可列可加性}：对两两不相交的 \(\{A_n\}_{n=1}^\infty \subset \Sigma\)，有 \(E\left(\bigcup_{n=1}^\infty A_n\right) = \sum_{n=1}^\infty E(A_n)\)（右端指强算子收敛，即对每个 \(x \in H\) 按范数收敛）。
\end{enumerate}

\textbf{引理 6.5（由定义立即导出的基本性质，p. 237）}：

\begin{itemize}
\tightlist
\item
  \begin{enumerate}
  \def\labelenumi{(\roman{enumi})}
  \tightlist
  \item
    \(E(\emptyset) = 0\)；
  \end{enumerate}
\item
  (ii) 有限可加性；
\item
  (iii) \textbf{正交性}：若 \(A \cap B = \emptyset\)，则 \(E(A)E(B) = E(B)E(A) = 0\)；
\item
  (iv) \(E(A \cup B) = E(A) + E(B) - E(A \cap B)\)；
\item
  (v) \textbf{乘交与交换性}：对任意 \(A, B \in \Sigma\)，均有 \(E(A)E(B) = E(A \cap B) = E(B)E(A)\)（整族谱投影互相可交换！）。
\end{itemize}

\begin{quote}
\textbf{连接普通测度的桥梁}：固定任意 \(x, y \in H\)，标量函数 \(\mu_{x,y}(\Delta) \triangleq (E(\Delta)x, y)\) 是一个普通的\textbf{复有限测度}；特别地，当 \(y = x\) 时，\(\mu_{x,x}(\Delta) = (E(\Delta)x, x) = \|E(\Delta)x\|^2\) 是总质量为 \(\|x\|^2\) 的\textbf{有限正测度}！
\end{quote}

\SubSubHeading{2. 有界可测函数的谱积分及其四大法则（定义 6.8--6.9、定理 6.15--6.16、推论 6.5--6.8，p. 239--243）}

记 \(B(\Omega, \Sigma)\) 为 \((\Omega, \Sigma)\) 上全体有界可测函数，配一致范数 \(\|f\|_\infty = \sup_{\omega \in \Omega}|f(\omega)|\)（教材简记 \(\|f\|\)）。

\begin{itemize}
\tightlist
\item
  \textbf{一致谱积分（定义 6.8，p. 239）}：把实值函数 \(f\) 的值域划分成小区间 \([y_{j-1}, y_j)\)，作 Riemann-Stieltjes 型算子和 \(S(D) = \sum_{j=1}^n \xi_j E(\{y_{j-1} \le f(\omega) < y_j\})\)，若分划细度 \(\delta(D) \to 0\) 时 \(S(D)\) 按\textbf{算子范数}收敛，则极限记为 \(\int_\Omega f(\omega)\,\mathrm{d}E(\omega)\)。
\item
  \textbf{弱谱积分（定义 6.9，p. 239）}：满足对一切 \(x, y \in H\) 都有 
\FitDisplayMath{\left(\left(\mathrm{w\text{-}}\int_\Omega f(\omega)\,\mathrm{d}E(\omega)\right)x, y\right) = \int_\Omega f(\omega)\,\mathrm{d}(E(\omega)x, y)}
 的有界线性算子（右端是关于普通复测度 \(\mu_{x,y}\) 的 Lebesgue 积分）。
\item
  \textbf{存在唯一与等价定理（定理 6.15--6.16，p. 239--241）}：对任意 \(f \in B(\Omega, \Sigma)\)，弱谱积分存在唯一，且\textbf{弱谱积分与一致谱积分完全相等}，今后统称为\textbf{谱积分} \(\int_\Omega f(\omega)\,\mathrm{d}E(\omega)\)。
\end{itemize}

\textbf{期末必背！谱积分的四大运算律（推论 6.5--6.8，p. 240--243）}：对任意 \(f, g \in B(\Omega, \Sigma)\) 及 \(\alpha, \beta \in \mathbb{C}\)，

{\def\LTcaptype{none} % do not increment counter
{\small\renewcommand{\arraystretch}{1.38}\setlength{\tabcolsep}{5pt}
\begin{longtable}[]{@{}
  >{\raggedright\arraybackslash}p{(\linewidth - 4\tabcolsep) * \real{0.2097}}
  >{\raggedright\arraybackslash}p{(\linewidth - 4\tabcolsep) * \real{0.5645}}
  >{\centering\arraybackslash}p{(\linewidth - 4\tabcolsep) * \real{0.2258}}@{}}
\toprule\noalign{}
\rowcolor{AccentSoft}
\begin{minipage}[b]{\linewidth}\raggedright
性质
\end{minipage} & \begin{minipage}[b]{\linewidth}\raggedright
公式
\end{minipage} & \begin{minipage}[b]{\linewidth}\centering
教材出处
\end{minipage} \\
\midrule\noalign{}
\endhead
\bottomrule\noalign{}
\endlastfoot
\textbf{线性性} & \(\displaystyle \int_\Omega (\alpha f + \beta g)\,\mathrm{d}E = \alpha \int_\Omega f\,\mathrm{d}E + \beta \int_\Omega g\,\mathrm{d}E\)，且 \(\displaystyle \int_\Omega \chi_A \,\mathrm{d}E = E(A)\) & 推论 6.5(i)(iii) \\
\textbf{共轭 / Hermite 性} & \(\displaystyle \left(\int_\Omega f(\omega)\,\mathrm{d}E(\omega)\right)^* = \int_\Omega \overline{f(\omega)}\,\mathrm{d}E(\omega)\)（实值 \(f \implies\) 自伴算子） & 推论 6.5(ii) \\
\textbf{乘法同态（最神奇！）} & \(\displaystyle \int_\Omega f g\,\mathrm{d}E = \left(\int_\Omega f\,\mathrm{d}E\right)\left(\int_\Omega g\,\mathrm{d}E\right)\)（且互相可交换） & 推论 6.6--6.7 \\
\textbf{范数估计} & \(\displaystyle \left\lVert\int_\Omega f(\omega)\,\mathrm{d}E(\omega)\right\rVert \le \sup_{\omega \in \Omega} \lvert f(\omega) \rvert\) & 推论 6.8 \\
\end{longtable}
}
}

\begin{quote}
\textbf{为什么``乘法同态''对普通测度不成立，对谱测度却成立？（注记 6.3）} 因为普通数字乘法没有``不交相乘为零''，但谱测度的投影满足 \(E(A)E(B) = E(A \cap B)\)！对简单函数相乘，交叉项自动坍缩成交集投影，取一致极限就得到了``积的积分 = 积分的积''。另外由推论 6.7 的式 (6.59) 立即有极好用的范数等式： 
\FitDisplayMath{\left\|\left(\int_\Omega f\,\mathrm{d}E\right)x\right\|^2 = \int_\Omega |f(\omega)|^2 \,\mathrm{d}\|E(\omega)x\|^2.}

\end{quote}

\SubSubHeading{3. 谱系与谱测度的一一对应（定义 6.10、定理 6.19，p. 243--247）}

正如实轴上概率分布函数 \(F(x) = \mu((-\infty, x])\) 与概率测度 \(\mu\) 一一对应，\textbf{谱系}就是投影值版的``分布函数''：

\begin{itemize}
\tightlist
\item
  \textbf{定义 6.10（谱系，p. 243）}：实参数正交投影族 \(\{E_\lambda\}_{\lambda \in \mathbb{R}}\) 称为 \(H\) 上的\textbf{谱系}，若满足：

  \begin{enumerate}
  \def\labelenumi{\arabic{enumi}.}
  \tightlist
  \item
    \textbf{单调性}：\(\lambda \ge \mu \implies E_\lambda \ge E_\mu\)（即 \(E_\lambda E_\mu = E_\mu E_\lambda = E_\mu\)）；
  \item
    \textbf{强右连续性}：对任意 \(\lambda_0 \in \mathbb{R}\)，\((\text{强})\lim_{\lambda \to \lambda_0 + 0} E_\lambda = E_{\lambda_0}\)；
  \item
    \textbf{两端极限}：\((\text{强})\lim_{\lambda \to -\infty} E_\lambda = 0\)，\((\text{强})\lim_{\lambda \to +\infty} E_\lambda = I\)。
  \end{enumerate}
\item
  \textbf{对应关系（定理 6.19，p. 247）}：由谱系 \(\{E_\lambda\}\) 可唯一导出 \((\mathbb{R}, \mathcal{B}_{\mathbb{R}})\) 上的谱测度 \(E\)，满足 \(E((a, b]) = E_b - E_a\)、\(E((a, b)) = E_{b-0} - E_a\)，从而 \(\int_{\mathbb{R}} f(\lambda)\,\mathrm{d}E(\lambda)\) 也常写成 \(\int_{\mathbb{R}} f(\lambda)\,\mathrm{d}E_\lambda\)。
\end{itemize}

\SoftDivider

\BookSubSection{sec:6-5}{6.5--6.6 酉算子与有界自伴算子的谱分解及 Borel 函数演算（p. 247--265）}{6.5–6.6 酉算子与有界自伴算子的谱分解及 Borel 函数演算（p. 247–265）}{\TagMust{闭卷必背}}

\SubSubHeading{1. 保距算子与酉算子（定义 6.11--6.12，p. 247）}

\begin{itemize}
\tightlist
\item
  \textbf{定义 6.11（保距算子）}：\(U \in \mathcal{L}(H_1, H_2)\) 满足对一切 \(x \in H_1\) 有 \(\|Ux\| = \|x\|\)。
\item
  \textbf{定义 6.12（酉算子）}：\(U\) 是 \(H_1\) 到 \(H_2\) 的\textbf{保距满射}（\(H_1 = H_2 = H\) 时简称 \(H\) 上的酉算子）。
\item
  \textbf{等价刻画（期末最常用！）}：\(U \in \mathcal{L}(H)\) 是酉算子 \(\iff U^* U = U U^* = I \iff U\) 可逆且 \(U^{-1} = U^* \iff U\) 满射且保内积 \((Ux, Uy) = (x, y)\)。
\item
  \textbf{酉算子的谱特征}：\(\sigma(U)\) 必包含在复平面单位圆周 \(\{z \in \mathbb{C} : |z| = 1\}\) 上（推论 6.10）。
\item
  \textbf{定理 6.21（酉算子的谱分解定理，p. 256）}：对复 Hilbert 空间 \(H\) 上任意酉算子 \(U\)，存在 \((\sigma(U), \mathcal{B}_{\sigma(U)})\) 上\textbf{唯一}的谱测度 \(F\)，使得 
\FitDisplayMath{U = \int_{\sigma(U)} \lambda \,\mathrm{d}F(\lambda),}
 且与 \(U\) 可交换的任意有界算子都与所有谱投影 \(F(M)\) 可交换。由定理 6.22，\(f(U) = \int_{\sigma(U)} f(\lambda)\,\mathrm{d}F(\lambda)\) 构成 \(U\) 的算子演算。
\end{itemize}

\SubSubHeading{2. Cayley 变换与有界自伴算子的谱分解（定义 6.15、引理 6.8--6.10、定理 6.24--6.25，p. 259--263）}

\textbf{如何把``单位圆上的酉算子谱定理''搬运到``实轴上的自伴算子''？} 复变函数里的分式线性变换 \(\varphi(z) = \frac{z - \mathrm{i}}{z + \mathrm{i}}\) 把整条实轴 \(\mathbb{R}\) 完美地一一双射到单位圆周去掉点 \(1\)：\(\{w : |w|=1, w \neq 1\}\)！

\begin{itemize}
\tightlist
\item
  \textbf{定义 6.15（Cayley 变换，p. 259）}：设 \(A \in \mathcal{L}(H)\) 是有界自伴算子，定义 
\FitDisplayMath{U = (A - \mathrm{i}I)(A + \mathrm{i}I)^{-1}}
 为 \(A\) 的 Cayley 变换。则 \(U\) 是酉算子（习题 23 的 \(\lambda = -\mathrm{i}\) 特例！）。
\item
  \textbf{引理 6.8--6.9（p. 259, 261）}：\(1\) 不是 \(U\) 的特征值，反演公式为 \(A = \mathrm{i}(I + U)(I - U)^{-1}\)，且 \(\varphi(\sigma(A)) = \sigma(U) \setminus \{1\}\)。
\item
  \textbf{定理 6.24 \& 定理 6.25（有界自伴算子的谱分解定理，p. 259, 262）}：设 \(A\) 是复 Hilbert 空间 \(H\) 上的有界自伴算子，则存在唯一集中在 \(\sigma(A)\) 上（即 \(E(\sigma(A)) = I\)，且不能集中在比 \(\sigma(A)\) 更小的闭集上）的谱测度 \(E\)（对应谱系 \(\{E_\lambda\}\)），使得 
\FitDisplayMath{A = \int_{\mathbb{R}} \lambda \,\mathrm{d}E_\lambda = \int_{\sigma(A)} \lambda \,\mathrm{d}E_\lambda.}

\item
  \textbf{推论 6.14（可交换性传递，p. 263）}：若 \(B \in \mathcal{L}(H)\) 与 \(A\) 可交换（\(BA = AB\)），则 \(B\) 与 \(A\) 的所有谱投影 \(E(M)\)（\(M \in \mathcal{B}_{\sigma(A)}\)）都可交换。
\end{itemize}

\SubSubHeading{3. 有界自伴算子的 Borel 函数演算（推论 6.15--6.17，p. 264）}

设 \(A \in \mathcal{L}(H)\) 自伴，\(E\) 是其谱测度。对任意 \(\sigma(A)\) 上的有界 Borel 可测函数 \(f \in \mathcal{B}(\sigma(A))\)，定义算子 
\FitDisplayMath{f(A) \triangleq \int_{\sigma(A)} f(\lambda)\,\mathrm{d}E(\lambda) = \int_{\sigma(A)} f(\lambda)\,\mathrm{d}E_\lambda.}


\begin{quote}
\textbf{记号防混指南：上面推论 6.14 里的 \(\mathcal{B}_{\sigma(A)}\) 和这里的 \(\mathcal{B}(\sigma(A))\) 有什么区别？} 两个符号长相几乎一样，但\textbf{一个装``集合''，一个装``函数''}：

\begin{itemize}
\tightlist
\item
  \textbf{带下标 \(\mathcal{B}_{\sigma(A)}\)（装集合 \(M\)）}：谱集 \(\sigma(A)\) 上的 \textbf{Borel \(\sigma\)-代数}——由 \(\sigma(A)\) 里的区间/开集经可数次并、交、补拼成的全部``正常子集''（例如对 \(\sigma(A)=\{2,5,9\}\) 就是全部 \(8\) 个子集；对 \([0,1]\) 包含任意区间 \([0.2,0.7]\)）。\(M \in \mathcal{B}_{\sigma(A)}\) 是喂给谱测度 \(E(M)\) 的\textbf{输入集合}。
\item
  \textbf{带括号 \(\mathcal{B}(\sigma(A))\)（装函数 \(f\)，亦记正体 \(B(\sigma(A))\)）}：\(\sigma(A)\) 上的全体\textbf{有界 Borel 可测函数}——必须同时\textbf{有界}（\(\sup_{\lambda\in\sigma(A)}|f(\lambda)|<\infty\)，否则算子无界）且 \textbf{Borel 可测}（几乎所有能写出来的有界函数都在内：多项式、幂级数、连续函数 \(\sqrt{\lambda}, \mathrm{e}^{\mathrm{i}\lambda}\)、分段函数和示性函数 \(\chi_M\) 全算；像 \(\frac{1}{\lambda}\) 在 \([0,1]\) 上无界就不在内）。\(f \in \mathcal{B}(\sigma(A))\) 是喂给谱积分 \(f(A)\) 的\textbf{被积函数}。
\end{itemize}
\end{quote}

\begin{itemize}
\item
  \textbf{推论 6.15（函数演算的代数同态性质，p. 264）}：

  \begin{enumerate}
  \def\labelenumi{\arabic{enumi}.}
  \tightlist
  \item
    \textbf{Hermite 性}：\(\bar{f}(A) = (f(A))^*\)（特别地，\(f\) 取实值时 \(f(A)\) 自伴）；
  \item
    \textbf{线性性}：\((\alpha f + \beta g)(A) = \alpha f(A) + \beta g(A)\)（注：教材 p. 264 印刷中误把等号左边的 \(\alpha\) 印成了 \(\lambda\)）；
  \item
    \textbf{可乘性}：\((fg)(A) = f(A)g(A)\)；
  \item
    \textbf{规范性}：\(f \equiv 1 \implies f(A) = I\)；\(f(\lambda) \equiv \lambda \implies f(A) = A\)（从而对多项式 \(p(\lambda) = \sum c_k \lambda^k\)，自然有 \(p(A) = \sum c_k A^k\)）。
  \end{enumerate}
\item
  \textbf{推论 6.16（范数估计，p. 264）}： 
\FitDisplayMath{\|f(A)\| \le \sup_{\lambda \in \sigma(A)} |f(\lambda)|.}


  \begin{quote}
  \textbf{极大威力}：如果函数列 \(f_n(\lambda)\) 在 \(\sigma(A)\) 上\textbf{一致收敛}于 \(f(\lambda)\)，那么 \(\|f_n(A) - f(A)\| \le \sup_{\lambda \in \sigma(A)} |f_n(\lambda) - f(\lambda)| \to 0\)，即算子列 \(f_n(A)\) 自动按\textbf{算子范数}收敛于 \(f(A)\)！（习题 45 就是直接用这一步！）
  \end{quote}
\item
  \textbf{推论 6.17（正算子唯一正平方根定理，p. 264）}：若 \(A \in \mathcal{L}(H)\) 满足 \(A \ge 0\)，则 \(\sigma(A) \subset [0, \infty)\)，且存在\textbf{唯一}的有界正算子 \(A_1 \ge 0\) 满足 \(A_1^2 = A\)，记作 \(A^{\frac{1}{2}} = \int_{\sigma(A)} \sqrt{\lambda}\,\mathrm{d}E_\lambda\)。
\end{itemize}

\SoftDivider

\BookSubSection{sec:6-tips}{考前速查：第 6 章期末大题``五大万能解题套路''}{考前速查：第 6 章期末大题“五大万能解题套路”}{\TagHigh{高频思路}}

\begin{enumerate}
\def\labelenumi{\arabic{enumi}.}
\tightlist
\item
  \textbf{证明 \(\lambda \in \rho(A)\)}：直接构造有界逆算子（如乘法算子做除法）；或借已知正则点变形为 \(T(I - B)\)，令 \(\|B\| < 1\) 用 Neumann 级数（第 2、5 题）。
\item
  \textbf{证明 \(\lambda \in \sigma(A)\)}：不必强求找特征向量！构造单位向量列 \(\{f_n\}\)（\(\|f_n\|=1\)）使 \(\|(\lambda I - A)f_n\| \to 0\)，直接推翻``有界逆必下有界''即可通吃点谱与连续谱（第 2 题）。
\item
  \textbf{求非自伴紧算子（如积分算子）的范数 \(\|A\|\)}：千万别去算 \(\Gamma(A)\)！牢记两步连招： 
\FitDisplayMath{\|A\|^2 = \|AA^*\| = \|A^*A\| = \Gamma(AA^*) = \text{紧正自伴算子 } AA^* \text{ 的最大特征值}}
 再把特征方程 \(AA^*f = \lambda f\) 求导化为常微分方程边值问题求解（第 24 题）。
\item
  \textbf{处理 \(|A| = (A^*A)^{1/2}\) 与 Hilbert-Schmidt 算子}：抓住黄金等式 \(\lVert\,|A|x\,\rVert = \lVert Ax\rVert\)，瞬间把一般算子 \(A\) 的范数/级数问题转化成\textbf{紧正自伴算子 \(|A|\)} 的特征值问题（第 31、37 题）。
\item
  \textbf{证明 Borel 函数演算的等式}：

  \begin{itemize}
  \tightlist
  \item
    涉及酉共轭/谱测度变换：先验证构造的新投影族满足谱测度公理（定义 6.7），再用弱积分内积核对，最后用谱测度唯一性收尾（第 39 题）；
  \item
    涉及级数/极限：先用多项式算清楚部分和 \(S_N(A)\)，再用范数估计 \(\|f(A)-S_N(A)\| \le \sup_{\sigma(A)}|f - S_N| \to 0\) 一步到位（第 45 题）。
  \end{itemize}
\end{enumerate}

\BookPart{}{第二部分：习题 6 课后作业逐题详解（共 9 题）}{第二部分：习题 6 课后作业逐题详解（共 9 题）}

\begin{quote}
\textbf{习题 6 统一约定（教材 p. 266）}：如无专门强调，以下用 \(H\) 和 \(X\) 分别表示复数域 \(\mathbb{C}\) 上的 Hilbert 空间和 Banach 空间。
\end{quote}

\SoftDivider

\BookProblem{prob:6-2}{第 2 题}{教材 p. 266}{第 2 题}

\SubPillPrimary{【原题】}

设 \(a \in C([0,1];\mathbb{R})\) 上实值连续函数\footnote{\textbf{教材表述注记}：原书``设 \(a \in C([0,1];\mathbb{R})\) 上实值连续函数''略有语病，意即``设 \(a \in C([0,1];\mathbb{R})\) 为 \([0,1]\) 上的实值连续函数''。}，定义 \(L^2(0,1)\) 上的算子如下： 
\FitDisplayMath{(Af)(x) = a(x)f(x), \quad \forall f \in L^2(0,1).}
 证明 \(A\) 是 \(L^2(0,1)\) 上自伴算子，并求 \(\sigma(A)\)。

\SubPillAmber{【零基础破题思路】}

这道题把乘法算子 \((Af)(x) = x f(x)\)（教材例 6.4）推广成了乘以任意实连续函数 \(a(x)\)，需要分三步走：

\begin{enumerate}
\def\labelenumi{\arabic{enumi}.}
\tightlist
\item
  \textbf{证有界自伴}：有界线性不能漏写！因为 \(a\) 在闭区间 \([0,1]\) 连续，所以有界（\(|a(x)| \le \|a\|_\infty\)），立刻得 \(A\) 有界；因为 \(a(x)\) 取实值，\(\overline{a(x)} = a(x)\)，直接把 \(a(x)\) 从内积第一项搬到第二项就证完自伴。
\item
  \textbf{猜出 \(\sigma(A)\)}：解方程 \((\lambda I - A)f = g\) 就是 \(f(x) = \frac{g(x)}{\lambda - a(x)}\)。什么时候除法会``炸掉''（分母为 \(0\)）？恰好是 \(\lambda\) 落在 \(a(x)\) 的值域 \(a([0,1])\) 里的时候！所以合理猜测 \(\sigma(A) = a([0,1])\)。
\item
  \textbf{双向证明 \(\sigma(A) = a([0,1])\)}：

  \begin{itemize}
  \tightlist
  \item
    当 \(\lambda \notin a([0,1])\) 时，因为连续函数把紧集 \([0,1]\) 映成紧集（闭集），\(\lambda\) 与 \(a([0,1])\) 保持正距离 \(\delta > 0\)，直接定义除法算子验证它是有界逆，即 \(\lambda \in \rho(A)\)；
  \item
    当 \(\lambda \in a([0,1])\) 时，\textbf{不要试图找特征函数}（如果 \(a(x)=x\)，根本没有特征函数！）。正确招式是：找一点 \(x_0\) 使 \(a(x_0)=\lambda\)，在 \(x_0\) 极小的邻域区间 \(I_n\) 上 \(|a(x)-\lambda| < \frac{1}{n}\)，把单位 \(L^2\)-范数的``窄脉冲'' \(f_n = |I_n|^{-1/2}\chi_{I_n}\) 送进 \(\lambda I - A\)，结果被压得趋于 \(0\)，这就彻底打破了``有界逆算子必然下有界''的要求，从而 \(\lambda \in \sigma(A)\)！
  \end{itemize}
\end{enumerate}

\SubPillMain{【规范解答】}

\textbf{第一步：证明 \(A\) 是 \(L^2(0,1)\) 上的有界自伴算子。}

显然 \(A\) 是线性算子。因为 \(a \in C([0,1];\mathbb{R})\) 在闭区间 \([0,1]\) 上连续，故有界，记 \(M \triangleq \max_{x \in [0,1]} |a(x)| < \infty\)。对任意 \(f \in L^2(0,1)\)，\(a f\) 可测且 
\FitDisplayMath{\|Af\|_2^2 = \int_0^1 |a(x)f(x)|^2 \,\mathrm{d}x \le M^2 \int_0^1 |f(x)|^2 \,\mathrm{d}x = M^2 \|f\|_2^2,}
 所以 \(Af \in L^2(0,1)\) 且 \(\|A\| \le M < \infty\)，即 \(A \in \mathcal{L}(L^2(0,1))\)。

又因为 \(a(x)\) 是实值函数，有 \(\overline{a(x)} = a(x)\)，从而对任意 \(f, g \in L^2(0,1)\)， 
\FitDisplayMath{(Af, g) = \int_0^1 a(x)f(x)\overline{g(x)}\,\mathrm{d}x = \int_0^1 f(x)\overline{a(x)g(x)}\,\mathrm{d}x = (f, Ag).}
 因此 \(A^* = A\)，即 \(A\) 是 \(L^2(0,1)\) 上的有界自伴算子。

\SoftDivider

\textbf{第二步：证明 \(\sigma(A) = a([0,1])\)。}

\textbf{(1) 证明 \(\sigma(A) \subset a([0,1])\)（等价于 \(\lambda \notin a([0,1]) \implies \lambda \in \rho(A)\)）。}

设 \(\lambda \in \mathbb{C} \setminus a([0,1])\)。由于 \([0,1]\) 是紧集且 \(a\) 连续，像集 \(a([0,1])\) 是 \(\mathbb{R}\)（从而也是 \(\mathbb{C}\)）中的紧集，故为闭集。因此 \(\lambda\) 到 \(a([0,1])\) 的距离为正，即存在 \(\delta > 0\) 使得 
\FitDisplayMath{|\lambda - a(x)| \ge \delta, \quad \forall x \in [0,1].}
 定义算子 \(B : L^2(0,1) \to L^2(0,1)\) 为 
\FitDisplayMath{(Bg)(x) \triangleq \frac{g(x)}{\lambda - a(x)}, \quad \forall g \in L^2(0,1).}
 由于 \(\left|\frac{1}{\lambda - a(x)}\right| \le \frac{1}{\delta}\)，对任意 \(g \in L^2(0,1)\) 有 \(\|Bg\|_2 \le \frac{1}{\delta}\|g\|_2 < \infty\)，故 \(B \in \mathcal{L}(L^2(0,1))\)。直接验证可知，对任意 \(f, g \in L^2(0,1)\) 及几乎处处的 \(x \in (0,1)\)， 
\FitDisplayMath{\bigl((\lambda I - A)Bg\bigr)(x) = (\lambda - a(x))\cdot \frac{g(x)}{\lambda - a(x)} = g(x), \qquad \bigl(B(\lambda I - A)f\bigr)(x) = \frac{(\lambda - a(x))f(x)}{\lambda - a(x)} = f(x).}
 即 \((\lambda I - A)B = B(\lambda I - A) = I\)，说明 \(\lambda I - A\) 是 \(L^2(0,1)\) 上的双射且 \((\lambda I - A)^{-1} = B \in \mathcal{L}(L^2(0,1))\)，故 \(\lambda \in \rho(A)\)。由此得 \(\sigma(A) \subset a([0,1])\)。

\textbf{(2) 证明 \(a([0,1]) \subset \sigma(A)\)（等价于 \(\lambda \in a([0,1]) \implies \lambda \in \sigma(A)\)）。}

任取 \(\lambda \in a([0,1])\)，则存在 \(x_0 \in [0,1]\) 使得 \(a(x_0) = \lambda\)。用反证法：假设 \(\lambda \in \rho(A)\)，由 Banach 逆算子定理（注记 6.1），\((\lambda I - A)^{-1} \in \mathcal{L}(L^2(0,1))\)。记 \(C \triangleq \|(\lambda I - A)^{-1}\| > 0\)，则对一切 \(f \in L^2(0,1)\) 必有 
\FitDisplayMath{\|f\|_2 = \bigl\|(\lambda I - A)^{-1}(\lambda I - A)f\bigr\|_2 \le C \|(\lambda I - A)f\|_2 \implies \|(\lambda I - A)f\|_2 \ge \frac{1}{C}\|f\|_2.}
 另一方面，由 \(a\) 在 \(x_0\) 处的连续性，对任意 \(n \in \mathbb{N}\)，存在 \(\delta_n > 0\)，使得对一切满足 \(x \in [0,1]\) 且 \(|x - x_0| < \delta_n\) 的 \(x\)，都有 
\FitDisplayMath{|\lambda - a(x)| = |a(x_0) - a(x)| < \frac{1}{n}.}
 令 \(I_n \triangleq (x_0 - \delta_n, x_0 + \delta_n) \cap (0, 1)\)。无论 \(x_0\) 位于 \((0,1)\) 内部还是端点 \(0, 1\)，\(I_n\) 都是 \((0,1)\) 中长度 \(|I_n| > 0\) 的非空开区间。构造归一化特征函数 
\FitDisplayMath{f_n(x) \triangleq \frac{1}{\sqrt{|I_n|}}\chi_{I_n}(x) \in L^2(0,1).}
 容易计算得 \(\|f_n\|_2^2 = \frac{1}{|I_n|}\int_{I_n} 1\,\mathrm{d}x = 1\)，且 
\FitDisplayMath{\|(\lambda I - A)f_n\|_2^2 = \frac{1}{|I_n|}\int_{I_n} |\lambda - a(x)|^2\,\mathrm{d}x \le \frac{1}{|I_n|}\cdot \frac{1}{n^2}\cdot |I_n| = \frac{1}{n^2},}
 即 \(\|(\lambda I - A)f_n\|_2 \le \frac{1}{n}\)。于是当 \(n \to \infty\) 时 \(\|(\lambda I - A)f_n\|_2 \to 0\)，与 \(\|(\lambda I - A)f_n\|_2 \ge \frac{1}{C}\|f_n\|_2 = \frac{1}{C} > 0\) 矛盾！

因此 \(\lambda \notin \rho(A)\)，即 \(\lambda \in \sigma(A)\)。结合 (1)(2) 得 
\FitDisplayMath{\sigma(A) = a([0,1]) = \left[\min_{x\in[0,1]} a(x),\; \max_{x\in[0,1]} a(x)\right]. \qquad \square}


\SubPillCoral{【考点复盘与避坑提示】}

\begin{itemize}
\tightlist
\item
  \textbf{常见失分点 1（漏证有界性）}：定义自伴算子前必须先说明 \(A \in \mathcal{L}(L^2(0,1))\)（即 \(\|Af\|_2 \le \|a\|_\infty \|f\|_2\)），否则不能直接套用有界自伴算子的概念。
\item
  \textbf{常见失分点 2（误把端点处 \(I_n\) 取错）}：写邻域时务必写成 \(I_n = (x_0 - \delta_n, x_0 + \delta_n) \cap (0,1)\)，不能直接假设 \((x_0 - \delta_n, x_0 + \delta_n) \subset (0,1)\)，因为 \(x_0\) 可能是端点 \(0\) 或 \(1\)！
\item
  \textbf{方法升华}：上面构造``单位列 \(f_n\) 使 \(\|(\lambda I - A)f_n\| \to 0\)''的方法叫\textbf{近似点谱法}，它同时覆盖点谱（当 \(a\) 在某正测度子集上取常值 \(\lambda\) 时）和连续谱（当 \(\{x: a(x)=\lambda\}\) 测度为 \(0\) 时），比分别讨论 \(\sigma_p\) 和 \(\sigma_c\) 简洁一倍且绝不漏情况。
\end{itemize}

\SoftDivider

\BookProblem{prob:6-5}{第 5 题}{教材 p. 266}{第 5 题}

\SubPillPrimary{【原题】}

设 \(A \in \mathcal{L}(X)\)，\(\lambda_0 \in \rho(A)\)。又设 \(\{A_n\}_{n=1}^\infty \subset \mathcal{L}(X)\) 满足 \(\lim_{n\to\infty}\|A - A_n\| = 0\)。证明：存在 \(N > 0\) 使得对所有 \(n > N\)，\(\lambda_0 \in \rho(A_n)\) 而且 
\FitDisplayMath{\lim_{n\to\infty} \|(\lambda_0 I - A_n)^{-1} - (\lambda_0 I - A)^{-1}\| = 0.}


\SubPillAmber{【零基础破题思路】}

这道题证明了``算子稍微动一点点，原来的正则点依然是正则点，而且逆算子也只动一点点''（即求逆映射在可逆算子开集上连续）。 怎么把未知的 \(\lambda_0 I - A_n\) 和已知可逆的 \(\lambda_0 I - A\) 联系起来？\textbf{强行提取公因式}！ 
\FitDisplayMath{\lambda_0 I - A_n = (\lambda_0 I - A) - (A_n - A) = (\lambda_0 I - A)\bigl[I - (\lambda_0 I - A)^{-1}(A_n - A)\bigr].}
 设 \(B_n = (\lambda_0 I - A)^{-1}(A_n - A)\)。只要 \(n\) 足够大，\(\|A_n - A\|\) 足够小，就能把 \(\|B_n\|\) 压到严格小于 \(1\)。一旦 \(\|B_n\| < 1\)，立刻祭出 \textbf{Neumann 级数（定理 6.1）} 算出 \((I - B_n)^{-1} = \sum_{k=0}^\infty B_n^k\)，后面的范数放缩全是水到渠成！

\SubPillMain{【规范解答】}

若 \(X = \{0\}\)，结论平凡成立；下设 \(X \neq \{0\}\)。

因为 \(\lambda_0 \in \rho(A)\)，由定义 6.1 与注记 6.1，预解算子 \(R_0 \triangleq (\lambda_0 I - A)^{-1} \in \mathcal{L}(X)\) 存在且有界，且由 \(R_0(\lambda_0 I - A) = I\) 知 \(\|R_0\| > 0\)。

对每个 \(n \in \mathbb{N}\)，作恒等分解： 
\FitDisplayMath{\lambda_0 I - A_n = (\lambda_0 I - A) - (A_n - A) = (\lambda_0 I - A)\bigl[I - R_0(A_n - A)\bigr].}
 记 \(B_n \triangleq R_0(A_n - A) \in \mathcal{L}(X)\)，则 
\FitDisplayMath{\|B_n\| \le \|R_0\|\,\|A_n - A\|.}


由于 \(\lim_{n\to\infty}\|A - A_n\| = 0\)，存在 \(N > 0\) 使得对所有 \(n > N\) 均有 
\FitDisplayMath{\|A_n - A\| < \frac{1}{\|R_0\|},}
 从而对所有 \(n > N\) 都有 
\FitDisplayMath{\|B_n\| \le \|R_0\|\,\|A_n - A\| < 1.}
 因此 \(\lim_{k\to\infty} \sqrt[k]{\|B_n^k\|} \le \|B_n\| < 1\)。由定理 6.1（Neumann 级数，取 \(\lambda = 1\)），当 \(n > N\) 时算子 \(I - B_n\) 可逆，其逆算子在 \(\mathcal{L}(X)\) 中按算子范数收敛于级数 
\FitDisplayMath{(I - B_n)^{-1} = \sum_{k=0}^\infty B_n^k \in \mathcal{L}(X),}
 且满足范数估计 
\FitDisplayMath{\|(I - B_n)^{-1}\| \le \sum_{k=0}^\infty \|B_n\|^k = \frac{1}{1 - \|B_n\|}.}


于是对所有 \(n > N\)，\(\lambda_0 I - A_n = (\lambda_0 I - A)(I - B_n)\) 是两个双射的复合，因而是 \(X\) 到 \(X\) 的双射，故 \(\lambda_0 \in \rho(A_n)\)，并且 
\FitDisplayMath{(\lambda_0 I - A_n)^{-1} = (I - B_n)^{-1}(\lambda_0 I - A)^{-1} = (I - B_n)^{-1}R_0.}


进一步，对 \(n > N\) 计算差值： 
\FitDisplayMath{(\lambda_0 I - A_n)^{-1} - (\lambda_0 I - A)^{-1} = \bigl[(I - B_n)^{-1} - I\bigr]R_0 = \left(\sum_{k=1}^\infty B_n^k\right) R_0 = B_n (I - B_n)^{-1} R_0.}
 取算子范数并利用函数 \(t \mapsto \frac{t}{1 - t}\) 在 \([0, 1)\) 上单调递增（因 \(0 \le \|B_n\| \le \|R_0\|\|A_n - A\| < 1\)），得 
\FitDisplayMath{\begin{aligned}
\|(\lambda_0 I - A_n)^{-1} - (\lambda_0 I - A)^{-1}\|
&\le \|B_n\|\,\|(I - B_n)^{-1}\|\,\|R_0\| \\
&\le \frac{\|B_n\|\,\|R_0\|}{1 - \|B_n\|}
\le \frac{\|R_0\|^2\,\|A_n - A\|}{1 - \|R_0\|\,\|A_n - A\|}.
\end{aligned}}
 令 \(n \to \infty\)，由于 \(\|A_n - A\| \to 0\)，右端趋于 \(0\)，因此 
\FitDisplayMath{\lim_{n\to\infty} \|(\lambda_0 I - A_n)^{-1} - (\lambda_0 I - A)^{-1}\| = 0. \qquad \square}


\SubPillCoral{【考点复盘与避坑提示】}

\begin{itemize}
\tightlist
\item
  \textbf{乘法顺序不要写反}：\((ST)^{-1} = T^{-1}S^{-1}\)！上面分解为 \(\lambda_0 I - A_n = (\lambda_0 I - A)(I - B_n)\)，求逆后 \((I - B_n)^{-1}\) 在\textbf{左边}，\(R_0 = (\lambda_0 I - A)^{-1}\) 在\textbf{右边}。如果习惯把 \(R_0\) 放右边，也可以一开始就右提公因式：\(\lambda_0 I - A_n = [I - (A_n - A)R_0](\lambda_0 I - A)\)，两种完全等价。
\item
  \textbf{对比教材定理 6.2}：定理 6.2 是固定 \(A\) 扰动标量 \(\lambda\)；本题（以及教材第 14 题，与本题内容重复）则是固定 \(\lambda_0\) 扰动算子 \(A_n\)，底层招式完全一样——全是提取已知可逆因子后上 Neumann 级数！
\end{itemize}

\SoftDivider

\BookProblem{prob:6-13}{第 13 题}{教材 p. 267}{第 13 题}

\SubPillPrimary{【原题】}

设 \(A \in \mathcal{L}(X)\)，\(\alpha \in \rho(A)\)，\(B = R(\alpha, A)\)。证明：

\begin{enumerate}
\def\labelenumi{(\roman{enumi})}
\tightlist
\item
  如果 \(\lambda, \mu \in \mathbb{C}\) 满足 \(\mu(\alpha - \lambda) = 1\)，则 \(\mu \in \sigma(B)\) 当且仅当 \(\lambda \in \sigma(A)\)。
\item
  如果 \(\mu \in \rho(B)\) 且 \(\mu(\alpha - \lambda) = 1\)，则 
\FitDisplayMath{R(\mu, B) = \frac{1}{\mu} + \frac{1}{\mu^2} R(\lambda, A).}
 （注：原书公式中 \(\frac{1}{\mu}\) 系 \(\frac{1}{\mu}I\) 的简写。）
\end{enumerate}

\SubPillAmber{【零基础破题思路】}

先把算子想象成普通数字！\(B = \frac{1}{\alpha - A}\)，那 \(\mu - B\) 是什么？做一下通分： 
\FitDisplayMath{\mu - \frac{1}{\alpha - A} = \frac{\mu(\alpha - A) - 1}{\alpha - A} = \frac{\mu\left(\alpha - \frac{1}{\mu} - A\right)}{\alpha - A} = \mu\cdot \frac{\lambda - A}{\alpha - A}!}
 再看分式取倒数（算子求逆）后的部分分式拆分： 
\FitDisplayMath{\frac{1}{\mu - B} = \frac{1}{\mu}\cdot \frac{\alpha - A}{\lambda - A} = \frac{1}{\mu}\cdot \frac{(\lambda - A) + (\alpha - \lambda)}{\lambda - A} = \frac{1}{\mu} + \frac{\alpha - \lambda}{\mu}\cdot \frac{1}{\lambda - A} = \frac{1}{\mu} + \frac{1}{\mu^2}\cdot\frac{1}{\lambda - A}!}
 因为这里出现的所有算子 \(\alpha I - A\)、\((\alpha I - A)^{-1}\)、\(\lambda I - A\) 全部\textbf{互相可交换}，标量通分公式搬到算子上一模一样成立！

\SubPillMain{【规范解答】}

由 \(\mu(\alpha - \lambda) = 1\) 立即得知 \(\mu \neq 0\) 且 \(\alpha - \lambda = \frac{1}{\mu} \neq 0\)，从而 \(\lambda = \alpha - \frac{1}{\mu}\)。

\textbf{核心恒等式}：因为 \(B = R(\alpha, A) = (\alpha I - A)^{-1} \in \mathcal{L}(X)\)，有 
\FitDisplayMath{\begin{aligned}
\mu I - B
&= \mu I - (\alpha I - A)^{-1} \\
&= (\alpha I - A)^{-1}\bigl[\mu(\alpha I - A) - I\bigr] \\
&= \mu(\alpha I - A)^{-1}\left[\left(\alpha - \frac{1}{\mu}\right)I - A\right] \\
&= \mu(\alpha I - A)^{-1}(\lambda I - A).
\end{aligned}}
 （另外，由于 \(\lambda I - A = (\alpha I - A) - \frac{1}{\mu}I\) 与 \(\alpha I - A\) 可交换，它也与 \((\alpha I - A)^{-1}\) 可交换，故上式右端两算子因子可交换。）

\SoftDivider

\textbf{(i) 证明 \(\mu \in \sigma(B) \iff \lambda \in \sigma(A)\)。}

等价地，我们证明 \(\mu \in \rho(B) \iff \lambda \in \rho(A)\)：

\begin{itemize}
\tightlist
\item
  \textbf{（\(\Leftarrow\)）} 设 \(\lambda \in \rho(A)\)，则 \(\lambda I - A\) 是 \(X\) 到 \(X\) 的双射。又 \(\mu \neq 0\) 且 \((\alpha I - A)^{-1}\) 是 \(X\) 到 \(X\) 的双射，故由核心恒等式，它们的复合 \(\mu I - B = \mu(\alpha I - A)^{-1}(\lambda I - A)\) 也是 \(X\) 到 \(X\) 的双射，即 \(\mu \in \rho(B)\)。
\item
  \textbf{（\(\Rightarrow\)）} 设 \(\mu \in \rho(B)\)，则 \(\mu I - B\) 是 \(X\) 到 \(X\) 的双射。对核心恒等式两端左乘 \(\frac{1}{\mu}(\alpha I - A)\)，得 
\FitDisplayMath{\lambda I - A = \frac{1}{\mu}(\alpha I - A)(\mu I - B).}
 由于 \(\frac{1}{\mu} \neq 0\)，且 \(\alpha I - A\) 与 \(\mu I - B\) 均是 \(X\) 到 \(X\) 的双射，故 \(\lambda I - A\) 也是双射，即 \(\lambda \in \rho(A)\)。
\end{itemize}

综上，\(\mu \in \rho(B) \iff \lambda \in \rho(A)\)，两边取补集即得 \(\mu \in \sigma(B) \iff \lambda \in \sigma(A)\)。

\SoftDivider

\textbf{(ii) 证明预解式关系 \(R(\mu, B) = \frac{1}{\mu}I + \frac{1}{\mu^2}R(\lambda, A)\)。}

设 \(\mu \in \rho(B)\) 且 \(\mu(\alpha - \lambda) = 1\)。由 (i) 的证明过程知 \(\lambda \in \rho(A)\)，故 \(R(\lambda, A) = (\lambda I - A)^{-1} \in \mathcal{L}(X)\) 存在。

对核心恒等式 \(\mu I - B = \mu(\alpha I - A)^{-1}(\lambda I - A)\) 两端求逆（利用 \((ST)^{-1} = T^{-1}S^{-1}\)），得 
\FitDisplayMath{\begin{aligned}
R(\mu, B) = (\mu I - B)^{-1}
&= \frac{1}{\mu}(\lambda I - A)^{-1}(\alpha I - A) \\
&= \frac{1}{\mu}R(\lambda, A)\bigl[(\lambda I - A) + (\alpha - \lambda)I\bigr] \\
&= \frac{1}{\mu}R(\lambda, A)(\lambda I - A) + \frac{\alpha - \lambda}{\mu}R(\lambda, A) \\
&= \frac{1}{\mu}I + \frac{1}{\mu^2}R(\lambda, A),
\end{aligned}}
 其中最后一步代入了 \(\alpha - \lambda = \frac{1}{\mu}\)。证毕。 \(\square\)

\SubPillCoral{【考点复盘与避坑提示】}

\begin{itemize}
\tightlist
\item
  \textbf{本质认识}：(i) 实质上是\textbf{谱映射定理}在分式线性函数 \(g(\lambda) = (\alpha - \lambda)^{-1}\) 上的特例，说明 \(\sigma(R(\alpha, A)) \setminus \{0\} = \{(\alpha - \lambda)^{-1} : \lambda \in \sigma(A)\}\)。
\item
  \textbf{顺便思考 \(\mu = 0\) 在哪里？} \(\mu = 0\) 不满足 \(\mu(\alpha-\lambda) = 1\)；但对有界算子 \(A\) 而言，\(B = (\alpha I - A)^{-1}\) 本身就有有界逆 \(B^{-1} = \alpha I - A \in \mathcal{L}(X)\)，所以 \(0 \in \rho(B)\) 恒成立，即实际上 \(\sigma(B) = \{(\alpha - \lambda)^{-1} : \lambda \in \sigma(A)\}\) 完整相等！（在无界闭算子中 \(0\) 才可能落进 \(\sigma(B)\)。）
\end{itemize}

\SoftDivider

\BookProblem{prob:6-23}{第 23 题}{教材 p. 268}{第 23 题}

\SubPillPrimary{【原题】}

设 \(A\) 为 \(H\) 上的有界自伴算子，\(\lambda \notin \mathbb{R}\)。证明：算子 \((A + \lambda I)(A + \bar{\lambda}I)^{-1}\) 为酉算子。

\SubPillAmber{【零基础破题思路】}

把算子 \(A\) 当作实数 \(x \in \mathbb{R}\)，把 \(\lambda\) 当作虚部非零的复数 \(a + \mathrm{i}b\)（\(b \neq 0\)）。那么分数 
\FitDisplayMath{u = \frac{x + \lambda}{x + \bar{\lambda}} = \frac{(x + a) + \mathrm{i}b}{(x + a) - \mathrm{i}b}}
 的分子和分母正好互为共轭复数！因此它的分母绝对不为 \(0\)（因为虚部 \(-b \neq 0\)），而且模长 \(|u| = \frac{\sqrt{(x+a)^2 + b^2}}{\sqrt{(x+a)^2 + b^2}} = 1\)！ 算子世界里，模长为 \(1\) 的复数对应的就是\textbf{酉算子}（\(U^*U = UU^* = I\)，保距满射）。证明分两步走：

\begin{enumerate}
\def\labelenumi{\arabic{enumi}.}
\tightlist
\item
  \textbf{先说明 \((A + \bar{\lambda}I)^{-1}\) 真的存在且有界}（别直接假定它存在！）：计算 \(\|(A + \bar{\lambda}I)x\|^2\) 会奇妙地展开成勾股定理 \(\|(A + aI)x\|^2 + b^2\|x\|^2 \ge b^2\|x\|^2\)，由下有界性一步证明它既是单射、值域又闭又稠密，从而满射可逆！（而且同一公式还能顺手证明 \(\|(A+\lambda I)x\| = \|(A+\bar{\lambda}I)x\|\)！）
\item
  \textbf{验证酉性}：用伴随运算法则 \((ST)^* = T^*S^*\)、\((T^{-1})^* = (T^*)^{-1}\) 与多项式算子可交换性，验证 \(U^*U = UU^* = I\)，即得保距满射（酉算子）。
\end{enumerate}

\SubPillMain{【规范解答】}

记 \(\lambda = a + \mathrm{i}b\)，其中 \(a, b \in \mathbb{R}\)。因为 \(\lambda \notin \mathbb{R}\)，所以 \(b \neq 0\)，且 \(\bar{\lambda} = a - \mathrm{i}b\)。

\textbf{第一步：证明 \((A + \bar{\lambda}I)^{-1}, (A + \lambda I)^{-1} \in \mathcal{L}(H)\) 均存在（从而 \(U\) 良定义）。}

因为 \(A\) 是自伴算子且 \(a \in \mathbb{R}\)，故 \(A + aI\) 也是自伴算子，从而对任意 \(x \in H\)，有 \(((A + aI)x, x) = (x, (A + aI)x) \in \mathbb{R}\)。于是对任意 \(x \in H\)，展开范数平方得 
\FitDisplayMath{\begin{aligned}
\|(A + \bar{\lambda}I)x\|^2
&= \|(A + aI)x - \mathrm{i}bx\|^2 \\
&= \|(A + aI)x\|^2 - ((A + aI)x, \mathrm{i}bx) - (\mathrm{i}bx, (A + aI)x) + b^2\|x\|^2 \\
&= \|(A + aI)x\|^2 + \mathrm{i}b((A + aI)x, x) - \mathrm{i}b(x, (A + aI)x) + b^2\|x\|^2 \\
&= \|(A + aI)x\|^2 + b^2\|x\|^2 \ge b^2\|x\|^2.
\end{aligned}}
 将 \(b\) 换为 \(-b\)，完全同理可得 
\FitDisplayMath{\|(A + \lambda I)x\|^2 = \|(A + aI)x\|^2 + b^2\|x\|^2 \ge b^2\|x\|^2, \quad \forall x \in H. \qquad (*)}


由不等式 \((*)\) 及 \(b \neq 0\) 立即推出以下三条：

\begin{enumerate}
\def\labelenumi{\arabic{enumi}.}
\tightlist
\item
  \textbf{单射性}：若 \((A + \bar{\lambda}I)x = 0\)，则 \(|b|\|x\| \le 0 \implies x = 0\)，故 \(A + \bar{\lambda}I\) 为单射；同理 \(A + \lambda I\) 也是单射。
\item
  \textbf{值域闭性}：设 \(y_n = (A + \bar{\lambda}I)x_n \to y\)（\(n \to \infty\)），则 
\FitDisplayMath{\|x_n - x_m\| \le \frac{1}{|b|}\|(A + \bar{\lambda}I)(x_n - x_m)\| = \frac{1}{|b|}\|y_n - y_m\| \to 0 \quad (n, m \to \infty).}
 故 \(\{x_n\}\) 是 \(H\) 中的 Cauchy 列。由 \(H\) 完备性，存在 \(x \in H\) 使 \(x_n \to x\)，再由 \(A + \bar{\lambda}I\) 连续性得 \(y = \lim_{n\to\infty}(A + \bar{\lambda}I)x_n = (A + \bar{\lambda}I)x \in \mathcal{R}(A + \bar{\lambda}I)\)，故 \(\mathcal{R}(A + \bar{\lambda}I)\) 是闭子空间。
\item
  \textbf{值域稠密性}：任取 \(z \in \mathcal{R}(A + \bar{\lambda}I)^\perp\)，则对一切 \(x \in H\)， 
\FitDisplayMath{0 = ((A + \bar{\lambda}I)x, z) = (x, (A + \bar{\lambda}I)^* z) = (x, (A + \lambda I)z).}
 由 \(x\) 的任意性得 \((A + \lambda I)z = 0\)；又 \(A + \lambda I\) 为单射，故 \(z = 0\)，即 \(\mathcal{R}(A + \bar{\lambda}I)^\perp = \{0\}\)。
\end{enumerate}

结合第 2、3 条知 \(\mathcal{R}(A + \bar{\lambda}I) = \overline{\mathcal{R}(A + \bar{\lambda}I)} = H\)。因此 \(A + \bar{\lambda}I\) 是 \(H\) 上的连续双射，由 Banach 逆算子定理，\((A + \bar{\lambda}I)^{-1} \in \mathcal{L}(H)\) 存在。同理，\((A + \lambda I)^{-1} \in \mathcal{L}(H)\) 也存在。

\SoftDivider

\textbf{第二步：证明 \(U \triangleq (A + \lambda I)(A + \bar{\lambda}I)^{-1}\) 是酉算子。}

首先，因 \((A + \bar{\lambda}I)^* = A^* + \lambda I = A + \lambda I\)，且对任意可逆算子 \(T \in \mathcal{L}(H)\) 均有 \((T^{-1})^* = (T^*)^{-1}\)（验证：\(T^*(T^{-1})^* = (T^{-1}T)^* = I\)，\((T^{-1})^*T^* = (TT^{-1})^* = I\)），故 
\FitDisplayMath{U^* = \bigl[(A + \lambda I)(A + \bar{\lambda}I)^{-1}\bigr]^* = \bigl[(A + \bar{\lambda}I)^{-1}\bigr]^*(A + \lambda I)^* = (A + \lambda I)^{-1}(A + \bar{\lambda}I).}


其次，\((A + \lambda I)\) 与 \((A + \bar{\lambda}I)\) 显然可交换： 
\FitDisplayMath{(A + \lambda I)(A + \bar{\lambda}I) = A^2 + (\lambda + \bar{\lambda})A + |\lambda|^2 I = (A + \bar{\lambda}I)(A + \lambda I).}
 在两边同时左乘、右乘 \((A + \bar{\lambda}I)^{-1}\)，得 
\FitDisplayMath{(A + \bar{\lambda}I)^{-1}(A + \lambda I) = (A + \lambda I)(A + \bar{\lambda}I)^{-1}.}
 同理 \((A + \lambda I)^{-1}\) 与 \((A + \bar{\lambda}I)\) 及 \((A + \bar{\lambda}I)^{-1}\) 均两两可交换。直接计算乘积： 
\FitDisplayMath{U^*U = (A + \lambda I)^{-1}(A + \bar{\lambda}I)(A + \lambda I)(A + \bar{\lambda}I)^{-1} = (A + \lambda I)^{-1}(A + \lambda I)(A + \bar{\lambda}I)(A + \bar{\lambda}I)^{-1} = I,}
 
\FitDisplayMath{UU^* = (A + \lambda I)(A + \bar{\lambda}I)^{-1}(A + \lambda I)^{-1}(A + \bar{\lambda}I) = (A + \lambda I)(A + \lambda I)^{-1}(A + \bar{\lambda}I)^{-1}(A + \bar{\lambda}I) = I.}


由此立得：对一切 \(x \in H\)，\(\|Ux\|^2 = (Ux, Ux) = (U^*Ux, x) = (x, x) = \|x\|^2\)（保距）；且对任意 \(y \in H\)，令 \(x = U^*y\) 即有 \(Ux = UU^*y = y\)（满射）。由定义 6.12，\(U\) 是 \(H\) 上的酉算子。 \(\square\)

\SubPillCoral{【考点复盘与避坑提示】}

\begin{itemize}
\tightlist
\item
  \textbf{更快捷的``两行保距+满射''验证（备选思路）}：做完第一步后其实根本不用算 \(U^*\)！由 \((*)\) 式已经得到 \(\|(A + \lambda I)y\| = \|(A + \bar{\lambda}I)y\|\) 对一切 \(y \in H\) 成立。对任意 \(x \in H\) 令 \(y = (A + \bar{\lambda}I)^{-1}x\)，立即得 \(\|Ux\| = \|(A+\lambda I)y\| = \|(A+\bar{\lambda}I)y\| = \|x\|\)（保距）；又 \(A + \bar{\lambda}I\) 与 \(A + \lambda I\) 都是满射，故 \(\mathcal{R}(U) = (A+\lambda I)(H) = H\)（满射），直接套定义 6.12 秒杀！
\item
  \textbf{与教材 6.6 节的联系}：本题取 \(\lambda = -\mathrm{i}\) 就是教材定义 6.15 的 \textbf{Cayley 变换} \(U = (A - \mathrm{i}I)(A + \mathrm{i}I)^{-1}\)！如果期末考``证明 Cayley 变换是酉算子''，直接套用本题一字不改即可拿满分！
\item
  \textbf{快捷补充}：若已可引用``自伴算子的谱全在实轴上 \(\sigma(A) \subset \mathbb{R}\)''，则由 \(-\bar{\lambda}, -\lambda \notin \mathbb{R}\) 也可直接推得 \(-\bar{\lambda}, -\lambda \in \rho(A)\)，从而快速得出第一步的可逆性；但上面的展开法完全不依赖后面的谱理论，更加稳妥自包含。
\end{itemize}

\SoftDivider

\BookProblem{prob:6-24}{第 24 题}{教材 p. 268}{第 24 题}

\SubPillPrimary{【原题】}

定义 \(L^2(0,1)\) 上的算子： 
\FitDisplayMath{Af(x) \triangleq \int_0^x f(t)\,\mathrm{d}t.}


\begin{enumerate}
\def\labelenumi{(\roman{enumi})}
\tightlist
\item
  求 \(\|A\|\)；
\item
  证明：\(A^*f(x) = \int_x^1 f(t)\,\mathrm{d}t,\; f \in L^2(0,1)\)；
\item
  计算 \(AA^*\) 的谱半径。
\end{enumerate}

\SubPillAmber{【零基础破题思路】}

\textbf{千万别按 (i)\(\to\)(ii)\(\to\)(iii) 的顺序死算！} 直接用定义求 \(\|A\| = \sup_{\|f\|_2=1} \|Af\|_2\) 是个很棘手的变分问题（其实等价于 Wirtinger 不等式）。命题人把 (ii) 和 (iii) 放在后面，其实是给了你一根绝妙的``倒推拐杖''！ 正确的\textbf{黄金做题顺序是 (ii) \(\to\) (iii) \(\to\) (i)}：

\begin{enumerate}
\def\labelenumi{\arabic{enumi}.}
\tightlist
\item
  \textbf{先做 (ii)}：用微积分的 \textbf{Fubini 定理}把积分区域 \(\{0 \le t \le x \le 1\}\) 交换积分次序，两行就求出 \(A^*\)！
\item
  \textbf{再做 (iii)}：观察 \(A\) 是一个平方可积核的积分算子（Hilbert-Schmidt 算子），所以 \(A\) 是\textbf{紧算子}，进而 \(AA^*\) 是非零\textbf{正自伴紧算子}！由 6.3 节核心定理（定理 6.12、推论 6.2），正自伴紧算子的谱半径 \(\Gamma(AA^*)\) 就等于它的\textbf{最大正特征值}！把特征方程 \(AA^*f = \lambda f\) 连求两次导数，瞬间变成大一高数就能手算的常微分方程 \(f'' + \frac{1}{\lambda}f = 0\)（配边值 \(f(0)=0, f'(1)=0\)），解出最大特征值 \(\lambda_1 = \frac{4}{\pi^2}\)！
\item
  \textbf{最后一步得出 (i)}：由 \(C^*\)-恒等式 \(\|A\|^2 = \|AA^*\| = \Gamma(AA^*) = \frac{4}{\pi^2}\)，开个根号立得 \(\|A\| = \frac{2}{\pi}\)！
\end{enumerate}

\SubPillMain{【规范解答】}

\begin{quote}
为逻辑顺畅，下面先证 \(A\) 为有界紧算子并证明 (ii)，再求 (iii)，最后由 (iii) 立即得出 (i)。
\end{quote}

\textbf{预备：说明 \(A\) 是 \(L^2(0,1)\) 上的有界线性紧算子。}

显然 \(A\) 线性，且可写作积分算子 \((Af)(x) = \int_0^1 K(x, t)f(t)\,\mathrm{d}t\)，其中核函数 
\FitDisplayMath{K(x, t) = \begin{cases} 1, & 0 \le t \le x \le 1, \\ 0, & 0 \le x < t \le 1. \end{cases}}
 直接计算核函数的 \(L^2\)-范数： 
\FitDisplayMath{\int_0^1 \int_0^1 |K(x, t)|^2 \,\mathrm{d}t\,\mathrm{d}x = \int_0^1 x\,\mathrm{d}x = \frac{1}{2} < \infty.}
 故 \(K \in L^2((0,1)\times(0,1))\)。由第 5 章（例 5.4、例 5.7 与定理 5.12），\(A\) 是 \(L^2(0,1)\) 上的 Hilbert-Schmidt 算子，从而是\textbf{有界线性紧算子}（\(A \in \mathcal{C}(L^2(0,1))\)）。

\SoftDivider

\textbf{(ii) 证明 \((A^*f)(x) = \int_x^1 f(t)\,\mathrm{d}t\)。}

任取 \(f, g \in L^2(0,1)\)。由 Cauchy-Schwarz 不等式， 
\FitDisplayMath{\int_0^1 \int_0^x |f(t)\overline{g(x)}|\,\mathrm{d}t\,\mathrm{d}x \le \left(\int_0^1 |f(t)|\,\mathrm{d}t\right)\left(\int_0^1 |g(x)|\,\mathrm{d}x\right) \le \|f\|_2 \|g\|_2 < \infty.}
 因此可以对三角形区域 \(\{(t, x) : 0 \le t \le x \le 1\}\) 应用 Fubini 定理交换积分次序： 
\FitDisplayMath{(Af, g) = \int_0^1 \left(\int_0^x f(t)\,\mathrm{d}t\right)\overline{g(x)}\,\mathrm{d}x = \int_0^1 f(t)\left(\int_t^1 \overline{g(x)}\,\mathrm{d}x\right)\mathrm{d}t = \int_0^1 f(t)\overline{\left(\int_t^1 g(x)\,\mathrm{d}x\right)}\,\mathrm{d}t.}
 由此知 \((Af, g) = (f, A^*g)\)，其中 \((A^*g)(t) = \int_t^1 g(x)\,\mathrm{d}x\)。换记变量名即得 
\FitDisplayMath{(A^*f)(x) = \int_x^1 f(t)\,\mathrm{d}t, \quad \forall f \in L^2(0,1).}


\SoftDivider

\textbf{(iii) 计算 \(AA^*\) 的谱半径 \(\Gamma(AA^*)\)。}

\textbf{步骤 1（确认 \(AA^*\) 的性质）}：因为 \(A\) 是紧算子、\(A^*\) 是有界算子，故 \(AA^*\) 是紧算子；又 \((AA^*)^* = (A^*)^*A^* = AA^*\)，且对一切 \(f \in L^2(0,1)\) 有 \((AA^*f, f) = \|A^*f\|_2^2 \ge 0\)，所以 \(AA^*\) 是\textbf{正自伴紧算子}（所有特征值均 \(\ge 0\)）。取 \(f \equiv 1\) 得 \((A^*1)(x) = 1 - x \neq 0\)，故 \(AA^* \neq 0\)。

由定理 6.12、推论 6.2 及正算子特征值非负，可知 \(\Gamma(AA^*) = \|AA^*\| > 0\)，且 \(\Gamma(AA^*)\) 恰好是 \(AA^*\) 的\textbf{最大正特征值}！

\textbf{步骤 2（将特征方程化为常微分方程边值问题）}：设 \(\lambda > 0\) 是 \(AA^*\) 的正特征值，\(f \in L^2(0,1)\) 是对应的非零特征函数，即 \(AA^*f = \lambda f\)： 
\FitDisplayMath{\lambda f(x) = (AA^*f)(x) = \int_0^x (A^*f)(s)\,\mathrm{d}s = \int_0^x \left(\int_s^1 f(t)\,\mathrm{d}t\right)\mathrm{d}s, \quad \text{a.e. } x \in (0,1). \qquad (\star)}
 我们先说明 \(f\) 可以（在修改零测集后）视为 \(C^2([0,1])\) 函数：

\begin{itemize}
\tightlist
\item
  因为 \(f \in L^2(0,1) \subset L^1(0,1)\)，函数 \(g(s) \triangleq \int_s^1 f(t)\,\mathrm{d}t\) 在 \([0,1]\) 上绝对连续（特别地连续）；
\item
  连续函数的变上限积分 \(h(x) \triangleq \int_0^x g(s)\,\mathrm{d}s\) 属于 \(C^1([0,1])\) 且 \(h'(x) = g(x)\)；
\item
  由 \((\star)\) 及 \(\lambda > 0\)，取 \(f = \frac{1}{\lambda}h\)，则 \(f \in C^1([0,1])\)；
\item
  既然 \(f\) 已连续，回代得 \(g(s) = \int_s^1 f(t)\,\mathrm{d}t \in C^1([0,1])\) 且 \(g'(s) = -f(s)\)，于是 \(f = \frac{1}{\lambda}h \in C^2([0,1])\)！
\end{itemize}

对光滑恒等式 \(\lambda f(x) = \int_0^x \left(\int_s^1 f(t)\,\mathrm{d}t\right)\mathrm{d}s\) 逐次求导并代入端点：

\begin{itemize}
\tightlist
\item
  令 \(x = 0\)，得 \(\lambda f(0) = 0 \implies f(0) = 0\)；
\item
  两边对 \(x\) 求一阶导，得 \(\lambda f'(x) = \int_x^1 f(t)\,\mathrm{d}t\)，令 \(x = 1\) 得 \(\lambda f'(1) = 0 \implies f'(1) = 0\)；
\item
  两边再对 \(x\) 求二阶导，得 \(\lambda f''(x) = -f(x)\)。
\end{itemize}

因此 \(f\) 满足边值问题： 
\FitDisplayMath{\begin{cases} f''(x) + \dfrac{1}{\lambda} f(x) = 0, & x \in [0, 1], \\ f(0) = 0, \quad f'(1) = 0. \end{cases}}


\textbf{步骤 3（解出全部正特征值并取最大者）}：记 \(\omega \triangleq \frac{1}{\sqrt{\lambda}} > 0\)，方程的通解为 
\FitDisplayMath{f(x) = C_1 \cos(\omega x) + C_2 \sin(\omega x).}
 由 \(f(0) = 0\) 得 \(C_1 = 0\)，故 \(f(x) = C_2 \sin(\omega x)\)（因 \(f \not\equiv 0\)，必有 \(C_2 \neq 0\)）。 求导得 \(f'(x) = C_2 \omega \cos(\omega x)\)，由 \(f'(1) = 0\) 及 \(C_2 \omega \neq 0\) 得 \(\cos(\omega) = 0\)。因为 \(\omega > 0\)，故 
\FitDisplayMath{\omega = \frac{(2n - 1)\pi}{2}, \quad n = 1, 2, 3, \dots}
 解得全部候选正特征值为 
\FitDisplayMath{\lambda_n = \frac{1}{\omega^2} = \frac{4}{(2n - 1)^2 \pi^2}, \quad n = 1, 2, 3, \dots}
 反向检验充分性：取 \(f_n(x) = \sin(\omega_n x)\)（\(\omega_n = \frac{(2n-1)\pi}{2}\)），直接积分得 
\FitDisplayMath{\int_0^x \left(\int_s^1 \sin(\omega_n t)\,\mathrm{d}t\right)\mathrm{d}s = \int_0^x \frac{\cos(\omega_n s) - \cos(\omega_n)}{\omega_n}\,\mathrm{d}s = \int_0^x \frac{\cos(\omega_n s)}{\omega_n}\,\mathrm{d}s = \frac{\sin(\omega_n x)}{\omega_n^2} = \lambda_n f_n(x),}
 故每个 \(\lambda_n\) 确实都是 \(AA^*\) 的正特征值！其中最大者在 \(n = 1\) 时取到，所以 \(AA^*\) 的谱半径为 
\FitDisplayMath{\Gamma(AA^*) = \|AA^*\| = \max_{n \ge 1} \lambda_n = \lambda_1 = \frac{4}{\pi^2}.}


\SoftDivider

\textbf{(i) 求算子范数 \(\|A\|\)。}

先证 Hilbert 空间上有界算子的通用恒等式 \(\|A\|^2 = \|AA^*\|\)：

\begin{itemize}
\tightlist
\item
  一方面，\(\|AA^*\| \le \|A\|\,\|A^*\| = \|A\|^2\)；
\item
  另一方面，对任意 \(x \in L^2(0,1)\) 有 \(\|A^*x\|_2^2 = (A^*x, A^*x) = (AA^*x, x) \le \|AA^*\|\,\|x\|_2^2\)，对 \(\|x\|_2 = 1\) 取上确界得 \(\|A\|^2 = \|A^*\|^2 \le \|AA^*\|\)。 从而 \(\|A\|^2 = \|AA^*\|\)。结合 (iii) 立刻得到 
\FitDisplayMath{\|A\| = \sqrt{\|AA^*\|} = \sqrt{\Gamma(AA^*)} = \sqrt{\frac{4}{\pi^2}} = \frac{2}{\pi}. \qquad \square}

\end{itemize}

\SubPillCoral{【考点复盘与避坑提示】}

\begin{itemize}
\tightlist
\item
  \textbf{超级经典大坑（切勿把 \(\Gamma(A)\) 和 \(\|A\|\) 混为一谈！）}：对于 Volterra 算子 \(A\) 本身，用归纳法可算得 \(\|A^n\| \le \frac{1}{n!}\)（参见教材例 6.5），所以 \(A\) 的谱半径 \(\Gamma(A) = \lim_{n\to\infty} \sqrt[n]{\|A^n\|} = 0\)！但它的范数 \(\|A\| = \frac{2}{\pi} \neq 0\)！原因是 \(A\) \textbf{不是自伴算子}，只有对自伴（或正规）算子才有 \(\Gamma(T) = \|T\|\)。所以求非自伴算子的范数，必须先构造自伴算子 \(AA^*\) 或 \(A^*A\)！
\item
  \textbf{容易漏写的得分点}：从 \(L^2\) 几乎处处相等到``可以两次求导''的光滑性自举（bootstrap）步骤，以及求完 \(\lambda_n\) 后的回代检验（验证充分性），在期末考试中往往占 2--3 分，务必简明写清！
\item
  \textbf{附：检验最大范数可达}：取 \(g_1(x) = \cos\left(\frac{\pi}{2}x\right)\)，则 \((Ag_1)(x) = \frac{2}{\pi}\sin\left(\frac{\pi}{2}x\right)\)，易算 \(\|g_1\|_2^2 = \frac{1}{2}\)，\(\|Ag_1\|_2^2 = \frac{4}{\pi^2}\cdot \frac{1}{2}\)，确实精确达到 \(\frac{\|Ag_1\|_2}{\|g_1\|_2} = \frac{2}{\pi}\)！
\end{itemize}

\SoftDivider

\BookProblem{prob:6-31}{第 31 题}{教材 p. 269}{第 31 题}

\SubPillPrimary{【原题】}

证明下列论断等价：

\begin{enumerate}
\def\labelenumi{(\roman{enumi})}
\tightlist
\item
  \(A\) 是 \(H\) 到复 Hilbert 空间 \(G\) 的 Hilbert-Schmidt 算子；
\item
  \((A^*A)^{\frac{1}{2}}\) 是 \(H\) 上的 Hilbert-Schmidt 算子；
\item
  存在非负数列 \(\{\lambda_n\}_{n=1}^\infty\)，\(\sum_{n=1}^\infty \lambda_n^2 < \infty\) 以及 \(H, G\) 中的标准正交基 \(\{x_n\}_{n=1}^\infty \subset H\)，\(\{y_n\}_{n=1}^\infty \subset G\)，使得对一切 \(x \in H\)， 
\FitDisplayMath{Ax = \sum_{n=1}^\infty \lambda_n (x, x_n) y_n.}

\end{enumerate}

\SubPillAmber{【零基础破题思路】}

这道题其实是线性代数里\textbf{矩阵奇异值分解（SVD）}的无穷维版！

\begin{itemize}
\tightlist
\item
  为什么 (i) \(\iff\) (ii) 几乎是一句话？因为记 \(T = (A^*A)^{\frac{1}{2}}\)（自伴正算子，满足 \(T^2 = A^*A\)），对\textbf{任意向量 \(x\)} 都有黄金等式： 
\FitDisplayMath{\|Tx\|^2 = (Tx, Tx) = (T^2x, x) = (A^*Ax, x) = \|Ax\|^2!}
 也就是说，\(T\) 和 \(A\) 把每个向量的长度拉成了\textbf{一模一样}！那按标准正交基算平方和 \(\sum \|Ae_n\|^2 = \sum \|Te_n\|^2\)，自然一边有限等价于另一边有限！
\item
  为什么 (ii) \(\Rightarrow\) (iii)？因为 \(T = (A^*A)^{\frac{1}{2}}\) 既是 Hilbert-Schmidt 算子（从而是\textbf{紧算子}，定理 5.12），又是\textbf{正自伴算子}，所以直接调用自伴紧算子的谱展开定理（定理 6.13），拿到 \(H\) 中由 \(T\) 的特征向量组成的标准正交基 \(\{x_n\}\)（\(Tx_n = \lambda_n x_n \ge 0\)）。然后对 \(\lambda_n > 0\) 的那些项，把 \(Ax_n\) 除以长度 \(\lambda_n\) 归一化定义 \(y_n = \frac{1}{\lambda_n}Ax_n\)，一算内积发现它们居然在 \(G\) 里自动两两正交且长度为 \(1\)！这就凑出了奇异值展开！
\item
  为什么 (iii) \(\Rightarrow\) (i)？最简单：把 \(x = x_k\) 直接代进公式 (iii)，立刻得 \(Ax_k = \lambda_k y_k\)，所以 \(\sum \|Ax_k\|^2 = \sum \lambda_k^2 < \infty\)，秒证！
\end{itemize}

\begin{quote}
\textbf{前置严谨提醒}：依第 5.2 节背景（p. 196）与题干记号，默认 \(H, G\) 为无穷维可分复 Hilbert 空间（有限维时将无穷求和改为有限和即可），且 (ii) 中出现 \(A^*\) 隐含 \(A \in \mathcal{L}(H, G)\)（在 (i) 中由定理 5.12(2) 知 \(A \in \mathcal{L}^2(H,G)\) 自动有界）。关于 (iii) 中``标准正交基''在 \(\dim\operatorname{Ker}A \neq \dim\mathcal{R}(A)^\perp\) 时的细微拓扑细节，请务必阅读文末【严谨注记】！
\end{quote}

\SubPillMain{【规范解答】}

\textbf{预备引理（核心等距恒等式）}：对有界线性算子 \(A \in \mathcal{L}(H, G)\)，显然 \(A^*A \in \mathcal{L}(H)\) 满足 \((A^*A)^* = A^*A\) 且 \((A^*Ax, x) = \|Ax\|^2 \ge 0\)，即 \(A^*A \ge 0\)。由正算子平方根定理（推论 6.17），存在唯一的正自伴算子 \(T \triangleq (A^*A)^{\frac{1}{2}} \in \mathcal{L}(H)\) 满足 \(T \ge 0\) 且 \(T^2 = A^*A\)。对任意 \(x \in H\)，有 
\FitDisplayMath{\|Tx\|^2 = (Tx, Tx) = (T^2x, x) = (A^*Ax, x) = (Ax, Ax) = \|Ax\|^2,}
 即 
\FitDisplayMath{\|(A^*A)^{\frac{1}{2}}x\| = \|Ax\|, \quad \forall x \in H. \qquad (\dagger)}


\SoftDivider

\textbf{第一步：证明 (i) \(\iff\) (ii)。}

\begin{itemize}
\tightlist
\item
  若 (i) 成立，由定理 5.12(2) 知 \(\|A\| \le \|A\|_2 < \infty\)，故 \(A \in \mathcal{L}(H, G)\)，从而 \(T = (A^*A)^{\frac{1}{2}}\) 存在。
\item
  任取 \(H\) 的一组标准正交基 \(\{e_n\}_{n=1}^\infty\)。由等式 \((\dagger)\) 得 
\FitDisplayMath{\sum_{n=1}^\infty \|Ae_n\|^2 = \sum_{n=1}^\infty \|(A^*A)^{\frac{1}{2}}e_n\|^2.}
 根据定义 5.4，左端有限当且仅当右端有限，故 \(A \in \mathcal{L}^2(H, G) \iff (A^*A)^{\frac{1}{2}} \in \mathcal{L}^2(H)\)（且此时 \(\|A\|_2 = \|(A^*A)^{\frac{1}{2}}\|_2\)）。因此 (i) \(\iff\) (ii)。
\end{itemize}

\SoftDivider

\textbf{第二步：证明 (iii) \(\Rightarrow\) (i)。}

设 (iii) 成立。对任意 \(x \in H\)，由于 \(\{x_n\}_{n=1}^\infty\) 是 \(H\) 的标准正交基，由 Parseval 等式知 \(\sum_{n=1}^\infty |(x, x_n)|^2 = \|x\|^2 < \infty\)；又 \(\sum_{n=1}^\infty \lambda_n^2 < \infty\) 保证数列 \(\{\lambda_n\}\) 有界（记 \(M = \sup_n \lambda_n < \infty\)），故 
\FitDisplayMath{\sum_{n=1}^\infty |\lambda_n (x, x_n)|^2 \le M^2 \|x\|^2 < \infty.}
 因 \(\{y_n\}_{n=1}^\infty\) 是 \(G\) 中的标准正交基，级数 \(Ax = \sum_{n=1}^\infty \lambda_n (x, x_n) y_n\) 在 \(G\) 中按范数收敛，且定义了一个有界线性算子 \(A \in \mathcal{L}(H, G)\)。

对每个 \(k \in \mathbb{N}\)，代入 \(x = x_k\)，由 \((x_k, x_n) = \delta_{kn}\) 得 \(Ax_k = \lambda_k y_k\)，从而 
\FitDisplayMath{\|Ax_k\|^2 = \|\lambda_k y_k\|^2 = \lambda_k^2 \|y_k\|^2 = \lambda_k^2.}
 对 \(k\) 求和得 
\FitDisplayMath{\sum_{k=1}^\infty \|Ax_k\|^2 = \sum_{k=1}^\infty \lambda_k^2 < \infty.}
 由定义 5.4，\(A\) 是 \(H\) 到 \(G\) 的 Hilbert-Schmidt 算子，即 (i) 成立。

\SoftDivider

\textbf{第三步：证明 (ii) \(\Rightarrow\) (iii)。}

设 \(T = (A^*A)^{\frac{1}{2}} \in \mathcal{L}^2(H)\)。由定理 5.12(5)，\(T\) 是 \(H\) 上的紧算子；又 \(T\) 是正自伴算子。

\begin{enumerate}
\def\labelenumi{\arabic{enumi}.}
\item
  \textbf{构造 \(H\) 的标准正交基 \(\{x_n\}_{n=1}^\infty\) 与非负数列 \(\{\lambda_n\}_{n=1}^\infty\)}： 由自伴紧算子的 Hilbert-Schmidt 谱展开定理（定理 6.13 及证明，p. 223--224），\(T\) 的非零特征值（计重数）至多可数且均为正实数，对应标准正交特征向量构成 \(\overline{\mathcal{R}(T)} = (\operatorname{Ker} T)^\perp\) 的标准正交基；将它与 \(\operatorname{Ker} T\) 的一组标准正交基（对应特征值 \(0\)）合并，就得到整个可分无穷维空间 \(H\) 的一组由 \(T\) 的特征向量构成的标准正交基 \(\{x_n\}_{n=1}^\infty\)，满足 
\FitDisplayMath{Tx_n = \lambda_n x_n, \qquad \lambda_n \ge 0 \quad (n = 1, 2, \dots).}
 由 \(T \in \mathcal{L}^2(H)\) 及定义 5.4，有 
\FitDisplayMath{\sum_{n=1}^\infty \lambda_n^2 = \sum_{n=1}^\infty \|Tx_n\|^2 = \|T\|_2^2 < \infty.}

\item
  \textbf{在 \(G\) 中构造标准正交系 \(\{y_n\}\)}： 记 \(P \triangleq \{n \in \mathbb{N} : \lambda_n > 0\}\)，\(Z \triangleq \{n \in \mathbb{N} : \lambda_n = 0\}\)。

  \begin{itemize}
  \tightlist
  \item
    对每个 \(n \in P\)，定义 \(y_n \triangleq \frac{1}{\lambda_n} A x_n \in G\)。对任意 \(n, m \in P\)，利用 \(T^2 = A^*A\) 验证正交归一化： 
\FitDisplayMath{(y_n, y_m)_G = \frac{1}{\lambda_n \lambda_m}(Ax_n, Ax_m)_G = \frac{1}{\lambda_n \lambda_m}(A^*Ax_n, x_m)_H = \frac{1}{\lambda_n \lambda_m}(T^2 x_n, x_m)_H = \frac{\lambda_n^2}{\lambda_n \lambda_m}(x_n, x_m)_H = \delta_{nm}.}
 故 \(\{y_n\}_{n \in P}\) 是 \(G\) 中的标准正交系，且由定义有 \(Ax_n = \lambda_n y_n\)（\(n \in P\)），同时 \(\overline{\operatorname{span}\{y_n : n \in P\}} = \overline{\mathcal{R}(A)}\)。
  \item
    对每个 \(n \in Z\)，由 \((\dagger)\) 有 \(\|Ax_n\| = \|Tx_n\| = 0\)，即 \(Ax_n = 0\)。
  \item
    在 \(\mathcal{R}(A)^\perp\) 中取一组标准正交基将 \(\{y_n\}_{n \in P}\) 扩充为 \(G\) 的标准正交基（当 \(\dim(\operatorname{Ker} A) = \dim(\mathcal{R}(A)^\perp)\) 时，可将这组补全向量与下标集 \(Z\) 一一对应记作 \(\{y_n\}_{n \in Z}\)，于是 \(\{y_n\}_{n=1}^\infty\) 构成 \(G\) 的标准正交基；一般情形详见下方严谨注记）。无论补全的 \(\{y_n\}_{n \in Z}\) 为何，对一切 \(n \in \mathbb{N}\) 均有 \(Ax_n = \lambda_n y_n\)（当 \(n \in Z\) 时两端皆为 \(0\)）。
  \end{itemize}
\item
  \textbf{验证展开式}： 对任意 \(x \in H\)，因为 \(\{x_n\}_{n=1}^\infty\) 是 \(H\) 的标准正交基，有按 \(H\) 范数收敛的展开 \(x = \sum_{n=1}^\infty (x, x_n)x_n\)。由 (i) \(\iff\) (ii) 知 \(A \in \mathcal{L}(H, G)\) 有界连续，故可将 \(A\) 移入无穷求和号内，得 
\FitDisplayMath{Ax = \sum_{n=1}^\infty (x, x_n) Ax_n = \sum_{n=1}^\infty \lambda_n (x, x_n) y_n.}
 由此 (ii) \(\Rightarrow\) (iii) 证毕，三论断等价。 \(\square\)
\end{enumerate}

\SubPillCoral{【考点复盘与避坑提示】}

\begin{itemize}
\tightlist
\item
  \textbf{【严谨注记：原题 (iii) 表述的细微瑕疵】}：如果严格按字面要求 \(\{x_n\}_{n=1}^\infty\) 和 \(\{y_n\}_{n=1}^\infty\) \textbf{同时}分别是 \(H, G\) 的\textbf{完备标准正交基}，并且必须共用完全相同的下标 \(n = 1, 2, \dots\)，那么当 \(\dim(\operatorname{Ker} A) \neq \dim(\mathcal{R}(A)^\perp)\) 时，技术上存在一个``配对人数不等''的瑕疵！

  \begin{itemize}
  \tightlist
  \item
    \textbf{具体反例}：取 \(H = G = \ell^2\)，定义加权右移算子 \(A(a_1, a_2, \dots) = (0, a_1, \frac{a_2}{2}, \frac{a_3}{3}, \dots)\)。显然 \(A \in \mathcal{L}^2(\ell^2)\)，且 \(\operatorname{Ker} A = \{0\}\)（故所有 \(\lambda_n = \frac{1}{n} > 0\)，没有零项），但 \(\mathcal{R}(A)^\perp = \operatorname{span}\{e_1\}\) 是一维的。此时 \(\{y_n\}_{n=1}^\infty\) 若全由 \(Ax_n = \lambda_n y_n\) 生成，它们全在 \(\mathcal{R}(A)\) 中，永远正交于 \(e_1\)，不可能成为整个 \(\ell^2\) 的完备基；而若硬把 \(e_1\) 塞进 \(\{y_n\}\) 并配以 \(\lambda_k = 0\)，就意味着 \(Ax_k = 0\)，又与 \(\operatorname{Ker} A = \{0\}\) 矛盾！
  \item
    \textbf{完美修正}：泛函分析文献中，奇异值展开（Schmidt 表示）最精确的条件 (iii) 是把``标准正交基''改为\textbf{``标准正交系（orthonormal system）''}——只需 \(\{x_n\} \subset H\) 与 \(\{y_n\} \subset G\) 各自两两正交且为单位向量（无需完备），三个条件便对\textbf{一切} Hilbert-Schmidt 算子（甚至有限秩、非可分空间上的）无条件严格等价！考试作答时照上述规范解答写出完整构造，并简短附注这一点，既保证拿满分又体现深刻功底。
  \end{itemize}
\end{itemize}

\SoftDivider

\BookProblem{prob:6-37}{第 37 题}{教材 p. 270}{第 37 题}

\SubPillPrimary{【原题】}

设 \(H, G\) 为可分的 Hilbert 空间，\(A\) 是 \(H\) 到 \(G\) 上的 Hilbert-Schmidt 算子，\(\{\lambda_n\}_{n=1}^N\)（\(N \in \mathbb{N} \cup \{\infty\}\)）是 \((A^*A)^{\frac{1}{2}}\) 的特征值。证明： 
\FitDisplayMath{\|A\|_2 = \left(\sum_{n=1}^N \lambda_n^2\right)^{\frac{1}{2}}.}


\SubPillAmber{【零基础破题思路】}

这道题是第 31 题的直接定量版，说明了 \textbf{Hilbert-Schmidt 范数就是全部奇异值平方和的算术平方根}（就像线性代数里矩阵的 Frobenius 范数 \(\|A\|_F = \sqrt{\operatorname{tr}(A^*A)} = \sqrt{\sum \sigma_i^2}\)）。 思路非常清爽：

\begin{enumerate}
\def\labelenumi{\arabic{enumi}.}
\tightlist
\item
  定义 5.4（式 (5.8)）告诉我们：\(\|A\|_2^2 = \sum \|Ae_k\|^2\) \textbf{不管选哪组标准正交基来算，结果都一模一样}！
\item
  那就挑一组最爽的标准正交基——直接拿正自伴紧算子 \(T = (A^*A)^{\frac{1}{2}}\) 的\textbf{特征向量标准正交基}来算！
\item
  别忘了一个小坑：非零特征值对应的特征向量 \(\{u_n\}_{n=1}^N\) 只张成 \(\overline{\mathcal{R}(T)} = (\operatorname{Ker} T)^\perp\)，不一定张成整个 \(H\)！所以必须把 \(\operatorname{Ker} T\) 的标准正交基 \(\{v_k\}\) 补进去，凑齐整个 \(H\) 的标准正交基；好在 \(v_k\) 对应的 \(\|Av_k\| = \|Tv_k\| = 0\)，求和时全部为 \(0\)，剩下的正好就是 \(\sum_{n=1}^N \lambda_n^2\)！
\end{enumerate}

\SubPillMain{【规范解答】}

\begin{quote}
\textbf{记号说明}：题干中 \(\{\lambda_n\}_{n=1}^N\) 表示 \((A^*A)^{\frac{1}{2}}\) 按几何重数重复计数的全体非零特征值（若将特征值 \(0\) 也计入，因 \(0^2 = 0\)，平方和不变；当 \(A = 0\) 时无非零特征值，两端皆为 \(0\)，结论平凡）。
\end{quote}

设 \(A \in \mathcal{L}^2(H, G)\) 且 \(A \neq 0\)。由定理 5.12(2)，\(A \in \mathcal{L}(H, G)\)，从而 \(A^*A \in \mathcal{L}(H)\) 是正自伴算子。由推论 6.17，存在唯一的正自伴算子 \(T \triangleq (A^*A)^{\frac{1}{2}} \in \mathcal{L}(H)\) 满足 \(T \ge 0\) 且 \(T^2 = A^*A\)。

\textbf{第一步：证明 \(T\) 是紧正自伴算子。}

对任意 \(x \in H\)，有 
\FitDisplayMath{\|Tx\|^2 = (Tx, Tx) = (T^2 x, x) = (A^*Ax, x) = \|Ax\|^2. \qquad (\ddagger)}
 任取 \(H\) 的一组标准正交基 \(\{e_m\}\)，由 \((\ddagger)\) 得 
\FitDisplayMath{\sum_m \|Te_m\|^2 = \sum_m \|Ae_m\|^2 = \|A\|_2^2 < \infty.}
 故 \(T \in \mathcal{L}^2(H)\) 是 \(H\) 上的 Hilbert-Schmidt 算子。由定理 5.12(5)，\(T\) 是 \(H\) 上的紧算子，因而 \(T\) 是正自伴紧算子。

\textbf{第二步：选取由 \(T\) 的特征向量组成的 \(H\) 的标准正交基。}

由自伴紧算子的谱展开定理（定理 6.13 及证明，p. 223--224）：

\begin{itemize}
\tightlist
\item
  \(T\) 的全体非零特征值（按重数计）至多可数，记为 \(\{\lambda_n\}_{n=1}^N\)（\(\lambda_n > 0\)，\(N \in \mathbb{N} \cup \{\infty\}\)），可取对应的单位特征向量 \(\{u_n\}_{n=1}^N\)（满足 \(Tu_n = \lambda_n u_n\)），使之构成 \(\overline{\mathcal{R}(T)} = (\operatorname{Ker} T)^\perp\) 的标准正交基；
\item
  因为 \(H\) 可分，闭子空间 \(\operatorname{Ker} T\) 存在至多可数的标准正交基 \(\{v_k\}_{k \in K}\)（其中 \(K\) 可能为空、有限集或可数无穷集），满足 \(Tv_k = 0\)。
\end{itemize}

由于 \(H = (\operatorname{Ker} T)^\perp \oplus \operatorname{Ker} T\)，并集 
\FitDisplayMath{\mathcal{E} \triangleq \{u_n\}_{n=1}^N \cup \{v_k\}_{k \in K}}
 构成整个空间 \(H\) 的一组标准正交基。

\textbf{第三步：利用标准正交基无关性计算 \(\|A\|_2\)。}

由定义 5.4 及式 (5.8)，Hilbert-Schmidt 范数与 \(H\) 上标准正交基的选取无关。选用标准正交基 \(\mathcal{E}\) 并代入 \((\ddagger)\)：

\begin{itemize}
\tightlist
\item
  对每个 \(n\)，\(\|Au_n\|^2 = \|Tu_n\|^2 = \|\lambda_n u_n\|^2 = \lambda_n^2\)；
\item
  对每个 \(k \in K\)，\(\|Av_k\|^2 = \|Tv_k\|^2 = 0\)。
\end{itemize}

因此 
\FitDisplayMath{\|A\|_2^2 = \sum_{n=1}^N \|Au_n\|^2 + \sum_{k \in K} \|Av_k\|^2 = \sum_{n=1}^N \lambda_n^2 + 0 = \sum_{n=1}^N \lambda_n^2.}
 两边开算术平方根，即得 
\FitDisplayMath{\|A\|_2 = \left(\sum_{n=1}^N \lambda_n^2\right)^{\frac{1}{2}}. \qquad \square}


\SubPillCoral{【考点复盘与避坑提示】}

\begin{itemize}
\tightlist
\item
  \textbf{必避小坑 1（必须写清``按重数计数''）}：如果某个非零特征值 \(\lambda\) 的特征子空间是 \(m\) 维的（\(m\) 重特征值），在标准正交基里它贡献 \(m\) 个基向量，所以 \(\lambda^2\) 必须累加 \(m\) 次！
\item
  \textbf{必避小坑 2（别忘了补全 \(\operatorname{Ker} T\)）}：定理 6.13 得到的非零特征向量 \(\{u_n\}\) 只是 \(\overline{\mathcal{R}(T)}\) 的标准正交基，只有并上 \(\operatorname{Ker} T\) 的基 \(\{v_k\}\) 才合法构成 \(H\) 的完整标准正交基，才能代入定义 5.4。
\item
  \textbf{对比记忆（与 6.3.2 节迹类算子对照）}：

  \begin{itemize}
  \tightlist
  \item
    算子范数：\(\|A\| = \sup_n \lambda_n\)（\(\ell^\infty\) 范数）；
  \item
    Hilbert-Schmidt 范数：\(\|A\|_2 = \left(\sum_n \lambda_n^2\right)^{1/2}\)（\(\ell^2\) 范数，本题结论）；
  \item
    迹类范数：\(\|A\|_1 = \sum_n \lambda_n\)（\(\ell^1\) 范数，定理 6.14(i)）。
  \end{itemize}
\end{itemize}

\SoftDivider

\BookProblem{prob:6-39}{第 39 题}{教材 p. 270}{第 39 题}

\SubPillPrimary{【原题】}

设 \(A, U \in \mathcal{L}(H)\)，\(A\) 是自伴算子，\(U\) 是酉算子，\(\tilde{A} = U^{-1}AU\)。证明：对任意的有界 Borel 函数 \(f\) 都有 
\FitDisplayMath{f(\tilde{A}) = U^{-1}f(A)U,}
 这里 \(f(\tilde{A}), f(A)\) 分别表示 \(\tilde{A}\) 与 \(A\) 的函数演算。

\SubPillAmber{【零基础破题思路】}

在线性代数里，相似矩阵满足 \(f(P^{-1}AP) = P^{-1}f(A)P\)（对多项式显然成立，因为 \((P^{-1}AP)^k = P^{-1}A^k P\)）。但本题的 \(f\) 是\textbf{任意有界 Borel 函数}（可以完全不连续，比如特征函数 \(\chi_{\mathbb{Q}}\)），不能靠多项式一致逼近直接蒙混过关！ 正确的谱论王道思路是\textbf{从``谱测度''动手}：

\begin{enumerate}
\def\labelenumi{\arabic{enumi}.}
\tightlist
\item
  既然 \(A\) 有谱分解 \(A = \int_{\sigma(A)} \lambda \,\mathrm{d}E(\lambda)\)，那么 \(\tilde{A} = U^{-1}AU\) 的谱测度应该长什么样？猜都猜得到：就是 \(\tilde{E}(\Delta) \triangleq U^{-1}E(\Delta)U\)！
\item
  逐条验证两件事：

  \begin{itemize}
  \tightlist
  \item
    第一，\(\tilde{E}\) 确实是一个合法的\textbf{谱测度}（投影性、自伴性、规范性、强可列可加性，注意只有 \(U\) 是\textbf{酉算子}即 \(U^{-1} = U^*\) 时，\(\tilde{E}(\Delta)\) 才保持自伴正交投影！）；
  \item
    第二，通过弱谱积分定义，对任意 \(x, y \in H\)，发现复测度 \((\tilde{E}(\cdot)x, y)\) 和 \((E(\cdot)Ux, Uy)\) \textbf{完全相等}，于是对\textbf{所有}有界 Borel 函数 \(f\) 瞬间得到 \(\int f \,\mathrm{d}\tilde{E} = U^{-1}f(A)U\)！
  \end{itemize}
\item
  最后画龙点睛：取 \(f(\lambda) = \lambda\) 看到 \(\int \lambda \,\mathrm{d}\tilde{E}(\lambda) = \tilde{A}\)，再利用\textbf{谱测度唯一性}确认 \(\tilde{E}\) 就是 \(\tilde{A}\) 的谱测度，从而 \(\int f\,\mathrm{d}\tilde{E}\) 就是 \(f(\tilde{A})\)！完美收官！
\end{enumerate}

\SubPillMain{【规范解答】}

\textbf{第一步：验证 \(\tilde{A}\) 为有界自伴算子且 \(\sigma(\tilde{A}) = \sigma(A)\)。}

因为 \(U\) 是酉算子，所以 \(U^{-1} = U^*\)。于是 
\FitDisplayMath{\tilde{A}^* = (U^* A U)^* = U^* A^* (U^*)^* = U^{-1} A U = \tilde{A},}
 即 \(\tilde{A}\) 是 \(H\) 上的有界自伴算子，故其 Borel 函数演算 \(f(\tilde{A})\) 有意义。

对任意 \(\lambda \in \mathbb{C}\)，有 
\FitDisplayMath{\lambda I - \tilde{A} = U^{-1}(\lambda I - A)U.}
 因为 \(U, U^{-1}\) 均为双射，故 \(\lambda I - \tilde{A}\) 是双射当且仅当 \(\lambda I - A\) 是双射，即 \(\rho(\tilde{A}) = \rho(A)\)，从而 \(\sigma(\tilde{A}) = \sigma(A)\)。

\SoftDivider

\textbf{第二步：构造共轭投影族 \(\tilde{E}\) 并验证其为谱测度。}

设 \(E\) 是 \(A\) 的谱测度（由定理 6.25，集中在紧集 \(\sigma(A) \subset \mathbb{R}\) 上），记 \(\Sigma \triangleq \mathcal{B}_{\sigma(A)}\)。对任意 \(\Delta \in \Sigma\)，定义算子 
\FitDisplayMath{\tilde{E}(\Delta) \triangleq U^{-1} E(\Delta) U = U^* E(\Delta) U.}
 逐条验证定义 6.7：

\begin{enumerate}
\def\labelenumi{\arabic{enumi}.}
\tightlist
\item
  \textbf{正交投影性}：对任意 \(\Delta \in \Sigma\)，利用 \(E(\Delta)^2 = E(\Delta) = E(\Delta)^*\) 及 \(UU^{-1} = I\)、\(U^{-1} = U^*\)，有 
\FitDisplayMath{\tilde{E}(\Delta)^2 = U^{-1}E(\Delta)UU^{-1}E(\Delta)U = U^{-1}E(\Delta)^2 U = U^{-1}E(\Delta)U = \tilde{E}(\Delta),}
 
\FitDisplayMath{\tilde{E}(\Delta)^* = (U^* E(\Delta) U)^* = U^* E(\Delta)^* U = U^{-1} E(\Delta) U = \tilde{E}(\Delta).}
 故每个 \(\tilde{E}(\Delta)\) 都是 \(H\) 上的正交投影算子。
\item
  \textbf{规范性}：\(\tilde{E}(\sigma(A)) = U^{-1} E(\sigma(A)) U = U^{-1} I U = I\)（且 \(\tilde{E}(\emptyset) = 0\)）。
\item
  \textbf{可列可加性}：设 \(\{\Delta_n\}_{n=1}^\infty \subset \Sigma\) 两两不相交，\(\Delta = \bigcup_{n=1}^\infty \Delta_n\)。对任意 \(x \in H\)，因为 \(E\) 强可列可加，有 \(\lim_{N\to\infty}\sum_{n=1}^N E(\Delta_n)(Ux) = E(\Delta)(Ux)\)。再由 \(U^{-1} \in \mathcal{L}(H)\) 的连续性， 
\FitDisplayMath{\lim_{N\to\infty}\sum_{n=1}^N \tilde{E}(\Delta_n)x = \lim_{N\to\infty} U^{-1}\left(\sum_{n=1}^N E(\Delta_n)(Ux)\right) = U^{-1}E(\Delta)(Ux) = \tilde{E}(\Delta)x.}
 即 \(\tilde{E}\left(\bigcup_{n=1}^\infty \Delta_n\right) = \sum_{n=1}^\infty \tilde{E}(\Delta_n)\)（强收敛）。
\end{enumerate}

因此，\(\tilde{E}\) 是 \((\sigma(A), \mathcal{B}_{\sigma(A)})\) 上的谱测度。

\SoftDivider

\textbf{第三步：证明对任意有界 Borel 函数 \(f\)，均有 \(\int_{\sigma(A)} f(\lambda)\,\mathrm{d}\tilde{E}(\lambda) = U^{-1}f(A)U\)。}

任取 \(x, y \in H\)。对任意 \(\Delta \in \Sigma\)，由 \(U^{-1} = U^*\) 得 
\FitDisplayMath{(\tilde{E}(\Delta)x, y) = (U^* E(\Delta) U x, y) = (E(\Delta)Ux, Uy).}
 这说明 \((\sigma(A), \Sigma)\) 上的两个复有限 Borel 测度 \(\Delta \mapsto (\tilde{E}(\Delta)x, y)\) 与 \(\Delta \mapsto (E(\Delta)Ux, Uy)\) 完全相同！因此对任意有界 Borel 函数 \(f \in \mathcal{B}(\sigma(A))\)，由谱积分与弱谱积分的等价性（定义 6.9、定理 6.16）得 
\FitDisplayMath{\begin{aligned}
\left(\left(\int_{\sigma(A)} f(\lambda)\,\mathrm{d}\tilde{E}(\lambda)\right)x, y\right)
&= \int_{\sigma(A)} f(\lambda)\,\mathrm{d}(\tilde{E}(\lambda)x, y) \\
&= \int_{\sigma(A)} f(\lambda)\,\mathrm{d}(E(\lambda)Ux, Uy) \\
&= (f(A)Ux, Uy) = (U^* f(A) U x, y) = (U^{-1}f(A)Ux, y).
\end{aligned}}
 由 \(x, y \in H\) 的任意性，即得 
\FitDisplayMath{\int_{\sigma(A)} f(\lambda)\,\mathrm{d}\tilde{E}(\lambda) = U^{-1}f(A)U. \qquad (\sharp)}


\SoftDivider

\textbf{第四步：证明 \(\tilde{E}\) 就是 \(\tilde{A}\) 的唯一谱测度，从而完成证明。}

设 \(E_{\tilde{A}}\) 是自伴算子 \(\tilde{A}\) 对应的谱测度（集中在 \(\sigma(\tilde{A}) = \sigma(A)\) 上）。在 \((\sharp)\) 式中特别取 \(f(\lambda) = \lambda\)，得 
\FitDisplayMath{\int_{\sigma(\tilde{A})} \lambda \,\mathrm{d}\tilde{E}(\lambda) = U^{-1}AU = \tilde{A} = \int_{\sigma(\tilde{A})} \lambda \,\mathrm{d}E_{\tilde{A}}(\lambda).}
 我们来验证 \(\tilde{E} = E_{\tilde{A}}\)（利用 Cayley 变换与教材定理 6.21）：

\begin{itemize}
\tightlist
\item
  记 \(\varphi(\lambda) \triangleq \frac{\lambda - \mathrm{i}}{\lambda + \mathrm{i}}\)（在实紧集 \(\sigma(\tilde{A})\) 上连续有界）。由谱积分的乘法同态性（推论 6.7），\(\tilde{A}\) 的 Cayley 变换 \(\tilde{V} \triangleq (\tilde{A} - \mathrm{i}I)(\tilde{A} + \mathrm{i}I)^{-1}\) 满足 
\FitDisplayMath{\tilde{V} = \int_{\sigma(\tilde{A})} \varphi(\lambda)\,\mathrm{d}\tilde{E}(\lambda) = \int_{\sigma(\tilde{A})} \varphi(\lambda)\,\mathrm{d}E_{\tilde{A}}(\lambda).}

\item
  由引理 6.9，\(\varphi\) 是 \(\sigma(\tilde{A})\) 到 \(\sigma(\tilde{V}) \setminus \{1\}\) 的同胚。将谱测度 \(\tilde{E}\) 与 \(E_{\tilde{A}}\) 通过 \(\varphi\) 推前到 \((\sigma(\tilde{V}), \mathcal{B}_{\sigma(\tilde{V})})\) 上，分别记为 \(\tilde{F}(M) = \tilde{E}(\varphi^{-1}(M))\) 和 \(F_{\tilde{V}}(M) = E_{\tilde{A}}(\varphi^{-1}(M))\)（在点 \(\{1\}\) 处取零，因 \(1\) 非 \(\tilde{V}\) 特征值），则二者都满足 \(\tilde{V} = \int_{\sigma(\tilde{V})} w \,\mathrm{d}\tilde{F}(w) = \int_{\sigma(\tilde{V})} w \,\mathrm{d}F_{\tilde{V}}(w)\)。
\item
  由酉算子谱测度的唯一性（定理 6.21），必有 \(\tilde{F} = F_{\tilde{V}}\)，从而对任意 Borel 集 \(\Delta \in \mathcal{B}_{\sigma(\tilde{A})}\) 有 \(\tilde{E}(\Delta) = \tilde{F}(\varphi(\Delta)) = F_{\tilde{V}}(\varphi(\Delta)) = E_{\tilde{A}}(\Delta)\)！
\end{itemize}

\begin{quote}
\textbf{（第四步的极简替代证明）}：对任意多项式 \(p(\lambda)\)，由乘法同态与线性性有 \(\int_{\sigma(\tilde{A})} p(\lambda)\,\mathrm{d}\tilde{E}(\lambda) = p(\tilde{A}) = \int_{\sigma(\tilde{A})} p(\lambda)\,\mathrm{d}E_{\tilde{A}}(\lambda)\)；由 Weierstrass 逼近定理，多项式在紧集 \(\sigma(\tilde{A}) \subset \mathbb{R}\) 上的连续函数空间 \(C(\sigma(\tilde{A}))\) 中一致稠密，故对一切 \(g \in C(\sigma(\tilde{A}))\) 及任意 \(x, y \in H\) 有 \(\int g\,\mathrm{d}(\tilde{E}(\cdot)x, y) = \int g\,\mathrm{d}(E_{\tilde{A}}(\cdot)x, y)\)；由 Riesz 表示定理（紧度量空间上复 Borel 测度的唯一性）立得 \((\tilde{E}(\Delta)x, y) = (E_{\tilde{A}}(\Delta)x, y)\)，从而 \(\tilde{E} = E_{\tilde{A}}\)！
\end{quote}

将 \(\tilde{E} = E_{\tilde{A}}\) 代回 \((\sharp)\) 式，即得对任意有界 Borel 函数 \(f\)： 
\FitDisplayMath{f(\tilde{A}) = \int_{\sigma(\tilde{A})} f(\lambda)\,\mathrm{d}E_{\tilde{A}}(\lambda) = \int_{\sigma(A)} f(\lambda)\,\mathrm{d}\tilde{E}(\lambda) = U^{-1}f(A)U. \qquad \square}


\SubPillCoral{【考点复盘与避坑提示】}

\begin{itemize}
\tightlist
\item
  \textbf{常见大坑（千万不要用多项式一致逼近企图一步证明！）}：多项式只能一致逼近\textbf{连续函数}，不能一致逼近一般有界 Borel 函数！所以必须通过\textbf{谱测度共轭 + 唯一性}来做，这样对所有有界可测函数一次性成立！
\item
  \textbf{为什么必须是酉算子而不是一般可逆算子？} 因为如果 \(U\) 只是可逆算子（\(U^{-1} \neq U^*\)），那么 \(\tilde{A} = U^{-1}AU\) 连自伴都不是，\(\tilde{E}(\Delta) = U^{-1}E(\Delta)U\) 也不是正交投影（不自伴），整个自伴算子谱理论前提会直接失效。
\end{itemize}

\SoftDivider

\BookProblem{prob:6-45}{第 45 题}{教材 p. 271}{第 45 题}

\SubPillPrimary{【原题】}

设 \(A\) 是 \(H\) 中的有界自伴算子，\(E\) 是 \(A\) 的谱测度，\(f(\lambda) = \sum_{n=1}^\infty \alpha_n \lambda_n\)\footnote{\textbf{教材印刷笔误勘误}：原书 p. 271 误将幂级数通项印成了 \(\alpha_n \lambda_n\)（把上标 \(n\) 印成了下标 \(n\)），正确应为幂级数 \(f(\lambda) = \sum_{n=1}^\infty \alpha_n \lambda^n\)。} 是 \(|\lambda| < \Gamma(A) + \varepsilon,\; \varepsilon > 0\) 上的解析函数。证明： 
\FitDisplayMath{f(A) = \sum_{n=1}^\infty \alpha_n A^n = \int_{\sigma(A)} f(\lambda)\,\mathrm{d}E_\lambda.}


\SubPillAmber{【零基础破题思路】}

这道题把``微积分/复变里的泰勒幂级数代入算子'' \(\sum_{n=1}^\infty \alpha_n A^n\) 和``第 6 章高大上的谱积分函数演算'' \(\int_{\sigma(A)} f(\lambda)\,\mathrm{d}E_\lambda\) 完美地统一到了一起！ 整道题的核心逻辑就三句话，极其优雅：

\begin{enumerate}
\def\labelenumi{\arabic{enumi}.}
\tightlist
\item
  \textbf{收敛半径大于谱半径 \(\implies\) 在谱集 \(\sigma(A)\) 上一致收敛}：因为 \(\sigma(A)\) 全都装在闭圆盘 \(|\lambda| \le \Gamma(A)\) 里，而幂级数在\textbf{更大}的开圆盘 \(|\lambda| < \Gamma(A) + \varepsilon\) 里解析，所以部分和多项式 \(S_N(\lambda) = \sum_{n=1}^N \alpha_n \lambda^n\) 在 \(\sigma(A)\) 上\textbf{一致收敛}到 \(f(\lambda)\)！
\item
  \textbf{对有限多项式，两种定义完全一样}：由谱积分的线性性（推论 6.5 / 6.15(ii)）和乘法同态（推论 6.7 / 6.15(iii)），\(\int_{\sigma(A)} \lambda^n \,\mathrm{d}E_\lambda = A^n\)，从而 \(\int_{\sigma(A)} S_N(\lambda)\,\mathrm{d}E_\lambda = \sum_{n=1}^N \alpha_n A^n\)！
\item
  \textbf{套用谱积分范数估计，极限误差归零}： 
\FitDisplayMath{\left\|\int_{\sigma(A)} f(\lambda)\,\mathrm{d}E_\lambda - \sum_{n=1}^N \alpha_n A^n\right\| = \left\|\int_{\sigma(A)} \bigl(f(\lambda) - S_N(\lambda)\bigr)\,\mathrm{d}E_\lambda\right\| \le \sup_{\lambda \in \sigma(A)} |f(\lambda) - S_N(\lambda)| \xrightarrow{N\to\infty} 0!}
 这一行不等式同时干了两件事：既证明了算子级数 \(\sum_{n=1}^\infty \alpha_n A^n\) 按算子范数收敛，又证明了它的和恰好等于 \(\int_{\sigma(A)} f(\lambda)\,\mathrm{d}E_\lambda\)！
\end{enumerate}

\SubPillMain{【规范解答】}

\textbf{第一步：证明幂级数在 \(\sigma(A)\) 上一致收敛，且 \(f(A)\) 良定义。}

由谱半径定义（定理 6.5），对任意 \(\lambda \in \sigma(A)\) 都有 \(|\lambda| \le \Gamma(A)\)（且因 \(A\) 自伴，\(\sigma(A) \subset [-\Gamma(A), \Gamma(A)]\)）。因此 
\FitDisplayMath{\sigma(A) \subset D \triangleq \{\lambda \in \mathbb{C} : |\lambda| \le \Gamma(A)\}.}


因为幂级数 \(f(\lambda) = \sum_{n=1}^\infty \alpha_n \lambda^n\) 在开圆盘 \(|\lambda| < \Gamma(A) + \varepsilon\)（\(\varepsilon > 0\)）内解析，其收敛半径 \(R \ge \Gamma(A) + \varepsilon > \Gamma(A)\)。由复变函数幂级数理论，幂级数在收敛圆内的闭圆盘 \(D = \{|\lambda| \le \Gamma(A)\}\) 上绝对且一致收敛。

特别地，记第 \(N\) 个部分和多项式为 
\FitDisplayMath{S_N(\lambda) \triangleq \sum_{n=1}^N \alpha_n \lambda^n \quad (N = 1, 2, \dots),}
 则 \(\{S_N\}\) 在紧集 \(\sigma(A)\) 上一致收敛于 \(f\)，即 
\FitDisplayMath{\lim_{N\to\infty} \sup_{\lambda \in \sigma(A)} |f(\lambda) - S_N(\lambda)| = 0.}
 同时 \(f\) 在 \(\sigma(A)\) 上连续有界，属于有界 Borel 函数类 \(\mathcal{B}(\sigma(A))\)，故由推论 6.15，函数演算 \(f(A) = \int_{\sigma(A)} f(\lambda)\,\mathrm{d}E_\lambda \in \mathcal{L}(H)\) 良定义。

\SoftDivider

\textbf{第二步：证明对部分和多项式 \(S_N\) 有 \(S_N(A) = \sum_{n=1}^N \alpha_n A^n\)。}

由有界自伴算子的谱分解定理（定理 6.24、定理 6.25），有 
\FitDisplayMath{A = \int_{\sigma(A)} \lambda \,\mathrm{d}E_\lambda.}
 利用谱积分的可乘性（推论 6.7 / 推论 6.15(iii)）作数学归纳法：

\begin{itemize}
\tightlist
\item
  当 \(n = 1\) 时，\(\int_{\sigma(A)} \lambda^1 \,\mathrm{d}E_\lambda = A\) 成立；
\item
  假设对 \(n\) 成立，则对 \(n + 1\) 有 
\FitDisplayMath{\int_{\sigma(A)} \lambda^{n+1}\,\mathrm{d}E_\lambda = \left(\int_{\sigma(A)} \lambda^n \,\mathrm{d}E_\lambda\right)\left(\int_{\sigma(A)} \lambda \,\mathrm{d}E_\lambda\right) = A^n \cdot A = A^{n+1}.}

\end{itemize}

从而对一切 \(n \in \mathbb{N}\) 均有 \(\int_{\sigma(A)} \lambda^n \,\mathrm{d}E_\lambda = A^n\)。再由谱积分的线性性（推论 6.5(i) / 推论 6.15(ii)），对任意 \(N \in \mathbb{N}\) 有 
\FitDisplayMath{\int_{\sigma(A)} S_N(\lambda)\,\mathrm{d}E_\lambda = \sum_{n=1}^N \alpha_n \int_{\sigma(A)} \lambda^n \,\mathrm{d}E_\lambda = \sum_{n=1}^N \alpha_n A^n.}


\SoftDivider

\textbf{第三步：利用范数估计取极限，完成证明。}

由谱积分的线性性与范数估计（推论 6.8 / 推论 6.16），对任意 \(N \in \mathbb{N}\) 有 
\FitDisplayMath{\begin{aligned}
\left\|\int_{\sigma(A)} f(\lambda)\,\mathrm{d}E_\lambda - \sum_{n=1}^N \alpha_n A^n\right\|
&= \left\|\int_{\sigma(A)} \bigl(f(\lambda) - S_N(\lambda)\bigr)\,\mathrm{d}E_\lambda\right\| \\
&\le \sup_{\lambda \in \sigma(A)} |f(\lambda) - S_N(\lambda)|.
\end{aligned}}
 令 \(N \to \infty\)，由第一步的一致收敛性知右端趋于 \(0\)。因此，算子级数 \(\sum_{n=1}^\infty \alpha_n A^n\) 在 \(\mathcal{L}(H)\) 中按算子范数收敛，且其和等于 \(\int_{\sigma(A)} f(\lambda)\,\mathrm{d}E_\lambda = f(A)\)，即 
\FitDisplayMath{f(A) = \sum_{n=1}^\infty \alpha_n A^n = \int_{\sigma(A)} f(\lambda)\,\mathrm{d}E_\lambda. \qquad \square}


\SubPillCoral{【考点复盘与避坑提示】}

\begin{itemize}
\tightlist
\item
  \textbf{为什么条件里一定要写 \(|\lambda| < \Gamma(A) + \varepsilon\)（多出个 \(\varepsilon > 0\)）？} 因为如果幂级数只在开圆盘 \(|\lambda| < \Gamma(A)\) 内解析，在边界点 \(|\lambda| = \Gamma(A)\) 处（而自伴算子必有谱点落在 \(\pm\Gamma(A)\) 上！）幂级数可能发散或不保证一致收敛！多出一个 \(\varepsilon > 0\)，才保证了闭圆盘 \(|\lambda| \le \Gamma(A)\) 严格落在收敛圆内部，从而绝对且一致收敛！
\item
  \textbf{含常数项的推广}：若 \(f(\lambda) = \sum_{n=0}^\infty \alpha_n \lambda^n\) 含常数项 \(\alpha_0\)，由于 \(\int_{\sigma(A)} 1\,\mathrm{d}E_\lambda = E(\sigma(A)) = I\)，约定 \(A^0 = I\) 后完全同理有 \(f(A) = \sum_{n=0}^\infty \alpha_n A^n = \int_{\sigma(A)} f(\lambda)\,\mathrm{d}E_\lambda\)。例如指数算子 \(\mathrm{e}^{tA} = \sum_{n=0}^\infty \frac{t^n}{n!} A^n = \int_{\sigma(A)} \mathrm{e}^{t\lambda}\,\mathrm{d}E_\lambda\) 对任意有界自伴算子 \(A\) 恒成立！
\item
  \textbf{另一角度验证收敛性}：因 \(A\) 自伴，有 \(\|A\| = \Gamma(A)\)，而 \(\sum_{n=1}^\infty |\alpha_n| \|A\|^n = \sum_{n=1}^\infty |\alpha_n|\Gamma(A)^n < \infty\)，所以在 Banach 空间 \(\mathcal{L}(H)\) 中级数 \(\sum_{n=1}^\infty \alpha_n A^n\) 本身也绝对收敛（上面第三步直接连收敛性和极限值一起证完了，更紧凑）。
\end{itemize}

\BookPart{}{考前 5 分钟自检清单}{考前 5 分钟自检清单}

\begin{enumerate}
\def\labelenumi{\arabic{enumi}.}
\tightlist
\item
  \textbf{谱的三大分类}：能一口气说出 \(\sigma_p\)（非单射）、\(\sigma_c\)（单射 + 真稠密值域）、\(\sigma_r\)（单射 + 不稠密值域），并记住 \(L^2(0,1)\) 乘法算子只有连续谱 \([0,1]\)。
\item
  \textbf{Neumann 级数与谱半径}：记住 \(\|B\| < 1 \implies (I - B)^{-1} = \sum_{k=0}^\infty B^k\)；Gelfand 公式 \(\Gamma(A) = \lim_{n\to\infty}\sqrt[n]{\|A^n\|} \le \|A\|\)；\textbf{自伴算子}才有 \(\Gamma(A) = \|A\|\)。
\item
  \textbf{紧算子 vs 自伴紧算子}：紧算子非零谱点必是有限重特征值、至多可数、只可能聚于 \(0\)；加上自伴后有正交基展开 \(Ax = \sum \lambda_n(x, e_n)e_n\) 且 \(\pm\|A\|\) 必为特征值。
\item
  \textbf{求积分算子范数的必杀链}：先求 \(A^*\) \(\to\) 证 \(AA^*\) 紧正自伴 \(\to\) 解 ODE 求最大特征值 \(\Gamma(AA^*)\) \(\to\) \(\|A\| = \sqrt{\Gamma(AA^*)}\)。
\item
  \textbf{绝对值算子黄金等式}：\(|A| = (A^*A)^{1/2} \implies \lVert\,|A|x\,\rVert = \lVert Ax\rVert\)，瞬间连通 \(A\) 与正自伴算子 \(|A|\)。
\item
  \textbf{谱积分与 Borel 函数演算四大法则}：线性、共轭（\(f \mapsto \bar{f}\)）、\textbf{乘法同态}（积的积分 = 积分的积）、范数估计 \(\|f(A)\| \le \sup_{\sigma(A)} |f|\)；以及正算子唯一正平方根 \(A^{1/2}\)。
\end{enumerate}


% =========================================================================
% 第 7 章 · Hilbert 空间上的无界算子
% =========================================================================
\SwitchChapterTheme{7}
\ChapterBanner{07}{ch:7}{第 7 章 · Hilbert 空间上的无界算子}{稠定与共轭 · 对称 vs 自伴 · 闭算子与自伴判别 · 无界谱分解 · 习题 7（4 题）}
\begin{quote}
\textbf{使用说明}：本资料面向\textbf{零泛函分析基础}的本科生期末急救与刷题，内容、记号与定理编号严格对齐教材第 7 章（p. 272--294）。所有作业解答均自包含、步骤严谨、零跳步，可直接作为期末考试规范作答模板。
\end{quote}

\BookPart{}{零基础导读：3 分钟看懂第 7 章到底在学什么}{零基础导读：3 分钟看懂第 7 章到底在学什么}

前六章里，我们一直生活在``温室''里：算子 \(A \in \mathcal{L}(H)\) 都是\textbf{有界}的，而且\textbf{定义在整个空间 \(H\) 上}。但只要走出纯数学温室、踏进\textbf{微分方程}和\textbf{量子力学}的真实世界，你遇到的核心算子几乎全都是\textbf{无界算子}——它们根本躲不掉！

先看两个让你一秒恍然大悟的``主角''：

\begin{enumerate}
\def\labelenumi{\arabic{enumi}.}
\tightlist
\item
  \textbf{求导算子（微分算子 / 量子力学中的动量算子）}：\(u(x) \mapsto u'(x)\)（或 \(-\mathrm{i} u'(x)\)）。

  \begin{itemize}
  \tightlist
  \item
    取一串振幅始终为 \(1\) 的高频正弦波 \(u_n(x) = \sin(nx)\)（\(x \in [0, 2\pi]\)），它们的``大小''一致固定（\(\|u_n\|_{L^2} = \sqrt{\pi}\)），但求导后 \(u_n'(x) = n\cos(nx)\)，大小被放大了 \textbf{\(n\) 倍}！随着频率 \(n \to \infty\)，放大倍数冲向 \(+\infty\)！
  \item
    而且，不是每个 \(L^2\) 函数（甚至不是每个连续函数）都可导——所以求导算子\textbf{不可能定义在全空间上}！
  \end{itemize}
\item
  \textbf{位置算子（乘法算子，教材例 7.1）}：在 \(L^2(\mathbb{R})\) 上定义 \((Af)(t) = t f(t)\)。

  \begin{itemize}
  \tightlist
  \item
    把一个宽度为 \(1\)、高度为 \(1\) 的方波脉冲 \(f_n(t) = \chi_{[n, n+1]}(t)\) 往远处平移到 \([n, n+1]\)，脉冲本身能量不变（\(\|f_n\|_2 = 1\)），但乘上 \(t\) 之后，幅度暴涨到至少 \(n\) 倍（\(\|Af_n\|_2 \ge n \to +\infty\)）！
  \item
    而且，像 \(f(t) = \frac{1}{1+|t|}\) 这样的函数明明属于 \(L^2(\mathbb{R})\)，乘上 \(t\) 变成 \(\frac{t}{1+|t|}\) 却\textbf{掉出了 \(L^2(\mathbb{R})\)}！所以 \(A\) 的定义域最大也只能取 \(\mathcal{D}(A) = \{f \in L^2(\mathbb{R}) \mid \int_{\mathbb{R}} t^2|f(t)|^2\,\mathrm{d}t < \infty\}\)，绝不可能是整个 \(L^2(\mathbb{R})\)！
  \end{itemize}
\end{enumerate}

进入无界算子的世界后，有三件颠覆直觉、也是期末最爱考的大事：

\begin{enumerate}
\def\labelenumi{\arabic{enumi}.}
\tightlist
\item
  \textbf{定义域 \(\mathcal{D}(A)\) 是算子的灵魂（不再只是背景板）}：同样的运算公式（比如 \(f \mapsto \mathrm{i} f'\)），给它配不同的边界条件（即不同的定义域），它可能\textbf{连对称都不是}，也可能\textbf{仅仅对称但不自伴}，也可能是完美的\textbf{自伴算子}！在无界算子里，\textbf{必须把一对 \((A, \mathcal{D}(A))\) 当成一个整体}。
\item
  \textbf{``闭算子''接棒``连续算子''}：无界算子都不连续，直接取极限 \(\lim Ax_n = A(\lim x_n)\) 会失效；于是数学家退一步，寻找``如果 \(x_n \to x\) \textbf{并且} \(Ax_n\) 自己也收敛到 \(y\)，那 \(x\) 必在定义域内且 \(Ax=y\)''的\textbf{闭算子}——幸运的是，所有共轭算子 \(A^*\)（特别是自伴算子）自动都是闭算子！
\item
  \textbf{自伴算子依然能像矩阵一样做谱分解}：虽然无界自伴算子的谱不再有界（比如位置算子的谱是整条实轴 \(\mathbb{R}\)），但利用神奇的 \textbf{Cayley 变换} \(U = (A-\mathrm{i}I)(A+\mathrm{i}I)^{-1}\)，我们可以把整条无界实轴``压扁卷成''单位圆周，把无界自伴算子 \(A\) 变成第 6 章已经解决的\textbf{有界酉算子 \(U\)}，借它的谱分解反推回来，就得到了无界自伴算子的谱分解定理 \(A = \int_{\mathbb{R}} \lambda\,\mathrm{d}E_\lambda\)！
\end{enumerate}

\BookSubSection{}{有界算子（第 6 章）vs 无界算子（第 7 章）全景对照表}{有界算子（第 6 章）vs 无界算子（第 7 章）全景对照表}{}

{\def\LTcaptype{none} % do not increment counter
{\small\renewcommand{\arraystretch}{1.38}\setlength{\tabcolsep}{5pt}
\begin{longtable}[]{@{}
  >{\raggedright\arraybackslash}p{(\linewidth - 4\tabcolsep) * \real{0.3333}}
  >{\raggedright\arraybackslash}p{(\linewidth - 4\tabcolsep) * \real{0.3333}}
  >{\raggedright\arraybackslash}p{(\linewidth - 4\tabcolsep) * \real{0.3333}}@{}}
\toprule\noalign{}
\rowcolor{AccentSoft}
\begin{minipage}[b]{\linewidth}\raggedright
对比维度
\end{minipage} & \begin{minipage}[b]{\linewidth}\raggedright
有界线性算子（第 6 章）
\end{minipage} & \begin{minipage}[b]{\linewidth}\raggedright
无界线性算子（第 7 章）
\end{minipage} \\
\midrule\noalign{}
\endhead
\bottomrule\noalign{}
\endlastfoot
\textbf{定义域 \(\mathcal{D}(A)\)} & 默认是全空间 \(H\) & 通常是真子空间，基本要求在 \(H\) 中\textbf{稠密}（稠定算子） \\
\textbf{算子相等 \(A = B\)} & 只需对所有 \(x \in H\) 有 \(Ax = Bx\) & 必须同时满足 \textbf{\(\mathcal{D}(A) = \mathcal{D}(B)\)} 且作用公式相同 \\
\textbf{共轭算子 \(A^*\)} & 永远存在、唯一，且 \(\mathcal{D}(A^*)=H\) & 必须 \(A\) \textbf{稠定}才能唯一定义 \(A^*\)；\(\mathcal{D}(A^*)\) 甚至可能只有 \(\{0\}\) \\
\textbf{对称 vs 自伴} & 全空间对称自动有界且自伴（二者无区别！） & \textbf{天壤之别}：对称仅要求 \(A \subset A^*\)；自伴要求 \(A = A^*\)（含定义域相等！） \\
\textbf{连续性 / 闭性} & 有界 \(\iff\) 连续 \(\implies\) 闭算子 & 不连续，但核心研究对象（\(A^*\)、自伴算子）恒为\textbf{闭算子} \\
\textbf{二阶共轭 \(A^{**}\)} & 恒有 \(A^{**} = A\) & 需 \(\mathcal{D}(A^*)\) 稠密才存在；\(A\) 稠定闭 \(\iff A = A^{**}\) \\
\textbf{自伴算子的谱 \(\sigma(A)\)} & 实轴上的非空\textbf{有界闭集}（紧集） & 实轴上的闭集 \(\sigma(A) \subset \mathbb{R}\)，但\textbf{可以无界}（如 \(\mathbb{R}\)），且无剩余谱 \\
\textbf{谱分解定理} & \(A = \int_{\sigma(A)} \lambda\,\mathrm{d}E_\lambda\)（有界积分） & \(A = \int_{\mathbb{R}} \lambda\,\mathrm{d}E_\lambda\)，且定义域由 \(\int_{\mathbb{R}}\lambda^2\,\mathrm{d}\lVert E_\lambda x\rVert^2 < \infty\) 刻画 \\
\end{longtable}
}
}

\BookSubSection{}{第 7 章 4 道课后作业题考点快速索引}{第 7 章 4 道课后作业题考点快速索引}{}

{\def\LTcaptype{none} % do not increment counter
{\small\renewcommand{\arraystretch}{1.38}\setlength{\tabcolsep}{5pt}
\begin{longtable}[]{@{}
  >{\centering\arraybackslash}p{(\linewidth - 6\tabcolsep) * \real{0.2778}}
  >{\centering\arraybackslash}p{(\linewidth - 6\tabcolsep) * \real{0.2778}}
  >{\raggedright\arraybackslash}p{(\linewidth - 6\tabcolsep) * \real{0.2222}}
  >{\raggedright\arraybackslash}p{(\linewidth - 6\tabcolsep) * \real{0.2222}}@{}}
\toprule\noalign{}
\rowcolor{AccentSoft}
\begin{minipage}[b]{\linewidth}\centering
题号
\end{minipage} & \begin{minipage}[b]{\linewidth}\centering
教材页码
\end{minipage} & \begin{minipage}[b]{\linewidth}\raggedright
核心考点
\end{minipage} & \begin{minipage}[b]{\linewidth}\raggedright
所用关键定理 / 核心技巧
\end{minipage} \\
\midrule\noalign{}
\endhead
\bottomrule\noalign{}
\endlastfoot
\hyperlink{prob:7-3}{\textbf{第 3 题}} & p. 293 & Legendre 微分算子的全部特征值与特征向量 & 多项式空间 \(\mathcal{P}_n\) 上三角矩阵 + Rodrigues 公式 + 分部积分对称性 + Weierstrass 多项式逼近 \\
\hyperlink{prob:7-4}{\textbf{第 4 题}} & p. 293 & \(L^2(\mathbb{R})\) 上乘法（位置）算子的无界性与闭性 & 示性函数 \(\chi_{[n,n+1]}\) 证无界 + 两次抽几乎处处收敛子列证闭性 \\
\hyperlink{prob:7-6}{\textbf{第 6 题}} & p. 293 & 对称算子的闭性 / 自伴性与 \(A, A^{**}, A^*\) 的关系 & 前置验证 \(\mathcal{D}(A^*)\) 稠密 + 定理 7.1（共轭必闭）+ 定理 7.2（稠定闭 \(\implies A=A^{**}\)） \\
\hyperlink{prob:7-7}{\textbf{第 7 题}} & p. 294 & 非负对称算子（\((Ax,x) \ge 0\)）的范数估计与闭性充要条件 & 内积展开得双控不等式 + Cauchy 列与 \(H\) 完备性 + 极限唯一性 \\
\end{longtable}
}
}

\BookPart{}{第一部分：第 7 章核心知识零基础精讲}{第一部分：第 7 章核心知识零基础精讲}

\begin{quote}
\textbf{全章基础约定}：如无特别说明，本章中 \(H\) 始终表示\textbf{复数域 \(\mathbb{C}\) 上的 Hilbert 空间}（完备内积空间），内积 \((\cdot, \cdot)\) 对第一变元线性、对第二变元共轭线性（即 \((\alpha x, \beta y) = \alpha \bar{\beta}(x, y)\)）。一个无界线性算子永远由两样东西共同决定：\textbf{线性子空间 \(\mathcal{D}(A) \subset H\)（定义域）} 与 \textbf{线性映射 \(A: \mathcal{D}(A) \to H\)}。
\end{quote}

\SoftDivider

\BookSubSection{sec:7-1-1}{7.1.1 算子的定义域、扩张与稠定算子的共轭算子（p. 272--275）}{7.1.1 算子的定义域、扩张与稠定算子的共轭算子（p. 272–275）}{\TagMust{闭卷必背}}

\SubSubHeading{1. 带定义域的算子运算与``扩张''（定义 7.1，p. 272--273）}

设 \(A, B\) 是分别以线性子空间 \(\mathcal{D}(A), \mathcal{D}(B) \subset H\) 为定义域的线性算子，\(\alpha \in \mathbb{F}\)（\(\mathbb{R}\) 或 \(\mathbb{C}\)）：

\begin{itemize}
\tightlist
\item
  \textbf{(i) 数乘 \(\alpha A\)}：定义域为 \(\mathcal{D}(A)\)，作用 \(x \mapsto \alpha Ax\)。

  \begin{itemize}
  \tightlist
  \item
    \textbf{零基础必踩避坑点}：\(0A\)（数 \(0\) 乘算子 \(A\)）的定义域仍是 \(\mathcal{D}(A)\)，而零算子 \(0\) 的定义域是全空间 \(H\)！所以当 \(\mathcal{D}(A) \subsetneq H\) 时，\textbf{\(0A \neq 0\)}，只能写 \(0A \subset 0\)！（当然作用在零向量上恒有 \(A0 = 0\)。）
  \end{itemize}
\item
  \textbf{(ii) 和 \(A + B\)}：定义域只能取公共部分 \(\mathcal{D}(A+B) = \mathcal{D}(A) \cap \mathcal{D}(B)\)，作用 \(x \mapsto Ax + Bx\)。
\item
  \textbf{(iii) 复合（积）\(BA\)}：先算 \(Ax\)，再要求 \(Ax\) 落进 \(B\) 的定义域，故 \(\mathcal{D}(BA) = \{x \in \mathcal{D}(A) \mid Ax \in \mathcal{D}(B)\}\)，作用 \(x \mapsto B(Ax)\)。

  \begin{itemize}
  \tightlist
  \item
    特别地，例 7.1 中乘法算子 \((Af)(t)=tf(t)\) 的 \(n\) 次幂 \(A^n\) 的定义域是 \(\mathcal{D}(A^n) = \{f \in L^2(\mathbb{R}) \mid \int_{\mathbb{R}} |t^n f(t)|^2\,\mathrm{d}t < \infty\}\)。
  \end{itemize}
\item
  \textbf{(iv) 延拓 / 扩张 \(A \subset B\)}：如果 \textbf{\(\mathcal{D}(A) \subset \mathcal{D}(B)\)}，并且对所有 \(x \in \mathcal{D}(A)\) 都有 \textbf{\(Ax = Bx\)}，就称 \(B\) 是 \(A\) 的\textbf{延拓（或扩张）}，记作 \textbf{\(A \subset B\)}（或 \(B \supset A\)）。

  \begin{itemize}
  \tightlist
  \item
    \emph{通俗直觉}：\(B\) 干的活跟 \(A\) 一模一样，只是 \(B\) 的``管辖范围（定义域）''比 \(A\) 更大（或相等）。
  \end{itemize}
\item
  \textbf{(v) 算子相等 \(A = B\)}：\(A \subset B\) 且 \(B \subset A\)，即 \textbf{\(\mathcal{D}(A) = \mathcal{D}(B)\) 且对一切 \(x \in \mathcal{D}(A)\) 有 \(Ax = Bx\)}。
\end{itemize}

\SubSubHeading{2. 稠定算子与共轭算子 \(A^*\)（定义 7.2--7.3、引理 7.1，p. 273--274）}

在有界算子里，共轭算子 \(A^*\) 由公式 \((Ax, y) = (x, A^*y)\)（\(\forall x, y \in H\)）定义。现在 \(A\) 只定义在子空间 \(\mathcal{D}(A)\) 上，我们还能用同一招来定义 \(A^*\) 吗？

这里会遇到\textbf{两个必须解决的问题}：

\begin{enumerate}
\def\labelenumi{\arabic{enumi}.}
\tightlist
\item
  \textbf{唯一性问题}：如果已知对一切 \(x \in \mathcal{D}(A)\) 都有 \((Ax, y) = (x, y^*)\)，这个 \(y^*\) 会不会不唯一？

  \begin{itemize}
  \tightlist
  \item
    设 \(y_1^*, y_2^*\) 都满足上式，相减得 \((x, y_1^* - y_2^*) = 0\)（\(\forall x \in \mathcal{D}(A)\)），即 \(y_1^* - y_2^* \perp \mathcal{D}(A)\)。
  \item
    要逼出 \(y_1^* - y_2^* = 0\)，当且仅当 \(\mathcal{D}(A)^\perp = \{0\}\)，也就是 \textbf{\(\overline{\mathcal{D}(A)} = H\)（\(\mathcal{D}(A)\) 在 \(H\) 中稠密）}！
  \end{itemize}
\item
  \textbf{存在性问题}：即使 \(\mathcal{D}(A)\) 稠密，也不是每个 \(y \in H\) 都能找到对应的 \(y^*\)！

  \begin{itemize}
  \tightlist
  \item
    固定 \(y \in H\)，看线性泛函 \(\varphi_y(x) \triangleq (Ax, y)\)（\(x \in \mathcal{D}(A)\)）。
  \item
    要存在 \(y^* \in H\) 使 \(\varphi_y(x) = (x, y^*)\)，由 \textbf{Riesz 表示定理}（将连续线性泛函从稠密子空间唯一延拓到 \(H\)），充要条件是 \(\varphi_y\) 在 \(\mathcal{D}(A)\) 上\textbf{有界}，即存在常数 \(M_y > 0\)（依赖于 \(y\)）使得 
\FitDisplayMath{|(Ax, y)| \le M_y \|x\|, \qquad \forall x \in \mathcal{D}(A).}

  \item
    只有那些能让上式成立的``好向量 \(y\)''，才有资格进入 \(\mathcal{D}(A^*)\)！
  \end{itemize}
\end{enumerate}

这就精准引出了教材的\textbf{定义 7.2、定义 7.3 与引理 7.1}：

\begin{itemize}
\tightlist
\item
  \textbf{定义 7.2（稠定算子，p. 273）}：若线性子空间 \(\mathcal{D}(A)\) 在 \(H\) 中稠密（\(\overline{\mathcal{D}(A)} = H\)），则称 \(A: \mathcal{D}(A) \to H\) 是 \(H\) 上的\textbf{稠定算子}。
\item
  \textbf{定义 7.3（共轭算子，p. 273）}：设 \(A\) 是 \(H\) 上的稠定线性算子，令 
\FitDisplayMath{\mathcal{D}(A^*) \triangleq \left\{y \in H \;\middle|\; \text{存在 } y^* \in H \text{ 使得对任意 } x \in \mathcal{D}(A),\ (Ax, y) = (x, y^*)\right\},}
 定义算子 \(A^*: \mathcal{D}(A^*) \to H\) 为 \(A^* y = y^*\)，称 \(A^*\) 为 \(A\) 的\textbf{共轭算子}（又称伴随算子）。
\item
  \textbf{引理 7.1（唯一性的充要条件，p. 273--274）}：对 \(y \in H\)，满足 \((Ax, y) = (x, y^*)\)（\(\forall x \in \mathcal{D}(A)\)）的向量 \(y^* \in H\) \textbf{最多只有一个}的充要条件是 \textbf{\(A\) 稠定}。
\end{itemize}

\begin{quote}
\textbf{零基础警示}：千万别以为``\(A\) 稠定''就能推出``\(\mathcal{D}(A^*)\) 稠密''！事实上存在稠定算子满足 \(\mathcal{D}(A^*) = \{0\}\)（习题 7 第 1 题就指出：\(\mathcal{D}(A^*) = \{0\} \iff A\) 的图像在 \(H \times H\) 中稠密）。所以写 \(A^{**} = (A^*)^*\) 之前，\textbf{必须先验证 \(\mathcal{D}(A^*)\) 在 \(H\) 中稠密}！（这是\hyperlink{prob:7-6}{\textcolor{LinkColor}{习题 7 第 6 题}}的关键得分点！）
\end{quote}

\SubSubHeading{3. 共轭算子的基本运算性质与``扩张反号铁律''（引理 7.2，p. 274--275）}

\begin{itemize}
\tightlist
\item
  \textbf{(i) 线性性}：\(A^*\) 是线性算子（即 \(\mathcal{D}(A^*)\) 是 \(H\) 的线性子空间，且 \(A^*(\alpha y_1 + \beta y_2) = \alpha A^* y_1 + \beta A^* y_2\)）。
\item
  \textbf{(ii) 和的共轭}：若 \(A, B\) 及 \(A+B\) 均稠定（即 \(\mathcal{D}(A)\cap\mathcal{D}(B)\) 稠密），则 
\FitDisplayMath{A^* + B^* \subset (A + B)^*.}

\item
  \textbf{(iii) 扩张反号铁律（极其高频！）}：设 \(A, B\) 都是稠定线性算子，若 \textbf{\(A \subset B\)}，则 \textbf{\(B^* \subset A^*\)}（即 \(A^* \supset B^*\)）。

  \begin{itemize}
  \tightlist
  \item
    \textbf{一句话悟透``为什么会反号''}： 想想进入 \(\mathcal{D}(B^*)\) 的``入场考试''：你必须通过\textbf{所有 \(x \in \mathcal{D}(B)\)} 的考题 \((Bx, y) = (x, y^*)\)；而进入 \(\mathcal{D}(A^*)\) 只需通过\textbf{子集 \(\mathcal{D}(A)\)} 的考题 \((Ax, y) = (x, y^*)\)！ \textbf{考题越多（\(A \subset B\)，范围更大、约束越严），能通过考试的人反而越少（\(\mathcal{D}(B^*) \subset \mathcal{D}(A^*)\)）！} 记住这个直觉，这辈子都不会把包含号写反！
  \end{itemize}
\item
  \textbf{(iv) 积的共轭}：设 \(A, B\) 和 \(BA\) 均为稠定算子，则 
\FitDisplayMath{A^* B^* \subset (BA)^*.}
 特别地，如果还有 \(B \in \mathcal{L}(H)\)（外层算子是全空间有界算子），则等号成立：\textbf{\(A^* B^* = (BA)^*\)}。
\end{itemize}

\begin{quote}
\textbf{教材印刷勘误提示（p. 274 引理 7.2）}：教材 p. 274 引理 7.2 的 (ii) 陈述与末尾把包含符号印反成了 \((A+B)^* \subset A^* + B^*\)，但查看该页正文证明第一行``如果 \(y \in \mathcal{D}(A^*+B^*)\)，则 \(y \in \mathcal{D}((A+B)^*)\)''，实际证明出的正确方向就是 \(A^* + B^* \subset (A+B)^*\)；(iv) 亦存在 \(AB\) 与 \(BA\) 混印，正确形式为上文给出的 \(A^* B^* \subset (BA)^*\)（当 \(B \in \mathcal{L}(H)\) 时取等）。复习与答题请以正确结论为准！
\end{quote}

\SoftDivider

\BookSubSection{sec:7-1-2}{7.1.2 对称算子 vs 自伴算子：一字之差，命运天壤之别（p. 275--277）}{7.1.2 对称算子 vs 自伴算子：一字之差，命运天壤之别（p. 275–277）}{\TagMust{闭卷必背}}

\SubSubHeading{1. 精确定义与等价判别（定义 7.4、引理 7.3--7.4，p. 275--277）}

\begin{itemize}
\tightlist
\item
  \textbf{定义 7.4（p. 275）}：设 \(A\) 是 \(H\) 上的\textbf{稠定}线性算子：

  \begin{itemize}
  \tightlist
  \item
    若 \textbf{\(A \subset A^*\)}，称 \(A\) 为\textbf{对称算子}（Symmetric operator）；
  \item
    若 \textbf{\(A = A^*\)}，称 \(A\) 为\textbf{自伴算子}（Self-adjoint operator）。
  \end{itemize}
\item
  \textbf{引理 7.3（对称算子的实操检验公式，p. 277）}：稠定线性算子 \(A\) 是对称算子，当且仅当对任意 \(x, y \in \mathcal{D}(A)\)，都有 
\FitDisplayMath{(Ax, y) = (x, Ay).}
 \emph{证明直观}：只要对所有 \(x, y \in \mathcal{D}(A)\) 满足 \((Ax, y) = (x, Ay)\)，按定义 7.3 立即看出每个 \(y \in \mathcal{D}(A)\) 自动满足 \(y \in \mathcal{D}(A^*)\) 且 \(A^* y = Ay\)，即 \(A \subset A^*\)！但注意：这\textbf{完全无法保证反方向的 \(\mathcal{D}(A^*) \subset \mathcal{D}(A)\)}！
\item
  \textbf{引理 7.4（自伴算子的极大性，p. 277）}：设 \(A\) 是自伴算子，\(B\) 是对称算子，如果 \(A \subset B\)，那么必有 \textbf{\(A = B\)}！

  \begin{itemize}
  \tightlist
  \item
    \emph{一行神证明}：由 \(A \subset B\) 及扩张反号律得 \(B^* \subset A^*\)，于是 
\FitDisplayMath{A = A^* \supset B^* \supset B \supset A \implies A = B!}

  \item
    \emph{几何直觉}：对称算子一旦``长''到了自伴，就已经到达了对称性的天花板，再也不可能被真扩张成更大的对称算子！
  \end{itemize}
\end{itemize}

\SubSubHeading{2. 醒目对比表：对称算子 vs 自伴算子}

{\def\LTcaptype{none} % do not increment counter
{\small\renewcommand{\arraystretch}{1.38}\setlength{\tabcolsep}{5pt}
\begin{longtable}[]{@{}
  >{\raggedright\arraybackslash}p{(\linewidth - 4\tabcolsep) * \real{0.3333}}
  >{\raggedright\arraybackslash}p{(\linewidth - 4\tabcolsep) * \real{0.3333}}
  >{\raggedright\arraybackslash}p{(\linewidth - 4\tabcolsep) * \real{0.3333}}@{}}
\toprule\noalign{}
\rowcolor{AccentSoft}
\begin{minipage}[b]{\linewidth}\raggedright
对比项
\end{minipage} & \begin{minipage}[b]{\linewidth}\raggedright
对称算子（Symmetric）
\end{minipage} & \begin{minipage}[b]{\linewidth}\raggedright
自伴算子（Self-adjoint）
\end{minipage} \\
\midrule\noalign{}
\endhead
\bottomrule\noalign{}
\endlastfoot
\textbf{核心关系式} & \(A \subset A^*\) & \(A = A^*\) \\
\textbf{定义域关系} & 允许严格更小：\(\mathcal{D}(A) \subsetneq \mathcal{D}(A^*)\) & 必须完全相等：\textbf{\(\mathcal{D}(A) = \mathcal{D}(A^*)\)} \\
\textbf{验证难度} & \textbf{容易}：只需验证 \(\mathcal{D}(A)\) 稠密且 \((Ax,y)=(x,Ay)\) & \textbf{困难得多}：除对称外，必须证明 \(\mathcal{D}(A^*) \subset \mathcal{D}(A)\)（或用定理 7.4 判别） \\
\textbf{闭性} & 未必是闭算子（但必\textbf{可闭}，因 \(\mathcal{D}(A^*)\) 稠密） & \textbf{必然是闭算子}（定理 7.1） \\
\textbf{谱的性质} & 谱\textbf{未必}在实轴上！甚至可能是整个复平面 \(\mathbb{C}\)（见例 7.2 中的 \(A_3\)） & 谱必然全在实轴上 \textbf{\(\sigma(A) \subset \mathbb{R}\)}，且无剩余谱（推论 7.2） \\
\textbf{谱分解定理} & \textbf{不成立}（无法生成酉群 \(\mathrm{e}^{\mathrm{i}tA}\)，在量子力学中不能代表物理观测量） & \textbf{成立} \(A = \int_{\mathbb{R}} \lambda\,\mathrm{d}E_\lambda\)（定理 7.8） \\
\end{longtable}
}
}

\SubSubHeading{3. 教科书级活例子：边界条件如何决定算子的命运（例 7.2，p. 275--277）}

在 \(H = L^2(0, 1)\) 上，考察\textbf{完全相同的微分运算公式} \(A_k f = \mathrm{i} f'\)，只改变边界条件（\(AC[0,1]\) 表示 \([0,1]\) 上的绝对连续函数，可分部积分）：


\FitDisplayMath{\begin{aligned}
\mathcal{D}(A_1) &= \{f \in L^2(0,1) \mid f \in AC[0,1],\ f' \in L^2(0,1)\} && \text{（最松：无任何边界条件）}, \\
\mathcal{D}(A_2) &= \{f \in \mathcal{D}(A_1) \mid f(0) = f(1)\} && \text{（适中：周期边界条件）}, \\
\mathcal{D}(A_3) &= \{f \in \mathcal{D}(A_1) \mid f(0) = f(1) = 0\} && \text{（最紧：两端固定为零）}.
\end{aligned}}


显然 \(A_3 \subsetneq A_2 \subsetneq A_1\)，且三者都是稠定的。做一次分部积分看边界项： 
\FitDisplayMath{(A_k f, g) = \int_0^1 \mathrm{i} f'(x) \overline{g(x)}\,\mathrm{d}x = \underbrace{\mathrm{i}\left[f(1)\overline{g(1)} - f(0)\overline{g(0)}\right]}_{\text{必须消掉的边界项！}} + \int_0^1 f(x) \overline{\mathrm{i} g'(x)}\,\mathrm{d}x.}


为了消掉这个边界项，三个算子走上了截然不同的命运：

\begin{itemize}
\tightlist
\item
  \textbf{\(A_3\)（约束太紧）}：\(f(0)=f(1)=0\) 直接把边界项杀死了，所以对 \(g\) 的端点值\textbf{没有任何要求}！结果 \textbf{\(A_3^* = A_1 \supsetneq A_3\)} —— \(A_3\) 是\textbf{对称算子，但不是自伴算子}！（它的谱甚至是整个复平面 \(\sigma(A_3) = \mathbb{C}\)！）
\item
  \textbf{\(A_1\)（约束太松）}：\(f(0), f(1)\) 随便取，为了消掉边界项，只好强迫 \(g(0)=g(1)=0\)！结果 \textbf{\(A_1^* = A_3 \subsetneq A_1\)} —— \(A_1\) \textbf{连对称算子都不是}！
\item
  \textbf{\(A_2\)（松紧恰到好处）}：\(f(0)=f(1)\)，边界项变成 \(\mathrm{i} f(0)[\overline{g(1)} - \overline{g(0)}]\)，为使其对一切 \(f(0)\) 为零，恰好等价于要求 \(g(0)=g(1)\)！结果 \textbf{\(A_2^* = A_2\)} —— \textbf{\(A_2\) 是自伴算子}！（其谱恰好是离散实点集 \(\{2n\pi \mid n \in \mathbb{Z}\}\)！）
\end{itemize}

\begin{quote}
\textbf{灵魂总结}：结合马上要学的 \textbf{Hellinger-Toeplitz 定理（定理 7.3 (i)）}——``定义在全空间 \(H\) 上的对称算子必定有界''，我们立刻明白：\textbf{任何无界对称算子都绝不可能定义在整个 \(H\) 上}。为了成为自伴算子，定义域不能太小（太小则共轭太大，只是对称），也不能太大（太大则共轭太小，失去对称），必须精准调配边界条件达到 \(A = A^*\) 的平衡！
\end{quote}

\SubSubHeading{4. 两大核心自伴算子范例与酉等价（例 7.3、7.1.3 节引理 7.5、定义 7.5、例 7.4，p. 277--279）}

\begin{itemize}
\tightlist
\item
  \textbf{例 7.3（位置/乘法算子自伴，p. 277）}：\(L^2(\mathbb{R})\) 上的乘法算子 \((Af)(t) = tf(t)\)（定义域 \(\mathcal{D}(A) = \{f \in L^2(\mathbb{R}) \mid tf \in L^2(\mathbb{R})\}\)）不仅是闭的无界算子（\hyperlink{prob:7-4}{\textcolor{LinkColor}{习题 7 第 4 题}}），而且是\textbf{自伴算子}（\(A = A^*\)）。
\item
  \textbf{引理 7.5 \& 定义 7.5（酉等价，p. 277--278）}：若 \(U\) 是 \(H\) 上的酉算子（\(U^* = U^{-1}\)），称 \(B = UAU^{-1}\)（定义域 \(\mathcal{D}(B) = U\mathcal{D}(A)\)）与 \(A\) \textbf{酉等价}。\textbf{若 \(A\) 自伴，则 \(UAU^{-1}\) 依然自伴}（且谱完全相同！）。
\item
  \textbf{例 7.4（Fourier 变换把微分变成乘法，p. 278--279）}：在 \(L^2(\mathbb{R})\) 中，动量算子 \(A_1 f = \frac{1}{\mathrm{i}} f'\)（定义域为局部绝对连续且 \(f, f' \in L^2(\mathbb{R})\)）通过酉的 Fourier 变换 \(U\) 与乘法算子 \(A_2 f(t) = tf(t)\) 满足 \textbf{\(A_1 = U A_2 U^{-1}\)}！由此一行秒证全直线上的微分算子 \(\frac{1}{\mathrm{i}}\frac{\mathrm{d}}{\mathrm{d}t}\) 是自伴算子！
\end{itemize}

\SoftDivider

\BookSubSection{sec:7-1-4}{7.1.4 算子的图像、闭算子与可闭算子（p. 279--280）}{7.1.4 算子的图像、闭算子与可闭算子（p. 279–280）}{\TagMust{闭卷必背}}

\SubSubHeading{1. 什么是闭算子？——``退一步海阔天空''的连续性替身}

在 \(H \times H = \{\{a, b\} \mid a, b \in H\}\) 上赋予内积 \((\{a,b\}, \{c,d\}) = (a,c) + (b,d)\) 和范数 \(\|\{a,b\}\|^2 = \|a\|^2 + \|b\|^2\)，它仍是 Hilbert 空间。

\begin{itemize}
\tightlist
\item
  \textbf{算子的图像}：\(G(A) \triangleq \{\{x, Ax\} \mid x \in \mathcal{D}(A)\} \subset H \times H\)。
\item
  \textbf{闭算子的等价定义}：称 \(A\) 是\textbf{闭算子}，如果它的图像 \(G(A)\) 是 \(H \times H\) 中的闭子空间；用序列语言翻译，就是： 
\FitDisplayMath{\text{若 } \{x_n\} \subset \mathcal{D}(A) \text{ 满足 } x_n \to x \text{ 且 } Ax_n \to y \quad (x, y \in H),\quad \text{则必有 } x \in \mathcal{D}(A) \text{ 且 } Ax = y.}

\end{itemize}

\textbf{零基础秒懂：对比``连续算子''与``闭算子''}

\begin{itemize}
\tightlist
\item
  \textbf{连续算子（高配版）}：只要你喂给我 \(x_n \to x\)，我\textbf{打包票} \(Ax_n\) 一定收敛，而且收敛到 \(Ax\)。
\item
  \textbf{闭算子（务实版）}：我承认我是无界算子，你给 \(x_n \to x\) 时，\(Ax_n\) 完全可能发散到无穷大（比如 \(u_n(x)=\frac{1}{n}\sin(n^2 x) \to 0\)，但导数 \(u_n'(x)=n\cos(n^2 x)\) 发散）；\textbf{但}如果 \(x_n \to x\) \textbf{而且} \(Ax_n\) 碰巧也收敛到某个 \(y\)，那我担保两件事：\CircNum{1} \(x\) 绝不会逃出定义域 \(\mathcal{D}(A)\)；\CircNum{2} 极限值绝不撒谎，\(y\) 必定等于 \(Ax\)！
\end{itemize}

\SubSubHeading{2. 旋转算子 \(V\) 与图像正交补（引理 7.6、推论 7.1，p. 279--280）}

在 \(H \times H\) 上定义一个把坐标旋转 \(90^\circ\) 的酉算子： 
\FitDisplayMath{V\{a, b\} = \{-b, a\}, \qquad \text{显然 } V^2 = -I.}


观察 \((Ax, y) = (x, z)\) 等价于 \((-Ax, y) + (x, z) = 0\)，即 \((\{-Ax, x\}, \{y, z\}) = 0\)！也就是 \(\{y, z\} \perp V\{x, Ax\}\)！这一神来之笔直接给出整个 7.1.4 节的底层几何引擎：

\begin{itemize}
\tightlist
\item
  \textbf{引理 7.6（p. 279--280）}：设 \(A\) 是 \(H\) 上的稠定线性算子，则 
\FitDisplayMath{G(A^*) = [V G(A)]^\perp.}

\item
  \textbf{推论 7.1（p. 280）}：如果 \(A\) 是 \(H\) 上的\textbf{稠定闭算子}（从而 \(VG(A)\) 是闭子空间），则有正交直和分解 
\FitDisplayMath{H \times H = V G(A) \oplus G(A^*).}

\end{itemize}

\SubSubHeading{3. 两大核心柱石定理（定理 7.1--7.2，p. 280）}

由引理 7.6 和推论 7.1，几乎不费吹灰之力就能摘下本章最重要的两颗果实（\textbf{\hyperlink{prob:7-6}{\textcolor{LinkColor}{习题 7 第 6 题}}直接考这两条！}）：

\begin{itemize}
\tightlist
\item
  \textbf{定理 7.1（共轭算子恒闭，p. 280）}：设 \(A\) 是 \(H\) 上的稠定线性算子，则 \textbf{\(A^*\) 是闭算子}。特别地，\textbf{自伴算子必是闭算子}。

  \begin{itemize}
  \tightlist
  \item
    \emph{一句话秒记证明}：因为 \(G(A^*) = [VG(A)]^\perp\)，而任何集合的正交补 \(M^\perp\) 永远是闭子空间！
  \end{itemize}
\item
  \textbf{定理 7.2（稠定闭算子的二阶共轭还原，p. 280）}：设 \(A\) 是 \(H\) 上的\textbf{稠定闭算子}，则 \textbf{\(\mathcal{D}(A^*)\) 在 \(H\) 中稠密，且 \(A = A^{**}\)}。

  \begin{itemize}
  \tightlist
  \item
    \emph{直观理解}：在 \(H \times H = VG(A) \oplus G(A^*)\) 两边再乘一次 \(V\)（用 \(V^2 = -I\)），得 \(H \times H = G(A) \oplus VG(A^*)\)，这就迫使 \(G(A) = [VG(A^*)]^\perp = G(A^{**})\)，即 \(A = A^{**}\)！
  \end{itemize}
\end{itemize}

\SubSubHeading{4. 进阶串联：可闭算子、闭包 \(\bar{A} = A^{**}\) 与对称算子的四级阶梯}

如果一个算子 \(A\) 还不是闭算子，我们能不能``把缺的极限点补进去''，把它扩张成一个闭算子？

\begin{itemize}
\tightlist
\item
  \textbf{可闭算子}：若 \(A\) 存在闭扩张，称 \(A\) \textbf{可闭}。等价判别：图像闭包 \(\overline{G(A)}\) 中不含形如 \(\{0, y\}\)（\(y \neq 0\)）的``一个输入对应两个输出''的坏点（即不存在 \(x_n \in \mathcal{D}(A)\) 使 \(x_n \to 0\) 而 \(Ax_n \to y \neq 0\)）。
\item
  \textbf{最小闭扩张（闭包 \(\bar{A}\)）}：当 \(A\) 可闭时，以 \(\overline{G(A)}\) 为图像的算子 \(\bar{A}\) 称为 \(A\) 的闭包（最小闭扩张）。
\item
  \textbf{稠定可闭算子的黄金公式}：稠定算子 \(A\) 可闭 \(\iff \mathcal{D}(A^*)\) 稠密，此时 \textbf{\(\bar{A} = A^{**}\)}（且 \(\bar{A}^* = A^*\)）！
\item
  \textbf{特别地，对任意对称算子 \(A\)}：因 \(A \subset A^*\) 自动保证 \(\mathcal{D}(A) \subset \mathcal{D}(A^*)\) 稠密，故\textbf{所有对称算子都可闭}，且恒满足\textbf{黄金包含链}： 
\FitDisplayMath{A \subset \bar{A} = A^{**} \subset A^*.}
 由此形成清晰的\textbf{四级阶梯}：

  \begin{enumerate}
  \def\labelenumi{\arabic{enumi}.}
  \tightlist
  \item
    \textbf{一般对称算子}：\(A \subsetneq \bar{A} = A^{**} \subsetneq A^*\)；
  \item
    \textbf{闭对称算子（非自伴）}：\(A = A^{**} \subsetneq A^*\)（\hyperlink{prob:7-6}{\textcolor{LinkColor}{习题 7 第 6 题}} (i)）；
  \item
    \textbf{本质自伴算子}（闭包自伴）：\(A \subsetneq \bar{A} = A^{**} = A^*\)；
  \item
    \textbf{自伴算子}：\(A = A^{**} = A^*\)（\hyperlink{prob:7-6}{\textcolor{LinkColor}{习题 7 第 6 题}} (ii)）。
  \end{enumerate}
\end{itemize}

\SoftDivider

\BookSubSection{sec:7-1-5}{7.1.5 对称算子成为自伴算子的判别准则（p. 280--282）}{7.1.5 对称算子成为自伴算子的判别准则（p. 280–282）}{\TagMust{闭卷必背}}

实际做题时，直接去算 \(\mathcal{D}(A^*)\) 常常非常麻烦。本节给出判断一个对称算子到底是不是自伴算子的几把``快刀''：

\SubSubHeading{1. 两个充分条件与 Hellinger-Toeplitz 定理（定理 7.3，p. 280--281）}

设 \(A\) 是 \(H\) 上的\textbf{对称算子}：

\begin{itemize}
\tightlist
\item
  \textbf{(i) Hellinger-Toeplitz 定理}：如果 \textbf{\(\mathcal{D}(A) = H\)}（到处定义），则 \textbf{\(A\) 是自伴的，且 \(A \in \mathcal{L}(H)\)（必为有界算子）}！

  \begin{itemize}
  \tightlist
  \item
    \emph{两行证明}：由 \(H = \mathcal{D}(A) \subset \mathcal{D}(A^*) \subset H\) 得 \(A = A^*\)（自伴）；由定理 7.1 知 \(A\) 是定义在全空间 \(H\) 上的闭算子，再由第 3 章\textbf{闭图像定理}立得 \(A\) 有界！
  \end{itemize}
\item
  \textbf{(ii) 满射判别法}：如果 \textbf{\(\mathcal{R}(A) = H\)}（值域为全空间），则 \textbf{\(A\) 是自伴的，且 \(A^{-1} \in \mathcal{L}(H)\)}！
\end{itemize}

\SubSubHeading{2. 虚数移位 \(A \pm \mathrm{i}I\) 的魔力与自伴基本判别准则（引理 7.7--7.9、定理 7.4，p. 281--282）}

对于任意\textbf{对称算子 \(A\)} 和任意 \(x \in \mathcal{D}(A)\)，把 \(\|(A \pm \mathrm{i}I)x\|^2\) 展开，交叉项因对称性精准抵消（\((Ax, \pm\mathrm{i}x) + (\pm\mathrm{i}x, Ax) = \mp\mathrm{i}(Ax,x) \pm \mathrm{i}(x,Ax) = 0\)），得到贯穿后半章的\textbf{勾股恒等式（式 7.9，p. 281）}： 
\FitDisplayMath{\|(A \pm \mathrm{i}I)x\|^2 = \|Ax\|^2 + \|x\|^2, \qquad \forall x \in \mathcal{D}(A).}


由式 (7.9) 立刻推倒一串多米诺骨牌：

\begin{itemize}
\tightlist
\item
  \textbf{引理 7.7（p. 281）}：若 \(A\) 对称，则 \textbf{\(\operatorname{Ker}(A \pm \mathrm{i}I) = \{0\}\)}（\(A \pm \mathrm{i}I\) 必单射）。
\item
  \textbf{引理 7.8（p. 281）}：若 \(A\) 是\textbf{闭对称算子}，则 \textbf{\(\mathcal{R}(A \pm \mathrm{i}I)\) 是闭子空间}。（\emph{对比\hyperlink{prob:7-7}{\textcolor{LinkColor}{习题 7 第 7 题}}：当对称算子额外非负 \((Ax,x)\ge 0\) 时，实数移位 \(A+I\) 扮演了完全相同的角色！}）
\item
  \textbf{引理 7.9（p. 281）}：对任意稠定算子 \(A\)，恒有正交关系 
\FitDisplayMath{\operatorname{Ker} A^* = (\mathcal{R}(A))^\perp, \qquad \overline{\mathcal{R}(A)} = (\operatorname{Ker} A^*)^\perp.}

\end{itemize}

将三块骨牌拼在一起，就得到了\textbf{自伴算子基本判别准则（定理 7.4，p. 281--282）}：

\begin{quote}
\textbf{定理 7.4}：设 \(A\) 是 Hilbert 空间 \(H\) 上的\textbf{对称算子}，则以下三个命题\textbf{等价}：

\begin{itemize}
\tightlist
\item
  \textbf{(i)} \(A\) 是\textbf{自伴算子}（\(A = A^*\)）；
\item
  \textbf{(ii)} \(A\) 是\textbf{闭算子}，且 \textbf{\(\operatorname{Ker}(A^* + \mathrm{i}I) = \{0\}\)} 与 \textbf{\(\operatorname{Ker}(A^* - \mathrm{i}I) = \{0\}\)}（合记为 \(\operatorname{Ker}(A^* \pm \mathrm{i}I) = \{0\}\)，即 \(A^*\) 没有特征值 \(+\mathrm{i}\) 和 \(-\mathrm{i}\)）；
\item
  \textbf{(iii)} \textbf{\(\mathcal{R}(A + \mathrm{i}I) = H\)} 且 \textbf{\(\mathcal{R}(A - \mathrm{i}I) = H\)}（合记为 \(\mathcal{R}(A \mp \mathrm{i}I) = H\)）。
\end{itemize}
\end{quote}

\textbf{零基础易错点}：(ii) 中\textbf{千万不能漏掉``\(A\) 是闭算子''}！如果 \(A\) 不闭，即使 \(\operatorname{Ker}(A^* \pm \mathrm{i}I) = \{0\}\)，也只能推出 \(\overline{\mathcal{R}(A \pm \mathrm{i}I)} = H\)（\(A\) 本质自伴，即 \(\bar{A}\) 自伴），不能推出 \(A\) 本身自伴！同时，(iii) 中\textbf{正负两个符号 \(+\mathrm{i}\) 和 \(-\mathrm{i}\) 缺一不可}！

\SoftDivider

\BookSubSection{sec:7-2}{7.2 无界自伴算子的谱、Cayley 变换与谱分解（p. 282--293）}{7.2 无界自伴算子的谱、Cayley 变换与谱分解（p. 282–293）}{\TagMust{闭卷必背}}

\SubSubHeading{1. 无界算子的预解集与谱（定义 7.6--7.7，p. 282--283）}

设 \(X\) 是复 Banach 空间，\(A: \mathcal{D}(A) \subset X \to X\) 是线性算子：

\begin{itemize}
\item
  \textbf{定义 7.6（预解集与谱，p. 282）}： 
\FitDisplayMath{\rho(A) = \left\{\lambda \in \mathbb{C} \;\middle|\; \lambda I - A \text{ 是 } \mathcal{D}(A) \to X \text{ 的双射，且 } (\lambda I - A)^{-1} \in \mathcal{L}(X)\right\}.}
 称 \(\rho(A)\) 为 \(A\) 的\textbf{预解集}（其中点称\textbf{正则点}），\(R(\lambda, A) \triangleq (\lambda I - A)^{-1}\) 为\textbf{预解式（预解算子）}，\(\sigma(A) \triangleq \mathbb{C} \setminus \rho(A)\) 为 \(A\) 的\textbf{谱集}。

  \begin{quote}
  \textbf{为什么这里必须手动补写 \((\lambda I - A)^{-1} \in \mathcal{L}(X)\)？} 第 6 章里 \(A\) 有界，Banach 逆算子定理自动保证逆算子有界；现在 \(A\) 可能既无界又不闭！如果 \(A\) 是\textbf{闭算子}，则 \(\lambda I - A\) 是闭算子，当它是双射时，逆算子 \((\lambda I - A)^{-1}\) 是定义在\textbf{全空间 \(X\)} 上的闭算子，由\textbf{闭图像定理}仍自动有界（属于 \(\mathcal{L}(X)\)）！
  \end{quote}
\item
  \textbf{定义 7.7（闭线性算子谱的分类，p. 282--283）}：对闭线性算子 \(A\)，每个 \(\lambda \in \sigma(A)\) 恰好落入以下三种互不相交的情形之一：

  \begin{itemize}
  \tightlist
  \item
    \textbf{(i) 点谱（特征值）\(\sigma_p(A)\)}：\(\lambda I - A\) 不是单射（即存在非零 \(x \in \mathcal{D}(A)\) 使 \(Ax = \lambda x\)）；
  \item
    \textbf{(ii) 连续谱 \(\sigma_c(A)\)}：\(\lambda I - A\) 是单射，\(\mathcal{R}(\lambda I - A) \neq X\) 但 \(\overline{\mathcal{R}(\lambda I - A)} = X\)；
  \item
    \textbf{(iii) 剩余谱 \(\sigma_r(A)\)}：\(\lambda I - A\) 是单射，且 \(\overline{\mathcal{R}(\lambda I - A)} \neq X\)。
  \end{itemize}
\end{itemize}

\SubSubHeading{2. 自伴算子的谱必在实轴上（定理 7.5、推论 7.2，p. 283--284）}

\begin{itemize}
\tightlist
\item
  \textbf{定理 7.5（下有界判别法 / Weyl 准则，p. 283）}：设 \(A\) 是 \(H\) 上的\textbf{自伴算子}，则 
\FitDisplayMath{\lambda \in \rho(A) \iff \text{存在常数 } C > 0 \text{ 使得 } \|(\lambda I - A)x\| \ge C \|x\|,\quad \forall x \in \mathcal{D}(A).}


  \begin{itemize}
  \tightlist
  \item
    \emph{逆否命题（求谱利器！）}：\(\lambda \in \sigma(A) \iff\) 存在单位向量列 \(\{x_n\} \subset \mathcal{D}(A)\)（\(\|x_n\|=1\)）使得 \(\|(\lambda I - A)x_n\| \to 0\)！（这就是证明命题 7.1 \(\sigma(A)=\mathbb{R}\) 的核心武器！）
  \end{itemize}
\item
  \textbf{推论 7.2（谱的实性，p. 283--284）}：设 \(A\) 是 \(H\) 上的自伴算子，则 
\FitDisplayMath{\sigma(A) \subset \mathbb{R}, \qquad \text{且 } \sigma_r(A) = \varnothing.}


  \begin{itemize}
  \tightlist
  \item
    \emph{证明直觉}：把 \(\lambda = \alpha + \mathrm{i}\beta\)（\(\beta \neq 0\)）代入展开，对称性使交叉项再次消去，得 
\FitDisplayMath{\|(\lambda I - A)x\|^2 = \|(\alpha I - A)x\|^2 + \beta^2 \|x\|^2 \ge |\beta|^2 \|x\|^2!}
 由定理 7.5 立得一切非实复数 \(\lambda\) 都属于 \(\rho(A)\)，从而 \(\sigma(A) \subset \mathbb{R}\)（且预解式满足估计 \(\|R(\lambda, A)\| \le \frac{1}{|\operatorname{Im}\lambda|}\)）！
  \end{itemize}
\end{itemize}

\SubSubHeading{3. Cayley 变换：把无界自伴算子``压''成有界酉算子（定理 7.6--7.7、定义 7.9、推论 7.3，p. 285--286）}

想把第 6 章的谱分解定理推广到无界自伴算子 \(A\)，最大障碍就是 \(A\) \textbf{无界}（谱跑到无穷远，做不了有界多项式逼近）。怎么化无界为有界？

复变函数里有一个神奇的映射： 
\FitDisplayMath{u(\lambda) = \frac{\lambda - \mathrm{i}}{\lambda + \mathrm{i}}.}
 当实数 \(\lambda\) 从 \(-\infty\) 跑到 \(+\infty\) 时，\(|u(\lambda)| = \frac{|\lambda - \mathrm{i}|}{|\lambda + \mathrm{i}|} = 1\) 始终落在\textbf{单位圆周}上，只是永远达不到点 \(1\)（\(\lambda \to \pm\infty\) 时趋近于 \(1\)）。

把这个标量公式直接换成算子，就是 \textbf{Cayley 变换（定义 7.9，p. 286）}： 
\FitDisplayMath{U = (A - \mathrm{i}I)(A + \mathrm{i}I)^{-1}.}


\begin{itemize}
\tightlist
\item
  \textbf{定理 7.6（正向 Cayley 变换，p. 285--286）}：设 \(A\) 是对称算子，令 \(U = (A - \mathrm{i}I)(A + \mathrm{i}I)^{-1}\)，则：

  \begin{itemize}
  \tightlist
  \item
    \begin{enumerate}
    \def\labelenumi{(\roman{enumi})}
    \tightlist
    \item
      \(U\) 是 \(\mathcal{R}(A+\mathrm{i}I) \to \mathcal{R}(A-\mathrm{i}I)\) 的\textbf{等距线性算子}（因式 7.9 保证 \(\|(A-\mathrm{i}I)x\| = \|(A+\mathrm{i}I)x\|\)）；
    \end{enumerate}
  \item
    (ii) \textbf{\(1 \notin \sigma_p(U)\)}（\(I - U\) 是单射）；
  \item
    (iii) 可以精确反解回 \(A\)：\textbf{\(A = \mathrm{i}(I+U)(I-U)^{-1}\)}，且 \(\mathcal{D}(A) = \mathcal{R}(I - U)\)；
  \item
    (iv) \(A\) 闭 \(\implies U\) 闭；\textbf{特别地，当 \(A\) 自伴时}，由定理 7.4 知 \(\mathcal{R}(A\pm\mathrm{i}I) = H\)，从而 \textbf{\(U\) 是全空间 \(H\) 上的（有界）酉算子}！
  \end{itemize}
\item
  \textbf{推论 7.3（p. 286）}：\(1 \in \rho(U) \iff A\) 是\textbf{有界}自伴算子。（因此无界自伴算子恰好对应 \(1 \in \sigma_c(U)\)！）
\item
  \textbf{定理 7.7（逆向 Cayley 变换，p. 286）}：反之，若 \(U\) 是等距算子且 \(\mathcal{R}(I-U)\) 稠密，则 \(1 \notin \sigma_p(U)\)，且 \(A = \mathrm{i}(I+U)(I-U)^{-1}\) 是定义在 \(\mathcal{R}(I-U)\) 上的对称算子；\textbf{特别地，若 \(U\) 是酉算子，则 \(A\) 是自伴算子}！
\end{itemize}

\SubSubHeading{4. 无界自伴算子的谱分解定理与乘法算子范例（定义 7.10、定理 7.8、命题 7.1--7.2，p. 290--293）}

有了 Cayley 变换，无界自伴算子的谱分解就水到渠成：

\begin{itemize}
\tightlist
\item
  \textbf{路线图}：无界自伴算子 \(A\) \(\xrightarrow{\text{Cayley变换}}\) 有界酉算子 \(U = \int_0^{2\pi} \mathrm{e}^{\mathrm{i}\theta}\,\mathrm{d}F_\theta\)（第 6 章推论 6.11，因 \(1 \notin \sigma_p(U)\) 故 \(F_0 = 0, \lim_{\theta\to 2\pi^-} F_\theta = I\)）\(\xrightarrow{\text{换元 }\lambda = -\cot(\theta/2),\ E_\lambda = F_\theta}\) 还原回实轴上的谱分解！
\item
  \textbf{定理 7.8（自伴算子谱分解定理，p. 290--292）}：设 \(A\) 是 \(H\) 上的自伴算子，则必有 \(H\) 上唯一的实轴谱系 \(\{E_\lambda \mid \lambda \in \mathbb{R}\}\)（单调递增、右连续、强极限 \(E_{-\infty}=0, E_{+\infty}=I\) 的正交投影族）使得 
\FitDisplayMath{A = \int_{\mathbb{R}} \lambda\,\mathrm{d}E_\lambda,}
 其定义域精确刻画为（按定义 7.10 无界函数谱积分）： 
\FitDisplayMath{\mathcal{D}(A) = \left\{x \in H \;\middle|\; \int_{\mathbb{R}} \lambda^2\,\mathrm{d}\|E_\lambda x\|^2 < +\infty\right\},}
 且对 \(x \in \mathcal{D}(A), y \in H\) 有 \(\|Ax\|^2 = \int_{\mathbb{R}} \lambda^2\,\mathrm{d}\|E_\lambda x\|^2\)，\((Ax, y) = \int_{\mathbb{R}} \lambda\,\mathrm{d}(E_\lambda x, y)\)。
\item
  \textbf{乘法算子的完整画像（命题 7.1--7.2，p. 292--293）}：对于例 7.1（即\hyperlink{prob:7-4}{\textcolor{LinkColor}{习题 7 第 4 题}}）中的 \(L^2(\mathbb{R})\) 位置算子 \((Ax)(t) = tx(t)\)：

  \begin{itemize}
  \tightlist
  \item
    \textbf{命题 7.1}：\(\sigma_p(A) = \sigma_r(A) = \varnothing\)，\(\sigma_c(A) = \sigma(A) = \mathbb{R}\)。（一个特征值都没有，整条实轴都是连续谱！）
  \item
    \textbf{命题 7.2}：它的谱系就是把函数在 \((-\infty, \lambda]\) 上截断的乘法投影： 
\FitDisplayMath{(E_\lambda f)(t) = \chi_{(-\infty, \lambda]}(t) f(t), \qquad A = \int_{\mathbb{R}} \lambda\,\mathrm{d}E_\lambda.}

  \end{itemize}
\end{itemize}

\SoftDivider

\BookSubSection{}{第 7 章考前速记总表}{第 7 章考前速记总表}{}

{\def\LTcaptype{none} % do not increment counter
{\small\renewcommand{\arraystretch}{1.38}\setlength{\tabcolsep}{5pt}
\begin{longtable}[]{@{}
  >{\raggedright\arraybackslash}p{(\linewidth - 6\tabcolsep) * \real{0.2500}}
  >{\raggedright\arraybackslash}p{(\linewidth - 6\tabcolsep) * \real{0.2500}}
  >{\raggedright\arraybackslash}p{(\linewidth - 6\tabcolsep) * \real{0.2500}}
  >{\raggedright\arraybackslash}p{(\linewidth - 6\tabcolsep) * \real{0.2500}}@{}}
\toprule\noalign{}
\rowcolor{AccentSoft}
\begin{minipage}[b]{\linewidth}\raggedright
考点模块
\end{minipage} & \begin{minipage}[b]{\linewidth}\raggedright
必背公式 / 核心结论
\end{minipage} & \begin{minipage}[b]{\linewidth}\raggedright
教材编号
\end{minipage} & \begin{minipage}[b]{\linewidth}\raggedright
考试最高频踩坑点
\end{minipage} \\
\midrule\noalign{}
\endhead
\bottomrule\noalign{}
\endlastfoot
\textbf{稠定与共轭} & \((Ax, y) = (x, A^*y)\)；\(A^*\) 唯一 \(\iff \overline{\mathcal{D}(A)} = H\) & 定义 7.2--7.3、引理 7.1 & 没证 \(\mathcal{D}(A^*)\) 稠密就直接写 \(A^{**}\) \\
\textbf{扩张反号律} & \(A \subset B \implies B^* \subset A^*\) & 引理 7.2 (iii) & 把包含方向写正了（记住：约束越多，满足者越少！） \\
\textbf{对称 vs 自伴} & 对称：\(A \subset A^*\)；自伴：\(A = A^*\)；\(A\) 自伴、\(B\) 对称、\(A \subset B \implies A = B\) & 定义 7.4、引理 7.3--7.4 & 只验 \((Ax,y)=(x,Ay)\) 就敢宣称 \(A\) 自伴，忘了核对定义域！ \\
\textbf{闭算子与 \(A^{**}\)} & \(G(A^*) = [VG(A)]^\perp\)；\(A^*\) 恒闭；\(A\) 稠定闭 \(\implies \mathcal{D}(A^*)\) 稠密且 \(A=A^{**}\) & 引理 7.6、定理 7.1--7.2 & 证闭性时只证了 \(Ax=y\) 却漏证 \(x \in \mathcal{D}(A)\) \\
\textbf{Hellinger-Toeplitz} & \(A\) 对称且 \(\mathcal{D}(A) = H \implies A\) 自伴且 \(A \in \mathcal{L}(H)\) & 定理 7.3 (i) & 误以为无界自伴算子可以定义在整个 \(H\) 上 \\
\textbf{自伴判别准则} & 对称算子 \(A\) 自伴 \(\iff A\) 闭且 \(\operatorname{Ker}(A^*\pm\mathrm{i}I)=\{0\} \iff \mathcal{R}(A\mp\mathrm{i}I)=H\) & 定理 7.4 & 漏掉``闭算子''条件，或只验了 \(+\mathrm{i}\) 漏了 \(-\mathrm{i}\) \\
\textbf{谱与下有界性} & \(A\) 自伴：\(\lambda \in \rho(A) \iff \exists C>0,\ \lVert (\lambda I-A)x\rVert \ge C\lVert x\rVert\)；\(\sigma(A)\subset\mathbb{R},\ \sigma_r(A)=\varnothing\) & 定义 7.6--7.7、定理 7.5、推论 7.2 & 无界自伴算子的谱可以\textbf{无界}（如 \(\sigma(A)=\mathbb{R}\)） \\
\textbf{Cayley 变换} & \(U = (A-\mathrm{i}I)(A+\mathrm{i}I)^{-1}\)，\(A = \mathrm{i}(I+U)(I-U)^{-1}\)；\(A\) 自伴 \(\iff U\) 酉且 \(1 \notin \sigma_p(U)\) & 定理 7.6--7.7、定义 7.9 & 正负号记混（记住分子是 \(-\mathrm{i}\)、分母是 \(+\mathrm{i}\)） \\
\textbf{谱分解定理} & \(A = \int_{\mathbb{R}} \lambda\,\mathrm{d}E_\lambda\)，\(\mathcal{D}(A) = \{x \colon \int_{\mathbb{R}} \lambda^2\,\mathrm{d}\lVert E_\lambda x\rVert^2 < \infty\}\) & 定理 7.8、命题 7.1--7.2 & 乘法算子谱系 \(E_\lambda f = \chi_{(-\infty,\lambda]} f\) \\
\end{longtable}
}
}

\BookPart{}{第二部分：第 7 章课后作业逐题详解（习题 7 第 3、4、6、7 题）}{第二部分：第 7 章课后作业逐题详解（习题 7 第 3、4、6、7 题）}

\begin{quote}
\textbf{全题前置约定}：按教材习题 7（p. 293）题首说明，以下各题如无特别说明，均假设 \(H\) 是复 Hilbert 空间。
\end{quote}

\SoftDivider

\BookProblem{prob:7-3}{习题 7 第 3 题}{教材 p. 293}{习题 7 第 3 题}

\SubPillPrimary{【原题】}

定义算子 \(A : C^2([-1, 1]) \to C([-1, 1])\) 如下： 
\FitDisplayMath{\mathcal{L}u(x) = -\frac{\mathrm{d}}{\mathrm{d}x}\left[(1 - x^2)\frac{\mathrm{d}u}{\mathrm{d}x}\right], \quad x \in [-1, 1].}
 求出 \(A\) 的特征值和特征向量。

\begin{quote}
\textbf{原题记号说明}：原题首句将算子记为 \(A\)，表达式中又沿用微分算子习惯记作 \(\mathcal{L}u(x)\)，二者指同一算子（即著名的 \textbf{Legendre 微分算子}）： 
\FitDisplayMath{(Au)(x) = \mathcal{L}u(x) = -\frac{\mathrm{d}}{\mathrm{d}x}\left[(1 - x^2)u'(x)\right] = -(1 - x^2)u''(x) + 2x u'(x), \qquad x \in [-1, 1].}

\end{quote}

\SubPillAmber{【零基础破题思路】}

这道题要找所有的标量 \(\lambda\) 和\textbf{非零}函数 \(u \in C^2([-1,1])\)，使得 \(Au = \lambda u\)，也就是解二阶变系数微分方程 
\FitDisplayMath{-(1-x^2)u''(x) + 2x u'(x) = \lambda u(x), \qquad x \in [-1,1].}


很多同学一看到变系数微分方程就懵了，其实这道题用\textbf{两步极其干净的泛函 + 线代组合拳}，连一丁点儿高深奇点理论都不需要，就能 100\% 严密拿下：

\begin{enumerate}
\def\labelenumi{\arabic{enumi}.}
\tightlist
\item
  \textbf{第一步（先找出``一批''特征值和特征函数）}：试试把最简单的单项式 \(x^k\) 喂给 \(A\)！算一下立刻惊喜地发现： 
\FitDisplayMath{A(x^k) = k(k+1)x^k - k(k-1)x^{k-2}.}
 次数不但没有变高，还只往低次掉！这说明对每个 \(n \ge 0\)，\(n\) 次及以下多项式空间 \(\mathcal{P}_n = \operatorname{span}\{1, x, \dots, x^n\}\) 是 \(A\) 的\textbf{有限维不变子空间}，且在基 \(\{1, x, \dots, x^n\}\) 下矩阵是\textbf{上三角矩阵}，对角线上明晃晃地写着 \(0, 2, 6, \dots, n(n+1)\)！ 于是每个 \(\lambda_n = n(n+1)\)（\(n = 0, 1, 2, \dots\)）铁定是特征值，而且对应一个唯一的 \(n\) 次多项式特征函数 \(P_n(x)\)（就是大名鼎鼎的 \textbf{Legendre 多项式}，既可以用递推系数写出来，也可以用 Rodrigues 公式写出来并 3 行验完！）。
\item
  \textbf{第二步（证明``全找齐了''——再无其他特征值，且每个特征空间只有 1 维）}：万一还有某个非多项式函数 \(u \in C^2([-1,1])\) 也是特征函数怎么办？ 这时\textbf{泛函分析的主角登场}：把 \(C^2([-1,1])\) 放进 Hilbert 空间 \(L^2(-1,1)\) 里！

  \begin{itemize}
  \tightlist
  \item
    做两次分部积分，发现由于权函数 \((1-x^2)\) 在端点 \(x = \pm 1\) \textbf{自己变成 \(0\)}，边界项自动消失，所以 \(A\) 在 \(L^2(-1,1)\) 上竟然\textbf{天然对称}：\((Au, v) = (u, Av)\)！
  \item
    于是任何满足 \(Au = \lambda u\) 的函数 \(u\) 都自动满足 \([\lambda - n(n+1)](u, P_n) = 0\)（\(\forall n \ge 0\)）。
  \item
    如果 \(\lambda \neq n(n+1)\)，那 \(u\) 就与\textbf{所有多项式 \(P_n\) 都正交}！而由 \textbf{Weierstrass 多项式逼近定理}，全体多项式在 \(C([-1,1])\)（从而在 \(L^2(-1,1)\)）里稠密——与稠密集正交的连续函数只能恒等于 \(0\)！
  \item
    同理，如果 \(\lambda = n_0(n_0+1)\)，把 \(u\) 减去它在 \(P_{n_0}\) 方向上的投影 \(c P_{n_0}\)，剩下的余项 \(w = u - c P_{n_0}\) 与所有 \(P_n\) 正交，立刻逼出 \(w \equiv 0\)，即 \(u = c P_{n_0}\)！
  \end{itemize}
\end{enumerate}

\SubPillMain{【规范解答】}

\textbf{结论}：算子 \(A\) 的全部特征值为 
\FitDisplayMath{\lambda_n = n(n + 1), \qquad n = 0, 1, 2, \dots,}
 对应于特征值 \(\lambda_n = n(n+1)\) 的全部特征向量（特征函数）为 
\FitDisplayMath{u(x) = c P_n(x), \qquad c \neq 0,}
 其中 \(P_n(x)\) 为 \(n\) 次 \textbf{Legendre 多项式}，其显式表达式（Rodrigues 公式）为 
\FitDisplayMath{P_n(x) = \frac{1}{2^n n!} \frac{\mathrm{d}^n}{\mathrm{d}x^n}\left[(x^2 - 1)^n\right], \qquad n = 0, 1, 2, \dots}
 （前四个为：\(P_0(x) = 1\)，\(P_1(x) = x\)，\(P_2(x) = \dfrac{1}{2}(3x^2 - 1)\)，\(P_3(x) = \dfrac{1}{2}(5x^3 - 3x)\)。）

下面分两步严格证明上述结论（空间默认在复数域 \(\mathbb{C}\) 上考虑；若限定实数域 \(\mathbb{R}\)，结论及证明完全相同）。

\SoftDivider

\textbf{第一步：证明每个 \(\lambda_n = n(n+1)\)（\(n = 0, 1, 2, \dots\)）都是 \(A\) 的特征值，且 \(P_n(x)\) 是对应的特征向量。}

先计算算子 \(A\) 作用在单项式 \(x^k \in C^2([-1,1])\)（\(k = 0, 1, 2, \dots\)）上的结果：

\begin{itemize}
\tightlist
\item
  当 \(k = 0\) 时，\(A(1) = 0 = 0 \cdot (0+1) \cdot 1\)；
\item
  当 \(k = 1\) 时，\(A(x) = -\dfrac{\mathrm{d}}{\mathrm{d}x}(1 - x^2) = 2x = 1 \cdot (1+1) x\)；
\item
  当 \(k \ge 2\) 时， 
\FitDisplayMath{A(x^k) = -\frac{\mathrm{d}}{\mathrm{d}x}\left[(1 - x^2) k x^{k-1}\right] = -\frac{\mathrm{d}}{\mathrm{d}x}\left(k x^{k-1} - k x^{k+1}\right) = k(k + 1)x^k - k(k - 1)x^{k-2}.}

\end{itemize}

对任意非负整数 \(n\)，记次数不超过 \(n\) 的多项式子空间为 \(\mathcal{P}_n = \operatorname{span}\{1, x, \dots, x^n\} \subset C^2([-1,1])\)。由上式可知 \(A(\mathcal{P}_n) \subset \mathcal{P}_n\)，即 \(\mathcal{P}_n\) 是 \(A\) 的 \((n+1)\) 维不变子空间。在基 \(\{1, x, x^2, \dots, x^n\}\) 下，限制算子 \(A|_{\mathcal{P}_n}\) 的表示矩阵为\textbf{上三角矩阵}（主对角线元素依次为 \(0, 2, 6, \dots, n(n+1)\)）： 
\FitDisplayMath{M_n = \begin{pmatrix} 0 & 0 & -2 & 0 & \cdots & \ast \\ 0 & 2 & 0 & -6 & \cdots & \ast \\ 0 & 0 & 6 & 0 & \cdots & \ast \\ 0 & 0 & 0 & 12 & \cdots & \ast \\ \vdots & \vdots & \vdots & \vdots & \ddots & \vdots \\ 0 & 0 & 0 & 0 & \cdots & n(n+1) \end{pmatrix}.}


由于主对角线上的 \((n+1)\) 个数 \(k(k+1)\)（\(k = 0, 1, \dots, n\)）两两互异，故对每个 \(n \ge 0\)，\(\lambda_n = n(n+1)\) 都是 \(A|_{\mathcal{P}_n}\)（从而也是 \(A\)）的特征值，且在 \(\mathcal{P}_n\) 中存在唯一（相差非零常数倍）的 \(n\) 次多项式特征向量。

下面直接验证 Rodrigues 公式给出的 \(n\) 次多项式 \(P_n(x) = \dfrac{1}{2^n n!} \dfrac{\mathrm{d}^n}{\mathrm{d}x^n}\left[(x^2-1)^n\right] \in C^2([-1,1])\) 确为对应于 \(\lambda_n = n(n+1)\) 的非零特征向量：

\begin{enumerate}
\def\labelenumi{\arabic{enumi}.}
\tightlist
\item
  \textbf{非零性与次数}：\((x^2 - 1)^n = x^{2n} + \text{低次项}\)，求 \(n\) 阶导数后首项为 \(\dfrac{(2n)!}{n!} x^n\)，故 \(P_n(x)\) 是精确的 \(n\) 次多项式（特别地 \(P_n \not\equiv 0\)；按 Leibniz 公式在 \(x=1\) 处计算还可得归一化 \(P_n(1)=1\)）。
\item
  \textbf{满足特征方程}：记 \(f(x) = (x^2 - 1)^n\)，直接求导得 \(f'(x) = 2nx(x^2-1)^{n-1}\)，两边同乘 \((x^2 - 1)\) 得恒等式 
\FitDisplayMath{(x^2 - 1)f'(x) = 2n x f(x).}
 对上式两边应用 Leibniz 乘积求导公式求 \((n+1)\) 阶导数：

  \begin{itemize}
  \tightlist
  \item
    左边： 
\FitDisplayMath{(x^2 - 1)f^{(n+2)}(x) + \binom{n+1}{1}(2x)f^{(n+1)}(x) + \binom{n+1}{2}(2)f^{(n)}(x) = (x^2 - 1)f^{(n+2)}(x) + 2(n+1)x f^{(n+1)}(x) + n(n+1)f^{(n)}(x);}

  \item
    右边： 
\FitDisplayMath{2n\left[x f^{(n+1)}(x) + \binom{n+1}{1} \cdot 1 \cdot f^{(n)}(x)\right] = 2nx f^{(n+1)}(x) + 2n(n+1)f^{(n)}(x).}
 两边相等并移项整理，得 
\FitDisplayMath{(x^2 - 1)f^{(n+2)}(x) + 2x f^{(n+1)}(x) - n(n+1)f^{(n)}(x) = 0.}
 两边同乘 \(\dfrac{1}{2^n n!}\) 并注意到 \(P_n(x) = \dfrac{1}{2^n n!} f^{(n)}(x)\)，立得 
\FitDisplayMath{-(1 - x^2)P_n''(x) + 2x P_n'(x) = n(n+1)P_n(x), \quad \text{即 } A P_n = n(n+1)P_n.}
 （注：亦可设 \(P_n(x) = \sum_{k=0}^n a_k x^k\) 代入 \(AP_n = n(n+1)P_n\) 比较系数，得到 \(a_{n-1}=0\) 及递推公式 \(a_k = -\dfrac{(k+2)(k+1)}{(n-k)(n+k+1)} a_{k+2}\) 唯一确定各项系数。）
  \end{itemize}
\end{enumerate}

\SoftDivider

\textbf{第二步：证明 \(A\) 再无其他特征值，且每个 \(\lambda_n\) 对应的特征空间恰为一维空间 \(\operatorname{span}\{P_n\}\)。}

将 \(C^2([-1,1])\) 与 \(C([-1,1])\) 嵌入复 Hilbert 空间 \(L^2(-1, 1)\)，内积与范数定义为 
\FitDisplayMath{(u, v) = \int_{-1}^{1} u(x)\overline{v(x)}\,\mathrm{d}x, \qquad \|u\|_{L^2} = \left(\int_{-1}^{1} |u(x)|^2\,\mathrm{d}x\right)^{1/2}.}


\begin{enumerate}
\def\labelenumi{\arabic{enumi}.}
\item
  \textbf{验证 \(A\) 的对称性及 \(\{P_n\}\) 的正交性}： 对任意 \(u, v \in C^2([-1,1])\)，由于 \((1 - x^2)\big|_{x=-1}^{x=1} = 0\)，连续分部积分两次得边界项全部为零： 
\FitDisplayMath{\begin{aligned}
  (Au, v) &= -\int_{-1}^{1} \left[(1 - x^2)u'(x)\right]' \overline{v(x)}\,\mathrm{d}x \\
  &= -\Big[(1 - x^2)u'(x)\overline{v(x)}\Big]_{-1}^{1} + \int_{-1}^{1} (1 - x^2)u'(x)\overline{v'(x)}\,\mathrm{d}x \\
  &= \int_{-1}^{1} (1 - x^2)u'(x)\overline{v'(x)}\,\mathrm{d}x \\
  &= \Big[(1 - x^2)u(x)\overline{v'(x)}\Big]_{-1}^{1} - \int_{-1}^{1} u(x)\overline{\left[(1 - x^2)v'(x)\right]'}\,\mathrm{d}x = (u, Av).
  \end{aligned}}
 特别地，对任意 \(m \neq n\)，由 \(m(m+1)(P_m, P_n) = (AP_m, P_n) = (P_m, AP_n) = n(n+1)(P_m, P_n)\) 及 \(m(m+1) \neq n(n+1)\)，得 
\FitDisplayMath{(P_m, P_n) = 0 \qquad (m \neq n),}
 且因为 \(P_n\) 连续非零，有 \(\|P_n\|_{L^2} > 0\)。
\item
  \textbf{证明多项式系 \(\{P_n\}_{n=0}^\infty\) 的完备性（唯一性引理）}： 设连续函数 \(w \in C([-1,1])\) 满足 \((w, P_n) = 0\) 对一切 \(n = 0, 1, 2, \dots\) 成立。 因为每个 \(P_n\) 恰为 \(n\) 次多项式，线性张成 \(\operatorname{span}\{P_0, P_1, \dots\}\) 恰好是 \([-1,1]\) 上的全体多项式，故对一切多项式 \(p(x)\) 均有 \((w, p) = 0\)。 由 \textbf{Weierstrass 多项式逼近定理}（对复值函数 \(w\)，分别对其连续的实部 \(\operatorname{Re} w\) 与虚部 \(\operatorname{Im} w\) 作实多项式一致逼近再合并），存在多项式列 \(\{p_k\}\) 在 \([-1,1]\) 上一致收敛于 \(w\)。于是 
\FitDisplayMath{\|w - p_k\|_{L^2}^2 = \int_{-1}^{1} |w(x) - p_k(x)|^2\,\mathrm{d}x \le 2 \max_{x \in [-1,1]} |w(x) - p_k(x)|^2 \to 0 \qquad (k \to \infty).}
 由内积的连续性得 
\FitDisplayMath{\|w\|_{L^2}^2 = (w, w) = \lim_{k\to\infty} (w, p_k) = 0.}
 因 \(w \in C([-1,1])\) 连续且 \(\int_{-1}^1 |w(x)|^2\,\mathrm{d}x = 0\)，必有 \(w(x) \equiv 0\)（对一切 \(x \in [-1,1]\)）。
\item
  \textbf{确定全部特征值与特征空间}： 现设 \(\lambda \in \mathbb{C}\) 与 \(u \in C^2([-1,1])\) 满足 \(Au = \lambda u\)。对一切 \(n = 0, 1, 2, \dots\)，利用 \(A\) 的对称性以及 \(n(n+1) \in \mathbb{R}\)，有 
\FitDisplayMath{\lambda (u, P_n) = (\lambda u, P_n) = (Au, P_n) = (u, A P_n) = \left(u, n(n+1)P_n\right) = n(n+1)(u, P_n),}
 整理即得 
\FitDisplayMath{\left[\lambda - n(n+1)\right](u, P_n) = 0, \qquad \forall n = 0, 1, 2, \dots}
 分两种情形讨论：

  \begin{itemize}
  \tightlist
  \item
    \textbf{情形 A}：\(\lambda \notin \{n(n+1) \mid n = 0, 1, 2, \dots\}\)。 此时对一切 \(n \ge 0\) 都有 \(\lambda - n(n+1) \neq 0\)，从而 \((u, P_n) = 0\)（\(\forall n \ge 0\)）。由上述完备性引理立得 \(u(x) \equiv 0\)。故 \(\lambda\) 不是特征值。
  \item
    \textbf{情形 B}：\(\lambda = n_0(n_0+1)\)（其中 \(n_0\) 为某个非负整数）。 此时对一切 \(n \neq n_0\) 都有 \((u, P_n) = 0\)。令 \(c = \dfrac{(u, P_{n_0})}{\|P_{n_0}\|_{L^2}^2}\)，构造 \(w = u - c P_{n_0} \in C^2([-1,1])\)。利用正交性 \((P_{n_0}, P_n) = 0\)（\(n \neq n_0\)）：

    \begin{itemize}
    \tightlist
    \item
      当 \(n \neq n_0\) 时，\((w, P_n) = (u, P_n) - c(P_{n_0}, P_n) = 0 - 0 = 0\)；
    \item
      当 \(n = n_0\) 时，\((w, P_{n_0}) = (u, P_{n_0}) - \dfrac{(u, P_{n_0})}{\|P_{n_0}\|_{L^2}^2}(P_{n_0}, P_{n_0}) = 0\)。 故 \((w, P_n) = 0\) 对一切 \(n \ge 0\) 成立。再由上述完备性引理得 \(w(x) \equiv 0\)，即 \(u(x) = c P_{n_0}(x)\)。当 \(u\) 为特征向量时要求 \(u \not\equiv 0\)，即 \(c \neq 0\)。
    \end{itemize}
  \end{itemize}
\end{enumerate}

综上，\(A\) 的全部特征值为 \(\lambda_n = n(n+1)\)（\(n = 0, 1, 2, \dots\)），对应的全部特征向量为 \(c P_n(x)\)（\(c \neq 0\)）。\(\blacksquare\)

\SubPillCoral{【考点复盘与避坑提示】}

\begin{itemize}
\tightlist
\item
  \textbf{为什么题目没有给任何边界条件（如 \(u(\pm 1)=0\)），\(A\) 却自动对称且选出了离散的特征值？} 因为系数 \((1-x^2)\) 在端点 \(x = \pm 1\) 自动退化为零（数理方程里称为\textbf{奇异 Sturm-Liouville 问题}），只要函数要求落在 \(C^2([-1,1])\)（端点有限且二阶连续可微），边界项就被 \((1-x^2)\) 自动抹平了！
\item
  \textbf{与数理方程中``幂级数解法''的融会贯通}： 学过《数学物理方法》的同学可能记得另一种做法：把 \(u(x) = \sum_{k=0}^\infty a_k x^k\) 代入方程得两步递推 \(a_{k+2} = \frac{k(k+1)-\lambda}{(k+2)(k+1)} a_k\)；当 \(\lambda \neq n(n+1)\) 时，级数永不截断，由比值判别法可证其发散于端点 \(x = \pm 1\)，从而不在 \(C^2([-1,1])\) 中（实际上连 \(C([-1,1])\) 都不在）。但幂级数法在考试中需要细致证明任意 \(C^2([-1,1])\) 解在 \((-1,1)\) 上解析且在端点处发散，很容易漏掉严格性细节；相比之下，上面的\textbf{``多项式上三角 + \(L^2\) 对称性 + Weierstrass 稠密性''}纯用泛函工具，无任何死角，是最稳妥的满分写法！
\item
  \textbf{避坑提醒}：千万别只写完第一步``找到 \(P_n(x)\)''就收工！题干问的是``求出 \(A\) 的特征值和特征向量''，必须证明\textbf{没有遗漏其他特征值}，且\textbf{每个特征空间都没有多项式以外的其他特征函数}（即第二步不可或缺！）。
\end{itemize}

\SoftDivider

\BookProblem{prob:7-4}{习题 7 第 4 题}{教材 p. 293}{习题 7 第 4 题}

\SubPillPrimary{【原题】}

设 \(H = L^2(\mathbb{R})\)，算子 \(A\) 定义如下： 
\FitDisplayMath{\begin{cases}
\mathcal{D}(A) = \left\{x \in H \;\middle|\; \displaystyle\int_{\mathbb{R}} t^2 |x(t)|^2\,\mathrm{d}t < +\infty\right\}, \\
(Ax)(t) = t x(t), \quad \forall x \in \mathcal{D}(A).
\end{cases}}
 证明：\(A\) 是一个闭的无界算子。

\SubPillAmber{【零基础破题思路】}

这就是贯穿全章的明星算子——量子力学里的\textbf{位置算子}（教材例 7.1、例 7.3、命题 7.1--7.2 的主角）。要证明三件事（其实第一件是隐含前提）：

\begin{enumerate}
\def\labelenumi{\arabic{enumi}.}
\tightlist
\item
  \textbf{线性与良定义}（快速补全，显得严谨）：确认 \(\mathcal{D}(A)\) 是线性子空间且 \(Ax \in L^2(\mathbb{R})\)。
\item
  \textbf{无界性}：要推翻``存在常数 \(M\) 使 \(\|Ax\| \le M\|x\|\)''，只需造出一列``自己大小固定为 \(1\)，但被 \(A\) 作用后大小爆到 \(+\infty\)''的测试函数！

  \begin{itemize}
  \tightlist
  \item
    \(A\) 的本质是乘以坐标 \(t\)，哪里的 \(t\) 最大？当然是远处 \([n, n+1]\)！
  \item
    取示性函数 \(x_n = \chi_{[n, n+1]}\)（宽度 \(1\)、高度 \(1\) 的小方块），\(\|x_n\|_2 = 1\)，而乘上 \(t\) 后至少放大 \(n\) 倍（\(\|Ax_n\|_2 \ge n \to +\infty\)），瞬间秒杀！
  \end{itemize}
\item
  \textbf{闭性}：回忆闭算子的定义： 
\FitDisplayMath{\text{若 } x_n \in \mathcal{D}(A),\ x_n \xrightarrow{L^2} x,\ Ax_n \xrightarrow{L^2} y,\quad \text{需证 } x \in \mathcal{D}(A) \text{ 且 } Ax = y.}


  \begin{itemize}
  \tightlist
  \item
    怎么把``积分意义下的 \(L^2\) 收敛''变成``逐点等式 \(y(t) = tx(t)\)''？用实变函数最经典的利器：\textbf{\(L^2\) 收敛必含几乎处处（a.e.）逐点收敛的子列}！
  \item
    \textbf{细节杀手锏}：不要一次糊弄说``取子列同时 a.e. 收敛''，而要\textbf{显式分两次抽子列}——先从 \(\{x_n\}\) 抽子列 \(\{x_{n_k}\}\) 使 \(x_{n_k} \to x\) a.e.；再从对应的 \(\{Ax_{n_k}\}\) 中\textbf{进一步抽子列} \(\{Ax_{n_{k_j}}\}\) 使 \(t x_{n_{k_j}} \to y\) a.e.！同一子列同时满足 \(t x_{n_{k_j}} \to tx\) a.e. 和 \(t x_{n_{k_j}} \to y\) a.e.，立得 \(y(t) = t x(t)\) a.e.！
  \item
    \textbf{临门一脚千万别忘}：利用已知的 \(y \in L^2(\mathbb{R})\)，得 \(\int_{\mathbb{R}} t^2|x(t)|^2\,\mathrm{d}t = \int_{\mathbb{R}} |y(t)|^2\,\mathrm{d}t < +\infty\)，这才真正证完 \textbf{\(x \in \mathcal{D}(A)\)} 且 \(Ax = y\)！
  \end{itemize}
\end{enumerate}

\SubPillMain{【规范解答】}

\textbf{第一步：验证 \(A\) 是良定义的线性算子。}

显然 \(0 \in \mathcal{D}(A)\)。任取 \(x, y \in \mathcal{D}(A)\) 及 \(\alpha, \beta \in \mathbb{C}\)，利用基本不等式 \(|a + b|^2 \le 2(|a|^2 + |b|^2)\)，有 
\FitDisplayMath{\int_{\mathbb{R}} t^2 |\alpha x(t) + \beta y(t)|^2\,\mathrm{d}t \le 2|\alpha|^2 \int_{\mathbb{R}} t^2 |x(t)|^2\,\mathrm{d}t + 2|\beta|^2 \int_{\mathbb{R}} t^2 |y(t)|^2\,\mathrm{d}t < +\infty,}
 故 \(\alpha x + \beta y \in \mathcal{D}(A)\)，即 \(\mathcal{D}(A)\) 是 \(H = L^2(\mathbb{R})\) 的线性子空间（事实上 \(\mathcal{D}(A)\) 包含全体紧支集连续函数 \(C_c(\mathbb{R})\)，故还是稠密子空间）。又对任意 \(x \in \mathcal{D}(A)\)， 
\FitDisplayMath{\|Ax\|^2 = \int_{\mathbb{R}} |t x(t)|^2\,\mathrm{d}t = \int_{\mathbb{R}} t^2 |x(t)|^2\,\mathrm{d}t < +\infty,}
 故 \(Ax \in L^2(\mathbb{R})\)，且显然 \(A(\alpha x + \beta y) = \alpha Ax + \beta Ay\)，因此 \(A: \mathcal{D}(A) \to H\) 是良定义的线性算子。

\SoftDivider

\textbf{第二步：证明 \(A\) 是无界算子。}

对每个正整数 \(n = 1, 2, \dots\)，构造区间 \([n, n+1]\) 上的示性函数： 
\FitDisplayMath{x_n(t) = \chi_{[n, n+1]}(t) = \begin{cases} 1, & t \in [n, n+1], \\ 0, & t \notin [n, n+1]. \end{cases}}
 直接计算它们的 \(L^2\) 范数： 
\FitDisplayMath{\|x_n\|^2 = \int_{\mathbb{R}} |x_n(t)|^2\,\mathrm{d}t = \int_n^{n+1} 1\,\mathrm{d}t = 1 \implies \|x_n\| = 1,}
 
\FitDisplayMath{\int_{\mathbb{R}} t^2 |x_n(t)|^2\,\mathrm{d}t = \int_n^{n+1} t^2\,\mathrm{d}t = \frac{(n+1)^3 - n^3}{3} = n^2 + n + \frac{1}{3} < +\infty.}
 由此可知每个 \(x_n \in \mathcal{D}(A)\) 且 \(\|x_n\| = 1\)，但 
\FitDisplayMath{\|Ax_n\| = \sqrt{n^2 + n + \frac{1}{3}} \ge n \to +\infty \qquad (n \to \infty).}
 因此不存在常数 \(M > 0\) 使得 \(\|Ax\| \le M\|x\|\) 对一切 \(x \in \mathcal{D}(A)\) 成立，故 \(A\) 是\textbf{无界算子}。

\SoftDivider

\textbf{第三步：证明 \(A\) 是闭算子。}

任取序列 \(\{x_n\}_{n=1}^\infty \subset \mathcal{D}(A)\) 以及元素 \(x, y \in L^2(\mathbb{R})\)，使得在 \(L^2(\mathbb{R})\) 范数意义下 
\FitDisplayMath{\lim_{n\to\infty} \|x_n - x\| = 0, \qquad \lim_{n\to\infty} \|Ax_n - y\| = 0.}
 需要证明：\(x \in \mathcal{D}(A)\) 且 \(Ax = y\)。

由 \(L^2(\mathbb{R})\) 空间的基本性质（Riesz 定理：\(L^2\) 强收敛序列必存在几乎处处逐点收敛的子列），分两步抽取子列：

\begin{enumerate}
\def\labelenumi{\arabic{enumi}.}
\tightlist
\item
  因为 \(x_n \xrightarrow{L^2} x\)，存在子列 \(\{x_{n_k}\}_{k=1}^\infty\) 使得 
\FitDisplayMath{\lim_{k\to\infty} x_{n_k}(t) = x(t), \qquad \text{a.e. } t \in \mathbb{R}.}

\item
  因为原序列满足 \(Ax_n \xrightarrow{L^2} y\)，其子列 \(\{Ax_{n_k}\}_{k=1}^\infty\) 依然在 \(L^2(\mathbb{R})\) 中强收敛于 \(y\)，故可从 \(\{n_k\}\) 中\textbf{进一步抽取子列} \(\{n_{k_j}\}_{j=1}^\infty\) 使得 
\FitDisplayMath{\lim_{j\to\infty} (Ax_{n_{k_j}})(t) = \lim_{j\to\infty} t x_{n_{k_j}}(t) = y(t), \qquad \text{a.e. } t \in \mathbb{R}.}

\end{enumerate}

另一方面，作为 \(\{x_{n_k}\}\) 的子列，\(\{x_{n_{k_j}}\}\) 仍满足 \(\lim_{j\to\infty} x_{n_{k_j}}(t) = x(t)\)（a.e. \(t \in \mathbb{R}\)）。对每个使该式成立的有限实数 \(t \in \mathbb{R}\)，两端同乘 \(t\) 即得 
\FitDisplayMath{\lim_{j\to\infty} t x_{n_{k_j}}(t) = t x(t), \qquad \text{a.e. } t \in \mathbb{R}.}


两个零测集之并仍为零测集。由逐点极限的唯一性，得 
\FitDisplayMath{y(t) = t x(t), \qquad \text{a.e. } t \in \mathbb{R}.}


最后验证 \(x \in \mathcal{D}(A)\) 且 \(Ax = y\)：由于已知 \(y \in L^2(\mathbb{R})\)，代入上式有 
\FitDisplayMath{\int_{\mathbb{R}} t^2 |x(t)|^2\,\mathrm{d}t = \int_{\mathbb{R}} |t x(t)|^2\,\mathrm{d}t = \int_{\mathbb{R}} |y(t)|^2\,\mathrm{d}t = \|y\|^2 < +\infty.}
 这就证明了 \(x \in \mathcal{D}(A)\)，进而 \((Ax)(t) = t x(t) = y(t)\)（a.e. \(t \in \mathbb{R}\)），即在 \(L^2(\mathbb{R})\) 中 \(Ax = y\)。

综上所述，\(A\) 是 \(L^2(\mathbb{R})\) 上的一个闭的无界算子。\(\blacksquare\)

\SubPillCoral{【考点复盘与避坑提示】}

\begin{itemize}
\tightlist
\item
  \textbf{扣分高发区 1（抽子列不严谨）}：切忌写``\(x_n \to x\) a.e. 且 \(Ax_n \to y\) a.e.''——\(L^2\) 收敛本身\textbf{不能}保证原序列几乎处处收敛！也不能各自独立抽两个不同的子列 \(\{x_{n_k}\}\) 和 \(\{Ax_{m_k}\}\)！必须\textbf{嵌套两次抽子列}（从第一个子列里再抽子列），才能保证同一个下标序列 \(\{n_{k_j}\}\) 上两条几乎处处极限同时成立。
\item
  \textbf{扣分高发区 2（漏证 \(x \in \mathcal{D}(A)\)）}：证明闭算子有两个硬指标：\CircNum{1} \textbf{\(x \in \mathcal{D}(A)\)}；\CircNum{2} \textbf{\(Ax = y\)}。本题中恰恰因为 \(y \in L^2(\mathbb{R})\) 且 \(t x(t) = y(t)\) a.e.，才把 \(\int t^2|x(t)|^2\mathrm{d}t < +\infty\) 给锁死，千万不能跳过这一行！
\item
  \textbf{另解拓展（与教材例 7.3 + 定理 7.1 联动）}：也可以先证明 \(A\) 稠定且自伴（\(A = A^*\)，见教材例 7.3，p. 277），再由\textbf{定理 7.1}（自伴算子必是闭算子）直接推出 \(A\) 是闭算子；但直接用序列定义最简短自包含，考场首选直接法！
\end{itemize}

\SoftDivider

\BookProblem{prob:7-6}{习题 7 第 6 题}{教材 p. 293}{习题 7 第 6 题}

\SubPillPrimary{【原题】}

设 \(A\) 是 \(H\) 上的对称算子。证明：

\begin{itemize}
\tightlist
\item
  \begin{enumerate}
  \def\labelenumi{(\roman{enumi})}
  \tightlist
  \item
    \(A\) 是闭算子当且仅当 \(A = A^{**} \subset A^*\)；
  \end{enumerate}
\item
  (ii) \(A\) 是自伴算子当且仅当 \(A = A^{**} = A^*\)。
\end{itemize}

\SubPillAmber{【零基础破题思路】}

这道题篇幅很短，却是检验你是否真正搞懂第 7 章概念逻辑链的``照妖镜''！

看到题目里出现了二阶共轭 \(A^{**} = (A^*)^*\)，零基础最容易直接埋头套定理，却漏掉一个最要命的\textbf{前提检查}：

\begin{itemize}
\tightlist
\item
  \textbf{前置暗坑}：凭什么 \(A^{**}\) 一定存在？按定义 7.3，只有\textbf{稠定算子}才有共轭算子！所以必须先花两行证明：\(A\) 对称 \(\implies A\) 稠定且 \(A \subset A^* \implies \mathcal{D}(A) \subset \mathcal{D}(A^*) \implies \mathcal{D}(A^*)\) 也在 \(H\) 中稠密 \(\implies A^{**} = (A^*)^*\) 良定义存在！
\end{itemize}

迈过这个坑之后，(i) 和 (ii) 就像搭乐高一样，精准对接教材 \textbf{定理 7.1}（稠定算子的共轭必闭）和 \textbf{定理 7.2}（稠定闭算子满足 \(A = A^{**}\)）：

\begin{itemize}
\tightlist
\item
  \textbf{(i)}

  \begin{itemize}
  \tightlist
  \item
    \(\Rightarrow\)：\(A\) 是稠定闭算子 \(\xrightarrow{\text{定理 7.2}} A = A^{**}\)；又 \(A\) 对称 \(\implies A \subset A^*\)；拼起来就是 \(A = A^{**} \subset A^*\)！
  \item
    \(\Leftarrow\)：已知 \(A = A^{**}\)；由于 \(A^*\) 稠定，把\textbf{定理 7.1} 用在 \(A^*\) 身上，立刻得出 \((A^*)^* = A^{**}\) 是闭算子，所以 \(A\) 就是闭算子！
  \end{itemize}
\item
  \textbf{(ii)}

  \begin{itemize}
  \tightlist
  \item
    \(\Rightarrow\)：\(A\) 自伴即 \(A = A^*\) \(\xrightarrow{\text{定理 7.1}} A\) 是闭算子 \(\xrightarrow{\text{(i)}} A = A^{**}\)，合并即 \(A = A^{**} = A^*\)！
  \item
    \(\Leftarrow\)：题设直接送了你 \(A = A^*\)，按定义 7.4 这就是自伴算子！
  \end{itemize}
\end{itemize}

\SubPillMain{【规范解答】}

\textbf{前置验证（证明 \(A^{**}\) 良定义存在）：}

因为 \(A\) 是 \(H\) 上的对称算子，由定义 7.4（p. 275）知，\(A\) 是 \(H\) 上的稠定线性算子（即 \(\overline{\mathcal{D}(A)} = H\)），且满足 \(A \subset A^*\)。由算子扩张的定义（定义 7.1 (iv)）可知 \(\mathcal{D}(A) \subset \mathcal{D}(A^*)\)，于是 
\FitDisplayMath{H = \overline{\mathcal{D}(A)} \subset \overline{\mathcal{D}(A^*)} \subset H,}
 故 \(\overline{\mathcal{D}(A^*)} = H\)。结合引理 7.2 (i) 知 \(A^*\) 是 \(H\) 上的稠定线性算子。因此，按定义 7.3，\(A^*\) 的共轭算子 \(A^{**} \triangleq (A^*)^*\) 良定义且唯一存在。

\SoftDivider

\textbf{证明 (i)：\(A\) 是闭算子 \(\iff A = A^{**} \subset A^*\)。}

\begin{itemize}
\tightlist
\item
  \textbf{必要性（\(\Rightarrow\)）}： 设 \(A\) 是闭算子。因为 \(A\) 是对称算子，所以 \(A\) 是 \(H\) 上的稠定闭线性算子。由\textbf{定理 7.2}（p. 280）立得 \(A = A^{**}\)。又因为 \(A\) 是对称算子，由定义 7.4 知 \(A \subset A^*\)。两式合并即得 
\FitDisplayMath{A = A^{**} \subset A^*.}

\item
  \textbf{充分性（\(\Leftarrow\)）}： 设 \(A = A^{**} \subset A^*\)。由前置验证知 \(A^*\) 是 \(H\) 上的稠定线性算子，对 \(A^*\) 应用\textbf{定理 7.1}（p. 280，稠定线性算子的共轭算子必为闭算子），可知 \((A^*)^* = A^{**}\) 是闭算子。因为 \(A = A^{**}\)，所以 \(A\) 是闭算子。
\end{itemize}

\SoftDivider

\textbf{证明 (ii)：\(A\) 是自伴算子 \(\iff A = A^{**} = A^*\)。}

\begin{itemize}
\tightlist
\item
  \textbf{必要性（\(\Rightarrow\)）}： 设 \(A\) 是自伴算子，即按定义 7.4 有 \(A = A^*\)。由\textbf{定理 7.1}（p. 280）知，自伴算子必为闭算子。再由刚刚证明的 (i)（或直接对 \(A = A^*\) 两边取共轭得 \(A^{**} = (A^*)^* = A^* = A\)）立得 
\FitDisplayMath{A = A^{**} = A^*.}

\item
  \textbf{充分性（\(\Leftarrow\)）}： 设 \(A = A^{**} = A^*\)，则直接有 \(A = A^*\)。又已知 \(A\) 是稠定线性算子，按定义 7.4 即知 \(A\) 是自伴算子。\(\blacksquare\)
\end{itemize}

\SubPillCoral{【考点复盘与避坑提示】}

\begin{itemize}
\tightlist
\item
  \textbf{加分亮点（一定要写的前置步）}：很多同学一上来就用``对 \(A^*\) 用定理 7.1 知 \(A^{**}\) 闭''，却忘了先说明 \textbf{为什么 \(A^*\) 是稠定的}！对一般稠定算子 \(T\)，\(T^*\) 并不一定稠定；本题之所以能无障碍地写 \(A^{**}\)，全靠题设里 \textbf{\(A\) 是对称算子} 给出的 \(\mathcal{D}(A) \subset \mathcal{D}(A^*)\)！
\item
  \textbf{彻底打通对称算子的``黄金包含链''与四级阶梯}： 实际上，对\textbf{任意}对称算子 \(A\)（哪怕不闭），都自动成立 
\FitDisplayMath{A \subset \bar{A} = A^{**} \subset A^*.}
 （速记理由：\CircNum{1} 对任意 \(x \in \mathcal{D}(A), y \in \mathcal{D}(A^*)\)，有 \((A^* y, x) = \overline{(x, A^* y)} = \overline{(Ax, y)} = (y, Ax)\)，直接推出 \(A \subset A^{**}\)；\CircNum{2} 对 \(A \subset A^*\) 两边取共轭并应用扩张反号律引理 7.2 (iii)，立得 \(A^{**} \subset A^*\)！） 于是对称算子的全部类型一目了然：
\end{itemize}

{\def\LTcaptype{none} % do not increment counter
{\small\renewcommand{\arraystretch}{1.38}\setlength{\tabcolsep}{5pt}
\begin{longtable}[]{@{}
  >{\raggedright\arraybackslash}p{(\linewidth - 4\tabcolsep) * \real{0.3333}}
  >{\raggedright\arraybackslash}p{(\linewidth - 4\tabcolsep) * \real{0.3333}}
  >{\raggedright\arraybackslash}p{(\linewidth - 4\tabcolsep) * \real{0.3333}}@{}}
\toprule\noalign{}
\rowcolor{AccentSoft}
\begin{minipage}[b]{\linewidth}\raggedright
算子层级
\end{minipage} & \begin{minipage}[b]{\linewidth}\raggedright
包含链状态
\end{minipage} & \begin{minipage}[b]{\linewidth}\raggedright
典型教材例子
\end{minipage} \\
\midrule\noalign{}
\endhead
\bottomrule\noalign{}
\endlastfoot
\textbf{1. 一般（非闭）对称算子} & \(A \subsetneq \bar{A} = A^{**} \subsetneq A^*\) & 把例 7.2 的 \(A_3\) 再缩小到 \(C_0^\infty(0,1)\) \\
\textbf{2. 闭对称算子（非自伴）} & \(A = A^{**} \subsetneq A^*\) & 例 7.2 中的 \(A_3\)（\(A_3 = A_3^{**} \subsetneq A_1 = A_3^*\)） \\
\textbf{3. 本质自伴算子（闭包自伴）} & \(A \subsetneq \bar{A} = A^{**} = A^*\) & 把乘法算子 \(A\) 限制在紧支光滑函数 \(C_c^\infty(\mathbb{R})\) 上 \\
\textbf{4. 自伴算子} & \(A = A^{**} = A^*\) & 例 7.2 的 \(A_2\)、例 7.3 的位置算子、例 7.4 的动量算子 \\
\end{longtable}
}
}

\SoftDivider

\BookProblem{prob:7-7}{习题 7 第 7 题}{教材 p. 294}{习题 7 第 7 题}

\SubPillPrimary{【原题】}

设 \(A\) 是 \(H\) 上的对称算子，而且对任意 \(x \in \mathcal{D}(A)\) 均有 \((Ax, x) \geqslant 0\)。证明：

\begin{itemize}
\tightlist
\item
  \begin{enumerate}
  \def\labelenumi{(\roman{enumi})}
  \tightlist
  \item
    \(\| (A + I)x \|^2 \geqslant \|x\|^2 + \|Ax\|^2\)；
  \end{enumerate}
\item
  (ii) \(A\) 是闭算子的充要条件是 \(\mathcal{R}(A + I)\) 是闭集。
\end{itemize}

\SubPillAmber{【零基础破题思路】}

这道题是对教材 \textbf{式 (7.9) 与引理 7.8（p. 281）} 的超级经典``双胞胎变形题''！

\begin{enumerate}
\def\labelenumi{\arabic{enumi}.}
\item
  \textbf{先看第 (i) 问}：直接把左边的范数平方按内积展开： 
\FitDisplayMath{\|(A+I)x\|^2 = (Ax + x, Ax + x) = \|Ax\|^2 + (Ax, x) + (x, Ax) + \|x\|^2.}
 因为 \(A\) 对称，所以 \((x, Ax) = (Ax, x)\)；又题目给了\textbf{非负条件} \((Ax, x) \ge 0\)，于是中间两项合起来 \(2(Ax, x) \ge 0\)，直接扔掉就得到 
\FitDisplayMath{\|(A+I)x\|^2 \ge \|x\|^2 + \|Ax\|^2!}
 立刻收获两把万能扳手（\textbf{双控不等式}）： 
\FitDisplayMath{\|x\| \le \|(A+I)x\|, \qquad \|Ax\| \le \|(A+I)x\|.}
 这两把扳手的威力是：\textbf{只要 \((A+I)x_n\) 乖乖收敛（是 Cauchy 列），\(x_n\) 和 \(Ax_n\) 就都被死死按住，同时变成 Cauchy 列！}
\item
  \textbf{再看第 (ii) 问（充要条件两个方向）}：

  \begin{itemize}
  \tightlist
  \item
    \textbf{必要性（\(A\) 闭 \(\implies \mathcal{R}(A+I)\) 闭）}：套路极稳！取 \(y_n = (A+I)x_n \to y\)，目标证明 \(y \in \mathcal{R}(A+I)\)。

    \begin{itemize}
    \tightlist
    \item
      \(\{y_n\}\) 收敛 \(\implies \{y_n\}\) 是 Cauchy 列；
    \item
      把 (i) 用在差向量 \(x_m - x_n\) 上：\(\|x_m - x_n\|^2 + \|Ax_m - Ax_n\|^2 \le \|y_m - y_n\|^2 \to 0\)；
    \item
      立刻推出 \(\{x_n\}\) 和 \(\{Ax_n\}\) 都是 \(H\) 中的 Cauchy 列；
    \item
      由 Hilbert 空间 \(H\) 完备，\(x_n \to x\) 且 \(Ax_n \to z\)；
    \item
      由 \(A\) 是闭算子，\(x \in \mathcal{D}(A)\) 且 \(Ax = z\)；
    \item
      于是 \(y = \lim(Ax_n + x_n) = Ax + x = (A+I)x \in \mathcal{R}(A+I)\)！收工！
    \end{itemize}
  \item
    \textbf{充分性（\(\mathcal{R}(A+I)\) 闭 \(\implies A\) 闭）}：设 \(x_n \in \mathcal{D}(A)\)，\(x_n \to x\) 且 \(Ax_n \to z\)，目标证明 \(x \in \mathcal{D}(A)\) 且 \(Ax = z\)。

    \begin{itemize}
    \tightlist
    \item
      相加立得 \((A+I)x_n = Ax_n + x_n \to z + x\)；
    \item
      因为 \(\mathcal{R}(A+I)\) 是闭集，极限 \(z+x\) 必落回 \(\mathcal{R}(A+I)\)，也就是\textbf{存在某个 \(x_0 \in \mathcal{D}(A)\) 使 \((A+I)x_0 = z + x\)}！
    \item
      现在怎么把陌生的 \(x_0\) 和目标 \(x\) 画上等号？再次祭出扳手 \(\|u\| \le \|(A+I)u\|\)，用在 \(x_n - x_0\) 上： 
\FitDisplayMath{\|x_n - x_0\| \le \|(A+I)(x_n - x_0)\| = \|(A+I)x_n - (z+x)\| \to 0!}

    \item
      这说明 \(x_n \to x_0\)！而题目本来就有 \(x_n \to x\)，由极限唯一性立得 \textbf{\(x = x_0 \in \mathcal{D}(A)\)}！
    \item
      最后顺手算 \(Ax = Ax_0 = (A+I)x_0 - x_0 = (z+x) - x = z\)，大功告成！
    \end{itemize}
  \end{itemize}
\end{enumerate}

\SubPillMain{【规范解答】}

注意 \(\mathcal{D}(A+I) = \mathcal{D}(A) \cap H = \mathcal{D}(A)\)。

\SoftDivider

\textbf{证明 (i)：}

对任意 \(x \in \mathcal{D}(A)\)，将范数平方按内积展开： 
\FitDisplayMath{\|(A + I)x\|^2 = (Ax + x, Ax + x) = (Ax, Ax) + (Ax, x) + (x, Ax) + (x, x) = \|Ax\|^2 + (Ax, x) + (x, Ax) + \|x\|^2.}
 因为 \(A\) 是对称算子，由引理 7.3（p. 277）得 \((x, Ax) = (Ax, x)\)；又由题设对任意 \(x \in \mathcal{D}(A)\) 有 \((Ax, x) \ge 0\)，因此 
\FitDisplayMath{\|(A + I)x\|^2 = \|Ax\|^2 + 2(Ax, x) + \|x\|^2 \ge \|x\|^2 + \|Ax\|^2.}


特别地，由 (i) 立即得到对任意 \(x \in \mathcal{D}(A)\) 均成立两条估计： 
\FitTaggedMath{\|x\| \le \|(A + I)x\|, \qquad \|Ax\| \le \|(A + I)x\|.}{$\ast$}


\SoftDivider

\textbf{证明 (ii)：\(A\) 是闭算子 \(\iff \mathcal{R}(A + I)\) 是闭集。}

\begin{itemize}
\item
  \textbf{必要性（\(\Rightarrow\)）：设 \(A\) 是闭算子，证明 \(\mathcal{R}(A+I)\) 是闭集。}

  任取序列 \(\{y_n\}_{n=1}^\infty \subset \mathcal{R}(A+I)\) 及 \(y \in H\)，满足 \(y_n \to y\)（\(n \to \infty\)）。我们需要证明 \(y \in \mathcal{R}(A+I)\)。

  由 \(y_n \in \mathcal{R}(A+I)\)，对每个 \(n \ge 1\) 存在 \(x_n \in \mathcal{D}(A)\) 使得 \(y_n = (A+I)x_n\)。因为 \(\{y_n\}\) 在 \(H\) 中收敛，故 \(\{y_n\}\) 是 Cauchy 列。对任意 \(m, n \ge 1\)，由于 \(x_m - x_n \in \mathcal{D}(A)\)，将 (i) 应用于 \(x_m - x_n\) 得 
\FitDisplayMath{\|x_m - x_n\|^2 + \|Ax_m - Ax_n\|^2 \le \|(A+I)(x_m - x_n)\|^2 = \|y_m - y_n\|^2 \to 0 \qquad (m, n \to \infty).}
 这表明 \(\{x_n\}\) 与 \(\{Ax_n\}\) 均为 \(H\) 中的 Cauchy 列。由 Hilbert 空间 \(H\) 的完备性，存在 \(x, z \in H\) 使得 
\FitDisplayMath{\lim_{n\to\infty} x_n = x, \qquad \lim_{n\to\infty} Ax_n = z.}
 因为 \(\{x_n\} \subset \mathcal{D}(A)\) 且 \(A\) 是闭算子，由闭算子定义立得 \(x \in \mathcal{D}(A)\) 且 \(Ax = z\)。于是 
\FitDisplayMath{y = \lim_{n\to\infty} y_n = \lim_{n\to\infty} (Ax_n + x_n) = z + x = Ax + x = (A + I)x \in \mathcal{R}(A + I).}
 故 \(\mathcal{R}(A + I)\) 包含其一切极限点，是 \(H\) 中的闭集。
\item
  \textbf{充分性（\(\Leftarrow\)）：设 \(\mathcal{R}(A+I)\) 是闭集，证明 \(A\) 是闭算子。}

  任取序列 \(\{x_n\}_{n=1}^\infty \subset \mathcal{D}(A)\) 及 \(x, z \in H\)，满足 
\FitDisplayMath{\lim_{n\to\infty} x_n = x, \qquad \lim_{n\to\infty} Ax_n = z.}
 我们需要证明 \(x \in \mathcal{D}(A)\) 且 \(Ax = z\)。

  考察序列 \((A+I)x_n = Ax_n + x_n \in \mathcal{R}(A+I)\)，显然有 
\FitDisplayMath{\lim_{n\to\infty} (A + I)x_n = \lim_{n\to\infty} (Ax_n + x_n) = z + x.}
 因为 \(\mathcal{R}(A+I)\) 是闭集，所以极限点 \(z + x \in \mathcal{R}(A+I)\)，即存在 \(x_0 \in \mathcal{D}(A)\) 使得 
\FitDisplayMath{(A + I)x_0 = z + x.}
 对任意 \(n \ge 1\)，由于 \(x_n - x_0 \in \mathcal{D}(A)\)，应用 (\(\ast\)) 式的第一条不等式得 
\FitDisplayMath{\|x_n - x_0\| \le \|(A + I)(x_n - x_0)\| = \|(A + I)x_n - (A + I)x_0\| = \|(A + I)x_n - (z + x)\| \to 0 \qquad (n \to \infty).}
 因此 \(\lim_{n\to\infty} x_n = x_0\)。由极限的唯一性（已有 \(\lim_{n\to\infty} x_n = x\)），立得 
\FitDisplayMath{x = x_0 \in \mathcal{D}(A).}
 进而由 \((A+I)x_0 = z + x\) 计算得 
\FitDisplayMath{Ax = Ax_0 = (A + I)x_0 - x_0 = (z + x) - x = z.}
 故 \(A\) 是闭算子。\(\blacksquare\)
\end{itemize}

\SubPillCoral{【考点复盘与避坑提示】}

\begin{itemize}
\tightlist
\item
  \textbf{为什么教材引理 7.8 用 \(A \pm \mathrm{i}I\)，而这里能用 \(A + I\)？} 对一般对称算子，实数内积 \((Ax, x)\) 可正可负，因此 \((A+I)x\) 展开里的 \(2(Ax,x)\) 可能变成负数导致下界失效（甚至出现特征值 \(-1\) 使得 \(A+I\) 连单射都不是！）；所以一般对称算子必须靠虚数单位 \(\pm\mathrm{i}\) 让交叉项严格抵消（式 7.9：\(\|(A\pm\mathrm{i}I)x\|^2 = \|Ax\|^2 + \|x\|^2\)）。而本题多了\textbf{非负条件 \((Ax,x) \ge 0\)}，相当于把谱挡在了非负半轴 \([0, +\infty)\) 上，于是点 \(-1\) 自动远离谱集，实数移位 \(A+I\) 就拥有了同样强大的双控不等式 \(\|(A+I)x\|^2 \ge \|x\|^2 + \|Ax\|^2\)！
\item
  \textbf{充分性证明的避坑提示}：证明``\(\mathcal{R}(A+I)\) 闭 \(\implies A\) 闭''时，得到存在 \(x_0 \in \mathcal{D}(A)\) 使 \((A+I)x_0 = z + x\) 后，\textbf{绝不能直接假设 \(x_0 = x\)}（因为此时还不知道 \(x \in \mathcal{D}(A)\)，更不能写 \((A+I)x\)！），必须用不等式 \(\|x_n - x_0\| \le \|(A+I)(x_n - x_0)\| \to 0\) 证出 \(x_n \to x_0\)，再借极限唯一性把 \(x_0\) 和 \(x\) 焊接在一起！
\item
  \textbf{考场拓展（和定理 7.3 的强强联合）}：若非负对称算子 \(A\) 进一步满足 \(\mathcal{R}(A+I) = H\)，则因为 \(A+I\) 也是对称算子且满射，由\textbf{定理 7.3 (ii)} 立得 \(A+I\) 自伴（从而 \(A\) 自伴！）。这是判断偏微分方程中正椭圆算子（如 \(-\Delta\)）是否自伴的超级实用判据！
\end{itemize}


% =========================================================================
% 第 8 章 · 广义函数与 Sobolev 空间
% =========================================================================
\SwitchChapterTheme{8}
\ChapterBanner{08}{ch:8}{第 8 章 · 广义函数与 Sobolev 空间}{广义导数与例 8.5/8.6 · 本书 2π 约定 Fourier 变换 · Sobolev 空间与椭圆估计 · 8 道题}
/\# 第 8 章 广义函数与 Sobolev 空间：零基础知识梳理与作业详解

\begin{quote}
\textbf{定位与使用指南}：本文档面向\textbf{此前没有学过泛函分析或偏微分方程}的本科生，严格对齐教材第 8 章（p.303--343）的记号、定义与定理编号，将「知识零基础精讲」「高频公式速查」与「全部 8 道必做作业逐题详解」合为一体。复习时建议先花 15 分钟读完\textbf{两大核心约定}与第 1 部分的核心直觉，再对照第 2 部分逐题刷完并背熟「规范解答」的得分动作。
\end{quote}

\BookPart{}{\WarnBadge{} 必读前置：本书第 8 章两大核心约定（防网课/参考书混坑！）}{【注意】 必读前置：本书第 8 章两大核心约定（防网课/参考书混坑！）}

不同泛函分析与偏微分方程教材在这一章有\textbf{两处最容易打架的记号差异}。做作业和期末考试前，请先锁死本书约定：

\begin{enumerate}
\def\labelenumi{\arabic{enumi}.}
\tightlist
\item
  \textbf{常规广义函数的配对与乘法带复共轭（教材 p.304 例 8.2，p.305）}

  \begin{itemize}
  \tightlist
  \item
    对局部可积函数 \(g \in L^1_{\mathrm{loc}}(\Omega)\)，其诱导的广义函数 \(\lambda_g \in \mathcal{D}'(\Omega)\) 定义为 
\FitDisplayMath{\lambda_g(f) \triangleq \int_{\mathbb{R}^n} f(x)\overline{g(x)}\,\mathrm{d}x, \quad \forall f \in \mathcal{D}(\Omega).}
 （之所以在 \(g\) 上取共轭，是为了与复 Hilbert 空间 \(L^2(\Omega)\) 的内积 \(\langle f, g\rangle_{L^2} = \int f\bar{g}\,\mathrm{d}x\) 保持一致。当 \(g\) 是\textbf{实值函数}时，\(\overline{g} = g\)，直接退化为 \(\lambda_g(f) = \int g(x)f(x)\,\mathrm{d}x\)。）
  \item
    对光滑函数 \(h \in C^\infty(\Omega)\) 与广义函数 \(\lambda \in \mathcal{D}'(\Omega)\)，乘积 \(h\lambda\) 定义为 
\FitDisplayMath{(h\lambda)(f) \triangleq \lambda(\bar{h}f), \quad \forall f \in \mathcal{D}(\Omega).}
 这样对常规分布 \(\lambda_g\) 刚好满足 \((h\lambda_g)(f) = \int \bar{h}f\bar{g}\,\mathrm{d}x = \int f\overline{hg}\,\mathrm{d}x = \lambda_{hg}(f)\)，与普通函数乘法完全相容！
  \end{itemize}
\item
  \textbf{Fourier 变换在指数上带 \(2\pi\) 因子（教材 p.312 式 (8.9)）}

  \begin{itemize}
  \tightlist
  \item
    对 \(f \in L^1(\mathbb{R}^n)\)，本书定义 Fourier 变换为 
\FitDisplayMath{\hat{f}(\xi) = (\mathcal{F}f)(\xi) \triangleq \int_{\mathbb{R}^n} f(x)\mathrm{e}^{-2\pi\mathrm{i}\xi\cdot x}\,\mathrm{d}x.}

  \item
    由此联动产生\textbf{全部标准系数}（考试写错 \(2\pi\) 会直接扣结果分！）：

    \begin{itemize}
    \tightlist
    \item
      求导变乘法：\(\widehat{\partial_j f}(\xi) = 2\pi\mathrm{i}\xi_j \hat{f}(\xi)\)（教材 p.313 命题 8.15）；
    \item
      乘坐标变求导：\(\partial_{\xi_j}\hat{f}(\xi) = -2\pi\mathrm{i}\widehat{x_j f}(\xi)\)（教材 p.312 命题 8.14）；
    \item
      \textbf{反演公式与 Plancherel 等式均不带 \((2\pi)^{-n}\)}： 
\FitDisplayMath{f(x) = \int_{\mathbb{R}^n} \hat{f}(\xi)\mathrm{e}^{2\pi\mathrm{i}\xi\cdot x}\,\mathrm{d}\xi, \qquad \|\hat{f}\|_{L^2(\mathbb{R}^n)} = \|f\|_{L^2(\mathbb{R}^n)}.}

    \end{itemize}
  \item
    \textbf{换算秘籍}：如果网课或别的参考书采用不带 \(2\pi\) 的版本 \(\tilde{f}(\omega) = \int f(x)\mathrm{e}^{-\mathrm{i}\omega\cdot x}\,\mathrm{d}x\)，二者关系极简单——只需代入 \(\omega = 2\pi\xi\)，即 \(\tilde{f}(\omega) = \hat{f}\!\left(\frac{\omega}{2\pi}\right)\)。
  \end{itemize}
\end{enumerate}

\BookPart{}{第一部分：第 8 章核心知识零基础精讲}{第一部分：第 8 章核心知识零基础精讲}

\BookSubSection{}{一、为什么要学广义函数与 Sobolev 空间？（一句话看懂全章主线）}{一、为什么要学广义函数与 Sobolev 空间？（一句话看懂全章主线）}{}

传统微积分有两个致命痛点：

\begin{enumerate}
\def\labelenumi{\arabic{enumi}.}
\tightlist
\item
  \textbf{很多物理与工程对象「不能求导」}：比如阶跃函数 \(\operatorname{sgn}(x)\) 在 \(0\) 点跳变、绝对值 \(|x|\) 在 \(0\) 点有尖角、点电荷密度 \(\delta\) 甚至不是普通函数。
\item
  \textbf{偏微分方程（PDE）的经典解空间不完备}：用逼近法解方程时，一列光滑解的极限往往失去经典导数，掉出了 \(C^k\) 空间。
\end{enumerate}

全章就是用一套\textbf{分步升级的巧妙设计}解决这两个痛点：

\begin{itemize}
\tightlist
\item
  \textbf{第 1 步（8.2 节）}：先挑出一批「极其乖巧」的探针——紧支光滑函数 \(\mathcal{D}(\Omega)\)；
\item
  \textbf{第 2 步（8.3 节）}：不再问函数``在某一点等于几''，而改为问它``和每个探针做积分得到几''，于是诞生了\textbf{广义函数（分布）\(\mathcal{D}'(\Omega)\)}，任何广义函数都能无穷次求导！
\item
  \textbf{第 3 步（8.4 节）}：为了把求导变成代数乘法（Fourier 变换），把探针稍微放大成\textbf{速降函数 \(\mathcal{S}(\mathbb{R}^n)\)}，对偶得到\textbf{缓增分布 \(\mathcal{S}'(\mathbb{R}^n)\)}；
\item
  \textbf{第 4 步（8.6--8.9 节）}：从广义函数里挑出``本身和各阶弱导数都在 \(L^p\)（尤其 \(L^2\)）里''的优等生，建成完备的 \textbf{Sobolev 空间 \(W^{m,p}(\Omega)\) 与实指数空间 \(H^s(\mathbb{R}^n)\)}，成为现代偏微分方程的绝对舞台。
\end{itemize}

\SoftDivider

\BookSubSection{sec:8-2}{二、8.2 试验函数空间 \(\mathcal{D}(\Omega) = C_c^\infty(\Omega)\)（教材 p.303--304）}{二、8.2 试验函数空间 D(Ω) = C\_c\textasciicircum{}∞(Ω)（教材 p.303–304）}{\TagLight{基础理解}}

\SubSubHeading{1. 零基础直觉：什么是「超级探针」？}

试验函数 \(\varphi \in C_c^\infty(\Omega)\)（也记作 \(C_0^\infty(\Omega)\)）同时满足两大神级性质：

\begin{itemize}
\tightlist
\item
  \textbf{无穷次光滑（\(C^\infty\)）}：想对它求多少阶偏导都行，永远不会出错；
\item
  \textbf{紧支集（subscript \(c\) / \(0\)）}：\(\operatorname{supp}\varphi \triangleq \overline{\{x \in \Omega : \varphi(x) \neq 0\}}\) 是包含在开集 \(\Omega\) 内部的紧集（在 \(\mathbb{R}^n\) 中即有界闭集）。换句话说，\textbf{只要离开某个有限区域，\(\varphi\) 及其一切导数统统严格等于 \(0\)！}
\end{itemize}

\begin{quote}
\textbf{为什么这个设计如此天才？} 因为做分部积分时，\(\int_\Omega u\,\varphi'\,\mathrm{d}x = \underbrace{[u\varphi]_{\partial\Omega}}_{=\,0} - \int_\Omega u'\,\varphi\,\mathrm{d}x\)，\textbf{边界项永远自动归零}！
\end{quote}

\SubSubHeading{2. 严格收敛定义与核心性质}

\begin{itemize}
\tightlist
\item
  \textbf{定义 8.6（\(\mathcal{D}(\Omega)\) 中的收敛，教材 p.303）}：设 \(\{f_k\}_{k=1}^\infty \subset C_c^\infty(\Omega)\)，\(f \in C_c^\infty(\Omega)\)。称 \(f_k \to f\) 于 \(C_c^\infty(\Omega)\)（记为 \(\mathcal{D}(\Omega)\)），当且仅当同时满足两条：

  \begin{enumerate}
  \def\labelenumi{\arabic{enumi}.}
  \tightlist
  \item
    \textbf{共同紧支集}：存在\textbf{同一个}紧集 \(K \subset \Omega\)，使得所有 \(f_k\) 与 \(f\) 的支集都包含在 \(K\) 中；
  \item
    \textbf{任意阶导数一致收敛}：对每个非负整数 \(j = 0, 1, 2, \dots\)，\(\nabla^j f_k\) 在 \(K\) 上一致收敛于 \(\nabla^j f\)（即对一切多重指标 \(\alpha\)，\(\sup_{x \in K} |\partial^\alpha f_k(x) - \partial^\alpha f(x)| \to 0\)）。
  \end{enumerate}
\item
  \textbf{命题 8.2（p.303）}：对紧集 \(K \subset \mathbb{R}^n\) 与赋范空间 \(X\)，线性映射 \(T: C_c^\infty(K) \to X\) 连续 \(\iff \exists k \ge 0, C > 0\) 使得 \(\|Tf\|_X \le C\|f\|_{C^k(K)}\)。
\item
  \textbf{命题 8.4（p.303）}：对任意 \(1 \le p \le \infty\)，嵌入 \(C_c^\infty(\Omega) \hookrightarrow L^p(\Omega)\) 连续单射。
\item
  \textbf{命题 8.6（p.304）}：光滑系数的线性微分算子是 \(\mathcal{D}(\Omega)\) 上的连续线性算子。
\end{itemize}

\SoftDivider

\BookSubSection{sec:8-3}{三、8.3 广义函数（分布）\(\mathcal{D}'(\Omega)\)（教材 p.304--312）}{三、8.3 广义函数（分布）D'(Ω)（教材 p.304–312）}{\TagHigh{作业核心}}

\SubSubHeading{1. 基本定义与连续性判定}

\begin{itemize}
\tightlist
\item
  \textbf{定义 8.7（广义函数，p.304）}：称连续线性泛函 \(\lambda: \mathcal{D}(\Omega) \to \mathbb{C}\) 为 \(\Omega\) 上的\textbf{广义函数}（也称分布），全体记为 \(\mathcal{D}'(\Omega)\)。
\item
  \textbf{连续性实用判定（式 (8.3)，p.304）}：线性泛函 \(\lambda: \mathcal{D}(\mathbb{R}^n) \to \mathbb{C}\) 是广义函数 \(\iff\) 对每个紧集 \(K \subset \mathbb{R}^n\)，存在非负整数 \(k\)（可依赖 \(K\)）与常数 \(C > 0\)，使得 
\FitDisplayMath{|\lambda(\varphi)| \le C\|\varphi\|_{C^k(K)}, \quad \forall \varphi \in C_c^\infty(K).}


  \begin{itemize}
  \tightlist
  \item
    若 \(k\) 可不依赖紧集 \(K\) 统一选取，则称 \(\lambda\) 至多 \(k\) 阶（定义 8.12，p.311）：如 \(\delta\) 为 \(0\) 阶，\(\delta'\) 与 \(\operatorname{p.v.}\frac{1}{x}\) 为 \(1\) 阶。
  \end{itemize}
\item
  \textbf{定义 8.8（广义函数列的收敛，p.304）}：称 \(\lambda_k \to \lambda\) 于 \(\mathcal{D}'(\Omega)\)，当且仅当\textbf{对每个固定的试验函数} \(\varphi \in \mathcal{D}(\Omega)\)，数列满足 
\FitDisplayMath{\lim_{k\to\infty}\lambda_k(\varphi) = \lambda(\varphi).}

\item
  \textbf{收敛强弱关系（p.305，必考结论！）}： 
\FitDisplayMath{L^p \text{ 收敛或一致收敛} \implies L^1_{\mathrm{loc}}(\Omega) \text{ 收敛} \implies \mathcal{D}'(\Omega) \text{ 收敛}.}
 只要 \(g_k \to g\) 于 \(L^1_{\mathrm{loc}}(\Omega)\)（即在每个紧集 \(K \subset \Omega\) 上 \(\int_K |g_k - g|\,\mathrm{d}x \to 0\)），就自动有 \(\lambda_{g_k} \to \lambda_g\) 于 \(\mathcal{D}'(\Omega)\)！
\end{itemize}

\SubSubHeading{2. 五大必背经典广义函数（教材例 8.2--8.6，p.304--305）}

{\def\LTcaptype{none} % do not increment counter
{\small\renewcommand{\arraystretch}{1.38}\setlength{\tabcolsep}{5pt}
\begin{longtable}[]{@{}
  >{\raggedright\arraybackslash}p{(\linewidth - 6\tabcolsep) * \real{0.1233}}
  >{\raggedright\arraybackslash}p{(\linewidth - 6\tabcolsep) * \real{0.1781}}
  >{\raggedright\arraybackslash}p{(\linewidth - 6\tabcolsep) * \real{0.4384}}
  >{\raggedright\arraybackslash}p{(\linewidth - 6\tabcolsep) * \real{0.2603}}@{}}
\toprule\noalign{}
\rowcolor{AccentSoft}
\begin{minipage}[b]{\linewidth}\raggedright
名称
\end{minipage} & \begin{minipage}[b]{\linewidth}\raggedright
记号
\end{minipage} & \begin{minipage}[b]{\linewidth}\raggedright
对 \(\varphi \in \mathcal{D}(\Omega)\) 的作用公式
\end{minipage} & \begin{minipage}[b]{\linewidth}\raggedright
教材出处与备注
\end{minipage} \\
\midrule\noalign{}
\endhead
\bottomrule\noalign{}
\endlastfoot
常规广义函数 & \(\lambda_g\)（\(g \in L^1_{\mathrm{loc}}(\Omega)\)） & \(\lambda_g(\varphi) = \int_{\Omega} \varphi(x)\overline{g(x)}\,\mathrm{d}x\) & \textbf{例 8.2（p.304）}；实值 \(g\) 时为 \(\int g\varphi\,\mathrm{d}x\) \\
Dirac 脉冲函数 & \(\delta\)（Radon 测度特例） & \(\delta(\varphi) = \varphi(0)\) & \textbf{例 8.3（p.304）}；\(0\) 阶，非函数 \\
\(\delta\) 的导数 & \(\delta'\)（\(n=1\)） & \(\delta'(\varphi) = -\varphi'(0)\) & \textbf{例 8.4（p.305）}；\(1\) 阶，非 Radon 测度 \\
Cauchy 主值 & \(\operatorname{p.v.}\dfrac{1}{x}\)（\(n=1\)） & \(\left(\operatorname{p.v.}\dfrac{1}{x}\right)(\varphi) = \lim\limits_{\varepsilon\to 0^+}\displaystyle\int_{\lvert x\rvert>\varepsilon}\dfrac{\varphi(x)}{x}\,\mathrm{d}x\) & \textbf{例 8.5（p.305）}；因 \(\frac{1}{x} \notin L^1_{\mathrm{loc}}\) 须挖洞对称取极限 \\
例 8.6 广义函数 & \(\lambda_r\)（\(n=1, r>0\)） & \(\lambda_r(\varphi) = \displaystyle\int_{\lvert x\rvert<r}\dfrac{\varphi(x)-\varphi(0)}{\lvert x\rvert}\,\mathrm{d}x + \displaystyle\int_{\lvert x\rvert\ge r}\dfrac{\varphi(x)}{\lvert x\rvert}\,\mathrm{d}x\) & \textbf{例 8.6（p.305）}；\textbf{习题 8 第 11(3) 题直接考 \(r=1\) 即 \(\lambda_1\)！} \\
\end{longtable}
}
}

\begin{quote}
\textbf{零基础秒懂例 8.5（\(\operatorname{p.v.}\frac{1}{x}\)）与例 8.6（\(\lambda_r\)）}：

\begin{itemize}
\tightlist
\item
  \textbf{例 8.5（Cauchy 主值 \(\operatorname{p.v.}\frac{1}{x}\)，p.305）}：\(\frac{1}{x}\) 在原点附近不可积（\(\int_0^1 \frac{\mathrm{d}x}{x} = +\infty\)），不能直接按普通积分定义。但 \(\frac{1}{x}\) 是\textbf{奇函数}，左右两侧分别冲向 \(-\infty\) 和 \(+\infty\)；只要\textbf{对称地挖掉小区间 \((-\varepsilon, \varepsilon)\)}（只积 \(|x|>\varepsilon\)），两边的无穷大刚好精准抵消！设 \(\operatorname{supp}\varphi \subset [-R, R]\)，由奇函数对称积分为 \(0\) 得 \(\int_{\varepsilon<|x|\le R}\frac{\varphi(0)}{x}\,\mathrm{d}x = 0\)，于是 
\FitDisplayMath{\int_{|x|>\varepsilon}\frac{\varphi(x)}{x}\,\mathrm{d}x = \int_{\varepsilon<|x|\le R}\frac{\varphi(x)-\varphi(0)}{x}\,\mathrm{d}x \xrightarrow{\varepsilon\to 0^+} \int_{-R}^{R}\frac{\varphi(x)-\varphi(0)}{x}\,\mathrm{d}x,}
 由中值定理 \(\left|\frac{\varphi(x)-\varphi(0)}{x}\right| \le \|\varphi'\|_\infty\)，积分收敛且 \(\left|\left(\operatorname{p.v.}\frac{1}{x}\right)(\varphi)\right| \le 2R\|\varphi'\|_\infty\)（\(1\) 阶广义函数，既非局部可积函数也非 Radon 测度）。
\item
  \textbf{例 8.6（\(\lambda_r\)，p.305）}：函数 \(\frac{1}{|x|}\) 同样在原点不可积（且是偶函数，挖洞也不能抵消！）。因此在 \(|x|<r\) 内直接把 \(\varphi(x)\) 减去常数 \(\varphi(0)\)，由于 \(|\varphi(x)-\varphi(0)| \le |x|\|\varphi'\|_\infty\)，分子的 \(|x|\) 恰好消掉分母的 \(|x|\)，积分立刻收敛；外围 \(|x|\ge r\) 本来无奇点，直接积分即可。
\end{itemize}
\end{quote}

\SubSubHeading{3. 广义函数的核心运算与雷区}

\begin{enumerate}
\def\labelenumi{\arabic{enumi}.}
\tightlist
\item
  \textbf{广义导数（p.307--308）}： 
\FitDisplayMath{\frac{\partial}{\partial x_j}\lambda(\varphi) \triangleq -\lambda\!\left(\frac{\partial}{\partial x_j}\varphi\right), \qquad \partial^\alpha \lambda(\varphi) = (-1)^{|\alpha|}\lambda(\partial^\alpha \varphi).}


  \begin{itemize}
  \tightlist
  \item
    \textbf{万能记忆口诀}：「\textbf{求导甩给试验函数，每甩一次加一个负号！}」
  \item
    \textbf{例 8.8（p.308）}：若 \(f \in C^m(\mathbb{R}^n)\)，其广义导数与经典导数完全一致。
  \item
    \textbf{命题 8.9（p.310，式 (8.6)）}：任意广义函数的混合偏导数\textbf{无条件可交换}次序：\(\partial_{x_i}\partial_{x_j}\lambda = \partial_{x_j}\partial_{x_i}\lambda\)。
  \item
    \textbf{\HotBadge{} 超级大杀器——求导与分布极限无条件交换（p.307）}： 
\FitDisplayMath{\lambda_k \to \lambda \text{ 于 } \mathcal{D}'(\Omega) \implies \partial^\alpha \lambda_k \to \partial^\alpha \lambda \text{ 于 } \mathcal{D}'(\Omega).}
 \emph{证明只有一行}：\((\partial^\alpha\lambda_k)(\varphi) = (-1)^{|\alpha|}\lambda_k(\partial^\alpha\varphi) \to (-1)^{|\alpha|}\lambda(\partial^\alpha\varphi) = (\partial^\alpha\lambda)(\varphi)\)！ 结合「\(L^1_{\mathrm{loc}}\) 收敛 \(\implies \mathcal{D}'\) 收敛」，遇到难算极限的 \(f_\varepsilon\) 时，\textbf{先找它的原函数 \(F_\varepsilon\)，证 \(F_\varepsilon \to F\) 于 \(L^1_{\mathrm{loc}}\)，再两边求广义导数}，一招秒杀（见\hyperlink{prob:8-14}{\textcolor{LinkColor}{习题 8 第 14 题}}）！
  \item
    \textbf{定理 8.7（p.310）}：一维区间上任意广义函数必有原函数；导数为 \(0\) 的广义函数必为常数。
  \item
    \textbf{定理 8.8（p.310）}：一维区间上若 \(\lambda' \in L^1_{\mathrm{loc}}((a,b))\)，则 \(\lambda\) 必为绝对连续函数。
  \end{itemize}
\item
  \textbf{光滑函数乘法（p.305，p.309 命题 8.8）}：

  \begin{itemize}
  \tightlist
  \item
    对 \(h \in C^\infty(\Omega), \lambda \in \mathcal{D}'(\Omega)\)，\((h\lambda)(\varphi) = \lambda(\bar{h}\varphi)\)；特别有 \(x\delta = 0\)，以及 Leibniz 乘法法则 \(\partial_{x_j}(h\lambda) = h(\partial_{x_j}\lambda) + (\partial_{x_j}h)\lambda\)。
  \item
    \textbf{\WarnBadge{} 两大禁区（p.311--312，注记 8.7）}：

    \begin{enumerate}
    \def\labelenumii{\arabic{enumii}.}
    \tightlist
    \item
      不能对一般广义函数施加非线性算子（如取绝对值 \(|\lambda|\)）；
    \item
      \textbf{不能把两个一般广义函数相乘}（如 \(\delta^2\) 无意义），也\textbf{不能用不光滑的函数乘一般广义函数}（如 \(f \in L^1(\mathbb{R})\) 时 \(f\delta\) 无法良定义，见\hyperlink{prob:8-7}{\textcolor{LinkColor}{习题 8 第 7 题}}）。注记 8.7 指出：即使 \(\varphi_{1/n} \to \delta\)，平方后的 \(\varphi_{1/n}^2\) 却发散到无穷！
    \end{enumerate}
  \end{itemize}
\item
  \textbf{支集与稠密性（p.306--307）}：

  \begin{itemize}
  \tightlist
  \item
    定义 8.11 与注记 8.5：\(\operatorname{supp}\lambda \subset K\) 要求 \(\varphi\) 在 \(K\) 的某个\textbf{开邻域}内恒为 \(0\) 时 \(\lambda(\varphi)=0\)（不能只在 \(K\) 本身为 \(0\)，例如 \(\varphi(0)=0\) 时 \(\delta'(\varphi)=-\varphi'(0)\) 仍可非零）。
  \item
    引理 8.7：\(\mathcal{D}(\mathbb{R}^n)\) 在 \(\mathcal{D}'(\mathbb{R}^n)\) 中稠密。
  \end{itemize}
\end{enumerate}

\SoftDivider

\BookSubSection{sec:8-4}{四、8.4 Schwartz 空间 \(\mathcal{S}(\mathbb{R}^n)\)、缓增分布 \(\mathcal{S}'(\mathbb{R}^n)\) 与 Fourier 变换（教材 p.312--319）}{四、8.4 Schwartz 空间 S(R\textasciicircum{}n)、缓增分布 S'(R\textasciicircum{}n) 与 Fourier 变换（教材 p.312–319）}{\TagHigh{作业核心}}

\SubSubHeading{1. 为什么不能直接在 \(\mathcal{D}'(\mathbb{R}^n)\) 上做 Fourier 变换？}

一个非零紧支光滑函数 \(\varphi \in \mathcal{D}(\mathbb{R}^n)\)，它的 Fourier 变换 \(\hat{\varphi}\) 是解析函数，\textbf{不可能再有紧支集}（除非 \(\varphi \equiv 0\)）！也就是说，\(\mathcal{F}\) 把 \(\mathcal{D}(\mathbb{R}^n)\) 映出了自己，无法通过 \((\mathcal{F}\lambda)(\varphi) = \lambda(\mathcal{F}\varphi)\) 给一般分布定义 Fourier 变换。

\textbf{解决方案}：把探针空间从「紧支光滑 \(\mathcal{D}(\mathbb{R}^n)\)」适度放宽为「无穷远比任何多项式衰减得都快的光滑函数」，即 \textbf{Schwartz 空间 \(\mathcal{S}(\mathbb{R}^n)\)}！

\begin{itemize}
\tightlist
\item
  \textbf{定义 8.14（Schwartz 空间 \(\mathcal{S}(\mathbb{R}^n)\)，p.313）}：\(f \in C^\infty(\mathbb{R}^n; \mathbb{C})\) 称为速降函数，如果对任意多重指标 \(\alpha, \beta\) 都有 
\FitDisplayMath{\lim_{|x|\to\infty} x^\alpha \partial^\beta f(x) = 0.}
 最典型的非紧支速降函数是 Gauss 函数 \(\mathrm{e}^{-a|x|^2}\)（\(a > 0\)）。
\item
  \textbf{命题 8.16--8.17（p.314--316）}：\(\mathcal{S}(\mathbb{R}^n)\) 对求导、乘多项式、乘法、卷积封闭，在 \(L^p(\mathbb{R}^n)\)（\(1\le p<\infty\)）中稠密；最重要的是，\textbf{\(\mathcal{F}: \mathcal{S}(\mathbb{R}^n) \to \mathcal{S}(\mathbb{R}^n)\) 是双射同胚，且逆变换就是伴随变换 \(\mathcal{F}^{-1} = \mathcal{F}^*\)}！
\end{itemize}

\SubSubHeading{2. 缓增分布 \(\mathcal{S}'(\mathbb{R}^n)\)（教材 p.317--318）}

\begin{itemize}
\item
  \textbf{定义 8.15（p.317）}：\(\mathcal{S}(\mathbb{R}^n)\) 上的连续线性泛函全体称为\textbf{缓增分布空间} \(\mathcal{S}'(\mathbb{R}^n)\)。
\item
  \textbf{包含链}：因为 \(\mathcal{D}(\mathbb{R}^n) \hookrightarrow \mathcal{S}(\mathbb{R}^n)\) 连续且稠密，取对偶后方向反转： 
\FitDisplayMath{\mathcal{D}(\mathbb{R}^n) \subset \mathcal{S}(\mathbb{R}^n) \subset L^2(\mathbb{R}^n) \subset \mathcal{S}'(\mathbb{R}^n) \subset \mathcal{D}'(\mathbb{R}^n).}

\item
  \textbf{直观辨析（例 8.14 \& 注记 8.9，p.317）}：``缓增''粗略指在无穷远处\textbf{至多呈多项式增长}（如多项式、有界函数、\(L^p\) 函数、\(\delta\) 都是缓增分布）；但像 \(\mathrm{e}^x\) 那样指数爆炸的函数属于 \(\mathcal{D}'(\mathbb{R})\) 却\textbf{不属于} \(\mathcal{S}'(\mathbb{R})\)！
\item
  \textbf{定义 8.16（缓增分布的 Fourier 变换，p.317）}：对 \(\lambda \in \mathcal{S}'(\mathbb{R}^n)\)，定义 \((\mathcal{F}\lambda)(f) \triangleq \lambda(\mathcal{F}f)\)（\(\forall f \in \mathcal{S}(\mathbb{R}^n)\)）。

  \begin{quote}
  \emph{注记}：在实值或双线性配对 \(\langle \lambda, f\rangle = \int gf\) 下，上式与普通函数的 Fourier 变换完全一致；在复共轭配对 \(\lambda_g(f) = \int f\bar{g}\) 下，由 Parseval 等式 \(\langle f, \mathcal{F}g\rangle_{L^2} = \langle \mathcal{F}^* f, g\rangle_{L^2}\) 可知应对偶使用 \(\mathcal{F}^*\)。本章所有实值/偶函数及命题 8.21 的运算律均直接按教材公式使用，不会产生任何符号冲突。
  \end{quote}
\end{itemize}

\SubSubHeading{3. \FastBadge{} 本书 \(2\pi\) 约定下 Fourier 变换高频公式速查表（p.312--318）}

{\def\LTcaptype{none} % do not increment counter
{\small\renewcommand{\arraystretch}{1.38}\setlength{\tabcolsep}{5pt}
\begin{longtable}[]{@{}
  >{\raggedright\arraybackslash}p{(\linewidth - 6\tabcolsep) * \real{0.2500}}
  >{\raggedright\arraybackslash}p{(\linewidth - 6\tabcolsep) * \real{0.2500}}
  >{\raggedright\arraybackslash}p{(\linewidth - 6\tabcolsep) * \real{0.2500}}
  >{\raggedright\arraybackslash}p{(\linewidth - 6\tabcolsep) * \real{0.2500}}@{}}
\toprule\noalign{}
\rowcolor{AccentSoft}
\begin{minipage}[b]{\linewidth}\raggedright
运算 / 性质
\end{minipage} & \begin{minipage}[b]{\linewidth}\raggedright
物理空间（\(x\)）
\end{minipage} & \begin{minipage}[b]{\linewidth}\raggedright
频域空间（\(\xi\)）
\end{minipage} & \begin{minipage}[b]{\linewidth}\raggedright
教材编号
\end{minipage} \\
\midrule\noalign{}
\endhead
\bottomrule\noalign{}
\endlastfoot
定义（正变换） & \(f(x)\) & \(\hat{f}(\xi) = \displaystyle\int_{\mathbb{R}^n} f(x)\mathrm{e}^{-2\pi\mathrm{i}\xi\cdot x}\,\mathrm{d}x\) & 式 (8.9)（p.312） \\
反演公式 & \(f(x) = \displaystyle\int_{\mathbb{R}^n} \hat{f}(\xi)\mathrm{e}^{2\pi\mathrm{i}\xi\cdot x}\,\mathrm{d}\xi\) & \(\hat{f}(\xi)\) & 命题 8.17、式 (8.18) \\
Riemann-Lebesgue & \(f \in L^1(\mathbb{R}^n)\) & \(\lim\limits_{\lvert\xi\rvert\to\infty}\hat{f}(\xi) = 0\)（且连续） & 命题 8.13（p.312） \\
导数 \(\leftrightarrow\) 乘多项式 & \(\partial_{x_j} f\)（或 \(\partial^\alpha f\)） & \(2\pi\mathrm{i}\xi_j \hat{f}(\xi)\)（或 \((2\pi\mathrm{i}\xi)^\alpha \hat{f}(\xi)\)） & 命题 8.15、8.21(6) \\
乘坐标 \(\leftrightarrow\) 频域求导 & \(-2\pi\mathrm{i}x_j f(x)\) & \(\partial_{\xi_j}\hat{f}(\xi)\) & 命题 8.14、8.21(6) \\
卷积 \(\leftrightarrow\) 逐点乘积 & \((f * g)(x)\) & \(\hat{f}(\xi)\hat{g}(\xi)\) & 命题 8.21(2)（p.317） \\
乘积 \(\leftrightarrow\) 卷积 & \(f(x)g(x)\) & \((\hat{f} * \hat{g})(\xi)\) & 命题 8.21(2)（p.317） \\
平移 \(\leftrightarrow\) 相位调制 & \((\tau_{x_0}f)(x) = f(x - x_0)\) & \(\mathrm{e}^{-2\pi\mathrm{i}\xi\cdot x_0}\hat{f}(\xi)\) & 命题 8.21(3)（p.317） \\
线性伸缩 & \(f_L(x) = f(Lx)\)（\(L\) 可逆） & \(\dfrac{1}{\lvert\det L\rvert}\hat{f}\big((L^*)^{-1}\xi\big)\) & 命题 8.18、8.21(5) \\
Parseval / Plancherel & \(\langle f, g\rangle_{L^2},\ \lVert f\rVert_{L^2}\) & \(\langle \hat{f}, \hat{g}\rangle_{L^2},\ \lVert\hat{f}\rVert_{L^2}\)（酉算子！） & 定理 8.9（p.316） \\
\(\delta\) 与常数 \(1\) & \(\delta\) ；常数 \(1\) & \(\hat{\delta} = 1\) ；\(\hat{1} = \delta\) & 命题 8.21(7)（p.318） \\
Gauss 核（\(a>0\)） & \(\mathrm{e}^{-ax^2}\)（\(1\) 维） & \(\sqrt{\dfrac{\pi}{a}}\,\mathrm{e}^{-\frac{\pi^2\xi^2}{a}}\) & 习题 8 第 17(3) 题 \\
Poisson 核（\(a>0\)） & \(\dfrac{1}{a^2+x^2}\)（\(1\) 维） & \(\dfrac{\pi}{a}\,\mathrm{e}^{-2\pi a\lvert\xi\rvert}\) & 习题 8 第 17(2) 题 \\
\end{longtable}
}
}

\SoftDivider

\BookSubSection{sec:8-6}{五、8.6--8.7 整数阶 Sobolev 空间 \(W^{m,p}(\Omega)\) 与嵌入定理（教材 p.321--335）}{五、8.6–8.7 整数阶 Sobolev 空间 W\textasciicircum{}m,p(Ω) 与嵌入定理（教材 p.321–335）}{\TagLight{基础理解}}

\SubSubHeading{1. 弱导数与 \(W^{m,p}(\Omega)\)}

\begin{itemize}
\tightlist
\item
  \textbf{定义 8.17（弱导数，p.321）}：设 \(f \in L^p(\Omega)\)，\(\alpha \in \mathbb{N}^n\)。若存在函数 \(g \in L^1(\Omega)\)（更一般地 \(L^1_{\mathrm{loc}}(\Omega)\)，见 p.309）满足对所有试验函数 \(\varphi \in C_0^\infty(\Omega)\) 均有 
\FitDisplayMath{\int_\Omega f\,\partial^\alpha\varphi\,\mathrm{d}x = (-1)^{|\alpha|}\int_\Omega g\,\varphi\,\mathrm{d}x,}
 则称 \(g\) 是 \(f\) 的 \(\alpha\) 阶\textbf{弱导数}，记作 \(\partial^\alpha f = g\)。
\item
  \textbf{\WarnBadge{} 弱导数 vs. 广义导数（注记 8.15，p.321，概念常考！）}：

  \begin{itemize}
  \tightlist
  \item
    任何局部可积函数都有\textbf{广义导数}（落在 \(\mathcal{D}'(\Omega)\) 里）；
  \item
    只有当这个广义导数\textbf{恰好还是一个局部可积函数}时，才叫\textbf{弱导数}！
  \item
    \emph{秒懂对比}：\(|x|\) 的广义导数是 \(\operatorname{sgn}(x) \in L^1_{\mathrm{loc}}(\mathbb{R})\)，所以 \(|x|\) \textbf{有弱导数}；而 \(\operatorname{sgn}(x)\) 的广义导数是 \(2\delta\)，不是函数，所以 \(\operatorname{sgn}(x)\) \textbf{没有弱导数}！
  \end{itemize}
\item
  \textbf{定义 8.18（Sobolev 空间，p.321--322）}：设 \(1 \le p \le \infty, m \in \mathbb{N}\)，定义 
\FitDisplayMath{W^{m,p}(\Omega) \triangleq \{f \in L^p(\Omega) : \partial^\alpha f \in L^p(\Omega),\ \forall |\alpha| \le m\}, \qquad \|f\|_{W^{m,p}(\Omega)} \triangleq \sum_{|\alpha|\le m}\|\partial^\alpha f\|_{L^p(\Omega)}.}


  \begin{itemize}
  \tightlist
  \item
    特别记 \(H^m(\Omega) \triangleq W^{m,2}(\Omega)\)（是 Hilbert 空间；注：教材 p.338 行首偶将 \(W^{m,2}\) 误印为 \(W^{2,m}\)，以定义 8.18 及 p.336 命题 8.28 为准）。
  \item
    \textbf{注记 8.16}：\(W^{0,p}(\Omega) = L^p(\Omega)\)，且 \(W^{m+1,p}(\Omega) \subset W^{m,p}(\Omega)\)。
  \item
    \textbf{引理 8.8（p.322）}：当 \(1 \le p < \infty\) 时，\(C_c^\infty(\mathbb{R}^n)\) 在 \(W^{k,p}(\mathbb{R}^n)\) 中稠密。
  \end{itemize}
\end{itemize}

\SubSubHeading{2. Sobolev 嵌入定理（p.327--335，直观口诀记忆！）}

\begin{quote}
\textbf{核心直觉}：「\textbf{导数阶数可以兑换可积指数（甚至连续性），维数 \(n\) 越高兑换成本越贵！}」记住核心标度指标 \(\frac{1}{p} - \frac{k}{n}\)：
\end{quote}

\begin{itemize}
\tightlist
\item
  \textbf{情形 A：低光滑 \(kp < n\)（可积性提升，定理 8.11 \& 8.13）}

  \begin{itemize}
  \tightlist
  \item
    一阶（定理 8.11，p.327）：\(1 \le p < n \implies W^{1,p}(\Omega) \hookrightarrow L^{p^*}(\Omega)\) 连续嵌入，其中 Sobolev 共轭指数满足 
\FitDisplayMath{\frac{1}{p^*} = \frac{1}{p} - \frac{1}{n} \iff p^* = \frac{np}{n - p}.}

  \item
    高阶（定理 8.13，p.331）：\(kp < n \implies W^{k,p}(\Omega) \hookrightarrow L^{\frac{np}{n-kp}}(\Omega)\)。
  \end{itemize}
\item
  \textbf{情形 B：高光滑 \(kp > n\)（跃升为连续/Hölder 光滑函数，定理 8.12 \& 8.13，p.331）}

  \begin{itemize}
  \tightlist
  \item
    \(W^{k,p}(\Omega) \hookrightarrow C^{m,\alpha}(\Omega)\)，其中 \(m\) 取小于 \(k - \frac{n}{p}\) 的非负整数（特别一阶 \(p > n\) 时嵌入到 \(C^{0, 1 - \frac{n}{p}}(\Omega)\)）。
  \end{itemize}
\item
  \textbf{紧嵌入（Rellich-Kondrachov 定理 8.14 \& 推论 8.3，p.333--335）}：

  \begin{itemize}
  \tightlist
  \item
    当 \(\Omega \subset \mathbb{R}^n\) 是\textbf{有界光滑区域}时，把连续嵌入的临界指数\textbf{严格往下降一点}（\(q < \frac{np}{n-p}\)；或 \(n \le p\) 时任意 \(q \in [1, \infty)\)），嵌入 \(W^{1,p}(\Omega) \hookrightarrow L^q(\Omega)\) 不仅是连续的，还是\textbf{紧算子}！推论 8.3：\(m_1 > m_2 \implies W^{m_1,p}(\Omega) \hookrightarrow W^{m_2,p}(\Omega)\) 紧。
  \item
    \textbf{注记 8.18（p.335）}：区域 \(\Omega\) 的\textbf{有界性绝对不能丢}，在无界区域（如 \(\mathbb{R}^n\)）上平移一列紧支函数即可破坏紧性。
  \end{itemize}
\end{itemize}

\SoftDivider

\BookSubSection{sec:8-8}{六、8.8--8.9 实指数 Sobolev 空间 \(H^s(\mathbb{R}^n)\)、迹定理、\(H^{-m}\) 与椭圆算子（教材 p.335--340）}{六、8.8–8.9 实指数 Sobolev 空间 H\textasciicircum{}s(R\textasciicircum{}n)、迹定理、H\textasciicircum{}-m 与椭圆算子（教材 p.335–340）}{\TagHigh{作业核心}}

\SubSubHeading{1. 用 Plancherel 定理把「求导」翻译到频域，解锁任意实数阶 \(s \in \mathbb{R}\)！}

对于整数 \(k \ge 0\) 与 \(f \in \mathcal{S}(\mathbb{R}^n)\)，由 Plancherel 定理和 \(\widehat{\partial_j f} = 2\pi\mathrm{i}\xi_j\hat{f}\)，有（教材 p.335--336 式 (8.47)--(8.48)）： 
\FitDisplayMath{\|f\|_{W^{k,2}(\mathbb{R}^n)}^2 = \int_{\mathbb{R}^n}\sum_{j=0}^k (2\pi|\xi|)^{2j}|\hat{f}(\xi)|^2\,\mathrm{d}\xi \asymp \int_{\mathbb{R}^n} (1 + |\xi|^2)^k |\hat{f}(\xi)|^2\,\mathrm{d}\xi.}
 引入极其方便的\textbf{日本括号（Japanese bracket）}记号： 
\FitDisplayMath{\langle\xi\rangle \triangleq (1 + |\xi|^2)^{1/2}.}
 于是整数 \(k\) 瞬间解放为\textbf{任意实数 \(s \in \mathbb{R}\)（正数、分数、负数统统适用！）}：

\begin{itemize}
\tightlist
\item
  \textbf{实指数 Sobolev 空间定义（教材 p.336）}：对任意 \(s \in \mathbb{R}\)， 
\FitDisplayMath{H^s(\mathbb{R}^n) \triangleq \big\{f \in \mathcal{S}'(\mathbb{R}^n) : \langle\xi\rangle^s \hat{f}(\xi) \in L^2(\mathbb{R}^n)\big\}, \qquad \|f\|_{H^s(\mathbb{R}^n)} \triangleq \|\langle\xi\rangle^s \hat{f}\|_{L^2(\mathbb{R}^n)}.}


  \begin{itemize}
  \tightlist
  \item
    \textbf{怎么凭直觉理解正负 \(s\)？}

    \begin{itemize}
    \tightlist
    \item
      \(s > 0\)（如 \(H^2\)）：要求 \(\hat{f}(\xi)\) 在高频 \(|\xi|\to\infty\) 衰减得更快 \(\implies\) 物理空间函数 \(f\) 更光滑；
    \item
      \(s = 0\)：\(H^0(\mathbb{R}^n) = L^2(\mathbb{R}^n)\)；
    \item
      \(s < 0\)（如 \(H^{-m}\)）：允许 \(\hat{f}(\xi)\) 在高频增长（只需乘上衰减权重 \(\langle\xi\rangle^{-m}\) 后落回 \(L^2\)）\(\implies\) 容纳了带奇点的分布（如 \(\delta \in H^{-s}(\mathbb{R}^n)\) 当且仅当 \(s > \frac{n}{2}\)！）。
    \end{itemize}
  \item
    \textbf{基本性质（p.336 命题 8.28）}：

    \begin{itemize}
    \tightlist
    \item
      \(s > s' \implies H^s(\mathbb{R}^n) \subset H^{s'}(\mathbb{R}^n)\)；
    \item
      \(\mathcal{S}(\mathbb{R}^n)\) 在每个 \(H^s(\mathbb{R}^n)\) 中稠密；非负整数 \(k\) 时 \(W^{k,2}(\mathbb{R}^n) = H^k(\mathbb{R}^n)\) 且范数等价；
    \item
      \(H^s(\mathbb{R}^n)\) 是复 Hilbert 空间，内积为 \(\langle f, g\rangle_{H^s} = \int_{\mathbb{R}^n}\langle\xi\rangle^{2s}\hat{f}(\xi)\overline{\hat{g}(\xi)}\,\mathrm{d}\xi\)（\emph{注：教材 p.336 内积公式漏印了帽子，应为 \(\hat{f},\hat{g}\)}）。
    \end{itemize}
  \end{itemize}
\item
  \textbf{对偶刻画（命题 8.29，p.336）}：\(H^s(\mathbb{R}^n)\) 上的连续线性泛函由唯一的 \(g \in H^{-s}(\mathbb{R}^n)\) 通过 \(\lambda(f) = \int_{\mathbb{R}^n}\hat{f}(\xi)\overline{\hat{g}(\xi)}\,\mathrm{d}\xi\) 表示，即 \((H^s(\mathbb{R}^n))^* \cong H^{-s}(\mathbb{R}^n)\)，且 
\FitDisplayMath{\|f\|_{H^s(\mathbb{R}^n)} = \sup\big\{|\langle f, g\rangle_{L^2(\mathbb{R}^n)}| : g \in \mathcal{S}(\mathbb{R}^n),\ \|g\|_{H^{-s}(\mathbb{R}^n)} \le 1\big\}.}

\item
  \textbf{\(H^s(\mathbb{R}^n)\) 的嵌入定理（定理 8.15--8.16，p.336--337）}：

  \begin{itemize}
  \tightlist
  \item
    \textbf{定理 8.15}：若 \(0 < s < \frac{n}{2}\) 且 \(\frac{n}{2} - s \le \frac{n}{q} \le \frac{n}{2}\)，则 \(H^s(\mathbb{R}^n) \hookrightarrow L^q(\mathbb{R}^n)\) 连续嵌入；
  \item
    \textbf{定理 8.16}：若 \(s > \frac{n}{2}\)，则对 \(0 < \alpha \le \min\!\left(s - \frac{n}{2}, 1\right)\)，\(H^s(\mathbb{R}^n) \hookrightarrow C^{0,\alpha}(\mathbb{R}^n)\) 连续嵌入（口诀：\textbf{「光滑度过半维数 \(n/2\)，立刻连续！」}）。
  \end{itemize}
\end{itemize}

\SubSubHeading{2. 常系数 \(m\) 阶椭圆算子（教材 p.337）}

\begin{itemize}
\tightlist
\item
  设 \(m \ge 0\)，\(L = \sum_{|\alpha|=m} c_\alpha \partial^\alpha\) 是 \(m\) 阶常系数齐次微分算子，定义其\textbf{象征（symbol）}为齐次多项式 \(l(\xi) \triangleq \sum_{|\alpha|=m} c_\alpha \xi^\alpha\)。
\item
  若存在常数 \(c > 0\)，使得对所有 \(\xi \in \mathbb{R}^n\) 均有 
\FitDisplayMath{|l(\xi)| \ge c|\xi|^m,}
 则称 \(L\) 为 \textbf{\(m\) 阶椭圆算子}。
\item
  \textbf{例 8.17（p.337）}：

  \begin{itemize}
  \tightlist
  \item
    \textbf{是椭圆算子}：Laplace 算子 \(\Delta = \sum_{j=1}^n \partial_j^2\)（\(2\) 阶，象征 \(l(\xi)=|\xi|^2\)）；\(\mathbb{R}^2\) 上的 Cauchy-Riemann 算子 \(\partial_{x_1} - \mathrm{i}\partial_{x_2}\)（\(1\) 阶，象征 \(l(\xi)=\xi_1-\mathrm{i}\xi_2\)，模长 \(|l(\xi)|=|\xi|\)）。
  \item
    \textbf{非椭圆算子}：热算子 \(\partial_t - \Delta\)、Schrödinger 算子 \(\mathrm{i}\partial_t + \Delta\)、波算子 \(-\partial_t^2 + \Delta\)（在 \(\mathbb{R}^{1+n}\) 上象征会在非零频率点处消失或阶数不齐）。
  \end{itemize}
\end{itemize}

\SubSubHeading{3. 8.9 节迹定理、\(H^{-m}(\Omega)\) 与 Poincaré 不等式（教材 p.338--340）}

\begin{itemize}
\tightlist
\item
  \textbf{迹定理（定理 8.17--8.19，p.338--339）}：\(L^2\) 函数在边界 \(\partial\Omega\)（零测集）上本来不能谈取值，但 \(H^1(\Omega)\) 的函数多了一阶弱导数后，限制映射 \(\gamma_0(\varphi) = \varphi|_{\partial\Omega}\) 可唯一延拓为有界线性算子 \(\gamma_0: H^1(\Omega) \to L^2(\partial\Omega)\)，从而严格赋予了 PDE 的边界值意义！将直到 \(m-1\) 阶法向迹全为 \(0\) 的子空间记为 \(H_0^m(\Omega)\)（\textbf{命题 8.30}：\(C_0^\infty(\Omega)\) 在 \(H_0^m(\Omega)\) 中稠密）。
\item
  \textbf{定理 8.20（\(H^{-m}(\Omega)\) 的结构，p.339）}：\(H^{-m}(\Omega) \triangleq (H_0^m(\Omega))'\) 恰好等于 \(\{\partial^\alpha f : |\alpha| \le m,\ f \in L^2(\Omega)\}\) 张成的广义函数子空间。
\item
  \textbf{定理 8.21（Poincaré 不等式，p.340）}：若 \(\Omega \subset \mathbb{R}^n\) 沿 \(x_1\) 方向有界（\(\sup_{x\in\Omega}|x_1| = K < \infty\)），则 \(\|\varphi\|_{L^2(\Omega)} \le 2K\|\partial_1\varphi\|_{L^2(\Omega)}\)（\(\forall \varphi \in H_0^1(\Omega)\)）。
\end{itemize}

\BookPart{}{第二部分：第 8 章全部 8 道作业逐题详解}{第二部分：第 8 章全部 8 道作业逐题详解}

\begin{quote}
\textbf{题目索引}：以下按顺序完整解答 \textbf{习题 8 第 7、9、11、14、17(2)(3)、20、28 题}（教材 p.341--343）以及 \textbf{教材 p.337 \hyperlink{prob:8-882}{\textcolor{LinkColor}{练习 8.8.2}}}，共 8 题。
\end{quote}

\SoftDivider

\BookProblem{prob:8-7}{题目 1：习题 8 第 7 题}{教材 p.341}{题目 1：习题 8 第 7 题}

\SubPillPrimary{【原题】}

设 \(f \in L^1(\mathbb{R})\)，能否定义 \(f\) 与 \(\delta\) 的乘积？

\SubPillAmber{【零基础破题思路】}

回想光滑函数 \(h \in C^\infty(\mathbb{R})\) 乘以 \(\delta\) 发生了什么： 
\FitDisplayMath{(h\delta)(\varphi) = \delta(\bar{h}\varphi) = \overline{h(0)}\varphi(0) = \big(\overline{h(0)}\delta\big)(\varphi)}
 （实值函数就是 \(h(0)\delta\)）。本质上，\textbf{乘以 \(\delta\) 会把 \(h\) 在单点 \(x = 0\) 处的精确取值``拎出来''}！

现在把光滑函数 \(h\) 换成 \(L^1(\mathbb{R})\) 里的函数 \(f\)，立刻撞上三堵死墙：

\begin{enumerate}
\def\labelenumi{\arabic{enumi}.}
\tightlist
\item
  \textbf{\(L^1\) 里没有``单点取值''}：\(L^1(\mathbb{R})\) 的元素是几乎处处相等的函数等价类，单点 \(\{0\}\) 测度为 \(0\)，随便改 \(f(0)\) 都不改变它在 \(L^1\) 中的身份；
\item
  \textbf{奇点可能发散}：\(L^1(\mathbb{R})\) 里有像 \(|x|^{-1/2}\chi_{(-1,1)}(x)\) 这样在原点趋于 \(+\infty\) 的函数，根本没法给 \(f(0)\) 赋有限值；
\item
  \textbf{哪怕想``先用光滑函数列在 \(L^1\) 里逼近 \(f\)，再对乘积取 \(\mathcal{D}'\) 极限''也不行}：我们可以构造出\textbf{都以 \(0\) 为 \(L^1\) 极限}的两列光滑函数（如 \(0\) 和越来越窄但顶点永远高度为 \(1\) 的高斯尖峰 \(\mathrm{e}^{-k^2 x^2}\)），它们乘上 \(\delta\) 后一个永远是 \(0\)，另一个永远是 \(\delta\)！极限不唯一，甚至还能构造出直接发散的序列！
\end{enumerate}

\SubPillMain{【规范解答】}

\textbf{结论：不能}（无法在 \(L^1(\mathbb{R})\) 上对所有 \(f\) 定义与光滑函数乘法相容且满足连续性的乘积 \(f\delta \in \mathcal{D}'(\mathbb{R})\)）。下面从三个层面严密说明：

\begin{enumerate}
\def\labelenumi{\arabic{enumi}.}
\tightlist
\item
  \textbf{直接沿用点值公式不良定义（等价类与无界奇点障碍）} 按教材 p.305 光滑函数与广义函数的乘积定义，对任意 \(h \in C^\infty(\mathbb{R})\) 及试验函数 \(\varphi \in \mathcal{D}(\mathbb{R})\)，有 
\FitDisplayMath{(h\delta)(\varphi) = \delta(\bar{h}\varphi) = \overline{h(0)}\varphi(0) = \big(\overline{h(0)}\delta\big)(\varphi).}
 若试图将该公式推广到 \(f \in L^1(\mathbb{R})\)，存在两个本质矛盾：

  \begin{itemize}
  \tightlist
  \item
    \textbf{等价类歧义}：\(L^1(\mathbb{R})\) 中的元素是几乎处处相等的可测函数等价类，单点集 \(\{0\}\) 是零测集。例如函数 \(f_1(x) \equiv 0\) 与 \(f_2(x) = \begin{cases} 1, & x = 0, \\ 0, & x \neq 0 \end{cases}\) 在 \(L^1(\mathbb{R})\) 中是\textbf{同一个元素}，但它们在 \(x = 0\) 处的值分别为 \(0\) 与 \(1\)，无法唯一确定 \(f(0)\)。
  \item
    \textbf{原点无有限值}：存在函数 \(f(x) = |x|^{-1/2}\chi_{(-1,1)}(x)\)，因 \(\int_{-1}^{1} |x|^{-1/2}\,\mathrm{d}x = 4 < \infty\) 故 \(f \in L^1(\mathbb{R})\)，但 \(\lim_{x\to 0} f(x) = +\infty\)，在原点处无法赋予任何有限值。
  \end{itemize}
\item
  \textbf{通过 \(L^1(\mathbb{R})\) 稠密逼近作连续延拓也不可能} 由于 \(C^\infty(\mathbb{R}) \cap L^1(\mathbb{R})\)（乃至 \(C_c^\infty(\mathbb{R})\)）在 \(L^1(\mathbb{R})\) 中稠密，考虑能否对 \(f \in L^1(\mathbb{R})\) 取光滑函数列 \(f_k \to f\)（于 \(L^1(\mathbb{R})\)），定义 \(f\delta \triangleq \lim_{k\to\infty} f_k\delta\)（于 \(\mathcal{D}'(\mathbb{R})\)）。以下反例证明该极限\textbf{既不唯一，也不一定存在}：

  \begin{itemize}
  \tightlist
  \item
    取 \(g_k(x) \equiv 0 \in C^\infty(\mathbb{R}) \cap L^1(\mathbb{R})\)，显然 \(g_k \to 0\) 于 \(L^1(\mathbb{R})\)，且对所有 \(k\) 有 \(g_k\delta = 0 \to 0\)（于 \(\mathcal{D}'(\mathbb{R})\)）。
  \item
    取 \(h_k(x) = \mathrm{e}^{-k^2 x^2} \in C^\infty(\mathbb{R}) \cap L^1(\mathbb{R})\)，计算其 \(L^1\) 范数： 
\FitDisplayMath{\|h_k - 0\|_{L^1(\mathbb{R})} = \int_{-\infty}^{\infty}\mathrm{e}^{-k^2 x^2}\,\mathrm{d}x = \frac{\sqrt{\pi}}{k} \xrightarrow{k\to\infty} 0,}
 即 \(h_k \to 0\) 于 \(L^1(\mathbb{R})\)。然而对任意 \(\varphi \in \mathcal{D}(\mathbb{R})\)， 
\FitDisplayMath{(h_k\delta)(\varphi) = \overline{h_k(0)}\varphi(0) = 1 \cdot \varphi(0) = \delta(\varphi) \implies \lim_{k\to\infty} h_k\delta = \delta \neq 0.}

  \item
    进一步取 \(u_k(x) = \sqrt{k}\,\mathrm{e}^{-k^2 x^2} \in C^\infty(\mathbb{R}) \cap L^1(\mathbb{R})\)，则 \(\|u_k\|_{L^1(\mathbb{R})} = \sqrt{\frac{\pi}{k}} \to 0\)，但对满足 \(\varphi(0) \neq 0\) 的试验函数，\((u_k\delta)(\varphi) = \sqrt{k}\varphi(0)\) 当 \(k \to \infty\) 时发散！
  \end{itemize}
\end{enumerate}

综上所述，映射 \(h \mapsto h\delta\) 不能从光滑函数空间合理延拓到 \(L^1(\mathbb{R})\)，因此\textbf{不能定义}一般 \(f \in L^1(\mathbb{R})\) 与 \(\delta\) 的乘积。

\SubPillCoral{【考点复盘与避坑提示】}

\begin{itemize}
\tightlist
\item
  \textbf{核心考点}：理解广义函数乘法的正则性门槛（教材 p.305、p.311--312 注记 8.7）以及 \(L^1\) 拓扑不控制单点值的本质。
\item
  \textbf{延伸思考}：如果把条件加强为 \(f \in C(\mathbb{R})\)（连续函数），那么\textbf{可以}定义 \(f\delta \triangleq \overline{f(0)}\delta\)！因为 \(\delta\) 是 \(0\) 阶分布（Radon 测度，命题 8.10），只需要乘子在原点连续（甚至只要求在原点逐点连续）就够了，不需要 \(C^\infty\)。但若是 \(f\delta'\)，则至少需要 \(f \in C^1(\mathbb{R})\)。
\end{itemize}

\SoftDivider

\BookProblem{prob:8-9}{题目 2：习题 8 第 9 题}{教材 p.341}{题目 2：习题 8 第 9 题}

\SubPillPrimary{【原题】}

在广义函数意义下求下列极限：

\begin{enumerate}
\def\labelenumi{(\arabic{enumi})}
\tightlist
\item
  \(\displaystyle\lim_{\varepsilon\to 0}\frac{\varepsilon}{(\varepsilon^2 + x^2)}\)；
\item
  \(\displaystyle\lim_{\varepsilon\to 0}\frac{1}{\sqrt{\pi\varepsilon}}\mathrm{e}^{-\frac{x^2}{\varepsilon}}\)。
\end{enumerate}

\SubPillAmber{【零基础破题思路】}

按定义 8.8，求一列常规函数 \(f_\varepsilon\) 在 \(\mathcal{D}'(\mathbb{R})\) 中的极限，\textbf{固定模板只有三步}：

\begin{enumerate}
\def\labelenumi{\arabic{enumi}.}
\tightlist
\item
  任取一个试验函数 \(\varphi \in \mathcal{D}(\mathbb{R})\)，写出积分 \(\int_{-\infty}^{\infty} f_\varepsilon(x)\varphi(x)\,\mathrm{d}x\)；
\item
  \textbf{变量替换``把 \(\varepsilon\) 从分母/指数里赶到 \(\varphi\) 的括号里''}：第 (1) 题令 \(x = \varepsilon y\)，第 (2) 题令 \(x = \sqrt{\varepsilon}\,y\)；
\item
  \textbf{用 Lebesgue 控制收敛定理把 \(\lim_{\varepsilon\to 0^+}\) 放进积分号里}（因为 \(\varphi\) 紧支连续有界，\(|\varphi(\cdot)| \le M\)，外层剩下的 \(\frac{1}{1+y^2}\) 和 \(\mathrm{e}^{-y^2}\) 都是经典可积函数），直接算出常数倍的 \(\varphi(0)\)！
\end{enumerate}

\begin{quote}
\textbf{特别提醒}：第 (1) 题的表达式关于 \(\varepsilon\) 是\textbf{奇函数}，当 \(\varepsilon \to 0^+\) 时极限是 \(\pi\delta\)，而当 \(\varepsilon \to 0^-\) 时极限是 \(-\pi\delta\)！作业与教材通常默认参数 \(\varepsilon > 0\)（即 \(\varepsilon \to 0^+\)），我们在解答中先给出 \(\varepsilon \to 0^+\) 的完整证明，再顺手补一句 \(\varepsilon \to 0^-\) 的结论，滴水不漏！
\end{quote}

\SubPillMain{【规范解答】}

任取试验函数 \(\varphi \in \mathcal{D}(\mathbb{R}) = C_c^\infty(\mathbb{R})\)。因为 \(\varphi\) 具有紧支集且连续，故有界，记 \(M \triangleq \sup_{x\in\mathbb{R}} |\varphi(x)| < \infty\)。两题中的函数均为实值光滑可积函数，按例 8.2 均定义 \(\mathcal{D}'(\mathbb{R})\) 中的常规广义函数。

\textbf{(1) 求解 \(\displaystyle\lim_{\varepsilon\to 0}\frac{\varepsilon}{\varepsilon^2 + x^2}\)}

\begin{itemize}
\item
  \textbf{情形一：\(\varepsilon \to 0^+\)（默认情形）} 对任意 \(\varepsilon > 0\)，作变量替换 \(x = \varepsilon y\)（则 \(\mathrm{d}x = \varepsilon\,\mathrm{d}y\)，积分限仍为 \((-\infty, +\infty)\)），得 
\FitDisplayMath{\int_{-\infty}^{\infty}\frac{\varepsilon}{\varepsilon^2 + x^2}\varphi(x)\,\mathrm{d}x = \int_{-\infty}^{\infty}\frac{\varepsilon}{\varepsilon^2(1 + y^2)}\varphi(\varepsilon y)\cdot\varepsilon\,\mathrm{d}y = \int_{-\infty}^{\infty}\frac{\varphi(\varepsilon y)}{1 + y^2}\,\mathrm{d}y.}
 验证 Lebesgue 控制收敛定理的条件：

  \begin{enumerate}
  \def\labelenumi{\arabic{enumi}.}
  \tightlist
  \item
    对每个固定的 \(y \in \mathbb{R}\)，由 \(\varphi\) 的连续性，有 \(\lim_{\varepsilon\to 0^+} \frac{\varphi(\varepsilon y)}{1 + y^2} = \frac{\varphi(0)}{1 + y^2}\)；
  \item
    对所有 \(\varepsilon > 0\) 及 \(y \in \mathbb{R}\)，有统一控制估计 
\FitDisplayMath{\left|\frac{\varphi(\varepsilon y)}{1 + y^2}\right| \le \frac{M}{1 + y^2},}
 且控制函数可积：\(\int_{-\infty}^{\infty}\frac{M}{1 + y^2}\,\mathrm{d}y = M\big[\arctan y\big]_{-\infty}^{+\infty} = \pi M < \infty\)。
  \end{enumerate}

  由 Lebesgue 控制收敛定理，得 
\FitDisplayMath{\lim_{\varepsilon\to 0^+}\int_{-\infty}^{\infty}\frac{\varepsilon}{\varepsilon^2 + x^2}\varphi(x)\,\mathrm{d}x = \int_{-\infty}^{\infty}\frac{\varphi(0)}{1 + y^2}\,\mathrm{d}y = \pi\varphi(0) = (\pi\delta)(\varphi).}
 故当 \(\varepsilon \to 0^+\) 时，在 \(\mathcal{D}'(\mathbb{R})\) 意义下， 
\FitDisplayMath{\lim_{\varepsilon\to 0^+}\frac{\varepsilon}{\varepsilon^2 + x^2} = \pi\delta.}

\item
  \textbf{情形二：\(\varepsilon \to 0^-\)（单侧极限补充）} 令 \(\eta = -\varepsilon > 0\)，则 \(\frac{\varepsilon}{\varepsilon^2 + x^2} = -\frac{\eta}{\eta^2 + x^2}\)。由情形一立刻得到，当 \(\varepsilon \to 0^-\) 时在 \(\mathcal{D}'(\mathbb{R})\) 中极限为 \(-\pi\delta\)。 （因此严格的双侧极限 \(\varepsilon \to 0\) 不存在；按教材默认 \(\varepsilon > 0\) 理解，答案为 \(\pi\delta\)。）
\end{itemize}

\SoftDivider

\textbf{(2) 求解 \(\displaystyle\lim_{\varepsilon\to 0}\frac{1}{\sqrt{\pi\varepsilon}}\mathrm{e}^{-\frac{x^2}{\varepsilon}}\)}

由表达式中的 \(\sqrt{\pi\varepsilon}\) 可知必有 \(\varepsilon > 0\)，即考虑 \(\varepsilon \to 0^+\)。作变量替换 \(x = \sqrt{\varepsilon}\,y\)（则 \(\mathrm{d}x = \sqrt{\varepsilon}\,\mathrm{d}y\)，积分限仍为 \((-\infty, +\infty)\)），得 
\FitDisplayMath{\int_{-\infty}^{\infty}\frac{1}{\sqrt{\pi\varepsilon}}\mathrm{e}^{-\frac{x^2}{\varepsilon}}\varphi(x)\,\mathrm{d}x = \frac{1}{\sqrt{\pi}}\int_{-\infty}^{\infty}\mathrm{e}^{-y^2}\varphi(\sqrt{\varepsilon}\,y)\,\mathrm{d}y.}
 验证 Lebesgue 控制收敛定理的条件：

\begin{enumerate}
\def\labelenumi{\arabic{enumi}.}
\tightlist
\item
  对每个固定的 \(y \in \mathbb{R}\)，由 \(\varphi\) 的连续性，有 \(\lim_{\varepsilon\to 0^+}\mathrm{e}^{-y^2}\varphi(\sqrt{\varepsilon}\,y) = \mathrm{e}^{-y^2}\varphi(0)\)；
\item
  对所有 \(\varepsilon > 0\) 及 \(y \in \mathbb{R}\)，有统一控制估计 
\FitDisplayMath{\left|\mathrm{e}^{-y^2}\varphi(\sqrt{\varepsilon}\,y)\right| \le M\mathrm{e}^{-y^2},}
 由熟知的 Gauss 积分 \(\int_{-\infty}^{\infty}\mathrm{e}^{-y^2}\,\mathrm{d}y = \sqrt{\pi} < \infty\)，控制函数属于 \(L^1(\mathbb{R})\)。
\end{enumerate}

由 Lebesgue 控制收敛定理，得 
\FitDisplayMath{\lim_{\varepsilon\to 0^+}\int_{-\infty}^{\infty}\frac{1}{\sqrt{\pi\varepsilon}}\mathrm{e}^{-\frac{x^2}{\varepsilon}}\varphi(x)\,\mathrm{d}x = \frac{\varphi(0)}{\sqrt{\pi}}\int_{-\infty}^{\infty}\mathrm{e}^{-y^2}\,\mathrm{d}y = \varphi(0) = \delta(\varphi).}
 故在 \(\mathcal{D}'(\mathbb{R})\) 意义下， 
\FitDisplayMath{\lim_{\varepsilon\to 0^+}\frac{1}{\sqrt{\pi\varepsilon}}\mathrm{e}^{-\frac{x^2}{\varepsilon}} = \delta.}


\SubPillCoral{【考点复盘与避坑提示】}

\begin{itemize}
\tightlist
\item
  \textbf{核心考点}：近似恒等元（Approximate Identity / \(\delta\) 序列）与 Lebesgue 控制收敛定理的规范用法。其实这两题藏着 PDE 最著名的两大基本核：

  \begin{itemize}
  \tightlist
  \item
    第 (1) 题除以 \(\pi\) 就是\textbf{上半平面 Poisson 核} \(P_\varepsilon(x) = \frac{1}{\pi}\frac{\varepsilon}{\varepsilon^2+x^2}\)（直接预告\textbf{第 20 题}！）；
  \item
    第 (2) 题令 \(\varepsilon = 4t\) 就是\textbf{一维热核} \(K_t(x) = \frac{1}{\sqrt{4\pi t}}\mathrm{e}^{-x^2/(4t)}\)（教材 p.319 式 (8.23)）。
  \end{itemize}
\item
  \textbf{避坑提示}：千万不要只写``把 \(\lim\) 放进积分号''而不写\textbf{可积控制函数}（\(M/(1+y^2)\) 和 \(M\mathrm{e}^{-y^2}\)）；也别忘了检查第 (1) 题 \(\varepsilon\) 正负号导致积分换元是否反转上下限！
\end{itemize}

\SoftDivider

\BookProblem{prob:8-11}{题目 3：习题 8 第 11 题}{教材 p.341}{题目 3：习题 8 第 11 题}

\SubPillPrimary{【原题】}

设 \(n = 1\)。证明下列结论：

\begin{enumerate}
\def\labelenumi{(\arabic{enumi})}
\tightlist
\item
  符号函数 \(\operatorname{sgn}(x)\) 的导数等于 \(2\delta\)。
\item
  局部可积函数 \(\ln|x|\) 的导数等于 \(\operatorname{p.v.}\dfrac{1}{x}\)。
\item
  局部可积函数 \(\ln|x|\operatorname{sgn}(x)\) 的导数等于例 8.6 中的广义函数 \(\lambda_1\)。
\item
  \(|x|\) 的导数是 \(\operatorname{sgn}(x)\)。
\end{enumerate}

\SubPillAmber{【零基础破题思路】}

四道小题全部围绕同一个口诀展开——\textbf{「求导甩给 \(\varphi\) 加负号，挖掉奇点分部积，边界消零看中值！」}

\begin{itemize}
\tightlist
\item
  第 (1)、(4) 题没有发散奇点，只在 \(x=0\) 处有跳变或折点：直接把 \((-\infty, \infty)\) 拆成 \((-\infty, 0)\) 和 \((0, \infty)\) 两段，逐段分部积分即可。
\item
  第 (2) 题 \(\ln|x|\) 在原点发散到 \(-\infty\)（虽然它是局部可积的！）：先用控制收敛定理挖掉 \((-\varepsilon, \varepsilon)\)，在 \((-\infty, -\varepsilon)\) 和 \((\varepsilon, \infty)\) 上分别分部积分。神奇的事发生了：边界项凑成了 \(\ln\varepsilon(\varphi(\varepsilon)-\varphi(-\varepsilon))\)，靠 \(\varphi\) 的光滑性（\(|\varphi(\varepsilon)-\varphi(-\varepsilon)| = O(\varepsilon)\)）压死对数奇性（\(\varepsilon\ln\varepsilon \to 0\)）！
\item
  第 (3) 题是\textbf{全题最精彩的一小问}！先抄出教材 p.305 例 8.6（取 \(r=1\)）： 
\FitDisplayMath{\lambda_1(\varphi) = \int_{|x|<1}\frac{\varphi(x)-\varphi(0)}{|x|}\,\mathrm{d}x + \int_{|x|\ge 1}\frac{\varphi(x)}{|x|}\,\mathrm{d}x.}
 观察到 \(\lambda_1\) 是在 \(|x|=1\) 处分段的，而且在内段 \(|x|<1\) 里藏着 \(\varphi(x)-\varphi(0)\)！这就直接告诉了我们\textbf{必胜绝招}：

  \begin{itemize}
  \tightlist
  \item
    把积分拆成四段 \((-R, -1), (-1, -\varepsilon), (\varepsilon, 1), (1, R)\)；
  \item
    在内侧两段 \((-1, -\varepsilon)\) 和 \((\varepsilon, 1)\) 上，把 \(\varphi'(x)\) 改写成 \([\varphi(x)-\varphi(0)]'\) 再分部积分——这样原点端 \(\pm\varepsilon\) 处出现 \(\ln\varepsilon(\varphi(\pm\varepsilon)-\varphi(0)) = O(\varepsilon\ln\varepsilon) \to 0\)，而接缝点 \(\pm 1\) 处恰好因为 \(\ln 1 = 0\)，\textbf{所有边界项全部清零}！
  \end{itemize}
\end{itemize}

\SubPillMain{【规范解答】}

本题中四个函数均为实值局部可积函数（对 (2)(3)，因 \(\int_0^R |\ln x|\,\mathrm{d}x = [x\ln x - x]_0^R < \infty\)，故 \(\ln|x|\) 与 \(\ln|x|\operatorname{sgn}(x)\) 均属于 \(L^1_{\mathrm{loc}}(\mathbb{R})\)），由例 8.2 知其诱导广义函数满足 \(\lambda_g(\varphi) = \int_{-\infty}^{\infty} g(x)\varphi(x)\,\mathrm{d}x\)。

以下统一任取 \(\varphi \in \mathcal{D}(\mathbb{R})\)，设 \(\operatorname{supp}\varphi \subset [-R, R]\)（取足够大的 \(R > 1\)），从而当 \(|x| \ge R\) 时 \(\varphi(x) = 0\)。

\SoftDivider

\textbf{(1) 证明 \((\operatorname{sgn} x)' = 2\delta\)}

符号函数定义为 \(\operatorname{sgn}(x) = \begin{cases} 1, & x > 0, \\ 0, & x = 0, \\ -1, & x < 0. \end{cases}\) 按广义导数定义与 Newton-Leibniz 公式，有 
\FitDisplayMath{\begin{aligned}
(\operatorname{sgn})'(\varphi) &= -\int_{-\infty}^{\infty}\operatorname{sgn}(x)\varphi'(x)\,\mathrm{d}x \\
&= -\left(-\int_{-R}^{0}\varphi'(x)\,\mathrm{d}x + \int_{0}^{R}\varphi'(x)\,\mathrm{d}x\right) \\
&= \big[\varphi(0) - \varphi(-R)\big] - \big[\varphi(R) - \varphi(0)\big] \\
&= 2\varphi(0) = (2\delta)(\varphi).
\end{aligned}}
 由 \(\varphi \in \mathcal{D}(\mathbb{R})\) 的任意性，得 \((\operatorname{sgn} x)' = 2\delta\)。

\SoftDivider

\textbf{(2) 证明 \((\ln|x|)' = \operatorname{p.v.}\dfrac{1}{x}\)}

因为 \(\ln|x| \in L^1_{\mathrm{loc}}(\mathbb{R})\) 且 \(\varphi' \in C_c^\infty(\mathbb{R})\) 有界紧支，故 \(\ln|x|\varphi'(x) \in L^1(\mathbb{R})\)。利用指示函数 \(\chi_{|x|>\varepsilon}(x) \to 1\)（a.e.）及控制函数 \(|\ln|x|\varphi'(x)| \in L^1(\mathbb{R})\)，由 Lebesgue 控制收敛定理得 
\FitDisplayMath{(\ln|x|)'(\varphi) = -\int_{-R}^{R}\ln|x|\,\varphi'(x)\,\mathrm{d}x = -\lim_{\varepsilon\to 0^+}\left(\int_{-R}^{-\varepsilon}\ln(-x)\,\varphi'(x)\,\mathrm{d}x + \int_{\varepsilon}^{R}\ln x\,\varphi'(x)\,\mathrm{d}x\right).}
 注意到当 \(x < 0\) 时 \(\frac{\mathrm{d}}{\mathrm{d}x}\ln(-x) = \frac{-1}{-x} = \frac{1}{x}\)。对两段分别分部积分（利用 \(\varphi(\pm R) = 0\)）： 
\FitDisplayMath{\begin{aligned}
\int_{-R}^{-\varepsilon}\ln(-x)\,\varphi'(x)\,\mathrm{d}x &= \Big[\ln(-x)\varphi(x)\Big]_{-R}^{-\varepsilon} - \int_{-R}^{-\varepsilon}\frac{\varphi(x)}{x}\,\mathrm{d}x = \ln\varepsilon\,\varphi(-\varepsilon) - \int_{-R}^{-\varepsilon}\frac{\varphi(x)}{x}\,\mathrm{d}x, \\
\int_{\varepsilon}^{R}\ln x\,\varphi'(x)\,\mathrm{d}x &= \Big[\ln x\,\varphi(x)\Big]_{\varepsilon}^{R} - \int_{\varepsilon}^{R}\frac{\varphi(x)}{x}\,\mathrm{d}x = -\ln\varepsilon\,\varphi(\varepsilon) - \int_{\varepsilon}^{R}\frac{\varphi(x)}{x}\,\mathrm{d}x.
\end{aligned}}
 两式相加并取负号，得 
\FitDisplayMath{-\left(\int_{-R}^{-\varepsilon} + \int_{\varepsilon}^{R}\right)\ln|x|\,\varphi'(x)\,\mathrm{d}x = \ln\varepsilon\big(\varphi(\varepsilon) - \varphi(-\varepsilon)\big) + \int_{|x|>\varepsilon}\frac{\varphi(x)}{x}\,\mathrm{d}x.}
 由 Lagrange 中值定理，\(|\varphi(\varepsilon) - \varphi(-\varepsilon)| \le 2\varepsilon\|\varphi'\|_\infty\)，从而 
\FitDisplayMath{\big|\ln\varepsilon\,\big(\varphi(\varepsilon) - \varphi(-\varepsilon)\big)\big| \le 2\|\varphi'\|_\infty \cdot \varepsilon|\ln\varepsilon| \xrightarrow{\varepsilon\to 0^+} 0.}
 令 \(\varepsilon \to 0^+\)，由例 8.5 得 
\FitDisplayMath{(\ln|x|)'(\varphi) = \lim_{\varepsilon\to 0^+}\int_{|x|>\varepsilon}\frac{\varphi(x)}{x}\,\mathrm{d}x = \left(\operatorname{p.v.}\frac{1}{x}\right)(\varphi).}
 故 \((\ln|x|)' = \operatorname{p.v.}\dfrac{1}{x}\)。

\SoftDivider

\textbf{(3) 证明 \((\ln|x|\operatorname{sgn}(x))' = \lambda_1\)}

回忆教材 p.305 例 8.6（取 \(r = 1\)）： 
\FitDisplayMath{\lambda_1(\varphi) = \int_{|x|<1}\frac{\varphi(x) - \varphi(0)}{|x|}\,\mathrm{d}x + \int_{|x|\ge 1}\frac{\varphi(x)}{|x|}\,\mathrm{d}x.}
 记 \(g(x) = \ln|x|\operatorname{sgn}(x)\)，则当 \(x > 0\) 时 \(g(x) = \ln x\)，当 \(x < 0\) 时 \(g(x) = -\ln(-x)\)；特别在 \(x \neq 0\) 处恒有 
\FitDisplayMath{g'(x) = \frac{1}{|x|}, \qquad \text{且 } g(\pm 1) = \pm\ln 1 = 0.}
 同 (2)，因 \(g(x)\varphi'(x) \in L^1(\mathbb{R})\)，由控制收敛定理挖去 \((-\varepsilon, \varepsilon)\)（\(0 < \varepsilon < 1 < R\)），并拆为四段积分： 
\FitDisplayMath{g'(\varphi) = -\lim_{\varepsilon\to 0^+}\big(I_1 + I_2 + I_3 + I_4\big),}
 其中四段分别作如下分部积分（内侧两段 \(I_2, I_3\) 改写 \(\varphi'(x) = \big(\varphi(x)-\varphi(0)\big)'\)，且所有端点利用 \(g(\pm 1)=0\) 及 \(\varphi(\pm R)=0\)）：

\begin{enumerate}
\def\labelenumi{\arabic{enumi}.}
\tightlist
\item
  \(I_1 = \displaystyle\int_{-R}^{-1} \big(-\ln(-x)\big)\varphi'(x)\,\mathrm{d}x = \Big[-\ln(-x)\varphi(x)\Big]_{-R}^{-1} - \int_{-R}^{-1}\frac{\varphi(x)}{|x|}\,\mathrm{d}x = -\int_{-R}^{-1}\frac{\varphi(x)}{|x|}\,\mathrm{d}x\)；
\item
  \(I_2 = \displaystyle\int_{-1}^{-\varepsilon} \big(-\ln(-x)\big)\big(\varphi(x)-\varphi(0)\big)'\,\mathrm{d}x = \Big[-\ln(-x)\big(\varphi(x)-\varphi(0)\big)\Big]_{-1}^{-\varepsilon} - \int_{-1}^{-\varepsilon}\frac{\varphi(x)-\varphi(0)}{|x|}\,\mathrm{d}x\) \(= -\ln\varepsilon\big(\varphi(-\varepsilon)-\varphi(0)\big) - \displaystyle\int_{-1}^{-\varepsilon}\frac{\varphi(x)-\varphi(0)}{|x|}\,\mathrm{d}x\)；
\item
  \(I_3 = \displaystyle\int_{\varepsilon}^{1} \ln x\,\big(\varphi(x)-\varphi(0)\big)'\,\mathrm{d}x = \Big[\ln x\,\big(\varphi(x)-\varphi(0)\big)\Big]_{\varepsilon}^{1} - \int_{\varepsilon}^{1}\frac{\varphi(x)-\varphi(0)}{|x|}\,\mathrm{d}x\) \(= -\ln\varepsilon\big(\varphi(\varepsilon)-\varphi(0)\big) - \displaystyle\int_{\varepsilon}^{1}\frac{\varphi(x)-\varphi(0)}{|x|}\,\mathrm{d}x\)；
\item
  \(I_4 = \displaystyle\int_{1}^{R} \ln x\,\varphi'(x)\,\mathrm{d}x = \Big[\ln x\,\varphi(x)\Big]_{1}^{R} - \int_{1}^{R}\frac{\varphi(x)}{|x|}\,\mathrm{d}x = -\int_{1}^{R}\frac{\varphi(x)}{|x|}\,\mathrm{d}x\)。
\end{enumerate}

将四式相加并取负号，得 
\FitDisplayMath{-(I_1 + I_2 + I_3 + I_4) = \ln\varepsilon\big(\varphi(-\varepsilon)-\varphi(0)\big) + \ln\varepsilon\big(\varphi(\varepsilon)-\varphi(0)\big) + \int_{\varepsilon < |x| < 1}\frac{\varphi(x)-\varphi(0)}{|x|}\,\mathrm{d}x + \int_{|x|\ge 1}\frac{\varphi(x)}{|x|}\,\mathrm{d}x.}
 由中值定理，\(|\varphi(\pm\varepsilon) - \varphi(0)| \le \varepsilon\|\varphi'\|_\infty\)，故边界项之和满足 
\FitDisplayMath{\Big|\ln\varepsilon\big(\varphi(-\varepsilon)-\varphi(0)\big) + \ln\varepsilon\big(\varphi(\varepsilon)-\varphi(0)\big)\Big| \le 2\|\varphi'\|_\infty \cdot \varepsilon|\ln\varepsilon| \xrightarrow{\varepsilon\to 0^+} 0.}
 又因在 \((-1, 1)\setminus\{0\}\) 上 \(\left|\frac{\varphi(x)-\varphi(0)}{|x|}\right| \le \|\varphi'\|_\infty \in L^1((-1,1))\)，令 \(\varepsilon \to 0^+\) 并由控制收敛定理得 
\FitDisplayMath{g'(\varphi) = \int_{|x|<1}\frac{\varphi(x)-\varphi(0)}{|x|}\,\mathrm{d}x + \int_{|x|\ge 1}\frac{\varphi(x)}{|x|}\,\mathrm{d}x = \lambda_1(\varphi).}
 故 \((\ln|x|\operatorname{sgn}(x))' = \lambda_1\)。

\SoftDivider

\textbf{(4) 证明 \((|x|)' = \operatorname{sgn}(x)\)}

因 \(|x|\) 连续（从而局部可积），拆分为 \([-R, 0]\) 与 \([0, R]\) 两段分部积分（注意 \(\varphi(\pm R)=0\) 且原点处 \(0\cdot\varphi(0)=0\)）： 
\FitDisplayMath{\begin{aligned}
(|x|)'(\varphi) &= -\int_{-R}^{R}|x|\varphi'(x)\,\mathrm{d}x = -\left(\int_{-R}^{0}(-x)\varphi'(x)\,\mathrm{d}x + \int_{0}^{R}x\varphi'(x)\,\mathrm{d}x\right) \\
&= -\left(\Big[-x\varphi(x)\Big]_{-R}^{0} + \int_{-R}^{0}\varphi(x)\,\mathrm{d}x + \Big[x\varphi(x)\Big]_{0}^{R} - \int_{0}^{R}\varphi(x)\,\mathrm{d}x\right) \\
&= -\int_{-R}^{0}\varphi(x)\,\mathrm{d}x + \int_{0}^{R}\varphi(x)\,\mathrm{d}x = \int_{-\infty}^{\infty}\operatorname{sgn}(x)\varphi(x)\,\mathrm{d}x = \operatorname{sgn}(\varphi).
\end{aligned}}
 故 \((|x|)' = \operatorname{sgn}(x)\)。

\SubPillCoral{【考点复盘与避坑提示】}

\begin{itemize}
\tightlist
\item
  \textbf{核心考点}：广义导数的定义、挖洞分部积分法、利用 \(\varphi(x)-\varphi(0)\) 消去边界发散项的经典技巧。
\item
  \textbf{第 (3) 题为什么必须是 \(\lambda_1\) 而不是 \(\lambda_r\)（\(r \neq 1\)）？} 因为只有切在 \(|x|=1\) 处，\(\ln 1 = 0\) 才刚好把内外交界处的边界项消成 \(0\)；如果硬切在 \(\pm r\) 处，会多出一个 \(2(\ln r)\varphi(0)\) 的边界项，实际上 \(\lambda_1 = \lambda_r + 2(\ln r)\delta\)！
\item
  \textbf{弱导数对照}：顺手记牢，由 (4) 知 \(|x|\) 有弱导数 \(\operatorname{sgn}(x)\)，由 (1) 知 \(\operatorname{sgn}(x)\) 无弱导数（因 \(2\delta\) 不是局部可积函数）。
\end{itemize}

\SoftDivider

\BookProblem{prob:8-14}{题目 4：习题 8 第 14 题}{教材 p.341}{题目 4：习题 8 第 14 题}

\SubPillPrimary{【原题】}

设 \(f_\varepsilon(x) = (x+\mathrm{i}\varepsilon)^{-1} + (x-\mathrm{i}\varepsilon)^{-1}\)。证明：极限 \(\lim_{\varepsilon\to 0} f_\varepsilon\) 属于 \(\mathcal{D}'(\mathbb{R}^n)\)。

\begin{quote}
\emph{教材印刷勘误提示}：原题公式中的变量 \(x\) 是一维实变量，题末的 \(\mathcal{D}'(\mathbb{R}^n)\)（原书印作 \(D'(\mathbb{R}^n)\)）系 \(\mathcal{D}'(\mathbb{R})\) 的笔误。
\end{quote}

\SubPillAmber{【零基础破题思路】}

第一步先做初中通分，立刻看穿纸老虎： 
\FitDisplayMath{f_\varepsilon(x) = \frac{1}{x+\mathrm{i}\varepsilon} + \frac{1}{x-\mathrm{i}\varepsilon} = \frac{(x-\mathrm{i}\varepsilon) + (x+\mathrm{i}\varepsilon)}{x^2 + \varepsilon^2} = \frac{2x}{x^2 + \varepsilon^2}.}
 注意两个惊喜：\CircNum{1} 它是\textbf{实值光滑函数}；\CircNum{2} 它只含 \(\varepsilon^2\)，所以关于 \(\varepsilon\) 是偶函数，\textbf{\(\varepsilon \to 0\) 的双侧极限天然相同！}

接下来怎么证明它在 \(\mathcal{D}'(\mathbb{R})\) 中收敛（并且明确算出极限是 \(2\operatorname{p.v.}\frac{1}{x}\)）？提供\textbf{两种满分证法}：

\begin{itemize}
\tightlist
\item
  \textbf{方法一（最优雅，3 步秒杀）}：一眼认出 \(\frac{2x}{x^2+\varepsilon^2} = \frac{\mathrm{d}}{\mathrm{d}x}\ln(x^2 + \varepsilon^2)\)！只要用控制收敛定理证原函数 \(\ln(x^2+\varepsilon^2) \to 2\ln|x|\) 于 \(L^1_{\mathrm{loc}}(\mathbb{R})\)，立刻推得在 \(\mathcal{D}'(\mathbb{R})\) 中收敛，再用「广义导数与分布极限可交换」及第 11(2) 题，一行写出极限 \(= (2\ln|x|)' = 2\operatorname{p.v.}\frac{1}{x} \in \mathcal{D}'(\mathbb{R})\)！
\item
  \textbf{方法二（完全自包含，不依赖第 11 题）}：利用奇函数积分 \(\int_{-R}^R \frac{2x}{x^2+\varepsilon^2}\varphi(0)\,\mathrm{d}x = 0\)，把 \(\varphi(x)\) 换成 \(\varphi(x)-\varphi(0)\)，瞬间用中值定理控住 \(\frac{2x}{x^2+\varepsilon^2}(\varphi(x)-\varphi(0))\)，直接控制收敛取极限并按式 (8.3) 验证极限泛函属于 \(\mathcal{D}'(\mathbb{R})\)！
\end{itemize}

\SubPillMain{【规范解答】}

先化简 \(f_\varepsilon(x)\)（\(\varepsilon \in \mathbb{R}\setminus\{0\}\)）： 
\FitDisplayMath{f_\varepsilon(x) = \frac{1}{x+\mathrm{i}\varepsilon} + \frac{1}{x-\mathrm{i}\varepsilon} = \frac{2x}{x^2 + \varepsilon^2}.}
 显然对任意 \(\varepsilon \neq 0\)，\(f_\varepsilon\) 为实值光滑函数，故定义 \(\mathcal{D}'(\mathbb{R})\) 中的常规广义函数；又 \(f_{-\varepsilon} = f_\varepsilon\)，故 \(\varepsilon \to 0\) 的双侧极限等价于 \(\varepsilon \to 0^+\)。下面给出两种严密证明（考试任选其一即可）。

\SoftDivider

\StepHeading{方法一：原函数 + 广义导数的连续性（最简洁）}

\begin{enumerate}
\def\labelenumi{\arabic{enumi}.}
\item
  \textbf{构造原函数}：对 \(\varepsilon \neq 0\)，令 \(F_\varepsilon(x) \triangleq \ln(x^2 + \varepsilon^2) \in C^\infty(\mathbb{R})\)。由经典求导与例 8.8，在 \(\mathcal{D}'(\mathbb{R})\) 中有 
\FitDisplayMath{F_\varepsilon' = \frac{2x}{x^2 + \varepsilon^2} = f_\varepsilon.}

\item
  \textbf{证明 \(F_\varepsilon \to 2\ln|x|\) 于 \(L^1_{\mathrm{loc}}(\mathbb{R})\)}： 任取紧区间 \([-R, R]\)（\(R > 0\)），考虑 \(0 < |\varepsilon| \le 1\)。

  \begin{itemize}
  \tightlist
  \item
    对每个 \(x \in [-R, R]\setminus\{0\}\)，有 \(\lim_{\varepsilon\to 0}\ln(x^2 + \varepsilon^2) = \ln(x^2) = 2\ln|x|\)（几乎处处收敛）。
  \item
    对任意 \(x \in [-R, R]\setminus\{0\}\) 及 \(0 < |\varepsilon| \le 1\)，有 \(x^2 \le x^2 + \varepsilon^2 \le R^2 + 1\)。因 \(\ln t\) 单调递增，故 \(\ln(x^2) \le \ln(x^2+\varepsilon^2) \le \ln(R^2+1)\)，从而 
\FitDisplayMath{\big|\ln(x^2 + \varepsilon^2)\big| \le 2|\ln|x|| + \ln(R^2 + 1) \in L^1([-R, R]).}
 由 Lebesgue 控制收敛定理，得 
\FitDisplayMath{\lim_{\varepsilon\to 0}\int_{-R}^{R}\big|F_\varepsilon(x) - 2\ln|x|\big|\,\mathrm{d}x = 0,}
 即 \(F_\varepsilon \to 2\ln|x|\) 于 \(L^1_{\mathrm{loc}}(\mathbb{R})\)。按教材 p.305，\(L^1_{\mathrm{loc}}(\mathbb{R})\) 收敛蕴含 \(\mathcal{D}'(\mathbb{R})\) 收敛，因此 \(F_\varepsilon \to 2\ln|x|\) 于 \(\mathcal{D}'(\mathbb{R})\)。
  \end{itemize}
\item
  \textbf{两边取广义导数}：由广义导数在 \(\mathcal{D}'(\mathbb{R})\) 上的连续性（教材 p.307）及第 11(2) 题结论 \((\ln|x|)' = \operatorname{p.v.}\frac{1}{x}\)，对任意 \(\varphi \in \mathcal{D}(\mathbb{R})\) 有 
\FitDisplayMath{\lim_{\varepsilon\to 0} f_\varepsilon(\varphi) = \lim_{\varepsilon\to 0} F_\varepsilon'(\varphi) = -\lim_{\varepsilon\to 0} F_\varepsilon(\varphi') = -(2\ln|x|)(\varphi') = (2\ln|x|)'(\varphi) = \left(2\operatorname{p.v.}\frac{1}{x}\right)(\varphi).}
 由例 8.5 知 \(2\operatorname{p.v.}\frac{1}{x} \in \mathcal{D}'(\mathbb{R})\)，故 \(\lim_{\varepsilon\to 0} f_\varepsilon = 2\operatorname{p.v.}\frac{1}{x} \in \mathcal{D}'(\mathbb{R})\)。
\end{enumerate}

\SoftDivider

\StepHeading{方法二：奇函数减常值 + 直接控制收敛（完全自包含）}

任取 \(\varphi \in \mathcal{D}(\mathbb{R})\)，设 \(\operatorname{supp}\varphi \subset [-R, R]\)（\(R > 0\)）。因 \(\frac{2x}{x^2+\varepsilon^2}\) 是奇函数，有 \(\int_{-R}^{R}\frac{2x}{x^2+\varepsilon^2}\varphi(0)\,\mathrm{d}x = 0\)，于是 
\FitDisplayMath{f_\varepsilon(\varphi) = \int_{-R}^{R}\frac{2x}{x^2+\varepsilon^2}\varphi(x)\,\mathrm{d}x = \int_{-R}^{R}\frac{2x}{x^2+\varepsilon^2}\big(\varphi(x) - \varphi(0)\big)\,\mathrm{d}x.}
 对所有 \(x \in [-R, R]\setminus\{0\}\) 及 \(\varepsilon \neq 0\)，由 \(\frac{|x|}{x^2+\varepsilon^2} \le \frac{1}{|x|}\) 与中值定理 \(|\varphi(x)-\varphi(0)| \le |x|\|\varphi'\|_\infty\)，得统一可积控制： 
\FitDisplayMath{\left|\frac{2x}{x^2+\varepsilon^2}\big(\varphi(x) - \varphi(0)\big)\right| \le \frac{2}{|x|}\cdot |x|\|\varphi'\|_\infty = 2\|\varphi'\|_\infty \in L^1([-R, R]).}
 且对每个 \(x \neq 0\)，被积函数逐点收敛于 \(2\frac{\varphi(x)-\varphi(0)}{x}\)。由 Lebesgue 控制收敛定理，得 
\FitDisplayMath{\lim_{\varepsilon\to 0} f_\varepsilon(\varphi) = 2\int_{-R}^{R}\frac{\varphi(x) - \varphi(0)}{x}\,\mathrm{d}x \triangleq \lambda(\varphi).}
 最后验证 \(\lambda \in \mathcal{D}'(\mathbb{R})\)：

\begin{itemize}
\tightlist
\item
  \(\lambda\) 显然是线性泛函；
\item
  对任意紧集 \(K \subset [-R, R]\) 及 \(\varphi \in C_c^\infty(K)\)，有 
\FitDisplayMath{|\lambda(\varphi)| \le 2\int_{-R}^{R}\frac{|\varphi(x)-\varphi(0)|}{|x|}\,\mathrm{d}x \le 2\int_{-R}^{R}\|\varphi'\|_\infty\,\mathrm{d}x = 4R\|\varphi'\|_\infty \le 4R\|\varphi\|_{C^1(K)},}
 由教材 p.304 式 (8.3) 立得 \(\lambda \in \mathcal{D}'(\mathbb{R})\)（且由对称挖洞即知 \(\lambda = 2\operatorname{p.v.}\frac{1}{x}\)）。证毕。
\end{itemize}

\SubPillCoral{【考点复盘与避坑提示】}

\begin{itemize}
\tightlist
\item
  \textbf{知识串联（Sokhotski-Plemelj 公式）}：把第 9(1) 题和本题拼在一起看，你会看到数学物理中最著名的公式之一： 
\FitDisplayMath{\frac{1}{x \pm \mathrm{i}\varepsilon} = \frac{x}{x^2+\varepsilon^2} \mp \mathrm{i}\frac{\varepsilon}{x^2+\varepsilon^2} \xrightarrow{\varepsilon\to 0^+} \operatorname{p.v.}\frac{1}{x} \mp \mathrm{i}\pi\delta.}
 本题把两者相加，虚部的 \(\mp\mathrm{i}\pi\delta\) 恰好抵消，只留下实部的 \(2\operatorname{p.v.}\frac{1}{x}\)！
\item
  \textbf{避坑提示}：用方法一控制 \(|\ln(x^2+\varepsilon^2)|\) 时，必须把\textbf{正负两端都控住}（下界用 \(\ln(x^2)\) 控原点附近，上界用 \(\ln(R^2+1)\) 控远离原点处），千万别直接错写成 \(|\ln(x^2+\varepsilon^2)| \le |\ln(x^2)|\)（当 \(x=1, \varepsilon=1\) 时左边是 \(\ln 2\)，右边是 \(0\)，不等式反了！）。
\end{itemize}

\SoftDivider

\BookProblem{prob:8-17}{题目 5：习题 8 第 17 题 (2)(3)}{教材 p.342}{题目 5：习题 8 第 17 题 (2)(3)}

\SubPillPrimary{【原题】}

计算下列函数的 Fourier 变换：

\begin{enumerate}
\def\labelenumi{(\arabic{enumi})}
\setcounter{enumi}{1}
\tightlist
\item
  \(\dfrac{1}{a^2 + x^2}\)；
\item
  \(\mathrm{e}^{-ax^2}\)；
\end{enumerate}

\SubPillAmber{【零基础破题思路】}

这两题都是\textbf{不依赖复变函数留数定理}、纯靠大一微积分就能优雅搞定的经典计算！全程牢记本书式 (8.9) 约定：\(\hat{f}(\xi) = \int_{-\infty}^{\infty} f(x)\mathrm{e}^{-2\pi\mathrm{i}x\xi}\,\mathrm{d}x\)。

\begin{itemize}
\tightlist
\item
  \textbf{(2) 借花献佛（先算指数衰减再反演）}：直接积分 \(\int \frac{\mathrm{e}^{-2\pi\mathrm{i}x\xi}}{a^2+x^2}\,\mathrm{d}x\) 需要留数法，但反过来算指数函数 \(g(t) = \mathrm{e}^{-2\pi |a||t|}\) 的 Fourier 变换却只是大一微积分的 \(\int \mathrm{e}^{\alpha t}\,\mathrm{d}t\)！一算出来居然恰好等于 \(\frac{|a|}{\pi(a^2+x^2)}\)，再用 Fourier 反演公式一翻身，立刻秒得答案！
\item
  \textbf{(3) 求导建立一阶常微分方程（ODE 大法）}：令 \(f(x) = \mathrm{e}^{-ax^2}\)（\(a>0\)），它满足最简单的微分关系 \(f'(x) = -2ax f(x)\)。两边做 Fourier 变换（命题 8.14 与 8.15），``乘 \(x\)''变成``对 \(\xi\) 求导''，``对 \(x\) 求导''变成``乘 \(2\pi\mathrm{i}\xi\)''，直接蹦出关于 \(\hat{f}(\xi)\) 的一阶齐次 ODE，10 秒解完再代初值 \(\hat{f}(0) = \sqrt{\frac{\pi}{a}}\)！
\end{itemize}

\SubPillMain{【规范解答】}

均采用教材 p.312 式 (8.9) 的 Fourier 变换约定：\(\hat{f}(\xi) = \int_{-\infty}^{\infty} f(x)\mathrm{e}^{-2\pi\mathrm{i}x\xi}\,\mathrm{d}x\)。

\SoftDivider

\textbf{(2) 计算 \(f(x) = \dfrac{1}{a^2 + x^2}\) 的 Fourier 变换（设实数 \(a \neq 0\)；注：\(a = 0\) 时 \(\frac{1}{x^2} \notin L^1_{\mathrm{loc}}(\mathbb{R})\)）}

构造辅助函数 \(g(t) \triangleq \mathrm{e}^{-2\pi |a||t|}\)。因 \(|a| > 0\)，显然 \(g \in L^1(\mathbb{R}) \cap L^2(\mathbb{R})\) 且连续。按定义分段直接计算 \(\hat{g}(x)\)： 
\FitDisplayMath{\begin{aligned}
\hat{g}(x) &= \int_{-\infty}^{\infty}\mathrm{e}^{-2\pi |a||t|}\mathrm{e}^{-2\pi\mathrm{i}tx}\,\mathrm{d}t \\
&= \int_{-\infty}^{0}\mathrm{e}^{2\pi(|a| - \mathrm{i}x)t}\,\mathrm{d}t + \int_{0}^{\infty}\mathrm{e}^{-2\pi(|a| + \mathrm{i}x)t}\,\mathrm{d}t \\
&= \left.\frac{\mathrm{e}^{2\pi(|a| - \mathrm{i}x)t}}{2\pi(|a| - \mathrm{i}x)}\right|_{-\infty}^{0} + \left.\frac{\mathrm{e}^{-2\pi(|a| + \mathrm{i}x)t}}{-2\pi(|a| + \mathrm{i}x)}\right|_{0}^{+\infty} \\
&= \frac{1}{2\pi(|a| - \mathrm{i}x)} + \frac{1}{2\pi(|a| + \mathrm{i}x)} \\
&= \frac{1}{2\pi}\cdot\frac{2|a|}{a^2 + x^2} = \frac{|a|}{\pi(a^2 + x^2)}.
\end{aligned}}
 注意到 \(\hat{g} \in L^1(\mathbb{R}) \cap L^2(\mathbb{R})\)。由 Plancherel 定理（定理 8.9）及 \(\mathcal{S}(\mathbb{R})\) 在 \(L^2(\mathbb{R})\) 中的稠密性（命题 8.16(6)），命题 8.17 中的反演关系 \(\mathcal{F}^*\mathcal{F} = I\) 延拓到整个 \(L^2(\mathbb{R})\)；又由注记 8.8，对 \(L^1 \cap L^2\) 函数 \(L^2\)-Fourier 变换与积分定义一致，故几乎处处成立 
\FitDisplayMath{g(\xi) = \int_{-\infty}^{\infty}\hat{g}(x)\mathrm{e}^{2\pi\mathrm{i}x\xi}\,\mathrm{d}x = \frac{|a|}{\pi}\int_{-\infty}^{\infty}\frac{1}{a^2 + x^2}\mathrm{e}^{2\pi\mathrm{i}x\xi}\,\mathrm{d}x.}
 因左右两端都是 \(\mathbb{R}\) 上的连续函数（右端由控制收敛定理连续），故上式对一切 \(\xi \in \mathbb{R}\) 逐点成立。将 \(\xi\) 替换为 \(-\xi\) 并利用 \(g(-\xi) = g(\xi)\)，即得 
\FitDisplayMath{\widehat{\frac{1}{a^2 + x^2}}(\xi) = \int_{-\infty}^{\infty}\frac{1}{a^2 + x^2}\mathrm{e}^{-2\pi\mathrm{i}x\xi}\,\mathrm{d}x = \frac{\pi}{|a|}\mathrm{e}^{-2\pi |a||\xi|}.}
 特别地，当 \(a > 0\) 时，结果为 \(\dfrac{\pi}{a}\mathrm{e}^{-2\pi a|\xi|}\)。

\SoftDivider

\textbf{(3) 计算 \(f(x) = \mathrm{e}^{-ax^2}\) 的 Fourier 变换（设 \(a > 0\)）}

\emph{参数说明}：若 \(a < 0\)，函数超指数增长，既不属于 \(L^1(\mathbb{R})\) 也不属于 \(\mathcal{S}'(\mathbb{R})\)；若 \(a = 0\)，常数函数 \(1\) 的缓增分布 Fourier 变换为 \(\delta\)（命题 8.21(7)）。以下求解标准情形 \(a > 0\)。

当 \(a > 0\) 时，\(f(x) = \mathrm{e}^{-ax^2} \in \mathcal{S}(\mathbb{R})\)，故 \(f, xf, f' \in L^1(\mathbb{R})\) 且 \(f'(x) = -2ax f(x)\)，即 \(x f(x) = -\frac{1}{2a}f'(x)\)。 依次应用教材 p.312 命题 8.14 与 p.313 命题 8.15，得 
\FitDisplayMath{\hat{f}'(\xi) = -2\pi\mathrm{i}\,\widehat{xf}(\xi) = -2\pi\mathrm{i}\left(-\frac{1}{2a}\right)\widehat{f'}(\xi) = \frac{\pi\mathrm{i}}{a}\cdot (2\pi\mathrm{i}\xi)\hat{f}(\xi) = -\frac{2\pi^2}{a}\xi\hat{f}(\xi).}
 （若直接分部积分，同样有 \(\hat{f}'(\xi) = \frac{\pi\mathrm{i}}{a}\int_{-\infty}^{\infty}(\mathrm{e}^{-ax^2})'\mathrm{e}^{-2\pi\mathrm{i}x\xi}\,\mathrm{d}x = -\frac{\pi\mathrm{i}}{a}(-2\pi\mathrm{i}\xi)\hat{f}(\xi) = -\frac{2\pi^2}{a}\xi\hat{f}(\xi)\)。）

由此得到一阶常微分方程 \(\hat{f}'(\xi) + \frac{2\pi^2}{a}\xi\hat{f}(\xi) = 0\)。两边同乘积分因子 \(\mathrm{e}^{\frac{\pi^2\xi^2}{a}}\)，得 
\FitDisplayMath{\frac{\mathrm{d}}{\mathrm{d}\xi}\left(\mathrm{e}^{\frac{\pi^2\xi^2}{a}}\hat{f}(\xi)\right) = 0 \implies \hat{f}(\xi) = \hat{f}(0)\mathrm{e}^{-\frac{\pi^2\xi^2}{a}}.}
 令 \(\xi = 0\) 计算初值（作换元 \(t = \sqrt{a}\,x\)）： 
\FitDisplayMath{\hat{f}(0) = \int_{-\infty}^{\infty}\mathrm{e}^{-ax^2}\,\mathrm{d}x = \frac{1}{\sqrt{a}}\int_{-\infty}^{\infty}\mathrm{e}^{-t^2}\,\mathrm{d}t = \sqrt{\frac{\pi}{a}}.}
 代入即得最终结果： 
\FitDisplayMath{\widehat{\mathrm{e}^{-ax^2}}(\xi) = \sqrt{\frac{\pi}{a}}\,\mathrm{e}^{-\frac{\pi^2\xi^2}{a}} \quad (a > 0).}


\SubPillCoral{【考点复盘与避坑提示】}

\begin{itemize}
\tightlist
\item
  \textbf{快速自检两绝招（考场防算错常数！）}：

  \begin{enumerate}
  \def\labelenumi{\arabic{enumi}.}
  \tightlist
  \item
    \textbf{代 \(\xi = 0\) 检验}：(2) 左边 \(\int_{-\infty}^\infty \frac{\mathrm{d}x}{a^2+x^2} = \frac{\pi}{|a|}\)，与右边 \(\frac{\pi}{|a|}\mathrm{e}^0\) 吻合；(3) 取 \(a = \pi\) 得 \(\widehat{\mathrm{e}^{-\pi x^2}}(\xi) = \mathrm{e}^{-\pi\xi^2}\)（Gauss 函数是 Fourier 变换的不动点，和教材 p.314 证明完全一致！）。
  \item
    \textbf{若对比其他不带 \(2\pi\) 的教材}（\(\tilde{f}(\omega) = \int f(x)\mathrm{e}^{-\mathrm{i}x\omega}\,\mathrm{d}x\)）：把结果中的 \(\xi\) 换成 \(\frac{\omega}{2\pi}\)，两题分别对应 \(\frac{\pi}{|a|}\mathrm{e}^{-|a||\omega|}\) 和 \(\sqrt{\frac{\pi}{a}}\mathrm{e}^{-\frac{\omega^2}{4a}}\)。\textbf{切忌把别的教材公式直接抄到本题试卷上！}
  \end{enumerate}
\end{itemize}

\SoftDivider

\BookProblem{prob:8-20}{题目 6：习题 8 第 20 题}{教材 p.342}{题目 6：习题 8 第 20 题}

\SubPillPrimary{【原题】}

利用 Fourier 变换求解以下方程： 
\FitTaggedMath{\begin{cases}
-\Delta u(x,y) = 0, & (x,y) \in \mathbb{R}_+^2 = (-\infty,\infty) \times (0,\infty), \\
u(x,0) = g(x), & x \in \mathbb{R}, \\
u(x,y) \to 0, & \text{若 } |(x,y)| \to \infty.
\end{cases}}{8.50}


\SubPillAmber{【零基础破题思路】}

这是偏微分方程里最经典的\textbf{「半空间偏 Fourier 变换法」}！记住一个原则：\textbf{哪个变量跑满整个 \((-\infty, +\infty)\)，就对哪个变量做 Fourier 变换！}

\begin{itemize}
\tightlist
\item
  这里 \(x \in (-\infty, \infty)\) 跑满整条实轴，而 \(y \in (0, \infty)\) 只有半条轴。
\item
  所以我们对 \(x\) 做一维 Fourier 变换（把 \(y\) 当普通参数看）：

  \begin{enumerate}
  \def\labelenumi{\arabic{enumi}.}
  \tightlist
  \item
    \(u_{xx}\) 瞬间变成代数项 \(-4\pi^2\xi^2 \hat{u}\)，原本的偏微分方程 \(u_{xx}+u_{yy}=0\) 立刻退化成关于 \(y\) 的大一常微分方程 \(\hat{u}_{yy} - (2\pi\xi)^2 \hat{u} = 0\)！
  \item
    写出通解 \(C_1\mathrm{e}^{-2\pi|\xi|y} + C_2\mathrm{e}^{2\pi|\xi|y}\)，靠 \(y\to+\infty\) 不爆炸砍掉 \(C_2\)，靠 \(y=0\) 边值定出 \(C_1 = \hat{g}(\xi)\)，得到 \(\hat{u}(\xi, y) = \hat{g}(\xi)\mathrm{e}^{-2\pi y|\xi|}\)；
  \item
    \textbf{梦幻联动第 17(2) 题}：因子 \(\mathrm{e}^{-2\pi y|\xi|}\) 刚好是 Poisson 核 \(P_y(x) = \frac{1}{\pi}\frac{y}{x^2+y^2}\) 的 Fourier 变换！频域相乘 = 物理空间卷积，直接写出上半平面 Poisson 积分公式！
  \item
    \textbf{再联动第 9(1) 题}：当 \(y \to 0^+\) 时 \(P_y \to \delta\)，所以 \(P_y * g \to g\) 自动还原边值！第 9(1)、17(2)、20 三题首尾呼应，形成完美闭环！
  \end{enumerate}
\end{itemize}

\SubPillMain{【规范解答】}

假定边值函数 \(g\) 满足适当的正则与衰减条件（例如 \(g \in \mathcal{S}(\mathbb{R})\)，或 \(g \in L^1(\mathbb{R}) \cap C_0(\mathbb{R})\)），且待求经典解 \(u(x,y)\) 及其关于 \(x\) 的偏导数充分衰减以保证 Fourier 变换与积分号下求导合法。

\begin{enumerate}
\def\labelenumi{\arabic{enumi}.}
\item
  \textbf{对水平变量 \(x\) 作一维 Fourier 变换（按本书式 (8.9) 约定）} 记 
\FitDisplayMath{\hat{u}(\xi, y) \triangleq \int_{-\infty}^{\infty} u(x, y)\mathrm{e}^{-2\pi\mathrm{i}x\xi}\,\mathrm{d}x, \qquad \xi \in \mathbb{R},\ y \ge 0.}
 方程 \(-\Delta u = 0\) 即 \(u_{xx}(x, y) + u_{yy}(x, y) = 0\)。两边关于 \(x\) 取 Fourier 变换，由命题 8.15 得 \(\widehat{u_{xx}}(\xi, y) = (2\pi\mathrm{i}\xi)^2 \hat{u}(\xi, y) = -4\pi^2\xi^2 \hat{u}(\xi, y)\)，且 \(\widehat{u_{yy}}(\xi, y) = \frac{\partial^2 \hat{u}}{\partial y^2}(\xi, y)\)，从而化为以 \(\xi\) 为参数的关于 \(y\) 的二阶常微分方程： 
\FitDisplayMath{\frac{\partial^2 \hat{u}}{\partial y^2}(\xi, y) - (2\pi\xi)^2 \hat{u}(\xi, y) = 0, \qquad y > 0.}

\item
  \textbf{求解常微分方程并利用边界与无穷远条件定常数} 对固定的 \(\xi \neq 0\)，特征方程 \(r^2 - (2\pi\xi)^2 = 0\) 的两根为 \(r = \pm 2\pi|\xi|\)，故通解为 
\FitDisplayMath{\hat{u}(\xi, y) = C_1(\xi)\mathrm{e}^{-2\pi|\xi|y} + C_2(\xi)\mathrm{e}^{2\pi|\xi|y}.}


  \begin{itemize}
  \tightlist
  \item
    当 \(y \to +\infty\) 时，\(\mathrm{e}^{2\pi|\xi|y} \to +\infty\)。由题设无穷远条件 \(u(x,y) \to 0\)（要求 \(\hat{u}(\xi, y)\) 在 \(y \to +\infty\) 时有界），必有 \(C_2(\xi) = 0\)。
  \item
    再由 \(y = 0\) 处的边值条件 \(u(x, 0) = g(x)\)，两边关于 \(x\) 取 Fourier 变换得 \(\hat{u}(\xi, 0) = \hat{g}(\xi)\)，从而 \(C_1(\xi) = \hat{g}(\xi)\)。
  \end{itemize}

  因此（由连续性包含 \(\xi = 0\)）， 
\FitDisplayMath{\hat{u}(\xi, y) = \hat{g}(\xi)\mathrm{e}^{-2\pi y|\xi|}, \qquad \xi \in \mathbb{R},\ y > 0.}

\item
  \textbf{Fourier 反演与卷积定理求显式解} 对任意固定的 \(y > 0\)，引入上半平面 \textbf{Poisson 核}： 
\FitDisplayMath{P_y(x) \triangleq \frac{1}{\pi}\frac{y}{x^2 + y^2}.}
 由习题 8 第 17(2) 题（令 \(a = y > 0\)），\(\widehat{\frac{1}{x^2+y^2}}(\xi) = \frac{\pi}{y}\mathrm{e}^{-2\pi y|\xi|}\)，因此恰好有 
\FitDisplayMath{\widehat{P_y}(\xi) = \frac{y}{\pi}\cdot\frac{\pi}{y}\mathrm{e}^{-2\pi y|\xi|} = \mathrm{e}^{-2\pi y|\xi|}.}
 于是 
\FitDisplayMath{\hat{u}(\xi, y) = \hat{g}(\xi)\widehat{P_y}(\xi) = \widehat{(P_y * g)}(\xi).}
 由 Fourier 反演公式，得边值问题 (8.50) 的解（\textbf{上半平面 Poisson 积分公式}）： 
\FitDisplayMath{u(x, y) = (P_y * g)(x) = \frac{y}{\pi}\int_{-\infty}^{\infty}\frac{g(t)}{(x - t)^2 + y^2}\,\mathrm{d}t, \qquad (x, y) \in \mathbb{R}_+^2.}

\item
  \textbf{解的合理性简明核验（对 \(g \in C_0(\mathbb{R})\)，即连续且 \(\lim_{|t|\to\infty} g(t) = 0\)）}

  \begin{itemize}
  \tightlist
  \item
    \textbf{调和性}：当 \(y > 0\) 时，直接求偏导易验证 \((\partial_{xx} + \partial_{yy})P_y(x-t) = 0\)，且各阶偏导在 \(y \ge y_0 > 0\) 上可积，积分号下求导即得 \(-\Delta u = 0\)。
  \item
    \textbf{边值还原}：当 \(y \to 0^+\) 时，如第 9(1) 题所示 \(P_y \to \delta\)（且 \(\int P_y = 1, P_y > 0\)），由 \(g\) 的连续有界性易证 \(u(x,y) \to g(x)\)。
  \item
    \textbf{无穷远趋于 \(0\)}：任给 \(\eta > 0\)，取 \(M > 0\) 使 \(|t| > M\) 时 \(|g(t)| < \eta\)。将积分拆成 \(|t| > M\) 与 \(|t| \le M\) 两部分：外段绝对值不超过 \(\eta\int_{\mathbb{R}} P_y(x-t)\,\mathrm{d}t = \eta\)；内段利用初等不等式 \(\frac{y}{(x-t)^2+y^2} \le \min\!\left\{\frac{1}{y},\ \frac{1}{2|x-t|}\right\}\) 可知当 \(|(x,y)| \to \infty\) 时内段一致趋于 \(0\)。故 \(\lim_{|(x,y)|\to\infty} u(x,y) = 0\)。
  \end{itemize}
\end{enumerate}

\SubPillCoral{【考点复盘与避坑提示】}

\begin{itemize}
\tightlist
\item
  \textbf{核心考点}：用 Fourier 变换求解常系数偏微分方程的标准四步（变换降维 \(\to\) 解 ODE \(\to\) 边界与衰减定常数 \(\to\) 卷积反演）。
\item
  \textbf{避坑提示}：写 \(\mathrm{e}^{-2\pi|\xi|y}\) 时\textbf{绝对不能漏掉 \(|\xi|\) 上的绝对值}！如果错写成 \(\mathrm{e}^{-2\pi\xi y}\)，当 \(\xi < 0\) 时它反而随 \(y \to +\infty\) 指数爆炸，全盘皆错！
\end{itemize}

\SoftDivider

\BookProblem{prob:8-28}{题目 7：习题 8 第 28 题}{教材 p.343}{题目 7：习题 8 第 28 题}

\SubPillPrimary{【原题】}

设 \(k \in \mathbb{N}\)。求使得 \(\dfrac{1}{|x|^k} \in H^{-m}(\mathbb{R}^n)\) 的最小自然数 \(m\)。

\SubPillAmber{【零基础破题思路】}

这道题是全章最能体现\textbf{实指数 Sobolev 空间威力}的压轴好题，也是最容易踩「漏查一半条件」大坑的题！请务必分清\textbf{两件事}：

\begin{enumerate}
\def\labelenumi{\arabic{enumi}.}
\item
  \textbf{先查入门资格}：函数 \(f_k(x) = \frac{1}{|x|^k}\) 要成为 \(\mathbb{R}^n\) 上合法的常规缓增分布（即局部可积 \(L^1_{\mathrm{loc}}(\mathbb{R}^n)\)），在原点用球坐标算一下： 
\FitDisplayMath{\int_{|x|<1}\frac{\mathrm{d}x}{|x|^k} = \omega_{n-1}\int_0^1 r^{n-1-k}\,\mathrm{d}r < \infty \iff k < n.}
 （若 \(k = 0\) 则 \(f_0 \equiv 1\)，\(\hat{1} = \delta\) 不是函数，不在任何 \(H^{-m}\) 里；若 \(k \ge n\)，它在原点不可积，作为普通函数根本不能定义常规广义函数，后面我们还会单独讨论其正则化情形！）
\item
  \textbf{用 \(H^{-m}(\mathbb{R}^n)\) 定义直击频域（教材 p.336）}： 
\FitDisplayMath{f_k \in H^{-m}(\mathbb{R}^n) \iff \int_{\mathbb{R}^n} (1 + |\xi|^2)^{-m} |\widehat{f_k}(\xi)|^2\,\mathrm{d}\xi < \infty.}
 关键是要算出 \(\widehat{|x|^{-k}}(\xi)\)！下面我们用第 17(3) 题的 Gauss 变换配合 \(\Gamma\) 函数\textbf{4 行严密推导出} \(\widehat{|x|^{-k}}(\xi) = c_{n,k}|\xi|^{k-n}\)（零黑箱！）。代入后转球坐标，积分变成： 
\FitDisplayMath{\int_0^\infty (1 + r^2)^{-m} r^{2k - n - 1}\,\mathrm{d}r < \infty.}
 \textbf{\WarnBadge{} 超级防坑关键}：这个积分在 \(r \to 0^+\) 和 \(r \to +\infty\) \textbf{两端都有可能发散}，必须两边同时检查！

  \begin{itemize}
  \tightlist
  \item
    \textbf{低频端 \(r \to 0^+\)}（对应物理空间无穷远 \(|x|\to\infty\) 的衰减）：需要 \(2k - n - 1 > -1 \iff \mathbf{k > \frac{n}{2}}\)！这一关和 \(m\) 无关——也就是说，\textbf{当 \(k \le \frac{n}{2}\) 时，不管 \(m\) 取多大都无解！}
  \item
    \textbf{高频端 \(r \to +\infty\)}（对应原点 \(x = 0\) 的奇性）：需要 \(2k - n - 1 - 2m < -1 \iff \mathbf{m > k - \frac{n}{2}}\)！由此立刻定出最小自然数 \(m = \left\lfloor k - \frac{n}{2}\right\rfloor + 1\)！
  \end{itemize}
\end{enumerate}

\SubPillMain{【规范解答】}

按教材惯例记自然数集 \(\mathbb{N} = \{0, 1, 2, \dots\}\)（下文可见，凡有解情形求得的最小 \(m\) 均满足 \(m \ge 1\)，故无论 \(\mathbb{N}\) 是否含 \(0\)，有解时的答案完全一致）。记 \(S^{n-1}\) 的表面积为 \(\omega_{n-1}\)。

\StepHeading{第一步：先排查 \(k = 0\) 与常规缓增分布的存在范围（\(1 \le k < n\)）}

\begin{itemize}
\tightlist
\item
  \textbf{若 \(k = 0\)}：此时 \(\frac{1}{|x|^0} \equiv 1\)。由命题 8.21(7) 知 \(\hat{1} = \delta\)，而对任意 \(m \in \mathbb{N}\)，\(\langle\xi\rangle^{-m}\delta = \delta\) 不是局部可积函数，更不属于 \(L^2(\mathbb{R}^n)\)。故 \(k = 0\) 时对任何 \(m \in \mathbb{N}\) 均有 \(1 \notin H^{-m}(\mathbb{R}^n)\)（无解）。
\item
  \textbf{检查局部可积性}：对正整数 \(k \ge 1\)，在单位球上用球坐标得 
\FitDisplayMath{\int_{|x|<1}\frac{\mathrm{d}x}{|x|^k} = \omega_{n-1}\int_0^1 r^{n-1-k}\,\mathrm{d}r < \infty \iff n - 1 - k > -1 \iff k < n.}
 而在 \(|x| \ge 1\) 上有 \(0 < |x|^{-k} \le 1\) 有界。因此当 \(1 \le k < n\) 时，\(f_k(x) \triangleq \frac{1}{|x|^k} \in L^1_{\mathrm{loc}}(\mathbb{R}^n) \cap \mathcal{S}'(\mathbb{R}^n)\) 是良定义的实值常规缓增分布。
\end{itemize}

\StepHeading{第二步：严格计算 \(\widehat{|x|^{-k}}(\xi)\)（\(1 \le k < n\)）}

利用 Gamma 函数定义 \(\Gamma(\alpha) = \int_0^\infty u^{\alpha-1}\mathrm{e}^{-u}\,\mathrm{d}u\)（\(\alpha > 0\)），取 \(\alpha = \frac{k}{2}\) 并作换元 \(u = \pi |x|^2 t\)（\(x \neq 0\)），得 Gauss 积分表示： 
\FitDisplayMath{\frac{1}{|x|^k} = \frac{\pi^{k/2}}{\Gamma(k/2)}\int_0^\infty t^{\frac{k}{2}-1}\mathrm{e}^{-\pi t|x|^2}\,\mathrm{d}t.}
 由第 17(3) 题（取 \(a = \pi t > 0\) 并对 \(n\) 个坐标连乘），\(n\) 维 Gauss 函数的 Fourier 变换为 
\FitDisplayMath{\int_{\mathbb{R}^n}\mathrm{e}^{-\pi t|x|^2}\mathrm{e}^{-2\pi\mathrm{i}x\cdot\xi}\,\mathrm{d}x = t^{-n/2}\mathrm{e}^{-\frac{\pi}{t}|\xi|^2}.}
 任取 \(\varphi \in \mathcal{S}(\mathbb{R}^n)\)。因 \(1 \le k < n\)，有 \(\int_{\mathbb{R}^n}|x|^{-k}|\hat{\varphi}(x)|\,\mathrm{d}x < \infty\) 且（对下述换元 \(s = 1/t\) 后）\(\int_{\mathbb{R}^n}|\xi|^{k-n}|\varphi(\xi)|\,\mathrm{d}\xi < \infty\)，两次应用 Fubini 定理得 
\FitDisplayMath{\begin{aligned}
(\mathcal{F}f_k)(\varphi) &= \int_{\mathbb{R}^n}\frac{1}{|x|^k}\hat{\varphi}(x)\,\mathrm{d}x \\
&= \frac{\pi^{k/2}}{\Gamma(k/2)}\int_0^\infty t^{\frac{k}{2}-1}\left(\int_{\mathbb{R}^n}\mathrm{e}^{-\pi t|x|^2}\hat{\varphi}(x)\,\mathrm{d}x\right)\mathrm{d}t \\
&= \frac{\pi^{k/2}}{\Gamma(k/2)}\int_0^\infty t^{\frac{k-n}{2}-1}\left(\int_{\mathbb{R}^n}\mathrm{e}^{-\frac{\pi}{t}|\xi|^2}\varphi(\xi)\,\mathrm{d}\xi\right)\mathrm{d}t \\
&= \frac{\pi^{k/2}}{\Gamma(k/2)}\int_{\mathbb{R}^n}\left(\int_0^\infty s^{\frac{n-k}{2}-1}\mathrm{e}^{-\pi s|\xi|^2}\,\mathrm{d}s\right)\varphi(\xi)\,\mathrm{d}\xi \\
&= c_{n,k}\int_{\mathbb{R}^n}|\xi|^{k-n}\varphi(\xi)\,\mathrm{d}\xi,
\end{aligned}}
 其中 \(s = \frac{1}{t}\)，正常数 
\FitDisplayMath{c_{n,k} = \pi^{k - \frac{n}{2}}\frac{\Gamma\!\left(\frac{n-k}{2}\right)}{\Gamma\!\left(\frac{k}{2}\right)} > 0.}
 于是，在 \(\mathcal{S}'(\mathbb{R}^n)\) 中有 
\FitDisplayMath{\widehat{|x|^{-k}}(\xi) = c_{n,k}|\xi|^{k-n}.}


\StepHeading{第三步：按 \(H^{-m}(\mathbb{R}^n)\) 定义检查两端的积分收敛性（\(1 \le k < n\)）}

按教材 p.336 定义，\(f_k \in H^{-m}(\mathbb{R}^n) \iff \langle\xi\rangle^{-m}\widehat{f_k}(\xi) \in L^2(\mathbb{R}^n)\)。因 \(c_{n,k} > 0\)，这等价于如下球坐标积分收敛： 
\FitDisplayMath{I \triangleq \int_{\mathbb{R}^n}(1 + |\xi|^2)^{-m}|\xi|^{2(k-n)}\,\mathrm{d}\xi = \omega_{n-1}\int_0^\infty (1 + r^2)^{-m} r^{2k - n - 1}\,\mathrm{d}r < \infty.}
 将 \((0, \infty)\) 拆分为 \((0, 1]\) 与 \([1, \infty)\) 两段分别检验：

\begin{enumerate}
\def\labelenumi{\arabic{enumi}.}
\tightlist
\item
  \textbf{低频端（\(r \to 0^+\)）}：在 \((0, 1]\) 上，\((1 + r^2)^{-m} \in [2^{-m}, 1]\)，故 \(\int_0^1 (1+r^2)^{-m} r^{2k-n-1}\,\mathrm{d}r < \infty\) 的充要条件为 
\FitDisplayMath{2k - n - 1 > -1 \iff k > \frac{n}{2}.}


  \begin{itemize}
  \tightlist
  \item
    特别地，当 \(1 \le k \le \frac{n}{2}\) 时，无论自然数 \(m\) 取多大，\(I = +\infty\)，故\textbf{不存在任何 \(m\) 使得 \(\frac{1}{|x|^k} \in H^{-m}(\mathbb{R}^n)\)}！
  \end{itemize}
\item
  \textbf{高频端（\(r \to +\infty\)）}：在 \([1, \infty)\) 上，\((1 + r^2)^{-m} \asymp r^{-2m}\)，故被积函数等价于 \(r^{2k - n - 1 - 2m}\)，积分 \(\int_1^\infty (1+r^2)^{-m} r^{2k-n-1}\,\mathrm{d}r < \infty\) 的充要条件为 
\FitDisplayMath{2k - n - 1 - 2m < -1 \iff 2m > 2k - n \iff m > k - \frac{n}{2}.}

\end{enumerate}

综合两端条件：当 \textbf{\(\dfrac{n}{2} < k < n\)} 时，\(m \in \mathbb{N}\) 满足 \(\frac{1}{|x|^k} \in H^{-m}(\mathbb{R}^n)\) 当且仅当 \(m > k - \frac{n}{2}\)，因此\textbf{最小自然数 \(m\)} 为： 
\FitDisplayMath{m = \left\lfloor k - \frac{n}{2}\right\rfloor + 1 = \begin{cases}
k - \dfrac{n}{2} + 1, & n \text{ 为偶数}, \\[8pt]
k - \dfrac{n - 1}{2}, & n \text{ 为奇数}.
\end{cases}}


\StepHeading{第四步：关于 \(k \ge n\) 的完整补全说明}

\begin{itemize}
\tightlist
\item
  \textbf{作为常规函数}：当 \(k \ge n\) 时，\(\frac{1}{|x|^k} \notin L^1_{\mathrm{loc}}(\mathbb{R}^n)\)，在原点附近不可积，故不能按例 8.2 直接视为 \(\mathcal{D}'(\mathbb{R}^n)\) 或 \(\mathcal{S}'(\mathbb{R}^n)\) 中的常规广义函数。
\item
  \textbf{若理解为正则化缓增分布（Hadamard 有限部 \(\operatorname{Pf}|x|^{-k} \in \mathcal{S}'(\mathbb{R}^n)\)）}：其 Fourier 变换在 \(k - n \notin 2\mathbb{N}\) 时为 \(c|\xi|^{k-n}\)，在 \(k - n \in 2\mathbb{N}\) 时为 \(c_1|\xi|^{k-n}\ln|\xi| + c_2|\xi|^{k-n}\)。此时因 \(k \ge n \implies 2k - n - 1 \ge n - 1 \ge 0 > -1\)，低频端自动收敛；高频端 \(\int_1^\infty (1+r^2)^{-m} r^{2k-n-1}(\ln r)^{2\text{ 或 }0}\,\mathrm{d}r < \infty\) 的充要条件依然是严格不等式 \(m > k - \frac{n}{2}\)（对数因子不影响临界严格不等式），故最小自然数仍为 
\FitDisplayMath{m = \left\lfloor k - \frac{n}{2}\right\rfloor + 1.}

\end{itemize}

\StepHeading{最终结论汇总}

\begin{itemize}
\tightlist
\item
  当 \(k \le \dfrac{n}{2}\)（含 \(k = 0\)）时：\textbf{不存在}满足条件的自然数 \(m\)（因无穷远处衰减不足，低频积分发散）；
\item
  当 \(\dfrac{n}{2} < k < n\) 时（常规缓增分布）：最小自然数为 
\FitDisplayMath{m = \left\lfloor k - \frac{n}{2}\right\rfloor + 1 = \begin{cases} k - \dfrac{n}{2} + 1, & n \text{ 为偶数}, \\ k - \dfrac{n-1}{2}, & n \text{ 为奇数}; \end{cases}}

\item
  当 \(k \ge n\) 时：作为普通函数不在 \(L^1_{\mathrm{loc}}(\mathbb{R}^n)\) 中；若理解为 Hadamard 有限部正则化缓增分布，最小自然数仍为 \(m = \left\lfloor k - \frac{n}{2}\right\rfloor + 1\)。
\end{itemize}

\SubPillCoral{【考点复盘与避坑提示】}

\begin{itemize}
\tightlist
\item
  \textbf{用具体数字秒验结论}：

  \begin{itemize}
  \tightlist
  \item
    \(n = 3, k = 1\)（三维牛顿势 \(\frac{1}{|x|}\)）：\(1 \le \frac{3}{2} \implies\) \textbf{无解}（\(\widehat{|x|^{-1}} = \frac{1}{\pi|\xi|^2}\)，在 \(\xi=0\) 处 \(|\xi|^{-4} \cdot r^2 = r^{-2}\) 发散！）；
  \item
    \(n = 3, k = 2\)（\(\frac{1}{|x|^2}\)）：\(\frac{3}{2} < 2 < 3 \implies m > 2 - \frac{3}{2} = 0.5 \implies\) 最小自然数 \(m = 1\)；
  \item
    \(n = 4, k = 3\)（\(\frac{1}{|x|^3}\)）：\(2 < 3 < 4 \implies m > 3 - 2 = 1 \implies\) 最小自然数 \(m = 2\)。
  \end{itemize}
\item
  \textbf{为什么 \(k \le \frac{n}{2}\) 无解？（物理空间直观验证）}：若取非负截断 \(\psi \in C_c^\infty(\{1 \le |x| \le 2\})\) 并放大为 \(\psi_R(x) = \psi(x/R)\)，则配对 \(\int |x|^{-k}\psi_R\,\mathrm{d}x \sim R^{n-k}\)，而 \(\|\psi_R\|_{H^m} \lesssim R^{n/2}\)；当 \(k < \frac{n}{2}\) 时比值 \(\sim R^{\frac{n}{2}-k} \to \infty\)，根本不是 \(H^m(\mathbb{R}^n)\) 上的有界线性泛函！
\item
  \textbf{Sobolev 嵌入对偶印证高频条件}：为什么原点奇性对应的恰好是 \(m > k - \frac{n}{2}\)？由定理 8.15，当 \(m < \frac{n}{2}\) 时 \(H^m \hookrightarrow L^q\)（\(\frac{1}{q} = \frac{1}{2} - \frac{m}{n}\)）；由 Hölder 不等式，要在局部与 \(L^q\) 配对，需要 \(|x|^{-k} \in L^{q'}(B_1)\)，即 \(kq' < n \iff \frac{k}{n} < \frac{1}{q'} = 1 - \frac{1}{q} = \frac{1}{2} + \frac{m}{n} \iff m > k - \frac{n}{2}\)！频域和物理空间完美印证！
\end{itemize}

\SoftDivider

\BookProblem{prob:8-882}{题目 8：教材第 337 页 练习 8.8.2}{}{题目 8：教材第 337 页 练习 8.8.2}

\SubPillPrimary{【原题】}

（教材 p.337 \textbf{练习 8.8.2}）如果 \(L\) 是一个 \(m\) 阶椭圆算子，\(f\) 是满足 \(f, Lf \in H^s(\mathbb{R}^n)\) 的缓增分布，则 \(f \in H^{s+m}(\mathbb{R}^n)\)，且有如下估计： 
\FitTaggedMath{\|f\|_{H^{s+m}(\mathbb{R}^n)} \le C\big(\|f\|_{H^s(\mathbb{R}^n)} + \|Lf\|_{H^s(\mathbb{R}^n)}\big),}{8.49}
 其中 \(C\) 依赖于 \(s, m, n, L\)。

\begin{quote}
\emph{教材 p.337 上方定义的两处印刷笔误提示}：

\begin{enumerate}
\def\labelenumi{\arabic{enumi}.}
\tightlist
\item
  原书将 \(m\) 阶常系数齐次微分算子误印为 \(\sum_{j_1+\dots+j_n=m} c_{j_1,\dots,j_n}\frac{\partial^n}{\partial x_{j_1}\cdots\partial x_{j_n}}\)，按标准多重指标记法应为 \(L = \sum_{|\alpha|=m} c_\alpha \partial^\alpha = \sum_{\alpha_1+\dots+\alpha_n=m} c_\alpha \frac{\partial^m}{\partial x_1^{\alpha_1}\cdots\partial x_n^{\alpha_n}}\)，其象征为 \(l(\xi) = \sum_{|\alpha|=m} c_\alpha \xi^\alpha\)；
\item
  原书将椭圆性下界印作 \(l(\xi) \ge c|\xi|^m\)，但结合紧随其后的\textbf{例 8.17}（Cauchy-Riemann 算子 \(\partial_{x_1} - \mathrm{i}\partial_{x_2}\) 的象征 \(l(\xi) = \xi_1 - \mathrm{i}\xi_2\) 为复值函数），正确定义应取绝对值（模长）：存在常数 \(c > 0\) 使得对所有 \(\xi \in \mathbb{R}^n\) 有 \(|l(\xi)| \ge c|\xi|^m\)。
\end{enumerate}
\end{quote}

\SubPillAmber{【零基础破题思路】}

千万不要被「\(m\) 阶椭圆算子」「实指数 \(H^s\)」「缓增分布」这堆吓人的名字唬住——\textbf{一转到频域，这道题其实是初中代数不等式！}

看清它的翻译词典：

\begin{itemize}
\tightlist
\item
  已知 \(f \in H^s\) \(\iff\) \(\langle\xi\rangle^s |\hat{f}(\xi)| \in L^2\)（已经控住了\textbf{低频} \(1 \cdot \langle\xi\rangle^s |\hat{f}|\)）；
\item
  已知 \(Lf \in H^s\) \(\iff\) \(\langle\xi\rangle^s |\widehat{Lf}(\xi)| = \langle\xi\rangle^s (2\pi)^m |l(\xi)|\,|\hat{f}(\xi)| \in L^2\)；
\item
  椭圆性条件 \(|l(\xi)| \ge c|\xi|^m\) \(\iff\) 用 \(|\widehat{Lf}(\xi)|\) 就能反过来压住\textbf{高频} \(|\xi|^m|\hat{f}(\xi)|\)！
\item
  目标 \(f \in H^{s+m}\) \(\iff\) 要控住 \(\langle\xi\rangle^{s+m}|\hat{f}| = \langle\xi\rangle^s (1 + |\xi|^2)^{m/2} |\hat{f}|\)！
\item
  而初等不等式 \((1 + |\xi|^2)^{m/2} \le 2^{m/2}(1 + |\xi|^m)\) 刚好把它一拆为二：\textbf{「1」交给 \(f\) 的 \(H^s\) 范数，「\(|\xi|^m\)」交给 \(Lf\) 的 \(H^s\) 范数}，两边取 \(L^2\) 范数用三角不等式，瞬间收官！
\end{itemize}

\SubPillMain{【规范解答】}

设 \(m\) 阶常系数齐次椭圆算子为 
\FitDisplayMath{L = \sum_{|\alpha|=m} c_\alpha \partial^\alpha \quad (m \ge 0),}
 其象征 \(l(\xi) = \sum_{|\alpha|=m} c_\alpha \xi^\alpha\) 满足椭圆性条件：存在常数 \(c > 0\)，使得 
\FitDisplayMath{|l(\xi)| \ge c|\xi|^m, \qquad \forall \xi \in \mathbb{R}^n.}


\begin{enumerate}
\def\labelenumi{\arabic{enumi}.}
\item
  \textbf{将假设翻译为频域可测函数} 由 \(H^s(\mathbb{R}^n)\) 的定义（教材 p.336），\(f \in H^s(\mathbb{R}^n)\) 意味着 \(f \in \mathcal{S}'(\mathbb{R}^n)\) 且 \(\langle\xi\rangle^s \hat{f}(\xi) \in L^2(\mathbb{R}^n)\)，其中 \(\langle\xi\rangle = (1 + |\xi|^2)^{1/2} \ge 1\)。由此 \(\hat{f}(\xi) = \langle\xi\rangle^{-s}\big(\langle\xi\rangle^s \hat{f}(\xi)\big)\) 是 \(\mathbb{R}^n\) 上的局部 \(L^2\)（从而局部可积）可测函数；同理，\(Lf \in H^s(\mathbb{R}^n)\) 意味着 \(\langle\xi\rangle^s \widehat{Lf}(\xi) \in L^2(\mathbb{R}^n)\) 且 \(\widehat{Lf}\) 是局部可积可测函数。
\item
  \textbf{计算 \(\widehat{Lf}(\xi)\) 并应用椭圆性下界} 在 \(\mathcal{S}'(\mathbb{R}^n)\) 中，反复应用教材 p.318 命题 8.21(6)（\(\mathcal{F}(\partial_{x_j}\lambda) = 2\pi\mathrm{i}\xi_j \mathcal{F}\lambda\)），对任意 \(|\alpha| = m\) 有 \(\widehat{\partial^\alpha f}(\xi) = (2\pi\mathrm{i})^m \xi^\alpha \hat{f}(\xi)\)。线性叠加得 
\FitDisplayMath{\widehat{Lf}(\xi) = \sum_{|\alpha|=m} c_\alpha (2\pi\mathrm{i})^m \xi^\alpha \hat{f}(\xi) = (2\pi\mathrm{i})^m l(\xi)\hat{f}(\xi).}
 因 \(l(\xi)\) 为多项式且 \(\hat{f}, \widehat{Lf} \in L^1_{\mathrm{loc}}(\mathbb{R}^n)\)，上式作为可测函数在 \(\mathbb{R}^n\) 上几乎处处成立。两边取模长并利用椭圆性条件 \(|l(\xi)| \ge c|\xi|^m\)，得 
\FitDisplayMath{|\widehat{Lf}(\xi)| = (2\pi)^m |l(\xi)|\,|\hat{f}(\xi)| \ge c(2\pi)^m |\xi|^m |\hat{f}(\xi)| \quad \text{a.e. } \xi \in \mathbb{R}^n,}
 即 
\FitDisplayMath{|\xi|^m |\hat{f}(\xi)| \le \frac{1}{c(2\pi)^m}|\widehat{Lf}(\xi)| \quad \text{a.e. } \xi \in \mathbb{R}^n.}

\item
  \textbf{利用初等不等式估计 \(\langle\xi\rangle^{s+m}|\hat{f}(\xi)|\)} 对一切 \(\xi \in \mathbb{R}^n\)，显然 \(1 + |\xi|^2 \le 2\max\{1, |\xi|^2\}\)。因 \(m \ge 0\)，两边取 \(\frac{m}{2}\) 次幂，得 
\FitDisplayMath{\langle\xi\rangle^m = (1 + |\xi|^2)^{m/2} \le 2^{m/2}\max\{1, |\xi|^m\} \le 2^{m/2}\big(1 + |\xi|^m\big).}
 （当 \(m = 0\) 时，左端 \(= 1\)，右端 \(= 2\)，不等式亦显然成立。）

  将上式与步骤 2 的估计相乘结合，对几乎处处的 \(\xi \in \mathbb{R}^n\) 有 
\FitDisplayMath{\begin{aligned}
  \langle\xi\rangle^{s+m}|\hat{f}(\xi)| &= \langle\xi\rangle^s \langle\xi\rangle^m |\hat{f}(\xi)| \\
  &\le 2^{m/2}\langle\xi\rangle^s |\hat{f}(\xi)| + 2^{m/2}\langle\xi\rangle^s |\xi|^m |\hat{f}(\xi)| \\
  &\le 2^{m/2}\langle\xi\rangle^s |\hat{f}(\xi)| + \frac{2^{m/2}}{c(2\pi)^m}\langle\xi\rangle^s |\widehat{Lf}(\xi)|.
  \end{aligned}}

\item
  \textbf{取 \(L^2(\mathbb{R}^n)\) 范数得出结论} 由已知 \(f, Lf \in H^s(\mathbb{R}^n)\)，右端两项均属于 \(L^2(\mathbb{R}^n)\)。由 \(L^2(\mathbb{R}^n)\) 范数的单调性与三角不等式（Minkowski 不等式），立得 \(\langle\xi\rangle^{s+m}\hat{f} \in L^2(\mathbb{R}^n)\)，即 \textbf{\(f \in H^{s+m}(\mathbb{R}^n)\)}，且满足估计 
\FitDisplayMath{\begin{aligned}
  \|f\|_{H^{s+m}(\mathbb{R}^n)} &= \|\langle\xi\rangle^{s+m}\hat{f}\|_{L^2(\mathbb{R}^n)} \\
  &\le 2^{m/2}\|\langle\xi\rangle^s \hat{f}\|_{L^2(\mathbb{R}^n)} + \frac{2^{m/2}}{c(2\pi)^m}\|\langle\xi\rangle^s \widehat{Lf}\|_{L^2(\mathbb{R}^n)} \\
  &= 2^{m/2}\|f\|_{H^s(\mathbb{R}^n)} + \frac{2^{m/2}}{c(2\pi)^m}\|Lf\|_{H^s(\mathbb{R}^n)} \\
  &\le C\big(\|f\|_{H^s(\mathbb{R}^n)} + \|Lf\|_{H^s(\mathbb{R}^n)}\big),
  \end{aligned}}
 其中常数可显式取为 
\FitDisplayMath{C = 2^{m/2}\max\left\{1,\ \frac{1}{c(2\pi)^m}\right\}.}
 （该常数只依赖于阶数 \(m\) 与 \(L\) 的椭圆常数 \(c\)，自然满足题设要求。）证毕。
\end{enumerate}

\SubPillCoral{【考点复盘与避坑提示】}

\begin{itemize}
\tightlist
\item
  \textbf{核心考点}：实指数 Sobolev 空间 \(H^s(\mathbb{R}^n)\) 的频域范数定义（p.336）、椭圆算子的象征下界与\textbf{全局椭圆正则性先验估计（式 (8.49)）}。
\item
  \textbf{为什么需要式中的 \(\|f\|_{H^s}\) 项，不能只用 \(\|Lf\|_{H^s}\) 控制？} 因为当 \(\xi \to 0\) 时，\(|\xi|^m \to 0\)，\(Lf\) 根本管不住低频附近的 \(\hat{f}(\xi)\)（例如常函数满足 \(\Delta f = 0\)，但 \(f \neq 0\)）！所以必须由 \(\|f\|_{H^s}\) 兜底低频的「\(1\)」，由 \(\|Lf\|_{H^s}\) 压住高频的「\(|\xi|^m\)」，二者缺一不可！
\item
  \textbf{为什么非椭圆算子不成立（教材 p.337 注记 8.19）？} 以波算子 \(L = -\partial_t^2 + \Delta\) 为例，其象征 \(l(\tau, \xi) = \tau^2 - |\xi|^2\) 在特征锥 \(|\tau| = |\xi| \neq 0\) 上恒等于 \(0\)！即使 \(Lf = 0\)，解 \(\hat{f}\) 也可以把任意高频率的能量集中堆在特征锥 \(|\tau|=|\xi|\) 上，完全无法由 \(Lf\) 控制，因此估计 (8.49) 失效。
\item
  \textbf{往深一步（衔接练习 8.8.3）}：若 \(Lf = 0\) 且 \(f \in H^s(\mathbb{R}^n)\)，反复套用本题结论可得 \(f \in H^{s+m} \implies H^{s+2m} \implies \cdots \implies \bigcap_{k\ge 1} H^{s+km}\)，再由 Sobolev 嵌入定理 8.16 立得 \(f \in C^\infty(\mathbb{R}^n)\)——这就是 PDE 中大名鼎鼎的\textbf{椭圆正则性自举（Bootstrap）思想}！
\end{itemize}

\BookPart{}{第三部分：期末考前 3 分钟冲刺速查清单}{第三部分：期末考前 3 分钟冲刺速查清单}

\begin{enumerate}
\def\labelenumi{\arabic{enumi}.}
\tightlist
\item
  \textbf{做广义函数求导题}：先写「任取 \(\varphi \in \mathcal{D}(\mathbb{R})\)，设 \(\operatorname{supp}\varphi \subset [-R, R]\)」\(\to\) 甩锅求导加负号 \(\to\) 奇点处挖去 \((-\varepsilon, \varepsilon)\) 用控制收敛定理 \(\to\) 分部积分用中值定理消掉 \(\varepsilon\ln\varepsilon\) 边界项。
\item
  \textbf{做分布极限题}：优先试两种大招——\CircNum{1} 换元 + 控制收敛定理（第 9 题）；\CircNum{2} 找原函数证 \(L^1_{\mathrm{loc}}\) 收敛再求广义导数（第 14 题）。
\item
  \textbf{做 Fourier 变换与 PDE 题}：时刻检查本书 \textbf{\(2\pi\) 约定}——求导乘 \(2\pi\mathrm{i}\xi\)，反演与 Plancherel \textbf{不带} \((2\pi)^{-n}\)；背熟 Gauss 核与 Poisson 核两对变换；半空间方程对无界变量做变换，写衰减通解别忘 \(|\xi|\) 的绝对值。
\item
  \textbf{做 \(H^s(\mathbb{R}^n)\) 题}：第一时间切到频域看 \(\langle\xi\rangle^s \hat{f}(\xi) \in L^2(\mathbb{R}^n)\)；查积分收敛时\textbf{低频 \(r \to 0^+\) 和高频 \(r \to +\infty\) 两端都必须查}（第 28 题）；证椭圆估计用 \(\langle\xi\rangle^m \le 2^{m/2}(1 + |\xi|^m)\) 一分为二（\hyperlink{prob:8-882}{\textcolor{LinkColor}{练习 8.8.2}}）。
\end{enumerate}


\end{document}
